\chapter{The C2 cases}
\label{ch:c2-short}
All source coefficients lie in the real closed field $K_c$, and all
testing measurements and states lie in the independently labelled real
closed field $K_t$. We use $K_c\otimes_{\mathbb R}K_t$ without an
algebraic-disjointness assumption. A retained positive square may have
unbounded pure-source coefficients. Only the matrix coefficients to
which the bounded strict-Gram lemma is applied must be bounded.
The native physical entry, bounded-coefficient theorem, and arbitrary-
scale noncommuting-residue theorem are the common preceding results.
Normalize $d=1$ by a pure-source square and write $P=F/4$. The common
four-state parent is
\begin{equation}
\begin{split}
F=4P={}&4nD_uD_vX^2+4mM^2+4CD_v(uW-Y)^2+4AD_uD_vY^2\\
&+4DJ^2+4\alpha D_u(vY-Z)^2\\
&+2\{J^2+D_u(vY-Z)^2+D_v(X+uZ)^2-D_uD_vZ^2\},\\
M={}&vW-X+uvY-uZ,\quad J=uW+uvX-Y-vZ,\quad
D_u=1+u^2,\quad D_v=1+v^2.
\end{split}\label{eq:c2-master-parent}
\end{equation}
The finite coefficient-extraction and monic-division operations are
those in Section~\ref{sec:c1-finite-operations}; they are field-independent.
For every displayed state substitution $s=s_0+A\xi$, every scalar,
complete first normal and hard coefficient below means respectively
$[1]\mathcal Q$, $[\xi_i]\mathcal Q$, and
$\frac12\partial_i\partial_j\mathcal Q$, where
$\mathcal Q=(F-\Pi)(s_0+A\xi)$ and $\Pi$ is the displayed sum of retained
positive squares. Quadratic state degree makes this the exact whole
polynomial. Coefficients need not be expanded to define them. All
support, divisibility, source-bound and residue verifications remain
part of the relevant calculation below.

The proof follows the source geometry. We first determine the stationary
coefficients and treat two weak complementary events by a weighted radial
completion. The remaining families are organized by the scales of the actual
angles. For one small angle the state parameter separates bounded and
unbounded source scales; for two small angles the source ratios define
R1--R10. Within each family the scalar calculation is completed by the
state-normal and complementary-testing arguments.

\section{Stationary sources and their classification}
\subsection{Source domains}
\label{component:pe351eb06b28b}
\label{pe351eb06b28b:sec:c2-complete-source-routing}
The native entry theorem supplies an actual physical stationary source
with $d=1$, $\lambda_{\rm event}=\beta=0$, $n\ge0$, and positive
$m,C,A,D,\alpha$. In the native coordinates $u_c=c$, $v_c=ac$ its
strict gaps are
\[
 c,a,b,t>0,\quad 0<r<1,\quad h=b-a^2t>0,\quad
 g=(1+c^2)b-(1-r+a^2c^2t)>0,
\]
\[
 j=1+c^2r-t>0,\qquad1-c^2h>0,\qquad b-a^2c^2r>0.
\]
These are the original source equations. If
$(1+n+m+C+A+D+\alpha)^{-1}$ is a unit, the bounded-coefficient
theorem applies. Henceforth that reciprocal is infinitesimal.
Call an event weak when its coefficient is infinitesimal and strong
otherwise; strong coefficients may be unbounded.

The exact reflection
\begin{equation}
 (u,v;W,X,Y,Z;n,m,C,A,D,\alpha)
 \longmapsto(1/v,1/u;Z,X,Y,W;n,D,\alpha,A,m,C)
 \label{eq:c2-short-reflection}
\end{equation}
satisfies $P=u^2v^2P'$ and is a homogeneous polynomial identity at
the literal endpoints. With $z=c^2$ its source map is
\[
 c'=(ac)^{-1},\ a'=a,\ b'=a^2/b,\ r'=a^2zr/b,
 \quad t'=(1-zh)/b,
\]
\[
 h'=a^2zh/b,\quad g'=a^2j/b,\quad j'=g/b,
 \quad\Delta'=\Delta/(b^2z^2),\quad N'=a^2N/b^2.
\]
The remaining three strict gaps are positive by the original gaps.
This map preserves the total coefficient scale. In particular it
orients an exactly-one-weak source to $C$ weak and $\alpha$ strong.

Write $\mathrm I,\mathrm U,\mathrm O$ for positive infinitesimal,
bounded-unit, and unbounded finite actual angles. The complete cases are
\[
\begin{array}{c|c|c}
 (u_c,v_c)&C,\alpha\text{ strong}&C\text{ weak},\alpha\text{ strong}\\\hline
 \mathrm{II}&\text{II calculation}&\varnothing\\
 \mathrm{IU}&\text{one small angle}&\varnothing\\
 \mathrm{IO}&\text{strong IH}&\text{weak IH}\\
 \mathrm{UI}&\text{strong UI}&\varnothing\\
 \mathrm{UU}&\text{five positive events}&\varnothing\\
 \mathrm{UO}&\text{reflected IU}&\text{reflected IU}\\
 \mathrm{OI}&\text{strong HI}&\varnothing\\
 \mathrm{OU}&\text{reflected UI}&\varnothing\\
 \mathrm{OO}&\text{reflected II}&\text{reflected weak II}.
\end{array}
\]
Each source domain is treated below for all independent testing angles and
states. The following bounds establish that these domains exhaust the source.

If $C,\alpha$ are both weak, the four event-cycle physicality tests
force $n,m,A,D$ unbounded. One test is
$4D\varepsilon^2-2\varepsilon+4(2+\varepsilon)C
+8(1+\varepsilon)\alpha$: a sufficiently small positive real
$\varepsilon$ excludes bounded $D$. The matrix-square core bound gives
$H\ge-2u_cv_c\|\psi_c\|^2$. Thus each matching event amplitude divided
by $\sqrt{u_cv_c}$ is infinitesimal. Compact source angles and the
reflection imply $c$ bounded, $a,b,t\sim1$, and
$r,h,g,j\in\mathrm I$. For $c\in\mathrm I$ the relative matching
equations give $v_c/u_c\sim1$ and $W_c/Z_c\sim1$ before scaled
stationarity is used. The radial calculations give
$N\asymp r+h+g+j$ and $4\alpha/(c^4N)\sim1$; they retain all four
state amplitudes and both state charts. For $c$ a unit,
$\alpha/C\asymp1$. The outer calculations in both cases use
$1/(Cn),1/(Cm),1/(CA),1/(CD)=O(1)$, with their reflected counterparts
when $c$ is a unit. This covers all two-weak sources.

The empty cells follow from the original source inequalities.
For II, $L<y(1+zr)$ and the positive source numerator show that weak
$C$ forces weak $\alpha$. For UI the same numerator and $n\ge0$
would force $r$ infinitesimal and $h$ a unit, contradicting weak $C$.
For IU, bounded $b$ obeys
\[
 r_0\{b_0+k_0^2(b_0+1)\}\le b_0,\quad
 r_0\ge1-b_0,\quad b_0\ge{k_0\over\sqrt{1+k_0^2}}>0,
 \qquad k=v_c,
\]
and gives strong $C$. For unbounded $b$, $C\asymp L^{-1}$ with
$L$ bounded, again giving strong $C$. For UU the positive numerator
and unit source denominator make $C,\alpha$ weak or strong together.
Finally, for OI or OU put $q=u_c^{-2}\in\mathrm I$ and $y=v_c^2$
bounded. The exact $n$ numerator is $-qP/(2rQ)$, $Q>0$, and
\[
 P\ge r^2(1+y)(H+yr)-Tqr(y+2)-q^2(T^2+T),\quad0<T<H+y.
\]
If $q/r$ were infinitesimal, $P>0$ would contradict $n\ge0$.
Thus $r/q$ is bounded, $H\sim1$, and the canceled $C$ numerator
is a unit with bounded positive denominator. These arguments include
equality $n=0$.

For IU, bounded $b$ is a unit by the preceding bound. Its exhaustive
opening cases are $g$ a unit; $g$ infinitesimal and $r$ a unit; and
$g,r$ both infinitesimal. The three one-angle calculations retain,
respectively, the ordinary hard normal, the full $m$ normal, and both
centered $n,m$ squares. If $b$ is unbounded, put
$x=u_c^2$, $B=xb$, $H=xh$, $L=x(b-h)$. The gaps imply
$0<H<1$, $0<L<k^2(1+xr)$ and $B=H+L$ bounded. Exact cancellation gives
\[
 1-r\asymp1,\quad m\asymp b^{-2},\quad C\asymp L^{-1},
 \quad A\asymp H^{-1},\quad D\asymp B/j,\quad
 \alpha\asymp B/(1-H).
\]
For infinitesimal $B$ use the homogeneous one-angle calculation; for
unit $B$ use the strong-event calculation, which includes unbounded
$C$ or $\alpha$. The reflection supplies every UO source.

On strong II the source equations give $B$ a unit, $H,1-r$ units, and
\[
 m/z^2,\quad A/y,\quad CL,\quad Dyj,\quad\alpha(1-H)\asymp1,
 \qquad nr=O(1).
\]
These are the coefficient assumptions of the shared II local and
outer calculations. For IO the weak IH source has $B$ unbounded and
$\rho=B/y$ with residue in $[0,1]$; its interior, two endpoints and
open-axis calculations exhaust that interval. The strong IH source
instead has $B$ a unit, $r=O(1/y)$ and $Dy^2\asymp1$ and uses the
same polynomial relations with its own coefficient bounds. Strong UI
has bounded $nr,m\mu,\alpha\pi,C\rho k^2,D\sigma k^2,Ct-Dj$,
which are the inputs of its three-normal calculation. On OI, direct
source cancellation gives $R=r/q$ bounded, $H\sim1$ and
\[
 mqK,\quad Cyt,\quad Dyj,\quad\alpha qE,\quad A/(qy)\asymp1.
\]
The HI circuit calculation includes every $n\ge0$, including $R_0=0$;
its multiplier $2(1+R_0)^2/(3+R_0)$ lies in $[2/3,2]$. On strong UU,
the source equations force $m,A,D$ strong, so the five-event proof
uses the positive real lower bound of $\min(1,m,C,A,D,\alpha)$.
The remaining reflected entries follow from
\eqref{eq:c2-short-reflection}. The directed weak OO entry is the
following weak-II calculation.

\subsubsection{The directed weak II source}
\label{component:pec6c617d3a3d}
\label{pec6c617d3a3d:sec:weak-c2-complete-join-rederived}
Use $z=u_c^2$, $y=v_c^2$, $H=v_cY_c$,
$L=v_c(u_cW_c-Y_c)$, $B=H+L$, $r=-X_c/u_c$.
The source equations are \eqref{eq:c2-ba-source}, with
$\delta=B\Delta$ and the positive numerator
\[
 M_s=B^2M=(1+z)(B+zr)H+yG+yzj+zrL(1+y),
 \qquad G=B+z(H+r-1).
\]
Since $L<y(1+zr)$, $L$ is
infinitesimal and $B$ bounded. If $B$ is infinitesimal, the exact bound
\[
 0<\alpha\le{(B+zr)(1+z)(H+y)\over2(1-H)}\in\mathrm I
\]
uses $\delta\ge B+z$. If $B$ is a unit, $H=B-L$, $\delta$ and the
positive numerator are units, giving strong $\alpha$. Thus weak
$\alpha$ is exactly $B\in\mathrm I$ in this chart.

Set $a=z/B$, $t=y/B$, $h=H/B$, $l=L/B=1-h$, $e=1-r$.
The exact gaps give
\[
 g=1-ae+hz>0,\quad j=1+zr-l/t>0,\quad ae<1+hz.
\]
Hence unit $e$ forces bounded $a$; unbounded $a$ forces infinitesimal
$e$; infinitesimal $a$ gives $g\sim1$; and infinitesimal $g$ gives
$ae\sim1$. The source formula gives
\[
 r(2n+1)\asymp1,\quad AH/y\asymp1,\qquad
 m={ (1+y)(1+zr)S_c\over2}{z\over\delta}
                       {1+zH/B-g\over g}.
\]
Thus $m$ is bounded for unit $g$, and $m\asymp1/g$ for
infinitesimal $g$. The positive source numerator also gives the
unbounded-coefficient equivalences
\[
\begin{array}{c|ccccc}
 \text{coefficient}&n&m&A&C&D\\\hline
 \text{condition}&r\in\mathrm I&g\in\mathrm I&H/y\in\mathrm I&
 L/(H+y)\in\mathrm I&t\in\mathrm I\text{ or }j\in\mathrm I.
\end{array}
\]
For $t$ unbounded, $j\sim1$ and $D=O(1/t)$.
These estimates concern the original source.

Split $t$ first. Infinitesimal $t$ is R1 and unbounded $t$ is R2,
proved below with their complete scalar and both-normal calculations.
For unit $t$ and bounded $a$, Proposition~\ref{prop:c2-bounded-a}
proves all cases at once.
For unit $t$ and unbounded $a$, necessarily $e$ is infinitesimal,
$r\sim1$. The remaining exhaustive cases are
\[
\begin{array}{c|l}
 h\text{ unit},\ l\in\mathrm I&\text{R3, the EN calculation}\\
 h,l\text{ units},\ j\in\mathrm I&\text{R4, the EN calculation}\\
 h,l,j\text{ units},\ g\in\mathrm I&\text{R6, RI or RB}\\
 h,g\in\mathrm I&\text{R7, RU}\\
 h,j\in\mathrm I,\ g\text{ unit}&\text{R8, the joint-D calculation}\\
 h\in\mathrm I,\ g,j\text{ units}&\text{R10, the weighted/thin calculation}.
\end{array}
\]
Indeed if $h,l,j,g$ are all units, all original coefficients are
bounded by the displayed identities, contradicting the assumed total
scale. Infinitesimal $r$ is impossible when $a$ is unbounded. No
other ordered alternative remains.

The R6 source inequality forces $b=a^2B$ bounded, and its two cases
$b$ unit or infinitesimal are RB and RI. RU treats both $b$ cases for
R7. In R8 and R10, split $n/e^2$ into noninfinitesimal, positive
infinitesimal, and zero. The first uses the weighted source Schur
bound; the second uses the explicitly paired positive coefficients and
the full contact-normal correction; equality uses the exact OP4
native C1 map. OP4 preserves the physical parent and the actual
contact, so only the native C1 small-$e$ or centered-$A$ theorem is
used at $n=0$. The positive-thin calculation never declares its
formal negative-$\beta$ tuple physical. These source splits and the
finite and homogeneous outer calculations below exhaust every
independent testing pair and state.

\subsection{Stationary coefficients and contact}\label{component:p6e5d41ed80b3}\label{app:c2-source}
In this appendix, $c$ denotes the actual half-angle and
$C$ the event weight previously denoted $c$. Set
\[
 z=c^2,\quad A_c=1+zr,\quad
 h=b-a^2t,\quad g=(1+z)b-(1-r+a^2zt),\quad j=1+zr-t,
\]
\begin{align*}
\Delta={}&(1+z)zb^2+2z(1+z)br+(1+z)(1+a^2z)b\\
 &+z^2(1-a^2)r^2+2z(1+a^2z)r+1+a^2z,\\
N={}&(1+z)(b+r)(b+a^2(1-t))-a^2(1+a^2z)(1-r)t.
\end{align*}
The unchanged $d=1$, $\lambda=\beta=0$ source is
\begin{equation}\begin{aligned}
n&=\frac{(b+r)A_c(A_c+a^2zt)-r\Delta}{2r\Delta},&
m&=\frac{(1+a^2z)(1-r)A_c(A_c+a^2zt)}{2g\Delta},\\
C&=\frac{A_cN}{2a^2t\Delta},&
A&=\frac{a^2A_c(b+r)(A_c+a^2zt)}{2h\Delta},\\
D&=\frac{(1+z)(b+r)(b-a^2zr)(A_c+a^2zt)}{2a^2j\Delta},&
\alpha&=\frac{z^2(b+r)N}{2(1-zh)\Delta}.
\end{aligned}\label{eq:c2-native-source}\end{equation}
Its exact contact is
\begin{equation}
 (u_c,v_c;W_c,X_c,Y_c,Z_c)
   =\left(c,ac;\frac ba,-cr,c\left(\frac ba-at\right),1\right).
\label{eq:c2-native-contact}\end{equation}
The positive chart requires $a,c,b,t>0$, $0<r<1$, $h,g,j>0$,
$1-zh>0$, $b-a^2zr>0$, and nonnegative original event weights.
The source-entry proof uses this positive chart without a shape-unit
assumption.

The numerator has the useful positive identity
\[
N=(1+c^2)(b+r)h+a^2\{g+j+rt(1+a^2c^2)\}>0.
\]
It follows by substituting the three source-gap definitions; every term
is positive in the domain used by the contact-entry theorem.
For the determinant calculation use the aliases
\[
 R_2=R_m,\quad R_3=uW-Y,\quad R_5=R_D,\quad R_\alpha=vY-Z,
 \qquad(e_1,e_2,e_3,e_4,e_5)=(n,m,C,A,D).
\]

\subsubsection{Uniqueness of the stationary coefficients}\label{component:p61209c7390b0}
\label{sec:c2-contact-determinant}

The six stationary coefficients are unique, as follows.
Let $M$ have rows indexed by $(W,X,Y,Z,u,v)$ and columns equal to the
derivatives of
\[
D_uD_vX^2,\quad R_2^2,\quad D_vR_3^2,\quad
D_uD_vY^2,\quad R_5^2,\quad D_uR_\alpha^2,
\]
evaluated at $(u,v;W,X,Y,Z)=(c,ac;b/a,-cr,c(b/a-at),1)$.
Write $z=c^2$, $D=1+z$, $E=1+a^2z$, $h=b-a^2t$,
$g=Db-(1-r+a^2zt)$, and $j=1+zr-t$.
Factoring the six columns by
\[
2cr,\quad2cg,\quad2act,\quad-2ch/a,\quad2acj,
\quad2(1-zh)
\]
and the first row by $c$ leaves the matrix
\[
B=\begin{pmatrix}
0&a&E&0&-1&0\\
-DE&-1&0&0&-az&0\\
0&az&-E&-DE&1&-acD\\
0&-c&0&0&ac&D\\
zrE&zh-1&(b/a)E&-zhE/a&azr-b/a&c(1-zh)\\
azrD&b/a+zh/a&a^2zt&-zDh&1+zr&-cDh/a
\end{pmatrix}.
\]
Use the pivots in row $2$, column $1$; row $1$, column $3$; and
row $4$, column $6$.  Eliminate those three columns, multiply the
remaining second and third rows by $D$ and $E$, and factor $-DE$
from their middle column.  The determinant reduces to
$\det B=D^2E^2\det T$, where
\[
T=\begin{pmatrix}
a&1&a^2z\\
-1-b-z(r+a^2t)&zh/a&az(r-1+zh)\\
b/a+az(h-r)&zh&E(1+zh)+z(r+a^2t)
\end{pmatrix}.
\]
Subtract $(zh/a)$ times the first row from the second and $zh$ times
the first row from the third.  Expanding along the middle column gives
\[
\det T=[1+b+z(b+r)][1+z(a^2+b+r)]
-z(1-r)(b-a^2zr)=\Delta.
\]
The last equality follows by expanding the two products and is exactly
the source definition of $\Delta$.  Restoring the factors proves
\begin{equation}
\det M=-64a c^6rgt h j(1-zh)D^2E^2\Delta.
\label{eq:c2-entry-determinant}
\end{equation}
Each source coefficient occurs with a factor $4$ in the parent, so the
stationary linear system in $(e_1,e_2,e_3,e_4,e_5,\alpha)$ has determinant
$4^6\det M$.  On the positive source domain this is nonzero.  Since the
six source formulas above satisfy all six stationary equations,
they are the unique stationary coefficients.

\section{Two weak complementary coefficients}
\subsection{Entry from two weak coefficients}
\label{component:p534f57e1da6f}

\label{p534f57e1da6f:sec:chart-free-c2}

Here weak means infinitesimal after division by the original positive
signed coefficient \(d\). All source calculations take place in one real
closed coefficient field. Testing fields remain independent.

\begin{theorem}[C2 entry without a contact-chart hypothesis]
\label{p534f57e1da6f:thm:chart-free-c2}
Let the physical canonical parent have \(\lambda=\beta=0\), \(d>0\),
and nonnegative original event coefficients. Suppose it has an actual
finite noncommuting stationary zero contact, and
\[
                    C/d,\ \alpha/d\quad\text{are infinitesimal}.
\]
Then its full parent is tensor SOS at every independent real or complex
projective test. No hypothesis on the other four event weights, on the
contact-angle ratio, or on the signs of the source gaps is imposed.

More precisely, normalize \(d=1\). After local real sign gauges and at
most the exact C2 reflection,
the actual contact enters the positive rational C2 source with
\[
 c\text{ bounded},\qquad a,b,t\longrightarrow1,\qquad
 r,h,g,j>0\text{ infinitesimal}.
\]
In particular all four matching weights are unbounded and both
complementary weights are strictly positive on this finite-contact domain.
\end{theorem}

The proof has three ingredients: a finite physical witness makes the
matching weights large; the quadratic CHSH excess controls the contact
amplitudes at the scale of the actual angles; stationarity fixes their
signs. This last step establishes the source chart rather than assuming it.

\subsubsection*{The physical bound at the correct angular scale}
For \(u,v>0\), put \(D_u=1+u^2,D_v=1+v^2\), and let \(M(u,v)\) be the
real symmetric matrix of the signed core
\[
H=2\{R_D^2+D_u(vY-Z)^2+D_v(X+uZ)^2-D_uD_vZ^2\},
\]
where
\[
R_m=vW-X+uvY-uZ,\qquad R_D=uW+uvX-Y-vZ.
\]
Write \(L=D_uD_v\). Direct matrix multiplication gives
\begin{equation}
(M-LI)^2=L^2I+4uvLJ,\qquad
J=\begin{pmatrix}0&0&0&1\\0&0&-1&0\\0&-1&0&0\\1&0&0&0\end{pmatrix},
\qquad J^2=I.                                      \label{p534f57e1da6f:eq:cf-core-square}
\end{equation}
For a unit real state \(\psi\), Cauchy--Schwarz therefore gives
\[
\psi^TM\psi\ge L-\sqrt{L^2+4uvL}\ge-2uv.
\]
The last inequality follows by squaring the positive quantities
\(L+2uv\) and \(\sqrt{L^2+4uvL}\); their squared difference is
\(4u^2v^2\). Homogeneity proves
\begin{equation}
                         H\ge-2uv\,\|\psi\|^2.        \label{p534f57e1da6f:eq:cf-core-bound}
\end{equation}
This is the quadratic CHSH bound in the present normalization.

At a normalized zero contact, every event term is nonnegative. Define
the four matching amplitudes, in the fixed order \((n,A,m,D)\), by
\[
\mathcal L=(\sqrt{D_uD_v}X,\sqrt{D_uD_v}Y,R_m,R_D).
\]
Equation~\eqref{p534f57e1da6f:eq:cf-core-bound} gives, for each matching event,
\begin{equation}
                   e_i\mathcal L_i^2\le uv/2.         \label{p534f57e1da6f:eq:cf-amplitude-bound}
\end{equation}
All square roots and inequalities in this paragraph are coefficient-side.

\subsubsection*{Matching weights and the compact contact geometry}
First every matching weight is unbounded. To see the mechanism for
\(D\), use the coefficient-field experiment
\[
 u=1,\quad v=\sqrt{1+\delta},\quad
 (W,X,Y,Z)=(1,0,0,\sqrt{1+\delta}),\qquad\delta>0.
\]
The \(n,A,m\) positive square terms vanish, and the full eight-event parent is exactly
\[
4D\delta^2-2\delta+8\lambda+4(2+\delta)C
 +8(1+\delta)\alpha+4(2+\delta)(1+\delta)\beta.
\]
If \(D\le B\) for a positive real \(B\), choose the positive real
\(\delta=1/[4(B+1)]\). The first two terms have a strictly negative
real upper bound, whereas the complementary terms are infinitesimal.
This contradicts physicality. The exact even rotations of the event
cycle give the same argument for \(n,A,m\), preserving the sum of the
four complementary weights.  Hence \(1/n,1/A,1/m,1/D\) are infinitesimal.

Normalize the actual state to norm one and use positive actual
half-angles, obtainable by local real basis gauges. The weaker
bound \(H/(D_uD_v)\ge-2\) shows that all normalized matching
probabilities vanish in residue. In particular \(X,Y\to0\).
Put
\[
c_u=(1+u^2)^{-1/2},\ s_u=uc_u,\qquad
c_v=(1+v^2)^{-1/2},\ s_v=vc_v.
\]
If \(\sigma=\operatorname{res}(c_us_v)\) and
\(\tau=\operatorname{res}(s_uc_v)\), the two remaining matching
equations read
\[
\sigma\overline W-\tau\overline Z=0,\qquad
\tau\overline W-\sigma\overline Z=0.
\]
At least one of \(\overline W,\overline Z\) is nonzero. Thus
\(\sigma=\tau\ge0\). If positive, the two actual half-angles are
bounded units with the same residue. If zero, both are infinitesimal
or both unbounded, by the two unit-circle identities. The exact C2
reflection sends \((u,v)\) to \((1/v,1/u)\) and exchanges
\((C,\alpha)\) and \((m,D)\), preserving the hypotheses. We may
therefore assume either ordinary bounded-unit angles or two
infinitesimal angles.

\subsubsection*{Relative coalescence when both angles vanish}
Suppose \(u,v\) are infinitesimal. From
\eqref{p534f57e1da6f:eq:cf-amplitude-bound}, all four
\(\mathcal L_i/\sqrt{uv}\) are infinitesimal. In particular
\(X/\sqrt{uv},Y/\sqrt{uv}\to0\). Set \(a=v/u>0\).
Dividing the two other matching amplitudes by \(\sqrt{uv}\) gives
\[
\sqrt a\,W-\frac Z{\sqrt a}\longrightarrow0,\qquad
\frac W{\sqrt a}-\sqrt a\,Z\longrightarrow0.
\]
The discarded terms are \(X/\sqrt{uv},Y/\sqrt{uv}\) and
\(\sqrt{uv}X,\sqrt{uv}Y\), all infinitesimal.
If \(a\) were unbounded, these equations would force both \(W,Z\)
infinitesimal; the same follows if \(a\) is infinitesimal.
Both contradict normalization. Thus \(a\) is a bounded unit.
Taking residues now yields \(a\to1\) and \(W/Z\to1\), with
both \(W,Z\) units. Consequently
\[
                  X/u,\ Y/u,\ R_m/u,\ R_D/u\longrightarrow0.
\]

Normalize \(W=1\) by a pure coefficient factor. Write
\[
x=X/u,\quad y=Y/u,\quad q=v/u,\qquad
(u,x,y,Z,q)\longrightarrow(0,0,0,1,1).
\]
The core divided by \(u^2\) is the polynomial
\begin{align}
\mathcal H:=H/u^2=2\{&
(1-y-qZ+u^2qx)^2+(1+u^2q^2)(x+Z)^2 \notag\\
&+(1+u^2)(u^2q^2y^2-2qyZ-q^2Z^2)\}. \label{p534f57e1da6f:eq:cf-scaled-core}
\end{align}
It has zero residue at this point. The nonnegative alpha positive square term and
the zero-contact equation give
\[
0\le4\frac{\alpha}{u^2}(1+u^2)(u^2qy-Z)^2
                   \le-\mathcal H.
\]
Its square factor is a positive unit. Hence
\begin{equation}
                         \alpha/u^2\longrightarrow0. \label{p534f57e1da6f:eq:cf-alpha-scale}
\end{equation}
This estimate is derived before the scaled stationary equations are used.

\subsubsection*{Stationarity fixes the source signs}
In the small-angle case introduce the four exact local coordinates
\[
\ell=\left(\sqrt{D_uD_v}\,x,\sqrt{D_uD_v}\,y,
 q-x+u^2qy-Z,\ 1+u^2qx-y-qZ\right).
\]
Their Jacobian with respect to \((x,y,Z,q)\) has residue
\[
J_0=\begin{pmatrix}
1&0&0&0\\0&1&0&0\\-1&0&-1&1\\0&-1&-1&-1
\end{pmatrix},\qquad \det J_0=2.
\]
Thus its inverse and the transformed derivatives are bounded at the
actual contact. Only this pointwise coefficient-field inverse is used.
Direct differentiation of \eqref{p534f57e1da6f:eq:cf-scaled-core} gives
\[
\operatorname{res}(\nabla_\ell\mathcal H)=2(1,-1,-1,1).
\]
The full parent divided by \(u^2\) is exactly
\[
4\sum_i e_i\ell_i^2
 +4C(1+u^2q^2)(1-y)^2
 +4(\alpha/u^2)(1+u^2)(u^2qy-Z)^2+\mathcal H.
\]
The two complementary-gradient terms have zero residue, by weak \(C\)
and \eqref{p534f57e1da6f:eq:cf-alpha-scale}. The original stationary equations imply
stationarity in \(\ell\), with \(u\) fixed. These equations first show
that every \(e_i\ell_i\) is bounded, so its residue is defined. Consequently
\begin{equation}
\operatorname{res}(n\ell_n,A\ell_A,m\ell_m,D\ell_D)
                         =\tfrac14(-1,1,1,-1).        \label{p534f57e1da6f:eq:cf-sign-residue}
\end{equation}
Every matching weight is positive, so \(X<0,Y>0,R_m>0,R_D<0\)
as exact source inequalities.

For ordinary angles the argument is shorter. Normalize \(W=1\);
the actual contact has residue
\((X,Y,Z,v)=(0,0,1,u)\), where \(u\) is a bounded unit.
Use the unscaled \(\mathcal L\) coordinates. Their Jacobian determinant
at that residue is \(2u(1+u^2)^2\), and
\[
\nabla_{\mathcal L}H=2u(1,-1,-1,1)
\]
there. All complementary gradients are bounded times \(C,\alpha\),
hence infinitesimal. Stationarity gives
\(\operatorname{res}(e_i\mathcal L_i)
 =\overline u\,(-1,1,1,-1)/4\), with the same exact signs.

\subsubsection*{The original source and all independent tests}
In either case \(W/Z\to1\), \(v/u\to1\), and \(X/u,Y/u\to0\).
Set, using the original actual contact,
\[
c=u,\quad a=v/u,\quad b=aW/Z,\quad
r=-X/(uZ),\quad t=(uW-Y)/(vZ).
\]
Then \(a,b,t\to1\) and
\[
h=b-a^2t=\frac{aY}{uZ}>0,\qquad
g=\frac{R_m}{uZ}>0,\qquad j=-\frac{R_D}{vZ}>0.
\]
All four \(r,h,g,j\) are infinitesimal, so \(0<r<1\).
The remaining chart boundaries are positive units:
\[
1-c^2h=\frac{Z-vY}{Z}>0,\qquad
b-a^2c^2r=\frac{a(W+vX)}Z>0.
\]
These definitions give exactly
\[
h=b-a^2t,\qquad
g=(1+c^2)b-(1-r+a^2c^2t),\qquad j=1+c^2r-t.
\]

Finally the stationary coefficient matrix has the nonzero determinant
\[
4^6[-64ac^6rgt\,hj(1-c^2h)(1+c^2)^2(1+a^2c^2)^2\Delta].
\]
The preceding three-by-three reduction gives this determinant.
For \(z=c^2\), its positive expansion is
\[
\Delta=1+(1+z)b+z(1+z)b^2+2z(1+z)br+z^2r^2+2zr
       +a^2z\{(1+z)b+1+zr(2-r)\}\ge1.
\]
The original coefficients
therefore equal the unique positive rational C2 source formulas.
In particular their positive numerator identity forces \(C,\alpha>0\).
If \(c\) is infinitesimal, apply the small-angle radial theorem and its
outer-test corollary; if \(c\) is a unit, apply the ordinary-angle radial
corollary. Reflect back when necessary and restore \(d\) by a pure-source
square root. This proves Theorem~\ref{p534f57e1da6f:thm:chart-free-c2}.

This covers the finite-contact two-weak branch. Weak \(C\) alone need
not imply weak \(\alpha\); the one-weak family below supplies the
complementary case.

\subsection{Matching the weak coefficients to the strong events}
\label{component:p8b53b68d8af6}

\label{p8b53b68d8af6:sec:c2-matching-strength}

A coefficient condition symmetric in the two weak positive square terms closes both
commuting angle corners at once.  It also gives a second exact radial
source family whose contact angles have ordinary positive residue.

\begin{theorem}[Two matched weak positive square terms]
\label{p8b53b68d8af6:thm:c2-matching-strength}
Let the original physical C2 parent have $d=1$, $\lambda=\beta=0$,
and event weights $(n,m,C,A,D,\alpha)$.  Suppose $C,\alpha>0$ are
infinitesimal, all four other weights are positive, and each of
\[
 Cn,\ Cm,\ CA,\ CD,\qquad
 \alpha n,\ \alpha m,\ \alpha A,\ \alpha D
\]
has bounded reciprocal in the coefficient field.  Then the full parent
is tensor SOS at every independent projective test and every real or
complex state.
\end{theorem}
\begin{proof}
Lemma~\ref{pc7dacaa3258e:lem:c2-four-strong-outer} applies directly to every commuting
residue test except the finite zero-zero corner.
For that corner use the exact reflection
\begin{equation}
 (u,v;W,X,Y,Z;n,m,C,A,D,\alpha)
 \longmapsto
 (1/v,1/u;Z,X,Y,W;n,D,\alpha,A,m,C).
\label{p8b53b68d8af6:eq:c2ms-reflection}
\end{equation}
More explicitly, if $F$ denotes the full parent with its indicated
arguments, then
\begin{equation}
 F(u,v;W,X,Y,Z;n,m,C,A,D,\alpha)
 =u^2v^2F(1/v,1/u;Z,X,Y,W;n,D,\alpha,A,m,C).
\label{p8b53b68d8af6:eq:c2ms-reflection-identity}
\end{equation}
This is a homogeneous polynomial identity, so its interpretation in
reciprocal charts includes literal zero and infinity endpoints.

For a short verification, the reflected $R_2,R_5$ amplitudes, after
multiplication by $uv$, are $-R_5,-R_2$.  The reflected $C$ and
$\alpha$ positive square terms exchange, while the $n,A$ positive square terms are fixed.
If $\mathcal H$ is the signed core, its common expanded form is
\begin{align*}
 \mathcal H/2={}&u^2(W^2+Z^2)+(1+v^2+u^2v^2)(X^2+Y^2)\\
 &+2u^2v(WX-YZ)+2u(XZ-WY)-2uv(WZ+XY).
\end{align*}
Multiplying its reflected expression by $u^2v^2$ fixes the first,
second and last terms, and exchanges the middle two.  This proves the
whole-parent identity, including its signed part.

The reflected weak weight is $\alpha$, its four strong weights are
$n,D,A,m$, and its other nonnegative weak positive square term is $C$.
Thus all hypotheses of Lemma~\ref{pc7dacaa3258e:lem:c2-four-strong-outer} hold for
the reflected parent.  Its infinity-infinity chart is precisely the
original zero-zero chart.  Restore the homogeneous pure-testing
factors to obtain the certificate there as well.

For noncommuting residue tests apply the physical-source theorem, using
positive pure-source normalization when a coefficient is unbounded.
Restore that factor and split real and imaginary state parts.
\end{proof}

\begin{corollary}[Radial C2 source with ordinary contact angles]
\label{p8b53b68d8af6:cor:c2-ordinary-angle-radial}
Use the same rational source \eqref{eq:c2-native-source}, now with source
half-angle $c$ a positive bounded unit.  Suppose $h=b-a^2t,r,g,j$ are
positive infinitesimals and $a,1-r,t$ are positive bounded units.
Then the full parent is tensor SOS at every independent projective test
and every real or complex state.
\end{corollary}
\begin{proof}
Put $z=c^2$ and continue to write $C=e_3$ for the weak event weight.
The exact source identities
\[
 b-1=-r+g-zh,\qquad t-1=zr-j,\qquad
 a-1=\frac{g+j-(1+z)(r+h)}{t(a+1)}
\]
remain valid for this ordinary source angle.  Thus $a,b,t$ have residue
one.  The positive source identity
\[
 N=(1+z)(b+r)h+a^2\{g+j+rt(1+a^2z)\}
\]
shows that $N>0$ is infinitesimal and $h/N,r/N,g/N,j/N$ are bounded:
each of the four positive coefficients in this sum is a bounded unit.

Here is the entire normalized coefficient calculation.  Write
$z_0=\overline z>0$.  Substitution of $a=b=t=1,r=0$ in the bounded
source numerators and unit denominators gives
\begin{gather}
 \overline\Delta=2(1+z_0)^2,\qquad
 \overline{C/N}=\frac1{4(1+z_0)^2},\qquad
 \overline{\alpha/C}=z_0^2,\notag\\
 \overline{rn}=\overline{hA}=\frac1{4(1+z_0)},\qquad
 \overline{gm}=\overline{jD}=\frac14.
\label{p8b53b68d8af6:eq:c2ms-ordinary-residues}
\end{gather}
For example, the exact quotients used in the first line are
\[
 \frac CN=\frac{1+zr}{2a^2t\Delta},\qquad
 \frac\alpha C=\frac{z^2(b+r)a^2t}{(1-zh)(1+zr)}.
\]
The four strong normalized products use the numerators of
\eqref{eq:c2-native-source} after multiplication by $r,h,g,j$, respectively;
no vanishing source gap is inverted in this calculation.
All six weights are therefore positive, $C,\alpha$ are infinitesimal,
and their normalized quotients displayed above are bounded units.
For instance
\[
 \frac1{Cn}=\frac{r/N}{(C/N)(rn)},\qquad
 \frac1{\alpha n}=\frac{C/\alpha}{Cn};
\]
the other six product bounds follow in the same way from $g/N,h/N,j/N$.
These are exactly the hypotheses of Theorem~\ref{p8b53b68d8af6:thm:c2-matching-strength}.

For completeness, the source is physical by the unchanged exact contact
argument.  The rational stationary identities
\eqref{eq:c2-native-contact} are valid at this value of $c$ as well.
At that noncommuting contact the state matrix has a nonzero kernel
vector.  Its restriction to $Z=0$ is positive definite: the signed core
is nonnegative there and the positive event positive square terms have zero common
kernel.  Hence the other three eigenvalues are positive.  The exact
contact and angular stationarity give the stationary zero of the
Laurent determinant; its constant-plus-positive-exponentials form makes
this its global minimum.  The determinant criterion and homogeneous
continuity prove physicality, exactly as for the small-angle source.
Theorem~\ref{p8b53b68d8af6:thm:c2-matching-strength} now applies.

The corollary applies to the stated radial subfamily of the ordinary-angle
cell.
\end{proof}

\subsection{Weighted radial completion}
\label{component:p9f1235511111}
\label{component:p53752031d16c}
\label{component:p163e0ebf59f6}
\label{p9f1235511111:sec:c2-unified}
\label{p53752031d16c:sec:c2-y-current}
\label{p163e0ebf59f6:sec:c2-radial-corner}
Normalize the original positive $d$ by a pure-source square. The
two-weak source reduction is applied after this normalization; it has
already given $c\in\mathrm I$,
$a,b,t\sim1$ and $r,h,g,j\in\mathrm I$, where $h=b-a^2t$.
Use the exact native coefficients in Section~\ref{app:c2-source}, with
$(e_1,e_2,e_3,e_4,e_5)=(n,m,C,A,D)$ and
\[
 (u_c,v_c;W_c,X_c,Y_c,Z_c)
 =(c,ac;b/a,-cr,c(b/a-at),1).
\]
The two source identities needed for all relative opening rates are
\begin{equation}
 N=(1+c^2)(b+r)h+a^2\{g+j+rt(1+a^2c^2)\}>0,
 \qquad N\asymp h+r+g+j.
 \label{p53752031d16c:eq:c2y-N-positive}
\end{equation}
Consequently
\[
 U_1=re_1,\quad U_2=ge_2,\quad U_4=he_4,\quad U_5=je_5
 \sim\tfrac14,\qquad e_3/N\sim\tfrac14,
 \qquad {4\alpha\over c^4N}\sim1.
\]
These follow directly from the native fractions after their respective
opening factors are canceled. In particular all four $e_i$ above are
unbounded, and $h/N,r/N,g/N,j/N$ are bounded. They also imply
$1/(Cn),1/(Cm),1/(CA),1/(CD)=O(1)$, the assumptions of the complete
radial outer calculation.

The same source reduction forces $n,m,A,D$ unbounded, so $n=0$ is
impossible in this case. No zero-coefficient face is being discarded.
\begin{proposition}
\label{prop:c2-weighted-radial}
\label{p163e0ebf59f6:thm:c2-radial-corner}
On this source domain the full parent is tensor SOS at every independent
projective testing pair and every real or complex state.
\end{proposition}
First let the independent tests $u,v$ be infinitesimal or zero and take
the pure-testing state chart $Z=1$, $W\ge0$ bounded. Define
\[
 R_m=vW-X+uvY-u,\qquad J=uW+uvX-Y-v,
\]
\[
 H_1=X+ru,\quad H_2=R_m-gu,\quad
 H_4=Y-(h/a)u,\quad H_5=J+jv.
\]
Retain the four whole positive squares
\[
 4e_1D_uD_vH_1^2+4e_2H_2^2+4e_4D_uD_vH_4^2+4e_5H_5^2.
\]
Their differences from the original terms are, respectively,
\[
 -8U_1D_uD_vuH_1+4rU_1D_uD_vu^2,
 \quad8U_2uH_2+4gU_2u^2,
\]
\[
 {8U_4\over a}D_uD_vuH_4+{4hU_4\over a^2}D_uD_vu^2,
 \quad-8U_5vH_5+4jU_5v^2.
\]
Thus the entire remaining polynomial has bounded coefficients,
including all its state normals. With $n_1=H_1$, $n_5=H_5$ the exact
inverse is
\[
 X=-ru+n_1,\qquad
 Y=uW-(1-j+ru^2)v+uvn_1-n_5.
\]
Define its scalar $f$, linear coefficients $L_1,L_5$, and quadratic
normal matrix by coefficient extraction in $(n_1,n_5)$. Quadratic
state degree proves that this is the whole remainder.

Put
\[
 R=au-v,\qquad S=u(W-W_c),\qquad T=u^2-c^2.
\]
Then $f\in(R,S,T)^2$ and $L_1,L_5\in(R,S,T)$ with bounded
coefficients. Here is the finite division proof. Parity makes
$f(u,au,W)=f_0(u^2,W)$ even and $f_v(u,au,W),L_i(u,au,W)$ odd.
At zero angles the value is $4\alpha$, independent of $W$, so both
nonconstant $W$ coefficients of $f_0$ are divisible by $u^2$.
Write $f_0=A(\xi)(W-W_c)^2+B(\xi)(W-W_c)+C(\xi)$.
The exact source contact gives
\[
 A=\xi A_1,\quad B=\xi(\xi-c^2)B_1,
 \quad C=(\xi-c^2)^2C_1.
\]
Monic division gives these polynomials without a factor $1/c$.
Likewise $f_v(u,au,W)=S\ell_1+uT\ell_2$. Finite Taylor division in
$v-au$ proves the scalar assertion. The identical odd-parity argument
proves the linear-normal assertions. All higher normal terms are
already quadratic by construction.

The exact current identity removes $S$. Put $K=1-j+ru^2$ and
\[
 Q(\xi)=aW_c-a^2(1-j)-a^2r(\xi+c^2).
\]
The identities $H_2=G-(1-u^2v^2)H_1-uvH_5$ and
$G=-[D_uW-Ku(v+au)]R+aD_uS+uTQ(u^2)$ give
\begin{equation}
 S={H_2+(1-u^2v^2)H_1+uvH_5
       +[D_uW-Ku(v+au)]R-uTQ(u^2)\over aD_u}.
 \label{p9f1235511111:eq:c2u-eliminate-S}
\end{equation}
Its denominator is a product of pure-source and pure-testing units.
After restoring the four retained squares this is an exact bounded
Gram on $K_i=\sqrt{e_i}H_i$ for $i=1,2,4,5$, together with $R,T$.

There is one further exact linear relation:
\begin{equation}
 \begin{split}
 \zeta={}&aD_uH_4-H_2-(1-u^2v^2+aD_uuv)H_1
 +(aD_u-uv)H_5\\
 &-[D_uW+(a-uv)K]R+u[Q(u^2)+a^2D_ur]T=0.
 \end{split}
 \label{p53752031d16c:eq:c2y-zero-relation}
\end{equation}
Add the exact zero
\begin{equation}
 {8U_4\over a^3}\zeta R+\zeta^2=0.
 \label{p53752031d16c:eq:c2y-zero-correction}
\end{equation}
The hard and $R$ residue can be computed directly on the present
all-small domain. Modulo $(h,r,g,j)$, followed by the joint real
residue $c=u=v=0$, the degree-two scalar is
\[
 2(u-v)(uW-v)=2R(R+S).
\]
Equation~\eqref{p9f1235511111:eq:c2u-eliminate-S} gives $S=WR$ on
that residue, while $\zeta=-(1+W)R$ and $U_4=1/4$.
The displayed zero changes $2(1+W)R^2$ into $(1+W)^2R^2$.
The four scaled hard diagonals are $4$ and their other couplings
vanish.

It remains to divide the radial column by $\sqrt N$. Let
$I=(h,r,g,j)$ in the bounded source coefficient ring. The exact
identities
\[
 b-1=-r+g-c^2h,\quad t-1=c^2r-j,\quad
 a-1={g+j-(1+c^2)(r+h)\over t(a+1)}
\]
show that subtraction of the $I=0$ specialization of any coefficient
leaves $I$ with bounded coquotients. Its denominators are units after
the four opening factors have been absorbed in $U_i$. Finite monic
division and the displayed current substitution preserve this fact.

The exact simplifications $Q(u^2)=h-a^2ru^2$ and
$Q(u^2)+a^2D_ur=h+a^2r$ belong to $I$. The only radial cross
coefficient requiring a further correction is the $RT$ coefficient.
For an explicit calculation before taking the testing residue, set
$D_c=1+c^2$, $Y_*=uW-v$, $G_*=D_uvW-D_vu$. At $I=0$ the whole
scalar remainder is
\[
 f_*=2uG_*+{2\over D_c}D_uD_vuY_*
       +2\{D_u(vY_*-1)^2+D_vu^2-D_uD_v\}.
\]
It satisfies
\[
 f_*(u,u,1)=0,\quad
 (f_*)_W(u,u,1)={2u^2(u^2-c^2)D_u\over D_c},\quad
 (f_*)_v(u,u,1)={-2u(u^2-c^2)D_u\over D_c}.
\]
The monic division order just specified therefore gives, after
eliminating $S$, the $RT$ coefficient
$2u(D_uW+1-uv)/D_c$. The earlier zero adds none at this
specialization. The $R$ coefficient of $\zeta$ is
$-(D_uW+1-uv)$, so add the further exact zero
\[
                         (2u/D_c)\zeta T=0.
\]
The resulting $RT$ coefficient belongs to $uI$.

Now use $T_N=\sqrt N T$. Every hard--$T_N$ coefficient is $u$
times bounded factors divided by $\sqrt{Ne_i}$, which is bounded
because $1/\sqrt{Ne_i}=\sqrt{(r,g,h,j)_i/N}/\sqrt{U_i}$.
Its residue is zero. The corrected $RT_N$ coefficient belongs to
$uI/\sqrt N$, again bounded with zero residue. The $T^2$ coefficient
belongs to $I$ because $f_*(u,u,1)=0$; every contribution other than
its value at $u=0$ has a further factor $u^2$. Its value at $u=0$ is
exactly $4\alpha/c^4$. Therefore its quotient by $N$ is bounded and
has residue $1$. The complete weighted residue is
\[
 4K_1^2+4K_2^2+4K_4^2+4K_5^2+(1+W_0)^2R^2+T_N^2\succ0.
\]
The separated strict-Gram lemma proves the full chart.

For $W=1$, $Z$ infinitesimal, retain the original four squares in
$X,R_m,Y,J$ and the whole $\alpha$ square. The exact testing identities
\[
 v={R_m+(1-u^2v^2)X+uD_vZ+uvJ\over D_u},
\]
\[
 u={D_uY+[Z(1-u^2v^2)-D_uuv]X+(D_u+uvZ)J+ZR_m
       \over D_u-D_vZ^2}
\]
have pure-testing unit denominators. They express the entire remaining
quadratic form on the four scaled retained amplitudes. To check this
without discarding a state term, the half-core has the exact completion
\[
 H=2\{C_*^2+D_vQ_*^2+D_uv^2Y^2
       +D_uu^2v^2(1-u^2v^2)Z^2/D_v\}-4D_uuvZ,
\]
where $C_*=u+uv(X+uZ)-Y$ and
$Q_*=X+uZ-D_uuv^2Z/D_v$. Substitute the two identities in every
term. All coefficients are bounded and the residue on the four scaled
amplitudes is $4I$; thus the same lemma applies.

The testing sign gauges allow $u,v\ge0$. The $X,Y$ block of half
the core has eigenvalues $L\pm uv>0$, $L=1+v^2+u^2v^2$, and Schur
complement
\[
 {uvD_uD_v\{uv(W^2+Z^2)-2LWZ\}\over L^2-u^2v^2}.
\]
For $WZ\le0$ the core is a nonnegative pure-testing value and all
events are already tensor squares. For $WZ>0$, the two preceding
state charts are exhaustive. Homogeneous restoration and real/imaginary
splitting prove every state. Section~\ref{component:pc7dacaa3258e}
then proves all other independent projective testing pairs.

\subsection{Complementary testing angles}
\label{component:pc7dacaa3258e}

\label{pc7dacaa3258e:sec:c2-radial-all-tests}

The radial corner has four strong event positive square terms and one weak positive square term.
The weak positive square term must be retained at its exact scale at the opposite
commuting corner.  An elementary conductor identity makes that scale
sufficient; it does not compare coefficient and testing infinitesimals.
Throughout this section $C$ denotes the event weight $e_3$, whereas $c$
continues to denote the source half-angle in \eqref{eq:c2-native-source}.

\begin{lemma}[Four strong positive square terms and a weak opposite-corner pivot]
\label{pc7dacaa3258e:lem:c2-four-strong-outer}
Consider the full parent \eqref{eq:c2-master-parent}, with
\[
 (e_1,e_2,e_3,e_4,e_5)=(n,m,C,A,D),\qquad \alpha\geq0.
\]
Suppose $C>0$ is infinitesimal and the four positive pure coefficient
quantities $Cn,Cm,CA,CD$ have bounded reciprocals.  Then the full parent
is tensor SOS at every commuting residue test outside the finite
zero-zero corner, for every real or complex projective state.
The four strong coefficients and the roots in the proof may be unbounded.
\end{lemma}

\begin{proof}
The hypotheses imply that $1/n,1/m,1/A,1/D$ are infinitesimal.
First consider the outer axes other than the infinity-infinity corner.
Choose $\mu=\min(1,n,m,A,D)$, a positive bounded unit.  Retain the
nonnegative excess of each strong positive square term, together with the whole
$C$ and $\alpha$ positive square terms.  The remaining parent consists of the signed
core and the four positive square terms with common weight $\mu$.

Here are the four kernel calculations, with positive factors
$1+u^2$ or $1+v^2$ omitted only from the positive square term amplitudes.
At $u=0$, finite nonzero $v$, the four amplitudes are
$(X,vW-X,Y,-Y-vZ)$; at $v=0$, finite nonzero $u$, they are
$(X,-X-uZ,Y,uW-Y)$.  Their common kernels are zero, and half the
signed core is, respectively,
\[
 D_v(X^2+Y^2),\qquad (uW-Y)^2+(X+uZ)^2.
\]
At $u=\infty$, finite $v$, the four homogenized amplitudes are
$(X,vY-Z,Y,W+vX)$; at $v=\infty$, finite $u$, they are
$(X,W+uY,Y,uX-Z)$.  These common kernels are zero even when the
finite angle is zero.  Half the homogenized signed core is, respectively,
\[
 (W+vX)^2+(vY-Z)^2,\qquad D_u(X^2+Y^2).
\]
Thus the selected state Gram has strict positive residue.  Its
coefficients are bounded in each indicated local testing chart, so the
strict-residue Gram lemma proves the whole chart.  This includes the
mixed zero-infinity endpoints and either sign of every finite angle.

It remains to treat $u=v=\infty$ at residue.  Put $p=1/u$, $q=1/v$,
using the homogeneous reciprocal coordinates at exact infinity, and
write $D_p=1+p^2$, $D_q=1+q^2$.  Thus $p,q$ are infinitesimal or zero.
The homogeneous full parent in this chart is exactly
\begin{align}
\widehat F={}&4nD_pD_qX^2+4mM^2+4CD_q(W-pY)^2
       +4AD_pD_qY^2+4DJ^2+4\alpha D_p(Y-qZ)^2+H,\notag\\
 M={}&pW-pqX+Y-qZ,\qquad
 J=qW+X-pqY-pZ,\notag\\
 H={}&2\{J^2+D_p(Y-qZ)^2+D_q(pX+Z)^2-D_pD_qZ^2\}.
\label{pc7dacaa3258e:eq:c2ro-parent}
\end{align}
For finite nonzero $u,v$ this is $p^2q^2F$; homogeneous interpretation
also gives the exact endpoint polynomial.

Set $L=1+p^2+p^2q^2$.  The needed conductor identity is
\begin{equation}
\frac H2=qWJ-qZM+L(X^2+Y^2)-2pqXY
             +q(W+pqZ)X-q(pqW+Z)Y.
\label{pc7dacaa3258e:eq:c2ro-core-conductor}
\end{equation}
For a direct check, put $M_0=pW-qZ$, $J_0=qW-pZ$.  The scalar part
of $H/2$ is
\[
 J_0^2+(q^2-p^2)Z^2=qWJ_0-qZM_0.
\]
Its linear $X,Y$ part is
$2q(W+pqZ)X-2q(pqW+Z)Y$, and its quadratic part is
$L(X^2+Y^2)-2pqXY$.  Substitution of
$M=M_0-pqX+Y$, $J=J_0+X-pqY$ gives
\eqref{pc7dacaa3258e:eq:c2ro-core-conductor}.

First normalize $W=1$ and assume $Z$ bounded.  No bound is required
on $X$ or $Y$.  Retain the $\alpha$ positive square term and use the five roots
\[
 K_1=\sqrt nX,\quad K_2=\sqrt mM,\quad
 K_4=\sqrt AY,\quad K_5=\sqrt DJ,\quad T_C=\sqrt C\,W.
\]
Here $T_C=\sqrt C$ in this state chart.  Equations
\eqref{pc7dacaa3258e:eq:c2ro-parent}--\eqref{pc7dacaa3258e:eq:c2ro-core-conductor} give the following
\emph{exact} Gram expression for the remainder:
\begin{align}
\widehat F-4\alpha D_p(Y-qZ)^2={}&
 4D_pD_q(K_1^2+K_4^2)+4K_2^2+4K_5^2
 +4D_q\left(T_C-p\sqrt{C/A}\,K_4\right)^2\notag\\
&+2L\left(\frac{K_1^2}{n}+\frac{K_4^2}{A}\right)
 -\frac{4pq}{\sqrt{nA}}K_1K_4\notag\\
&+2qT_C\left\{
 \frac{K_5}{\sqrt{CD}}-\frac{ZK_2}{\sqrt{Cm}}
 +\frac{(1+pqZ)K_1}{\sqrt{Cn}}
 -\frac{(pq+Z)K_4}{\sqrt{CA}}\right\}.
\label{pc7dacaa3258e:eq:c2ro-Wunit-Gram}
\end{align}
Every denominator is pure coefficient.  The four inverse square roots
$1/\sqrt{Ce_i}$ are bounded by hypothesis, and
$\sqrt{C/A}=C/\sqrt{CA}$ is infinitesimal.
All coefficients of this Gram are therefore in the bounded tensor
ring.  Every displayed correction has zero residue, whereas its five
positive square term pivots have residue $4$.  Thus its residue is exactly
\[
             4(K_1^2+K_2^2+K_4^2+K_5^2+T_C^2)\succ0.
\]
The strict-residue Gram lemma applies without any bound on the roots.
In particular it makes no comparison between $C$ and $p$ or $q$.

The other state chart is $Z=1$, $W$ infinitesimal or zero.  Here the
exact pure-testing conductor is a unit:
\begin{equation}
 q=\Gamma:=\frac{M+WJ+(pq-W)X+(Wpq-1)Y}{W^2-1}.
\label{pc7dacaa3258e:eq:c2ro-Zunit-conductor}
\end{equation}
Indeed $M+WJ=q(W^2-1)+(W-pq)X+(1-Wpq)Y$.
Substitution in \eqref{pc7dacaa3258e:eq:c2ro-core-conductor} gives
\begin{equation}
 \frac H2=L(X^2+Y^2)-2pqXY
 +\Gamma\{WJ-M+(W+pq)X-(pqW+1)Y\}.
\label{pc7dacaa3258e:eq:c2ro-Zunit-core}
\end{equation}
This is an exact quadratic form on $X,M,Y,J$, with bounded pure-testing
coefficients.  Retain the whole $C$ and $\alpha$ positive square terms separately.
Scale the four remaining roots as before.  Every coefficient of the
scaled core has zero residue, because each contains two inverse strong
positive square term square roots.  Adding the four strong positive square terms consequently
produces the strict residue
$4(K_1^2+K_2^2+K_4^2+K_5^2)$.  Again only the Gram coefficients,
not $X,Y$ or the positive square term roots, needed boundedness.

These two state charts are exhaustive when $(W,Z)\ne(0,0)$: if
$W\ne0$ and $Z/W$ is bounded use the first; otherwise $Z\ne0$ and
$W/Z$ is infinitesimal, giving the second.  Normalization and restoration
use pure-testing factors.  When $W=Z=0$,
\eqref{pc7dacaa3258e:eq:c2ro-core-conductor} is just the bounded quadratic form
$H/2=L(X^2+Y^2)-2pqXY$; after the same four-root scaling its
coefficients have zero residue, so the four retained strong positive square terms
again give a strict Gram.  This includes the zero state.  Splitting a
complex state into real and imaginary parts completes the proof.
\end{proof}

\begin{corollary}[The radial source has no missing outer commuting chart]
\label{pc7dacaa3258e:cor:c2-radial-outer}
Use the unchanged exact source \eqref{eq:c2-native-source} under the hypotheses
of Theorem~\ref{p163e0ebf59f6:thm:c2-radial-corner}: $h,r,g,j$ are positive
infinitesimals and $a,1-r,t$ are positive bounded units.
Then the full parent is tensor SOS at every commuting residue test and
every real or complex state.  Together with the physical-source
noncommuting-residue theorem, this proves all tests for this exact
radial source family.
\end{corollary}
\begin{proof}
The source identities and positive decomposition of $N$ proved in
Section~\ref{p163e0ebf59f6:sec:c2-radial-corner} give
\[
 N\asymp h+r+g+j,\qquad
 U_1=re_1,\ U_2=ge_2,\ U_4=he_4,\ U_5=je_5>0
 \quad\hbox{bounded units},
\]
where $N/(h+r+g+j)$ has residue $1$.  Moreover, directly from
\eqref{eq:c2-native-source},
\[
 \frac CN=\frac{A_c}{2a^2t\Delta}>0
 \quad\hbox{is a bounded unit}.
\]
For example,
\[
 \frac1{Ce_1}=\frac{r/N}{(C/N)U_1},
\]
and the three analogous formulas use $g/N,h/N,j/N$.
Thus all four product reciprocals in
Lemma~\ref{pc7dacaa3258e:lem:c2-four-strong-outer} are bounded, while $C$ is positive
infinitesimal.  That lemma covers every outer commuting chart, and
Theorem~\ref{p163e0ebf59f6:thm:c2-radial-corner} covers the zero-zero chart.

The final noncommuting join uses only the already-established
physicality of the exact source and the physical-source
noncommuting-residue theorem.  In particular the exact contact and
Laurent-determinant physicality argument used for the unit-$N$ source
requires $N>0$, not that $N$ have positive residue.  All five event
weights are positive.  On the three-dimensional state subspace $Z=0$
the signed core is nonnegative, and the positive event positive square terms have
zero common kernel at the noncommuting contact; hence the state
restriction is positive definite.  The exact stationary state is a
nonzero kernel vector, so the other three state eigenvalues are positive.
Determinant stationarity remains exact, and the same Laurent-determinant
argument proves physicality.
No coefficient-entry claim is made here for C2 sources outside the
stated radial family.
\end{proof}

\section{Strong sources with two small actual angles}
The following full-parent polynomial identities and normal divisions are
also used by the C1 coefficient domains. Their use there requires the
C1 source bounds and coefficient identities proved in the C1 chapter;
it does not use a C2 source conclusion with altered coefficients.

\subsection{Local strong-II Gram}
\label{component:pe16499ccc4c6}

\label{pe16499ccc4c6:sec:r7-local-strong-II}

\begin{theorem}[Strong-II local allocation]
Let the positive C2 source satisfy $z,y\in\Inf_{>0}$ and let
$B=H+L$ be a positive unit.  Keep the exact source coefficients of
(GE1), with $n\geq0$ and the remaining displayed weights positive.
Then the original C2 parent belongs to $\sum(K_c\otimes_{\mathbb R}
K_t)^2$ whenever both independent affine testing coordinates $u,v$
are infinitesimal, including either literal testing axis.
There is no comparison hypothesis between $u^2,z$, between $v^2,y$,
or between the two source rates.  All openings of $r$, $L/y$,
$1+zr-L/y$, and $1-H$ are retained, as is literal $n=0$.
\end{theorem}

This is the local full-parent calculation. Section~\ref{component:p590112f102bd} proves the other testing charts and the literal projective endpoints.
All identities below are polynomial in testing coordinates.  Inverses
are taken only in the pure source field, and the final step uses the
bounded Gram strict-real-residue lemma already used in the completed
local UI certificate.  In particular no ordering or inverse of the
mixed tensor is used.

\paragraph{Source units and notation.}
Write $U=u^2,V=v^2,k=\sqrt y$ and
\[
 a_c=1+zr,\quad t=L/y,\quad j=a_c-t,\quad s=a_c/B,
 \quad c=CL,\quad d=Dyj.
\]
The source entry gives bounded $t,j,r$, positive units $B,c,d,m/z^2,A/y$,
bounded $nr$ and $\alpha(1-H)$, and strong $\alpha$.
These facts include the arbitrary openings $t\to0$, $j\to0$,
$r\to0$, and $1-H\to0$.  They follow directly from the positive
source formulas and were proved in (DI7).  Let $h=\operatorname{st}B$;
then $0<h\leq1$ and, with $r_0=\operatorname{st}r$,
\[
\begin{gathered}
 c_0=d_0^{\rm pair}=\frac{h}{2(1+h)},\quad
 \operatorname{st}(m/z^2)=\frac{1-r_0}{2h^2(1+h)},\quad
 \operatorname{st}(A/y)=\frac1{2h(1+h)},\\
 \operatorname{st}(nr)=\frac1{2(1+h)}-\frac{r_0}{2},\qquad
 \operatorname{st}\{\alpha(1-H)\}=\frac{h^2}{2(1+h)}.
\end{gathered}
\tag{LII1}
\]
These residues follow directly from (GE1): $L\to0$, $H\to h$,
$a_c\to1$, $G\to h$, $\delta\to h(1+h)$, and the source
numerator $M\to h^2$.  Substitution in the six rational coefficient
formulas gives (LII1), including the identity
$\operatorname{st}(nr)=1/[2(1+h)]-r_0/2$ without division by $n$.
Define the source quantities (with $q,g$ bounded)
\[
 \ell=y(C+D)=c/t+d/j,\qquad
 q=\frac{D}{C+D}=\frac{dt}{cj+dt},\qquad
 g=\frac{yCD}{C+D}=\frac{cd}{cj+dt}.
\]
Here $0<q<1$, $g$ is a positive unit with residue $h/[2(1+h)]$,
and $\ell$ has a positive real lower bound, although it may be
unbounded.  Indeed $cj+dt$ is a unit because $t+j=a_c$ is a unit.
Choose positive real $\ell_*\leq\ell$ and
$\alpha_*\leq\alpha$.

\paragraph{The state roots and the exact C,D completion.}
Use the source-fixed parity graph
\[
 X=-\frac{zr v}{B}W+\xi,\qquad
 Y=\frac{Hu}{B}W+\eta,\qquad
 Z=\frac{uv}{B}W+\zeta.                         \tag{LII2}
\]
It agrees with the actual kernel at $U=z,V=y$ and is a change of all
four homogeneous state coordinates with determinant one.  Set
\[
 Q_n=u\xi-\zeta,\quad K=ukW,\quad
 R_h=\frac{vQ_n-su(V-y)W}{k},\qquad
 N=\frac\eta k-qR_h.                              \tag{LII3}
\]
Thus $\eta=k(N+qR_h)$.  Every root in (LII3) vanishes at the actual
contact.  If $Q=uW-Y$ and $J=Q+v(uX-Z)$, then
\[
 Q/k=(t/B)K-N-qR_h,\qquad
 J/k=(1-q)R_h-N-(j/B)K.
\]
Consequently the exact weighted completion is
\[
\begin{split}
 4CQ^2+4DJ^2={}&4\ell N^2+4gR_h^2
 +\frac{8(d-c)}BKN-8gsKR_h
 +\frac{4(ct+dj)}{B^2}K^2,\\
 4CVQ^2={}&4CV\eta^2
 +\frac{4ct}{B^2}VK^2-\frac{8c}BVK(N+qR_h).
                                                        \end{split}\tag{LII4}
\]
There is no estimate in these identities.  Retain the following
four nonnegative, contact-centered positive square terms outside the remainder:
\[
 4(\ell-\ell_*)N^2,\qquad 4CV\eta^2,\qquad
 4n(1+U)(1+V)\xi^2,\qquad
 4(\alpha-\alpha_*)(1+U)(v\eta-\zeta)^2.          \tag{LII5}
\]
Call the exact remainder $G$.  Its coefficients are bounded after
writing $\eta=k(N+qR_h)$ and using (LII4).  For completeness, its
entire formula is
\[
\begin{split}
G={}&4\ell_*N^2+4gR_h^2+\frac{8(d-c)}BKN-8gsKR_h\\
 &+\frac{4(ct+dj)}{B^2}K^2
 +\frac{4ct}{B^2}VK^2-\frac{8c}BVK(N+qR_h)\\
 &+4n(1+U)(1+V)
   \left\{\frac{z^2r^2V}{B^2}W^2-\frac{2zr v}B W\xi\right\}\\
 &+4\alpha_*(1+U)(v\eta-\zeta)^2
  +4\alpha(1+U)
   \left\{\frac{(1-H)^2UV}{B^2}W^2
        -\frac{2(1-H)uv}B W(v\eta-\zeta)\right\}\\
 &+4mM^2+4A(1+U)(1+V)Y^2\\
 &+2\{J^2+(1+U)(vY-Z)^2+(1+V)(X+uZ)^2
                      -(1+U)(1+V)Z^2\},        \end{split}\tag{LII6}
\]
where $M=vW-X+uvY-uZ$ and (LII2) is used in the last line.
Every potentially unbounded original coefficient occurs here only
through $nr$, $\alpha(1-H)$, or the bounded quantities in (LII1).
No division by a testing coordinate has been introduced.

\paragraph{One exact gauge removes the dependent-root first normal.}
In the formal bounded polynomial (LII6), the coefficient of $R_hW$
is $uk\rho(U,V)$ for a bounded source polynomial $\rho$.  Replace
this entire term using the exact identity
\[
 KR_h=uvWQ_n-sU(V-y)W^2.                         \tag{LII7}
\]
Call the resulting polynomial $\widetilde G$.  It is equal to $G$
on the roots (LII3), and has no $R_hW$ term.  This gauge is useful
because $R_h$ is a dependent root: it permits the state and angular
contact equations to be read as formal coefficient vanishings.

More precisely, by parity write uniquely
\[
 \widetilde G=Q_0(N,\xi,\zeta,R_h)
 +W\{vP_\xi(U,V)\xi+uvP_\zeta(U,V)\zeta
                         +ukP_N(U,V)N\}
 +\phi(U,V)W^2,                                  \tag{LII8}
\]
where all coefficients are bounded source polynomials and $Q_0$ is
quadratic in the four displayed roots.  Formula (LII6), followed by
the one replacement (LII7), is an exact coefficient-extraction
specification for every entry in (LII8).

The retained positive square terms (LII5) vanish together with their first jets at
the actual contact.  The exact original source state equations thus
give
\[
 P_\xi(z,y)=P_\zeta(z,y)=P_N(z,y)=0.              \tag{LII9}
\]
Indeed the physical $\eta$ derivative first forces the $N$ coefficient
to vanish; the physical $\xi,\zeta$ derivatives then force their
formal coefficients to vanish because the $R_hW$ coefficient was
removed.  The exact original contact and the two angular derivatives
now give
\[
 \phi(z,y)=\phi_U(z,y)=\phi_V(z,y)=0.              \tag{LII10}
\]
This uses the six exact original source substitutions already checked
in R6.  It imposes no new relation between source and testing fields.

\paragraph{Bounded division and the complete angular allocation.}
There are two additional elementary support facts in (LII8):
\[
 P_\xi\in(U,z),\qquad \deg_U\phi\leq2,\qquad
 \phi(0,V)=V(d_s+e_sV),                           \tag{LII11}
\]
where
\[
 d_s=\frac4{B^2}\{m(B+zr)^2+(n+\tfrac12)z^2r^2\},\qquad
 e_s=\frac4{B^2}(n+\tfrac12)z^2r^2.
\]
Both $d_s/z^2$ and $e_s/z^2$ are bounded.  The first support assertion
is literal bounded-coefficient membership: setting $U=z=0$ in the
extracted polynomial gives zero.  It is also clear before expansion:
all $\xi W$ terms not supplied by the gauge contain $U$, $m$, $nzr$,
or $zr$; those supplied by the gauge contain $U$.

Monic division using (LII9) therefore gives bounded polynomials
\[
\begin{split}
 P_\xi&=(U-z)A_\xi+z(V-y)B_\xi,\\
 P_\zeta&=(U-z)A_\zeta+(V-y)B_\zeta,\\
 P_N&=(U-z)A_N+(V-y)B_N.                         \end{split}\tag{LII12}
\]
For the first line, evaluate $P_\xi$ at $U=z$: the support assertion
supplies the factor $z$ with bounded quotient, and its value at $V=y$
is zero.  All other divisions are by the monic polynomials $U-z$
or $V-y$.  In particular none creates an inverse of a small source
or a testing difference.

The scalar has the exact bounded decomposition
\[
 \phi=\frac{d_s+e_sV}{z^2}V(U-z)^2
       +U(U-z)(V-y)A_s(V)+U(V-y)^2C_s(V).        \tag{LII13}
\]
Here is a constructive proof, including boundedness.  Write
$\phi=a(V)U^2+b(V)U+V(d_s+e_sV)$ and subtract the first term of
(LII13).  The remaining $U^2,U$ coefficients $a_1,b_1$ vanish at
$V=y$ by the value and $U$ derivative in (LII10).  Divide them by
$V-y$ to obtain $A_s,B_s$.  The $V$ derivative gives
$zA_s(y)+B_s(y)=0$, so set $C_s=(B_s+zA_s)/(V-y)$.
All three divisions are monic divisions of bounded polynomials.
This proves (LII13) exactly, not asymptotically.

Introduce the two legal angular roots
\[
 R=v(U-z)W,\qquad T=uk(U-z)W.                    \tag{LII14}
\]
The following identities allocate every term in (LII12)--(LII13):
\[
\begin{split}
 zv(V-y)W&=\frac{uVQ_n-uvkR_h}{s}-(V-y)R,\\
 uv(V-y)W&=\frac{VQ_n-vkR_h}{s},\\
 uk(V-y)W&=\frac{vkQ_n-yR_h}{s},\\
 U(U-z)(V-y)W^2&=-\frac{R_hT-uRQ_n}{s},\\
 U(V-y)^2W^2&=\frac{VQ_n^2-2vkR_hQ_n+yR_h^2}{s^2}.
                                                        \end{split}\tag{LII15}
\]
Together with $v(U-z)W=R$, $uv(U-z)W=uR$, and
$uk(U-z)W=T$, these are an explicit complete allocation of the
full remainder as a bounded symmetric Gram in
\[
                    (N,\xi,\zeta,R_h,R,T).       \tag{LII16}
\]
The original first state normals are all present.  The scalar
restriction has not been used as a substitute for them.

\paragraph{The finite block and the missing pivot.}
The residue of $Q_0$ on the first four roots is
\[
             \operatorname{diag}(4\ell_*,2,4\alpha_*,4g_0).
                                                        \tag{LII17}
\]
For example the signed core at $W=0$ is exactly
$2(\xi+u\zeta-uv\eta)^2+2(1+U)V\xi^2+2(1+V)\eta^2$;
all its other entries vanish in the present residue.  The extra
terms in the first three lines of (LII15) have infinitesimal Gram
coefficients except for $A_\xi\xi R$ and the allowed $A_NNT$.
Likewise the last line has only infinitesimal coefficients, and
the fourth line has one potentially finite $R_hT$ entry.  Thus,
apart from entries in the last row and column $T$, the sole finite
off-diagonal entry is between $\xi$ and $R$.

The exact coefficient extraction yields
\[
 \operatorname{st}A_\xi=\frac4{h(1+h)},\qquad
 \operatorname{st}(d_s/z^2)=\frac2{h^2(1+h)}.     \tag{LII18}
\]
For a short independent check of the first value, the unscaled
original first-normal coefficient on the $V=y$ slice is
\[
 -4nzr(1+y)+4m(1-H)-(4D+2)yj+2(1+y)
 \longrightarrow 2-4d_0^{\rm pair}=\frac2{1+h}.
\]
Its full multiplier is $2/B$.  Passing from the original normal
to (LII8) changes this $U$-quotient only by a bounded multiple of
$U$ or $z$, so its residue is unchanged.  The second value follows
from (LII1):
\[
 \operatorname{st}\{m(B+zr)^2/z^2+(n+\tfrac12)r^2\}
 =\frac{1-r_0}{2(1+h)}+\frac{r_0}{2(1+h)}
 =\frac1{2(1+h)}.
\]
Hence the finite $(\xi,R)$ Gram block is exactly
\[
 \begin{pmatrix}
  2&\dfrac2{h(1+h)}\\[2mm]
  \dfrac2{h(1+h)}&\dfrac2{h^2(1+h)}
 \end{pmatrix},\qquad
 \text{Schur complement}=\frac2{h(1+h)^2}>0.      \tag{LII19}
\]
This calculation retains $r_0=0$, $n=0$, and $h=1$.

Finally the exact toric conductor is the zero identity
\[
 \boxed{\quad
 T^2-UR^2-\frac{U-z}{s}R_hT
             +\frac{u(U-z)}s RQ_n=0.\quad}       \tag{LII20}
\]
Add any positive real multiple $\Lambda$ of this zero identity to
the Gram.  Its only nonzero residue contribution is $\Lambda T^2$.
The five-by-five block excluding $T$ is strictly positive by
(LII17)--(LII19); all entries incident to $T$ are bounded.  Choose
a sufficiently large real $\Lambda$ by the ordinary real Schur
complement.  The entire six-by-six residue Gram is then strictly
positive.  The bounded Gram lemma gives a tensor SOS for $G$, and
adding the four positive square terms (LII5) proves the theorem.

\paragraph{Quantifiers and endpoint discipline.}
No proof step divides by $u$, $v$, $U-z$, or $V-y$: parity quotients
and monic polynomial divisions are coefficient extractions whose
final identities are polynomial.  Thus the two literal affine axes
are included.  The source inverses $1/B$, $1/k$, and the source
positive square term square roots are legal even when their values are large;
all coefficients of the final Gram are bounded.  The source ratios
$t,j,r,1-H$ have not been bounded away from zero.  The new roots
raise testing degree, so the local certificate is not homogenized
to assert the original literal projective endpoints.  Those are
joined by the separate direct-parent remote certificate.

\paragraph{Source input.}
Source contact and stationarity follow from the native-entry identification in Chapter~\ref{ch:entry} and the exact source formulas in Section~\ref{component:p6e5d41ed80b3}. The scaled bounds used above are the II specialization of GE1--GE2 in Section~\ref{component:pf5726218aec8}, recorded there as DI7.

\subsection{Complementary strong-II tests}
\label{component:p590112f102bd}

\label{p590112f102bd:sec:r7-c2-remote}

Use the positive C2 source with $z,y$ infinitesimal and $B=H+L$ a
positive unit. We prove the full parent for every testing pair with at
least one noninfinitesimal angle, including literal axes and projective
infinities. The double-small testing chart was treated above.

\paragraph{Source constants, including all openings.}
Put $t=L/y$, $j=1+zr-t$, and write $h=\operatorname{res}H$.
Then $0<h\leq1$, $t+j=1+zr\sim1$, and the exact source formulas give
\[
 Cyt\sim Dyj\sim\frac{h}{2(1+h)},\qquad
 \frac Ay\sim\frac1{2h(1+h)}.
\]
Consequently $C,D$ are unbounded, and the positive pure-source quantity
\[
 \gamma=\frac{yCD}{C+D}
 =\frac{(Cyt)(Dyj)}{(Cyt)j+(Dyj)t}
 \sim\frac{h}{2(1+h)}.                                      \tag{RR1}
\]
These statements are uniform when either $t$ or $j$ tends to zero.
If $h<1$, then
\[
 \alpha\sim\frac{h^2}{2(1-h^2)};
\]
if $h=1$, then $\alpha$ is unbounded.  Thus, for every real
$0<h_0<h$, the positive real number
\[
 \beta_0=\frac{h_0^2}{2(1-h_0^2)}<\alpha                     \tag{RR2}
\]
is an admissible cap.  No lower bound on $r,j,t,1-H$ is imposed,
and $n=0$ is retained.

\paragraph{Homogeneous testing coordinates.}
For finite $u,v$ define
\[
 a=\frac{u}{\sqrt{1+u^2}},\quad b=\frac1{\sqrt{1+u^2}},\qquad
 c=\frac{v}{\sqrt{1+v^2}},\quad d=\frac1{\sqrt{1+v^2}}.
\]
Use the same homogeneous formulas at infinity, with $(a,b)=(1,0)$
or $(c,d)=(1,0)$.  These are the original parent's homogeneous
endpoint values.  In particular no raised-degree graph is being
specialized at an endpoint.  Put
\[
 E=aW-bY,\quad K=aX-bZ,\quad L=bX+aZ,
 \quad J=dE+cK,\quad Q=cY-dZ.
\]
The original parent divided by $4(1+u^2)(1+v^2)$ is exactly
\[
 nX^2+m\{bcW-bdX+acY-adZ\}^2+CE^2+AY^2
 +(D+\tfrac12)J^2+(\alpha+\tfrac12)Q^2
 +\tfrac12(L^2-Z^2).                                    \tag{RR3}
\]
At projective endpoints this equation means the original homogeneous
form, not a division by an infinite number.  All radicals and
normalizations used here belong to a single testing field.

\begin{proposition}[All tests with $v$ not infinitesimal]
For every such $v$, including $v=\infty$, and every $u$, including
$u=0,\infty$, the original full parent is tensor-SOS.
\end{proposition}
\begin{proof}
Here $c$ is a unit.  Set $P=bL$, $R_1=\sqrt C E$, and
$R_2=\sqrt D J$.  Orthogonality gives
\[
 L^2-Z^2=P^2+2aKP-b^2K^2,
 \qquad
 K=\frac{R_2}{c\sqrt D}-\frac{dR_1}{c\sqrt C}.            \tag{RR4}
\]
Both coefficients on the right are infinitesimal bounded tensors:
$1/c$ is bounded in the testing field, while $1/\sqrt C$ and
$1/\sqrt D$ are infinitesimal in the source field.  Therefore
\[
 R_1^2+R_2^2+\tfrac12P^2+aKP-\tfrac12b^2K^2
\]
has a bounded Gram with residue $\operatorname{diag}(1,1,1/2)$.
It is tensor-SOS by the strict bounded-Gram lemma.  Every other term
of (RR3), namely the retained $n,m,A,\alpha$ positive square terms and
$J^2/2+Q^2/2$, is already SOS.  This uses all four state coordinates.
\end{proof}

\paragraph{An exact harmonic identity for the remaining remote strip.}
Now let $v$ be infinitesimal, possibly zero, and let $u$ not be
infinitesimal, possibly infinity.  Thus $a$ is a unit.
Write $D_v=1+v^2$, and put
\[
 s=C+D,\quad\theta=D/s,\quad
 N=\sqrt{s}(E+\theta vK),\quad
 V=\frac{vK}{\sqrt y},\quad q=\sqrt yY.
\]
No comparison between $v$ and $\sqrt y$ is made.  The roots may be
unbounded; their Gram coefficients below are bounded.  The exact
pure-source harmonic identity is
\[
 CD_vE^2+D(E+vK)^2=N^2+\gamma V^2+Cv^2E^2.               \tag{RR5}
\]
Also set $x=bX$, $z_0=Z$, and $\ell=vY$.  The single toric relation
needed below is the polynomial identity
\[
 a x\ell-b^2z_0\ell=bVq.                                \tag{RR6}
\]
For a real cap $0<\beta_0<\alpha$, the original parent divided
by $4(1+u^2)$ is the sum of the explicitly positive square terms
\[
 \begin{split}
 &nD_vX^2+m\{b(vW-X)+a(vY-Z)\}^2+N^2+Cv^2E^2\\
 &\hspace{15mm}+\tfrac12(E+vK)^2+(\alpha-\beta_0)(\ell-z_0)^2
 \end{split}                                             \tag{RR7}
\]
and the following exact bounded quadratic form:
\[
 G=\gamma V^2+\frac AyD_vq^2
 +(\beta_0+\tfrac12)(\ell-z_0)^2
 +\frac{D_v}{2}(x^2+2axz_0-b^2z_0^2).                   \tag{RR8}
\]
This decomposition is valid also at $u=\infty$, by the original
homogeneous convention.  It has used no state graph or state inverse.

\begin{proposition}[All tests with $u$ not infinitesimal and $v$ infinitesimal]
The full original parent is tensor-SOS for all these tests, including
$v=0$ and $u=\infty$.
\end{proposition}
\begin{proof}
For a real constant $k$ define
\[
 M_k=(\beta_0+\tfrac12)(\ell-z_0)^2
 +\tfrac12(x^2+2axz_0-b^2z_0^2)
 +k(ax\ell-b^2z_0\ell).
\]
By (RR6), the exact form (RR8) equals
\[
 \boxed{G=D_vM_k+
 \gamma V^2+\frac AyD_vq^2-kD_vbVq
 -v^2(\beta_0+\tfrac12)(\ell-z_0)^2.}                   \tag{RR9}
\]
This is a full-parent identity; every state-normal term is present.

First suppose $b$ has positive residue $b_0$.  Since $a$ is a unit,
$0<b_0<1$.  Choose $0<h_0<h$ sufficiently close to $h$ that
\[
 b_0(1+h-h_0)<1,
\]
and take the admissible cap (RR2).  There exists a positive real
$\kappa$ satisfying
\[
 \frac1{1+b_0h_0}<\kappa<
 \min\left\{\frac1{1-b_0h_0},\frac1{b_0(1+h)}\right\}.
                                                               \tag{RR10}
\]
Set $k=-\kappa$.  The Gram of $M_k$ on $(x,z_0,\ell)$ has
first pivot $1/2$, second Schur pivot $\beta_0$, and determinant
\[
 \det M_k=\frac{2\beta_0 b^2k^2-(2\beta_0+1)(k+1)^2}{8}.
                                                               \tag{RR11}
\]
Its residue is strictly positive by the first two inequalities in
(RR10), because $h_0^2=2\beta_0/(2\beta_0+1)$.  The Gram on $(V,q)$
in (RR9) has positive diagonal residues and determinant residue
\[
 \frac14\left\{\frac1{(1+h)^2}-k^2b_0^2\right\}>0.       \tag{RR12}
\]
Thus both blocks have strictly positive residue.  The final term in
(RR9) is a bounded infinitesimal perturbation on $(x,z_0,\ell)$,
so the combined bounded Gram remains strict and is tensor-SOS.

If $b$ is infinitesimal, including $b=0$, choose any positive real
cap $\beta_0<\alpha$ and take $k=-1$.  There is now an exact
three-square identity
\[
 M_{-1}=\tfrac12\{x+a(z_0-\ell)\}^2
       +\beta_0(z_0-\ell)^2+\tfrac12(b\ell)^2.           \tag{RR13}
\]
On these three roots and $(V,q)$, (RR9) has positive diagonal
residue; the only off-block term is $D_vbVq$, whose coefficient
is infinitesimal.  The negative perturbation changes only the
$(z_0-\ell)^2$ coefficient by an infinitesimal.  Its Gram is again
strict.  Formula (RR7) restores all retained positive square terms.
The argument does not divide by $u,v,b$, or any state coordinate,
so the asserted literal boundaries are included.
\end{proof}

\paragraph{Scope and join.}
The two propositions cover precisely the complement of the region
where both testing angles are infinitesimal, for every strong-II
source with $B$ a unit.  They preserve independent source and testing
fields, arbitrary source openings, testing signs, and literal $n=0$.
Combined with a separately verified full-parent case of that
local region, they close the strong II/HH orbit by the exact C2
reflection.  They do not close the $B$-infinitesimal weak face.

\section{One small actual angle: bounded state scale}
\subsection{A lower bound on the state parameter}
\label{component:p7294aae40ee2}

\label{p7294aae40ee2:sec:c2-bounded-state-entry}

This short reduction concerns the same exact rational source, before
assuming that $b$ is a unit.  Suppose $x>0$ is infinitesimal, $k$ is
a positive bounded unit, $b>0$ is bounded, $0<r<1$, $0<h<b$,
$g=b-1+r+xh>0$, and $n\geq0$.  No further source physicality is
needed for the following necessary bound.

All quantities in
\[
             (b+r)A_cS_c\geq r\Delta
\]
are bounded, so taking residues gives
\[
 r_0\{b_0+k_0^2(b_0+1)\}\leq b_0.
\]
The brace is strictly positive even if $b_0=0$.  On the other hand
$g>0$ gives $r_0\geq1-b_0$.  If $b_0\leq1$, combining the two
inequalities yields
\[
 (1-b_0)\{b_0+k_0^2(b_0+1)\}\leq b_0,
 \quad\text{or equivalently}\quad
              b_0^2\geq\frac{k_0^2}{1+k_0^2}.
\]
For $b_0\geq1$ this lower bound is automatic.  Hence in every case
\[
                b_0\geq\frac{k_0}{\sqrt{1+k_0^2}}>0.
\]
Thus bounded $b$ is necessarily a unit.  In particular there is no
additional $b$-infinitesimal source domain in this actual-contact
scaling.  Combining the two new single-angle certificates, the
uncaptured part of this bounded-$b$ entry domain has exactly the
joint necessary degeneration $r_0=g_0=0$, hence $b_0=1$; the other
remaining scale is unbounded $b$.  This is a necessary source
reduction, not a certificate for the joint degeneration.

\subsection{A bounded-shape family at all tests}
\label{component:pec59340d644f}

\label{pec59340d644f:sec:c2-general-single-angle}

This calculation gives a source-domain extension of the fixed rational
family.  Its scalar uniqueness proof is an exact six-term Laurent
identity and a three-cycle arithmetic--geometric mean inequality.
The full-normal and outer-chart arguments include literal $n=0$.

Put $x>0$ infinitesimal in the coefficient field and $\rho=\sqrt{x}$.
Choose positive bounded coefficient units $k,b,r,h$ such that
\[
 1-r,\quad b-h,\quad g=b-1+r+xh,\quad b-k^2r
\]
are positive bounded units.  Set
\[
 c=\rho,\quad a=k/\rho,\quad t=x(b-h)/k^2,
 \quad j=1+xr-t,\quad a_2=k^2/x,
 \quad A_c=1+xr,\quad S_c=A_c+a_2xt.
\]
Thus $h=b-a_2t$, $g=(1+x)b-(1-r+a_2xt)$, and $j=1+xr-t$.
Use the exact original rational source
\begin{align*}
 \Delta={}&(1+x)xb^2+2x(1+x)br+(1+x)(1+a_2x)b
       +x^2(1-a_2)r^2+2x(1+a_2x)r+1+a_2x,\\
 N={}&(1+x)(b+r)(b+a_2(1-t))-a_2(1+a_2x)(1-r)t,\\
 n={}&\frac{(b+r)A_cS_c-r\Delta}{2r\Delta},\\
 m={}&\frac{(1+a_2x)(1-r)A_cS_c}{2g\Delta},\\
 C={}&\frac{A_cN}{2a_2t\Delta},\\
 A={}&\frac{a_2A_c(b+r)S_c}{2h\Delta},\\
 D={}&\frac{(1+x)(b+r)(b-a_2xr)S_c}{2a_2j\Delta},\\
 \alpha={}&\frac{x^2(b+r)N}{2(1-xh)\Delta}.
\end{align*}
Assume $n\geq0$ and physicality of this canonical source.  The latter
is the same inherited source-physicality hypothesis used for the
general noncommuting-residue exit; no conclusion below is predicated
on a new source enumeration.  Literal $n=0$ is allowed.

The parent, with $d=1$ and $\lambda=\beta=0$, is
\begin{align*}
 F={}&4nD_uD_vX^2+4mM^2+4CD_v(uW-Y)^2+4AD_uD_vY^2
                +4DJ^2+4\alpha D_u(vY-Z)^2+H,\\
 M={}&vW-X+uvY-uZ,\qquad J=uW+uvX-Y-vZ,\\
 H/2={}&J^2+D_u(vY-Z)^2+D_v(X+uZ)^2-D_uD_vZ^2,
 \qquad D_u=1+u^2,\quad D_v=1+v^2.
\end{align*}
Its exact stationary contact is
\[
 (u,v;W,X,Y,Z)=(\rho,k;\rho b/k,-\rho r,xh/k,1).
\]
Writing $\kappa=b/k$ and using a subscript zero for residues gives
\begin{align*}
 \Delta_0&=(1+k_0^2)(b_0+1),\qquad
 n_0=\frac{b_0+r_0-r_0\Delta_0}{2r_0\Delta_0}\geq0,\\
 m_0&=\frac{1-r_0}{2(b_0-1+r_0)(b_0+1)}>0,\\
 C_*:=(xC)_0&=\frac{k_0^2(b_0+r_0)}{2(b_0-h_0)\Delta_0},\qquad
 A_*:=(xA)_0=\frac{k_0^2(b_0+r_0)}{2h_0\Delta_0},\\
 D_*:=(D/x)_0&=\frac{(b_0+r_0)(b_0-k_0^2r_0)}{2k_0^2\Delta_0},\qquad
 E_*:=(\alpha/x)_0=\frac{k_0^2(b_0+r_0)^2}{2\Delta_0}.
\end{align*}
All starred quantities are positive.  Thus $n$ is bounded, $m$ is a
positive unit, $C,A$ have scale $x^{-1}$, and $D,\alpha$ have scale $x$.

\begin{theorem}[General bounded-shape single-angle source]
\label{pec59340d644f:thm:c2-general-single-angle}
Every source in the displayed domain has tensor SOS for the full
original parent at all independent real or complex projective testing
states and all projective testing angles.  No comparison between a
coefficient infinitesimal and a testing infinitesimal is needed.
The conclusion includes $n=0$.  Applying the exact C2 reflection gives
the corresponding original orientation with weak $C$ and strong
$\alpha$; this is a sufficient source domain, not a global one-weak
C2 classification.
\end{theorem}

\paragraph{The finite $u=0$ axis and its full normals.}
Set $a_0=4m_0$, $b_0'=a_0+2$, and
\[
 H_0(v)=(4n_0+2)(1+v^2)+a_0,\quad \eta=h/b,
 \quad \gamma=-r-\frac{a_0b-b_0'}{H_0(k)}.
\]
Here $\gamma$ is infinitesimal and $\eta_0=C_* /(C_*+A_*)$.
Use the exact contact graph
\[
 X_g=\frac{a_0vW-b_0'uZ}{H_0(v)}+\gamma uZ,
 \qquad Y_g=\eta uW.
\]
The denominator $H_0(v)$ is a pure testing unit and $H_0(k)$ is a
pure coefficient unit.  Let $N_X=X-X_g$, $N_Y=(Y-Y_g)/\rho$.
Writing $f=F|_g=f_{ww}W^2+f_{wz}WZ+f_{zz}Z^2$, exact monic divisions give
\[
 f=aP^2+b_s B^2+c_s V^2+eE^2+g_sBV,
 \qquad(P,B,V,E)=(uW/\rho,uZ,vW,\rho Z),
\]
where
\[
 \begin{gathered}
 a=x(f_{ww}-f_{ww}|_{u=0})/u^2,\qquad
 b_s=(f_{zz}-f_{zz}|_{u=0})/u^2,\\
 c_s=f_{ww}|_{u=0}/v^2,\qquad e=f_{zz}|_{u=0}/x,\qquad
 g_s=f_{wz}/(uv).
 \end{gathered}
\]
These displayed quotients are polynomial divisions, not inverses of
testing infinitesimals.  Parity gives the divisions by $u^2,v^2,uv$;
$f_{zz}|_{u=0}=4Dv^2+4\alpha$ gives the $x$ factor.  Only $C,A$ can
produce coefficient poles, and they occur in $f_{ww}$ with a factor
$u^2$; its prefactor $x$ removes them.

Put $Q=4C_*A_* /(C_*+A_*)>0$.  The bounded coefficient residues are
\begin{align*}
 q(v)=a_0^{\rm sc}&=Q D_v,&
 b(v)&=b_0'\frac{H_0-a_0-2}{H_0},\\
 \bar c(v)&=a_0\frac{H_0-a_0}{H_0},&
 g(v)&=-2b_0'\frac{H_0-a_0}{H_0},&
 e(v)&=4D_*v^2+4E_*.
\end{align*}
Thus $c(v):=v^2\bar c(v)$ and $d(v):=-vg(v)$ give the weighted
leading scalar
\[
 q(v)u^2W^2/x+b(v)u^2Z^2+c(v)W^2-d(v)uWZ+e(v)xZ^2.
\]
The full expansion is
\[
 F=f+N_XL_X+N_YL_Y+H_{XX}N_X^2+2H_{XY}N_XN_Y+H_{YY}N_Y^2.
\]
Its hard residue is $\operatorname{diag}(H_0(v),4D_v(C_*+A_*))$.
Direct differentiation gives
\[
 L_X=V\ell_w+B\ell_z,\qquad
 L_Y=P(B_2+xe_w)+B(\rho u e_z)+E e_*,
\]
where
\[
 B_2=8D_v\{(xA)\eta D_u-(xC)(1-\eta)\},\qquad
 e_*=8v(D-\alpha).
\]
All $\ell_w,\ell_z,e_w,e_z$ are bounded; the residues of
$\ell_w,\ell_z,B_2,e_*$ vanish identically in finite $v$.
For the second identity, after subtracting $uWB_2/x$ from $F_Y|_g$,
the remaining $W$ coefficient is $ue_w$ and the remaining $Z$
coefficient is $u^2e_z+e_*$.  The constant signed-core derivatives
cancel.  This proves full normal restoration wherever the scalar
Gram on these four roots is strict.

\paragraph{Global scalar strictness by an exact cycle identity.}
Let $\nu=2n_0+1\geq1$, $\mu=2m_0+1>1$, $z=v^2$,
$\mathcal H=\nu(1+z)+\mu-1$, $e_2=4D_*$ and $e_0=4E_*$.  Then
\[
 b=\frac{2\mu(\nu(1+z)-1)}{\mathcal H},\quad
 c=\frac{2(\mu-1)z\nu(1+z)}{\mathcal H},\quad
 d=\frac{4\mu v\nu(1+z)}{\mathcal H},\quad
 e=e_2z+e_0.
\]
For $s,z>0$ define
$\mathcal D(s,z)=(q s+c)(b s+e)-d^2s/4$.
The following exact identity is the key:
\begin{align*}
 \frac{\mathcal H\mathcal D(s,z)}{(1+z)sz}
 ={}&2Q\mu(\nu-1)s/z+2Q\mu\nu s
 +2(\mu-1)\nu e_2z/s+2(\mu-1)\nu e_0/s\\
 &+Qe_0(\nu+\mu-1)/z+Qe_2\nu z\\
 &+Q\{e_2(\nu+\mu-1)+e_0\nu\}-4\mu\nu.
\end{align*}
Every nonconstant coefficient is positive, except the first may be
zero when $n_0=0$.  Exact contact and stationarity, followed by taking
residues, give a stationary zero at $(s,z)=(1,k_0^2)$.  Here is an
algebraic proof that it is the unique zero on the whole positive
quadrant.  Rescale $z=k_0^2t$ and write the Laurent expression as
\[
 \mathcal L=a_\ell s/t+b_\ell s+c_\ell t/s+d_\ell/s
                      +e_\ell/t+f_\ell t-\Sigma,
 \qquad \Sigma=a_\ell+b_\ell+c_\ell+d_\ell+e_\ell+f_\ell.
\]
The constant term has this form by its zero value at $(1,1)$.
Stationarity gives
$a_\ell-c_\ell=d_\ell-b_\ell=f_\ell-e_\ell=:\tau$.
Put $p_\ell=\min(a_\ell,c_\ell)\geq0$,
$q_\ell=\min(b_\ell,d_\ell)>0$,
$r_\ell=\min(e_\ell,f_\ell)>0$.  Exactly,
\begin{align*}
 \mathcal L={}&p_\ell(s/t+t/s-2)+q_\ell(s+1/s-2)
                              +r_\ell(t+1/t-2)\\
 &+\begin{cases}
 \tau(s/t+1/s+t-3),&\tau\geq0,\\
 (-\tau)(t/s+s+1/t-3),&\tau\leq0.
 \end{cases}
\end{align*}
Each bracket is nonnegative by the two- or three-term
arithmetic--geometric mean inequality.  The two strictly positive
pair terms force $s=t=1$ at equality.  They also give a positive
definite Hessian there: the other brackets are nonnegative and
vanish at that point.  Thus both remote uniqueness and local
nondegeneracy are proved by the same finite algebraic identity, over
any real closed residue field.  No inference from a merely local
Hessian is being used.

Consequently, for finite nonzero $v$ with $v^2\ne k_0^2$,
\[
 |d(v)|<2\sqrt{b(v)c(v)}+2\sqrt{q(v)e(v)}.
\]
Indeed the minimum of $\mathcal D(s,z)/s$ over $s>0$ is
$(\sqrt{bc}+\sqrt{qe})^2-d^2/4$ and is attained, so positivity is
strict.  Split the exact cross term using $BV=vPE$ and the pure
testing function
\[
 \theta=\frac{|v|\sqrt{b(v)\bar c(v)}}
 {|v|\sqrt{b(v)\bar c(v)}+\sqrt{q(v)e(v)}}.
\]
The resulting residue blocks on $(B,V)$ and $(P,E)$ are strictly
positive.  The denominator is a pure testing unit.  At $v_0=0$ the
same conclusion holds when $n_0>0$, with zero cross residues.
The case $n_0=0,v_0=0$ is treated separately below.

\paragraph{The matched finite-angle residue.}
For $v_0=k_0$, normalize $Z=1$ when $Z\ne0$.  If $W$ is
infinitesimal, put
\[
 R=(uW-\kappa x)/\rho,\quad S=u-\kappa^{-1}W,
 \quad U=\rho(v-k),\quad T=u(v-k).
\]
The exact conductors are
\[
 u^2-x=uS+\kappa^{-1}\rho R,\qquad
 W^2-\kappa^2x=-\kappa WS+\kappa\rho R,
 \qquad uW-\kappa x=\rho R.
\]
Write $\xi=u^2,\omega=W^2,p=uW$ and
\[
 f=p^2Q_b/x+B_b,\qquad
 Q_b=4D_v\{(xC)(1-\eta)^2+(xA)\eta^2D_u\}.
\]
Here $B_b$ is bounded and polynomial in these three invariants with
only the separate unit denominators already specified.  Put
$G=\kappa^2xQ_b+B_b$ and take successive monic differences
\[
 G=G_0(v)+(\xi-x)G_\xi+(\omega-\kappa^2x)G_\omega
                    +(p-\kappa x)G_p.
\]
Exactly,
\[
 f=R^2Q_b+\kappa^2S^2G_\omega+\rho R\mathcal A+uS\mathcal B+G_0,
\]
where $\mathcal A=2\kappa Q_b+\kappa^{-1}G_\xi+
\kappa G_\omega+G_p$ and $\mathcal B=G_\xi-\kappa^2G_\omega$.
The contact Jacobian of $(R,S)$ with respect to $(u,W)$ has
determinant $-2$.  Exact stationarity therefore gives
$\mathcal A_c=\mathcal B_c=0$.  Cancelling $\rho$ here proves
vanishing and does not introduce a Gram coefficient inverse.
Monic contact division of these two functions and the conductors
above give a bounded Gram on $(R,S,U,T)$.
Also $G_0=xg_0(v)$ with bounded $g_0$, and exact contact gives a
double factor $(v-k)^2$, producing its bounded $U^2$ coefficient.

The full normals have zero-residue lifts on these roots.  For $L_X$,
write $L_X=uA_1+WB_1=-\kappa B_1S+u(A_1+\kappa B_1)$.
Both coefficient functions have identically zero leading part for
every $v$; exact normal contact makes $A_1+\kappa B_1$ vanish at
contact.  Monic division gives coefficients with factors $u\rho$,
$u^2$, or $uW$, and a $(v-k)$ quotient of an identically zero
leading function.  All have zero residue.  For $L_Y$, write
$L_Y=RB_2+\rho E_2$.  Its bounded even $E_2$ vanishes exactly at
contact, and both $B_2,E_2$ have identically zero leading part.
Monic division gives factors $x,\rho u,\rho W$, or the same
zero-leading-function quotient.  Thus the entire normal restoration
is bounded and has zero residue.

The soft residue has the form
\[
 q(k_0)R^2+\kappa_0^2c(k_0)S^2+A_1'RU+B_1'ST+E_1'U^2.
\]
It satisfies
\[
 E_1'>\frac{(A_1')^2}{4q(k_0)}+
          \frac{(B_1')^2}{4\kappa_0^2c(k_0)}.
\]
To prove this strict inequality, the leading scalar at
$u=\rho s,W=\rho w$ is
$\Phi=q s^2w^2+b s^2+c w^2-dsw+e$.
Complete its $w$ square:
\[
 \Phi=(qs^2+c)\left(w-\frac{ds}{2(qs^2+c)}\right)^2
           +\frac{\mathcal D(s^2,v^2)}{qs^2+c}.
\]
The global Laurent calculation gives a positive definite Hessian of
the second term in $(s,v)$ at $(1,k_0)$; both coordinate derivatives
are nonzero there.  The first square supplies the missing $w$
direction.  Hence the full three-variable Hessian is positive
definite.  On first variations $T=U$; the displayed inequality is
exactly its Schur complement after the invertible $(R,S)$ change.
Choose a real $\delta$ strictly between the second summand and
$E_1'-(A_1')^2/(4q(k_0))$.  Add $\delta$ times the exact zero
\[
 T^2-U^2-ST(v-k)-\kappa^{-1}RU(v-k)=0.
\]
The two blocks $(R,U)$ and $(S,T)$ are now strictly positive, and so
is the full six-root residue Gram.

If the ratio $W/Z$ is not infinitesimal, instead normalize $W=1$;
then $Z$ is bounded.  The earlier scalar is $aP^2$ plus a term of
positive residue $c(k_0)$; its normal lifts have zero residue on
$(P,1)$.  This is strict.  The case $Z=0$ is a literal sum of the
event squares and the three positive core squares.  The simultaneous
sign change $(u,v,X,Y)\mapsto(-u,-v,-X,-Y)$ covers $v_0=-k_0$.

\paragraph{The double-zero corner when $n_0=0$.}
Use instead the pure coefficient quantities
\[
 H_c=4n+4m+2,\quad \xi_c=4m/H_c,\quad
 \zeta=(4m+2)/H_c,\quad \eta_c=C/(C+A),
\]
and graph $X_g=\xi_cvW-\zeta uZ$, $Y_g=\eta_cuW$.
Put $\epsilon=4\alpha+4n(4m+2)/H_c>0$.
Then $\epsilon$ is infinitesimal, while $n/\epsilon$ and
$x/\epsilon$ are bounded.  Use the hard roots as before and scalar
roots
\[
 P=uW/\rho,\quad V=vW,\quad E=\rho Z,\quad
 B_\epsilon=\sqrt\epsilon\,uZ,\quad T=uvZ.
\]
Direct positive square term expansion gives exactly
\[
 f_{zz}=4\alpha+\epsilon u^2+4n\zeta^2u^4+4Dv^2
                         +u^2v^2 B_b,
\]
where $B_b$ is bounded with residue $2$.  Its first four terms
give coefficients $4\alpha/x$, $1+(4n\zeta^2/\epsilon)u^2$,
and $(4D/x)v^2$ on $E^2,B_\epsilon^2,E^2$ respectively.  The
$W^2$ contribution has bounded coefficients on $P^2,V^2$ with
positive residues $Q$ and $4m_0/(2m_0+1)$.  The $WZ$ contribution
is $uvWZ$ times a bounded function; use $uvWZ=vPE$ to give it zero
cross residue.  Thus the five scalar roots have strict diagonal
residue $(Q,4m_0/(2m_0+1),4E_*,1,2)$.

For completeness, $L_X$ has a zero-residue bounded coefficient on
$V$.  At $v=0$ its $Z$ coefficient is exactly $-8n\zeta u^3$;
on $B_\epsilon$ this is $-8n\zeta u^2/\sqrt\epsilon$, which is
infinitesimal because $n/\epsilon$ is bounded.  The remaining
$Z$ coefficient has a factor $uv^2$, giving a zero-residue multiple
of $T$.  For $L_Y$, the large-positive square term coefficient is now exactly
$B_2=8D_v(xA)\eta_cu^2$; the remaining terms lift as
\[
 L_Y=P(B_2+xe_w)+T(\rho u e_z)+E\,8v(D-\alpha),
\]
with bounded $e_w,e_z$.  All coefficients have zero residue.  The
hard residue is $\operatorname{diag}(4m_0+2,4(C_*+A_*))$.
This includes literal $n=0$ without deleting a coordinate.

\paragraph{The remaining commuting charts.}
The following exact core identities retain all unused event positive square terms.
Every retained coefficient is nonnegative; in particular the $n$
positive square term can be retained even when $n=0$.  They prove strict bounded
Grams without dividing by a testing infinitesimal.

If $u$ is a finite unit and $v\to0$, take
\[
 R=(uW-Y)/\rho,\quad L=Y/\rho,\quad B=X+uZ,\quad V=vZ.
\]
Then $M=-B+(\rho v/u)(R+D_uL)$, and exactly
\begin{align*}
 H/2={}&(1+D_uv^2)B^2+u^2D_uV^2-2uvD_uBV\\
 &+2\rho R(uvB-D_uV)+xR^2-2D_u\rho LV+D_uxv^2L^2.
\end{align*}
Retain $n,D,\alpha$.  The residue on $(B,V,R,L)$ is diagonal
with entries $(4m_0+2,2u_0^2D_{u,0},4C_*,4A_*D_{u,0})$.

If $u\to\infty$ and $v$ is finite, put $p=1/u$ and use $p^2F$.
Take $R=(W-pY)/\rho$, $L=Y/\rho$, $B=Z+pX$, $V=vX$.
The homogenized $m$ amplitude is
$-B+\rho pvR+\rho vD_pL$, and the exact core is
\begin{align*}
 \widehat H/2={}&B^2+D_pV^2+2\rho D_pRV-2\rho pvRB+xR^2\\
 &+2D_p\rho pLV-2D_p\rho vLB+D_pxv^2L^2.
\end{align*}
Retain $n,D,\alpha$.  The residue on $(B,V,R,L)$ is diagonal
with entries $(4m_0+2,2,4C_*D_{v,0},4A_*D_{v,0})$.

If $v\to\infty$ and $u$ is bounded, put $q=1/v$ and use $q^2F$.
Take $R=(uW-Y)/\rho$, $L=Y/\rho$, $T=uqZ$, $N_1=\rho Z$.
Exactly,
\begin{align*}
 \widehat H/2={}&(D_u+q^2)X^2+D_uxL^2+T^2+2qXT
       +q^2xR^2+2uq\rho XR-2qRN_1-2D_uqLN_1,\\
 4D\widehat J^2={}&4(D/x)(-N_1+\rho uX+xqR)^2.
\end{align*}
Retain $n,m,\alpha$.  The residue on $(X,R,L,T,N_1)$ is diagonal
with entries
\[ (2D_{u,0},4C_*,4A_*D_{u,0},2,4D_*). \]

Finally if both angles are infinite, use $p=1/u,q=1/v$ and
$p^2q^2F$.  Set $R=(W-pY)/\rho$, $L=Y/\rho$, $T=qZ$.  Then
\[
 \widehat H/2=(1+p^2D_q)X^2+xq^2R^2+xD_pL^2+T^2
       +2\rho qXR+2pqXT-2\rho pRT-2\rho D_pLT.
\]
Retain $n,m,D,\alpha$.  The residue on $(X,R,L,T)$ is
diagonal with entries $(2,4C_*,4A_*,2)$.

\paragraph{Exhaustion and independence.}
These charts cover every commuting angle residue: the whole finite
$u=0$ axis, the finite-unit $v=0$ axis, either infinite axis, and
their intersections.  When both angle residues are noncommuting,
apply the inherited physical noncommuting-residue theorem at arbitrary
coefficient scales.  Complex states require no additional chart:
if $Z\ne0$, multiply by a phase and normalize $Z=1$; the real parent
is the sum of its real-state value and its imaginary-state value,
the latter having $Z=0$ and hence being a literal square sum.
If $Z=0$ this is already true of both parts.  Homogeneity covers all
projective state charts.  Roots may be unbounded throughout.
All inverses or radicals used in coefficients are pure coefficient
or pure testing operations on units, except $\sqrt\epsilon$ in
roots and its controlled coefficient ratio $n/\sqrt\epsilon$.
No mixed inverse and no comparison of independent infinitesimals
occurs.  Applying the strict bounded-tensor Gram lemma completes
the proof.

\subsection{Three arbitrarily small source openings}
\label{component:p81893ad08bed}

\label{p81893ad08bed:sec:c2-three-openings}

This theorem removes three unit assumptions from the source of
Theorem~\ref{pec59340d644f:thm:c2-general-single-angle}.  Both positive openings $h$
and $b-h$ may approach zero, and $r$ may approach zero, making $n$
arbitrarily large.  The conclusions concern the entire original parent
and every independent test.  Source physicality remains an explicit
hypothesis; no new source classification is claimed.

\begin{theorem}[Three-opening single-angle source]
\label{p81893ad08bed:thm:c2-three-openings}
Use exactly the rational C2 source and contact of
Section~\ref{pec59340d644f:sec:c2-general-single-angle}.  Suppose $x>0$ is
infinitesimal, $k,b$ are positive bounded units,
\[
             0<r<1,\qquad 0<h<b,
\]
and $g=b-1+r+xh$ is a positive bounded unit.  Assume $n\geq0$
and physicality of the original source.  There is no lower bound on
$r,h,b-h$.  Then the entire original parent has tensor SOS at all
independent projective measurement pairs and all real or complex
projective states.  Literal $n=0$ is included.
\end{theorem}

Write $\rho=\sqrt x$, $\eta=h/b$, and retain the source quantities
$A_c,S_c,\Delta$ from that section.

\paragraph{Entry and the two complementary openings.}
First the former assumptions that $1-r$ and $b-k^2r$ are units follow
from $n\geq0$.  Indeed all quantities in
\[
                 (b+r)A_cS_c\geq r\Delta
\]
are bounded, and $\Delta_0=(1+k_0^2)(b_0+1)$.  Taking residues yields
\[
 r_0\leq\frac{b_0}{b_0+k_0^2(b_0+1)}<1,
 \qquad
 b_0-k_0^2r_0\geq
 \frac{b_0^2(1+k_0^2)}{b_0+k_0^2(b_0+1)}>0.
\]
Thus $m$ is a positive bounded unit, and $D/x,\alpha/x$ are positive
bounded units, regardless of whether $r$ has zero residue.

Set
\[
 T_C=xC(b-h),\qquad T_A=xAh,\qquad H=x(C+A).
\]
The original source formulas, with their apparent opening denominators
cancelled, show that $T_C,T_A$ are bounded positive units and have the
same residue
\[
 T_0=\frac{k_0^2(b_0+r_0)}{2\Delta_0}>0.
\]
Consequently $H$ is bounded below by a positive real, although it may
be unbounded.  The four identities
\begin{align*}
 xC(1-\eta)&=T_C/b,&xA\eta&=T_A/b,\\
 xC(1-\eta)^2&=T_C(b-h)/b^2,&
 xA\eta^2&=T_Ah/b^2                                      \tag{CO1}
\end{align*}
are the reason no new scalar singularity appears when either opening
collapses.  In particular their last two terms add to a bounded unit,
with residue $T_0/b_0$.  The effective scalar coefficient is
\[
                         Q=4T_0/b_0>0.                  \tag{CO2}
\]
It is independent of the limiting value of $h$.

\paragraph{A bounded $n$: restoration with the normalized hard root.}
If $r$ is a unit, the exact formula for $n$ makes $n$ bounded.  Use the
contact graph $X_g,Y_g=\eta uW$ of the prior theorem.  All of its scalar
coefficients remain bounded: the only potentially unbounded weights
are $C,A$, and on the graph they occur solely through the last two
products in (CO1), divided by $x$ and multiplied by $u^2W^2$.
Its leading scalar is exactly the scalar in that theorem with $Q$
replaced by (CO2).  The quantities $n_0,m_0,D_*,E_*$ have their former
values.  Hence the entire six-term Laurent identity, its balanced
AM--GM decomposition, global equality characterization, and positive
contact Hessian apply without change.

Replace the old hard root $N_Y=(Y-Y_g)/\rho$ by
\[
        \widehat N_Y=\sqrt H\,N_Y=\sqrt{C+A}(Y-Y_g).
                                                               \tag{CO3}
\]
Its exact hard coefficient from the $C,A$ positive square terms is
\[
             4D_v\frac{C+AD_u}{C+A}
       =4D_v\left(1+\frac{A}{C+A}u^2\right),             \tag{CO4}
\]
which is bounded and has residue $4D_{v,0}$ on the finite $u=0$
residue axis.  The other positive square term contributions to this coefficient
and to the mixed hard entry are bounded: $1/(C+A)=x/H$ and
$1/\sqrt{C+A}=\rho/\sqrt H$ are bounded and infinitesimal.

For clarity the complete first-normal verification is also retained.
The previous unnormalized $Y$ derivative has exactly the form
\[
 L_Y=P(B_2+xe_w)+B(\rho u e_z)+E e_*,\qquad
 B_2=8D_v\{(xA)\eta D_u-(xC)(1-\eta)\},
 \quad e_*=8v(D-\alpha).                                \tag{CO5}
\]
Here $P=uW/\rho$, $B=uZ$, $E=\rho Z$.  The functions $e_w,e_z$
come only from the bounded $n,m,D,\alpha$ positive square terms and the signed
core, so remain bounded.  Equation (CO1) makes $B_2$ bounded with
identically zero residue for every finite $v$ on $u=0$; $e_*$ has
the same property.  The coefficient of $\widehat N_Y$ is exactly
$L_Y/\sqrt H$, so its bounded zero-residue lifts are preserved.
The $X$ derivative contains neither $C$ nor $A$ and is unchanged.
Thus all four-root remote certificates restore both full normals.

For matched finite-angle residues the scalar has exactly the old
form $f=p^2Q_b/x+B_b$, with
\[
 Q_b=4D_v\{(xC)(1-\eta)^2+(xA)\eta^2D_u\}.
\]
Both $Q_b$ and $B_b$ have bounded coefficients by (CO1).  The monic
divisions in $u^2-x$, $W^2-\kappa^2x$, $uW-\kappa x$, and $v-k$,
where $\kappa=b/k$, therefore give precisely the same bounded Gram
on $R,S,U,T$.  Its residue is unchanged.  The old first-normal
contact divisions use only bounded $B_2,e_w,e_z$ and exact stationary
contact; multiplication by the pure coefficient $1/\sqrt H$
preserves their zero-residue lifts.  Hence the complete six-root
Gram has strictly positive residue.

If $n_0=0$ at the double-zero test, use the previous exact graph
with $\eta_c=C/(C+A)$.  Now
\[
 xC(1-\eta_c)=xA\eta_c=\frac{xCA}{C+A},
 \qquad
 \left(\frac{xCA}{C+A}\right)_0=T_0/b_0.                \tag{CO6}
\]
These quantities are bounded units even when $H$ is unbounded.
Thus the exact $\epsilon$ reserve and five scalar roots in that
proof remain valid.  Its $Y$ normal is again divided by $\sqrt H$,
and its hard $Y$ residue is now $4$ rather than $4(C_*+A_*)$.
This proves the literal $n=0$ case as well as its infinitesimal
neighborhood without deleting a state coordinate.

\paragraph{An infinitesimal $r$: the opposite event can be unbounded.}
Suppose now $r$ is infinitesimal.  Let $N=4n+2$.  The exact source
formula gives
\[
             Nr=\frac{2(b+r)A_cS_c}{\Delta}.             \tag{CO7}
\]
Thus $Nr$ is a bounded positive unit, $N$ is unbounded, and both
$nr$ and $nr^2$ are bounded.  Use the exact-contact graph
\[
 X_g=-ruZ,\qquad Y_g=\eta uW,
 \qquad
 \widehat N_X=\sqrt N(X+ruZ),\quad
 \widehat N_Y=\sqrt{C+A}(Y-\eta uW).                    \tag{CO8}
\]
At the original contact $u=\rho,v=k,W=\rho b/k,Z=1$ this graph
has exactly the required $X,Y$ values.  The potentially large $n$
positive square term contributes $4nr^2D_uD_vu^2Z^2$ to the graph scalar, and
$-8nrD_uD_vuZ$ to its unscaled $X$ derivative.  Equation (CO7)
therefore controls the entire scalar and every full first normal,
not only their values at contact.

The hard residue on $(\widehat N_X,\widehat N_Y)$ is
\[
                  \operatorname{diag}(D_{v,0},4D_{v,0}). \tag{CO9}
\]
Indeed $H_{XX}/N=D_uD_v+O(1/N)$, since the remaining source
weights that enter $H_{XX}$ are $m,D$; (CO4) controls $H_{YY}$.
The mixed entry is bounded times
$1/(\sqrt N\sqrt{C+A})$ and has zero residue.

The five exact scalar monic divisions on $P,B,V,E$ from the previous
proof still apply.  Their residues, with $z=v^2$, are now
\[
 q=Q(1+z),\quad b_s=2\mu,\quad c_s=2(\mu-1)z,
 \quad d_s=4\mu v,\quad e=E_2z+E_0,                    \tag{CO10}
\]
where $\mu=2m_0+1>1$, $E_2=4(D/x)_0>0$,
$E_0=4(\alpha/x)_0>0$, and $Q$ is given by (CO2).
Define
$\mathcal D(s,z)=(qs+c_s)(b_ss+e)-d_s^2s/4$.
Exactly,
\begin{align*}
 \frac{\mathcal D(s,z)}{sz}={}&
 2Q\mu s/z+2Q\mu s
 +2(\mu-1)E_2z/s+2(\mu-1)E_0/s\\
 &+QE_0/z+QE_2z+Q(E_0+E_2)-4\mu.                    \tag{CO11}
\end{align*}
All six nonconstant Laurent coefficients are positive.  Exact
stationary contact, applied after the bounded graph substitution
$u=\rho s_1,W=\rho w_1$, gives a stationary zero at
$(s,z)=(1,k_0^2)$.  Alternatively this follows by direct substitution
of
\[
 m_0=\frac1{2(b_0-1)(b_0+1)},\quad
 Q=\frac{2k_0^2}{\Delta_0},\quad
 E_2=\frac{2b_0^2}{k_0^2\Delta_0},\quad
 E_0=\frac{2k_0^2b_0^2}{\Delta_0}.
\]
Here $b_0>1$ follows from $g_0=b_0-1>0$.
Rescale $z=k_0^2t$ and use exactly the balanced reciprocal-pair and
three-cycle identity of the previous theorem.  Its $s+1/s-2$ and
$t+1/t-2$ coefficients are strictly positive.  Thus (CO11) has a
unique zero in the positive quadrant and a positive definite Hessian
there.  This proves remote scalar strictness and the matched-contact
Schur inequality, with no appeal to numerical evidence.

For the full $X$ normal, direct differentiation and parity give
\[
 \left.F_X\right|_g=V A_X+B B_X,
\]
where the complete coefficients are
\begin{align*}
 A_X={}&-8m(1+\eta u^2)+(8D+4)(1-\eta)u^2,\\
 B_X={}&-8nrD_uD_v+8m(1-r)
       -(8D+4)v^2(1+ru^2)+4D_v(1-r).
\end{align*}
These are bounded polynomials in the bounded testing coordinates.
Consequently the coefficient of $\widehat N_X$ has bounded
zero-residue coefficients on $V,B$ because $1/\sqrt N$ is
infinitesimal.  The full $Y$ normal is (CO5), with the bounded
graph (CO8); (CO1) gives all of its required bounds and zero residues.
Its complete residual coefficients are
\begin{align*}
 e_w={}&8mv^2(1+\eta u^2)-(8D+4)(1-\eta)
                     +(8\alpha+4)D_uv^2\eta,\\
 e_z={}&v\{-8m(1-r)+(8D+4)r-(8\alpha+4)\}.
\end{align*}
This restores the remote scalar certificates, including $v_0=0$,
where (CO10) has strictly positive coefficients on $P,B,V,E$ and
zero cross residues.  No special $n_0=0$ corner occurs in this case.

At a matched residue write, exactly,
\[
 \frac{F_X|_g}{\sqrt N}=uA_1+WB_1
   =-\kappa B_1S+u(A_1+\kappa B_1).
\]
Both $A_1,B_1$ have identically zero residue for every finite $v$;
their finite monic coquotients also retain the factor $1/\sqrt N$.
Exact normal contact makes $A_1+\kappa B_1$ vanish at the source
contact.  The same four conductors $R,S,U,T$ therefore give bounded
zero-residue coefficients, exactly as in the bounded-$n$ argument.
For the $Y$ normal the previous division
$L_Y=RB_2+\rho E_2'$ has bounded $B_2,E_2'$ with identically zero
leading functions; division by $\sqrt H$ preserves those properties.
All quadratic normal coefficients are already bounded by
(CO4),(CO7),(CO9).  Thus the matched-contact Gram for the entire
parent is strictly positive in residue.  The other state charts and
the simultaneous sign copy use precisely the previous proof.

\paragraph{Every other commuting test.}
Choose the pure coefficient
\[
                         \sigma=\min(1,xC,xA).
\]
It is a positive bounded unit, since $T_C,T_A$ are units while
$h,b-h$ are bounded.  Retain the nonnegative excess positive square terms
\[
 4(C-\sigma/x)D_v(uW-Y)^2,
 \qquad4(A-\sigma/x)D_uD_vY^2,
\]
and use the four outer core identities in the previous theorem with
$C=A=\sigma/x$.  These identities retain the entire $n$ positive square term;
they never use boundedness of $n$.  Their displayed hard entries
$4C_*,4A_*$ are simply replaced by $4\sigma_0$, while $m_0>0$
and $(D/x)_0>0$ are unchanged.  Every residue Gram stays strictly
positive.  This covers the finite-unit $u$ axis, both reciprocal
axes and their double corners, including literal projective ends.

The finite $u=0$ axis, those outer axes, and the inherited theorem for
noncommuting-residue tests exhaust all independent measurements.
Restore the retained positive square terms and positive homogeneous factors, and
split complex states into real and imaginary parts.  This proves
Theorem~\ref{p81893ad08bed:thm:c2-three-openings}.  In particular the newly proved
source scales permit either $A$ or $C$ to be arbitrarily larger than
$x^{-1}$, independently of an arbitrarily large $n$.  They cannot
both have unbounded ratios to this scale because $h+(b-h)=b$ is a unit.  There is no
prescribed comparison between these coefficient scales and any
testing infinitesimal.

\subsection{The remaining opening with bounded opposite event}
\label{component:pd7bd7f0d7fb6}

\label{pd7bd7f0d7fb6:sec:c2-last-opening}

\begin{theorem}[The unbounded-$m$ opening with bounded $n$]
\label{pd7bd7f0d7fb6:thm:c2-no-opening-bounds}
In Theorem~\ref{p81893ad08bed:thm:c2-three-openings}, replace the assumption that
$g=b-1+r+xh$ is a positive bounded unit by $g>0$ and the assumption that at least one of $r,g$ is a bounded
unit.  The all-tests tensor-SOS conclusion still holds.
\end{theorem}

The only new case is infinitesimal $g$; it is automatically bounded.
Our hypothesis now makes $r$ a unit.  Since $b_0+r_0=1$ and $b$ is
a positive unit, $r_0<1$.  Thus $n$ is bounded,
$0<b_0<1$, and $1-r$ is a unit.  The source formulas imply
\[
               mg=\frac{(1+k^2)(1-r)A_cS_c}{2\Delta},   \tag{CG1}
\]
which is a positive bounded unit.  Thus $m$ is unbounded.  The
two-opening bounds (CO1), positivity of $D/x,\alpha/x$, and the
consequence that $b-k^2r$ is a unit remain valid.

\paragraph{Use the original strong amplitude as the hard normal.}
Put $\eta=h/b$ and use the exact contact graph
\[
 Y_g=\eta uW,\qquad
 X_g=v(1+\eta u^2)W-(1+g)uZ.                            \tag{CG2}
\]
At contact this gives $X_g=\rho[b+xh-1-g]=-\rho r$.
The original $m$ amplitude satisfies exactly $M|_g=guZ$.
Normalize the full hard directions by
\[
 E_M=\sqrt m\{X-X_g-uv(Y-Y_g)\},\qquad
 E_Y=\sqrt{C+A}(Y-Y_g).                                 \tag{CG3}
\]
In these variables
\[
 M=guZ-E_M/\sqrt m,
 \quad
 X=X_g+E_M/\sqrt m+uvE_Y/\sqrt{C+A}.
\]
Consequently the unbounded positive square term has no $E_Y$ contribution at all.
This triangular change is essential: scaling $X$ alone would leave
the uncontrolled coefficient $\sqrt{m/(C+A)}uv$.
The entire scalar contribution of that positive square term is $4mg^2u^2Z^2$,
which is bounded by (CG1).  Its full $X$ derivative contributes
$-8mg\,uZ$, also bounded, before division by $\sqrt m$.

All other source weights entering the $X$ row are $n,D$ and the
unit signed weight, so the full $X$ derivative on (CG2) has bounded
coefficients on $V=vW,B=uZ$.  Hence the $E_M$ first normal has
zero-residue bounded coefficients on those roots.  The $E_Y$ first
normal is
\[
       \frac{(F_Y+uvF_X)|_g}{\sqrt{C+A}}.               \tag{CG4}
\]
The $m$ positive square term cancels identically in $F_Y+uvF_X$, before any graph
substitution.  The remaining large terms are exactly the $C,A$
terms controlled by (CO1); all other coefficients are bounded.
In particular the old (CO5) form on $P,B,E$ holds with bounded
coquotients and identically zero leading corrections.
The complete quadratic hard block is bounded and has residue
\[
                         \operatorname{diag}(4,4D_{v,0})  \tag{CG5}
\]
on the finite $u=0$ axis.  The $m$ positive square term gives precisely its first
diagonal entry, the $C,A$ positive square terms give the second, and all mixed
entries from bounded positive square terms contain $1/(\sqrt m\sqrt{C+A})$.

\paragraph{The remaining five Laurent terms already force strictness.}
The exact scalar divisions on $P,B,V,E$ are bounded for the same
reason: $C,A$ enter through (CO1), and $m$ only through $mg^2$.
Their leading weighted scalar is
\[
 q\,u^2W^2/x+b_su^2Z^2+c_sW^2-d_suWZ+exZ^2,
\]
where, for $z=v^2$ and $\nu=2n_0+1\geq1$,
\[
 \begin{gathered}
 q=Q(1+z),\quad b_s=2\nu(1+z)-2,\quad
 c_s=2z\nu(1+z),\quad d_s=4v\nu(1+z),\\
 e=E_2z+E_0,\qquad Q=4T_0/b_0,
 \quad E_2=4(D/x)_0>0,\quad E_0=4(\alpha/x)_0>0.
 \end{gathered}                                         \tag{CG6}
\]
Set $\mathcal D(s,z)=(qs+c_s)(b_ss+e)-d_s^2s/4$.
Its exact Laurent identity is
\begin{align*}
 \frac{\mathcal D(s,z)}{(1+z)sz}={}&
 2Q(\nu-1)s/z+2Q\nu s
 +2\nu E_2z/s+2\nu E_0/s\\
 &+QE_0/z+QE_2-4\nu.                                  \tag{CG7}
\end{align*}
Contact and stationarity give a zero at $(s,z)=(1,k_0^2)$.
After $z=k_0^2t$, write the five coefficients in order as
$a,b,c,d,e$.  Here $a\geq0$ and $b,c,d,e>0$.  Stationarity gives
$c=a+e$ and $b=d+e$.  The zero value fixes the constant, so exactly
\[
 \mathcal L=
 a(s/t+t/s-2)+d(s+1/s-2)+e(t/s+s+1/t-3).                 \tag{CG8}
\]
Every summand is nonnegative by AM--GM.  The last bracket alone
forces $s=t=1$ at equality: its three positive terms have product
one, and equality makes all three equal.  Its exponent vectors
$(-1,1),(1,0),(0,-1)$ span two dimensions; equivalently its Hessian
in $\log s,\log t$ at $(1,1)$ is
$\left(\begin{smallmatrix}2&-1\\-1&2\end{smallmatrix}\right)$.
Hence it already has a positive definite contact Hessian.
This proves global uniqueness and local nondegeneracy even when
$n_0=0$ and the first pair is absent.

For remote finite nonzero $v$ the same scalar cross-term split as
before gives a strict Gram.  The case $v_0=0,n_0>0$ is also strict
on $P,B,V,E$.  The full normals are restored by (CG3)--(CG5).
At matched finite-angle residues the whole scalar still has the
form $p^2Q_b/x+B_b$ with bounded $Q_b,B_b$.  The four monic
conductors $R,S,U,T$ of the prior theorem therefore give a bounded
Gram, and (CG8) proves its strict Schur inequality.

The full first normals have the required same conductor lifts.
For $E_M$, their coefficient is $(uA_1+WB_1)/\sqrt m$ with bounded
polynomial $A_1,B_1$.  Monic contact division retains the factor
$1/\sqrt m$ in every coquotient.  Exact contact makes the residual
coefficient $A_1+\kappa B_1$ vanish, giving the old four-root lift
with zero residue.  For $E_Y$, the cancellation in (CG4) leaves
precisely the old large-positive square term term $B_2$ plus bounded residual
terms; its exact contact and identically zero leading part give
the same $R B_2+\rho E_2'$ division, then the bounded factor
$1/\sqrt H$.  These are full normal identities, not just Hessian
checks.  Quadratic state degree and (CG5) finish the complete
matched-contact certificate.

\paragraph{The double-zero test when $n_0=0$.}
The graph (CG2) is replaced in this remote corner.  Put
\[
 H_c=4m+4n+2,\quad \xi=4m/H_c,\quad
 \zeta=(4m+2)/H_c,\quad \eta_c=C/(C+A).
\]
Use
\[
 Y_g=\eta_cuW,\quad
 X_g=\xi v(1+\eta_cu^2)W-\zeta uZ,
\]
and the triangular normals
\[
 E_X=\sqrt{H_c}\{X-X_g-uv(Y-Y_g)\},\qquad
 E_Y=\sqrt{C+A}(Y-Y_g).                                \tag{CG9}
\]
The entire graph amplitude is
\[
 M|_g=(1-\xi)v(1+\eta_cu^2)W+(\zeta-1)uZ.
\]
Both $m(1-\xi)$ and $m(\zeta-1)$ are bounded because $n$ is
bounded.  Thus the scalar and full first normals have bounded
coquotients after the same $x$ scaling as before.

Set $\epsilon=4\alpha+4n\zeta>0$.  It is infinitesimal, with
$n/\epsilon,x/\epsilon$ bounded.  The complete scalar $Z^2$
coefficient is exactly
\[
 f_{zz}=4\alpha+\epsilon u^2+4n\zeta^2u^4+4Dv^2
                            +u^2v^2 B_b,               \tag{CG10}
\]
where $B_b$ is bounded with residue $2$.  This identity follows
directly from the positive square terms and the exact relation
$4n\zeta^2+(4m+2)(1-\zeta)^2=4n\zeta$.
The $W^2$ coefficients on $P=uW/\rho,V=vW$ have positive residues
$Q,2$.  The cross term is $uvWZ$ times a bounded function and
can be written as a zero-residue multiple of $PE$ using
$uvWZ=vPE$.  Thus the five scalar roots
\[
           P,\quad V,\quad E=\rho Z,\quad
           B_\epsilon=\sqrt\epsilon\,uZ,\quad T=uvZ
\]
have strict diagonal residue $(Q,2,E_0,1,2)$.

For the full $X$ derivative, at $v=0$ the $Z$ coefficient is exactly
$-8n\zeta u^3$.  After division by $\sqrt{H_c}$ its coefficient
on $B_\epsilon$ is infinitesimal because
$n/\sqrt\epsilon=(n/\epsilon)\sqrt\epsilon$ is infinitesimal.
The remaining $Z$ coefficient has the factor $uv^2$ and bounded
coquotient; its coefficient on $T$ has zero residue.  The $W$
coefficient has a factor $v$ and bounded coquotient, multiplied by
$1/\sqrt{H_c}$.  For $E_Y$ the exact $m$ cancellation in
$F_Y+uvF_X$ again applies; (CO6) gives
$B_2=8D_v(xA)\eta_c u^2$, and the remaining lift is
\[
 P(B_2+xe_w)+T(\rho u e_z)+E\,8v(D-\alpha),
\]
divided by $\sqrt H$, with bounded $e_w,e_z$.
Every coefficient has zero residue.  The hard residue of (CG9) is
$\operatorname{diag}(1,4)$, and all mixed hard entries vanish in
residue.  This proves the full-parent corner including literal $n=0$.

\paragraph{All other tests.}
For the remaining commuting axes use the earlier outer identities,
cap $C,A$ at $\sigma/x$ as before, and also cap $m$ at
$\min(m,1)$.  Retain the entire nonnegative excess $m$ positive square term.
The selected $m$ weight is a unit and all previous displayed
strict residues apply.  The noncommuting-residue theorem covers
the other tests.  Restore all positive factors and split complex
states as before.  This proves the theorem for arbitrary positive
$g$, including the new unbounded-$m$ regime.

The remaining assumptions in this source family are now only that
$k,b$ are units, $x$ is positive infinitesimal, all original contact
openings are positive, at least one of $r,g$ is a unit, $n\geq0$,
and the source is physical.  The joint degeneration $r,g\to0$,
where both $n,m$ can be unbounded, is not covered.  The
result does not yet treat unbounded $b$ or a degenerate residue of
the finite actual angle $k$.  Section~\ref{p7294aae40ee2:sec:c2-bounded-state-entry}
excludes infinitesimal $b$ under these entry hypotheses.

\subsection{Joint opening of the two event coefficients}
\label{component:p8f912583feaa}

\label{p8f912583feaa:sec:c2-joint-centered}

\begin{theorem}[Both event openings may close]
Use the exact rational C2 source and original parent of
Section~\ref{component:pec59340d644f}.  Suppose $x>0$ is infinitesimal,
$k,b$ are positive bounded units, $0<r<1$, $0<h<b$, and
$g=b-1+r+xh>0$.  Suppose $r,g$ are infinitesimal, $n\geq0$, and
the source is physical.  Then the entire original parent is tensor
SOS at every independent real or complex projective testing state
and every independent projective pair of testing angles.  No lower
bound on either $h$ or $b-h$ is required.
\end{theorem}

Together with the three-opening and last-opening theorems this covers
the bounded-$k,b$ rational source family without requiring either $r$
or $g$ to be a unit. Nonunit $k$ and unbounded $b$ remain separate
source domains.

\paragraph{Source bounds, including both complementary openings.}
Write $\rho=\sqrt{x}$, $\eta=h/b$, $\kappa=b/k$, $K=k_0>0$,
and $D_K=1+K^2$.  Since $r,g$ are infinitesimal, $b_0=1$.
The source formulas give
\[
 (nr)_0=\frac1{4D_K},\qquad (mg)_0=\frac14,
 \qquad \Delta_0=2D_K.
\]
Thus $n,m$ are positive unbounded pure coefficient quantities,
$nr,mg$ are bounded units, and $nr^2,mg^2$ are infinitesimal.
In particular $m-1>0$.  The existing opening cancellation identities,
with $T_C=xC(b-h)$ and $T_A=xAh$, give bounded positive units
\[
 (T_C)_0=(T_A)_0=\frac{K^2}{4D_K},\qquad
 (D/x)_0=\frac1{4K^2D_K},\quad
 (\alpha/x)_0=\frac{K^2}{4D_K}.
\]
They hold even if $h$ or $b-h$ is infinitesimal.  Specifically,
\[
 xC(1-\eta)=T_C/b,\quad xA\eta=T_A/b,\quad
 xC(1-\eta)^2=T_C(b-h)/b^2,\quad
 xA\eta^2=T_Ah/b^2.                                  \tag{JC1}
\]
Put $H_c=x(C+A)$.  This quantity is bounded below by a positive
real, although it can be unbounded.  Its reciprocal and reciprocal
square root are bounded pure coefficient elements.

\paragraph{Center the excess positive square terms before taking normals.}
Retain exactly the original amplitudes
\[
 M=vW-X+uvY-uZ,\qquad J=uW+uvX-Y-vZ,
 \qquad D_u=1+u^2,\quad D_v=1+v^2.
\]
Let $F$ denote the full original parent, including its signed core
$H$.  Define $G$ by the following exact polynomial identity:
\[
 F=G+4nD_uD_v(X+ruZ)^2
             +4(m-1)(M-guZ)^2.                       \tag{JC2}
\]
An explicit expression, requiring no elimination or mixed inverse, is
\begin{align*}
G={}&4M^2+4CD_v(uW-Y)^2+4AD_uD_vY^2
       +4DJ^2+4\alpha D_u(vY-Z)^2+H\\
 &-8nrD_uD_vuZX-4nr^2D_uD_vu^2Z^2\\
 &+8(m-1)g\,uZM-4(m-1)g^2u^2Z^2.                  \tag{JC3}
\end{align*}
Every coefficient in (JC3) except $C,A$ is bounded.  At the exact
original contact
\[
 (u,v;W,X,Y,Z)=(\rho,k;\kappa\rho,-r\rho,xh/k,1)
\]
both subtracted squares and all of their first derivatives vanish.
Consequently $G$ has exactly the same zero and full stationary
contact as $F$.  This statement is essential: merely capping the
uncentered original positive square terms would not preserve the contact.

On the finite $u=0$ residue axis use
\begin{gather*}
 X_g=-ruZ,\qquad Y_g=\eta uW,\qquad E_X=\sqrt n(X+ruZ),\\
 E_Y=\sqrt{C+A}(Y-\eta uW),\qquad
 E_M=\sqrt{m-1}(M-guZ).                              \tag{JC4}
\end{gather*}
Expand only $G$ in $E_X,E_Y$ and retain $4E_M^2$ as written.
The root $E_M$ is allowed to be redundant.  No substitution for
$W$ is made.  In particular, the positive $C$ positive square term introduces
no coefficient of the form $Cu^2/n$ or $Cu^2/m$.

\paragraph{The complete graph scalar and its global leading inequality.}
Write $f=G|_g=f_{ww}W^2+f_{wz}WZ+f_{zz}Z^2$.  The same exact monic
divisions as in the prior single-angle proof give
\[
 f=aP^2+b_sB^2+c_sV^2+eE^2+g_sBV,
 \qquad (P,B,V,E)=(uW/\rho,uZ,vW,\rho Z),              \tag{JC5}
\]
where
\[
 a=x\frac{f_{ww}-f_{ww}|_{u=0}}{u^2},\quad
 b_s=\frac{f_{zz}-f_{zz}|_{u=0}}{u^2},\quad
 c_s=\frac{f_{ww}|_{u=0}}{v^2},\quad
 e=\frac{f_{zz}|_{u=0}}x,\quad g_s=\frac{f_{wz}}{uv}.
\]
These are polynomial coquotients, not inverses of testing
infinitesimals.  Their coefficients are bounded by (JC1)--(JC3).
The exact constant $f_{zz}|_{u=0}=4Dv^2+4\alpha$ supplies the $x$
factor.  For finite $v$ their residues are
\[
 q(v)=\frac{K^2(1+v^2)}{D_K},\qquad
 b_s=4,\quad c_s=4,\quad g_s=-10,\qquad
 e(v)=\frac{v^2/K^2+K^2}{D_K}.                        \tag{JC6}
\]
Thus the leading scalar, using the exact identity $BV=vPE$, is
\[
                  qP^2-2vPE+eE^2+4(V-B)^2.           \tag{JC7}
\]
The exact identity on the bounded full state graph is
\[
                   q(v)e(v)-v^2
                    =\frac{(v^2-K^2)^2}{D_K^2}.      \tag{JC8}
\]
If $v$ is a finite nonzero residue with $v^2\ne K^2$, then
$10|v|<8|v|+2\sqrt{qe}$.  Splitting the exact $BV=vPE$ cross
term with the pure testing function
\[
             \theta=\frac{4|v|}{4|v|+\sqrt{q(v)e(v)}}
\]
gives strictly positive residue blocks on $(B,V)$ and $(P,E)$.
At $v_0=0$, allocate the entire cross term using $BV=vPE$;
(JC5) then has strictly positive diagonal residue
$(q(0),4,4,e(0))$ and zero cross residues.  There is no exceptional
$n_0=0$ corner in the joint regime.

Here is also a global equality and Hessian check needed at the
matched angles.  For $z=v^2>0$ and $s>0$, set
\[
 \mathcal D(s,z)=(q s+4z)(4s+e)-25zs.
\]
Exactly,
\begin{align*}
 \frac{\mathcal D(s,z)}s
 &=4qs+4ez/s+qe-9z\\
 &=4\bigl(\sqrt{qs}-\sqrt{ez/s}\bigr)^2
       +(\sqrt{qe}-\sqrt z)(\sqrt{qe}+9\sqrt z).    \tag{JC9}
\end{align*}
The second term is nonnegative by (JC8); its equality forces
$z=K^2$.  The first then forces $s=1$.  Moreover the second term is
\[
 \frac{(z-K^2)^2}{D_K^2}
       \frac{\sqrt{qe}+9\sqrt z}{\sqrt{qe}+\sqrt z},
\]
whose quadratic coefficient at contact is strictly positive.
The first square supplies a nonzero $s$ derivative there.
Hence $\mathcal D$ has a positive definite Hessian at $(1,K^2)$.
These radicals are used only in the residue-field inequality; no
radical of an element mixing coefficient and testing fields is used.

\paragraph{Restore every normal, with bounded coefficients.}
From (JC2)--(JC4) the entire parent is exactly
\[
 F=f+E_XL_X+E_YL_Y+
 H_{XX}E_X^2+2H_{XY}E_XE_Y+H_{YY}E_Y^2+4E_M^2.        \tag{JC10}
\]
The first normals are $L_X=G_X|_g/\sqrt n$ and
$L_Y=G_Y|_g/\sqrt{C+A}$.  The only unbounded weights $C,A$ do not
enter the $X$ row.  Thus
\[
 L_X=V\ell_w+B\ell_z
\]
with bounded coefficient functions carrying the infinitesimal
pure coefficient factor $1/\sqrt n$.  Both have identically zero
residue on the entire finite $u=0$ axis, including $v=0$.

For the second normal the exact form is
\[
 L_Y=\frac1{\sqrt{H_c}}\left\{
 P(B_2+xe_w)+B(\rho u e_z)+E e_*\right\},            \tag{JC11}
\]
where
\[
 B_2=8D_v\{(xA)\eta D_u-(xC)(1-\eta)\},\qquad
 e_*=8v(D-\alpha).
\]
All $e_w,e_z$ are bounded: all new terms in (JC3) have bounded
coefficients, and their $Y$ derivative has an additional $u^2Z$
factor.  Equation (JC1) makes $B_2$ bounded with identically zero
residue, even at either closing complementary opening.  Thus every
coefficient in (JC11) has zero residue.

The entire hard block is bounded and has residue
\[
             \operatorname{diag}(4D_{v,0},4D_{v,0},4)
                   \quad\hbox{on }(E_X,E_Y,E_M).      \tag{JC12}
\]
The first entry comes from the retained centered $n$ positive square term;
the bounded $G$ contribution is divided by $n$.  The second is
\[
 4D_v\frac{C+AD_u}{C+A}
   =4D_v\left(1+\frac A{C+A}u^2\right)
\]
plus bounded terms divided by $C+A$.  The mixed entry contains
$1/(\sqrt n\sqrt{C+A})$ times a bounded function.  The third
entry is the unexpanded centered $m$ positive square term.  This proves the
full-parent assertion at all remote finite-axis tests, by the
strict Gram in (JC5) and the zero-residue first normals.

\paragraph{The matched finite-angle residues.}
At $v_0=K$, normalize $Z=1$ and suppose $W$ is infinitesimal.
Use exactly the prior contact conductors
\[
 R=(uW-\kappa x)/\rho,\quad S=u-\kappa^{-1}W,
 \quad U=\rho(v-k),\quad T=u(v-k),
\]
with exact identities
\[
 u^2-x=uS+\kappa^{-1}\rho R,\quad
 W^2-\kappa^2x=-\kappa WS+\kappa\rho R,
 \quad uW-\kappa x=\rho R.                           \tag{JC13}
\]
Here $f=p^2Q_b/x+B_b$, with bounded $B_b$ and
\[
 Q_b=4D_v\{(xC)(1-\eta)^2+(xA)\eta^2D_u\}.
\]
Boundedness follows from (JC1); the remaining terms use only
the bounded coefficients of (JC3).  The full exact stationary
contact of $G$ proved after (JC3) supplies every vanishing required
by the prior monic-contact-division argument.  Successive monic
differences in $u^2-x,W^2-\kappa^2x,uW-\kappa x,v-k$ therefore
give a bounded Gram on $R,S,U,T$.

For precision, its strict residue follows from the same explicit
Schur-complement step in that argument, now with
\[
 \Phi=q(v)s^2w^2+4s^2+4v^2w^2-10vsw+e(v).
\]
Completing the $w$ square gives
\[
 \Phi=(qs^2+4v^2)
       \left(w-\frac{5vs}{qs^2+4v^2}\right)^2
       +\frac{\mathcal D(s^2,v^2)}{qs^2+4v^2}.
\]
Equation (JC9) proves a positive definite three-variable Hessian
at $(s,w,v)=(1,1/K,K)$.  Consequently the residue before the
conductor-zero correction has form
\[
 q(K)R^2+4S^2+A_1RU+B_1ST+E_1U^2,
 \qquad E_1>A_1^2/(4q(K))+B_1^2/16.
\]
Choose a real $\delta$ strictly between $B_1^2/16$ and
$E_1-A_1^2/(4q(K))$, and add $\delta$ times the exact zero
\[
 T^2-U^2-ST(v-k)-\kappa^{-1}RU(v-k)=0.
\]
The two residue blocks $(R,U)$ and $(S,T)$ are strictly positive.

An explicit equivalent residue check can be used in place of this
Schur choice.  The capped current has the exact conductor lift
\[
 (M-gu)|_g=-v\kappa S+\kappa T
                  +\eta uv\rho R+\eta u\kappa\rho U.
\]
Retaining its $T^2$ term, rather than first moving that term to $U^2$,
gives the scalar residue
\[
 K^2R^2+\frac{2(K^2-1)}{D_K}RU+\frac{U^2}{K^2}
                  +4(S-T/K)^2.
\]
Adding $2/D_K^2$ times the same conductor zero leaves the $(R,U)$
Schur margin $2/D_K^2>0$, and leaves
$4(S-T/K)^2+(2/D_K^2)T^2$ on $(S,T)$, also strictly positive.

The full-normal lifts are also the ones actually required for
(JC10), rather than scalar-only lifts.  Write
$L_X=uA_1'+WB_1'=-\kappa B_1'S+u(A_1'+\kappa B_1')$.
Both functions and all their bounded monic coquotients retain
$1/\sqrt n$.  Exact $G_X=0$ at contact supplies vanishing of
$A_1'+\kappa B_1'$.  Using (JC13) gives zero-residue coefficients
on $R,S,U,T$.  For (JC11), write its numerator as $RB_2+\rho E_2$.
Both $B_2,E_2$ have identically zero leading functions and $E_2$
vanishes at contact.  The same monic divisions give factors
$x,\rho u,\rho W$, or a zero-leading-function $(v-k)$ coquotient;
division by $\sqrt{H_c}$ preserves boundedness and zero residue.
Thus the complete seven-root Gram on
$(R,S,U,T,E_X,E_Y,E_M)$ is strict in residue.

If $W/Z$ is not infinitesimal, normalize $W=1$ and keep $Z$
bounded.  Then the scalar has a strict residue on $(P,1)$:
its positive constant contribution is $4K^2$.  All normalized
first normals have zero residue as above.  For $Z=0$, the original
parent is already the literal sum of its event positive square terms and three
positive core squares.  The simultaneous sign copy gives $v_0=-K$.

\paragraph{Every other independent test.}
The outer-chart proof need not use the centered decomposition.
Use the original parent, cap $C,A$ at $\sigma/x$, where
$\sigma=\min(1,xC,xA)$ is a positive bounded unit, and cap $m$ at
$1$.  Retain the original nonnegative excess positive square terms and the
entire original $n$ positive square term.  The four exact outer-core identities
in the prior single-angle theorem apply without change: their
strict residues require only $\sigma_0>0$, the selected $m=1$,
and $(D/x)_0>0$.  They never bound $n$.  This covers the finite
unit $u$ axis, both reciprocal axes, all their intersections, and
literal projective ends.  The inherited physical
noncommuting-residue theorem covers the remaining tests.

These charts and the finite $u=0$ axis exhaust the projective measurement
pairs. Restore the retained squares and homogeneous factors, then apply
the real identities to the real and imaginary state parts.

\section{One small actual angle: unbounded state scale}
\subsection{Source reduction at unbounded state scale}
\label{component:pb60046a37191}

\label{pb60046a37191:sec:c2-unbounded-b-reduction}

Use the original positive rational source with $d=1$, actual
$u_c=\sqrt x$ positive infinitesimal and $v_c=k$ a positive bounded
unit. The global entry theorem supplies its positive chart inequalities.
The preceding rational theorems, including the joint-opening theorem,
now cover every bounded $b$: the inherited $n\ge0$ entry inequality
makes bounded $b$ a unit. We classify the remaining $b$ unbounded
without changing any original coefficient.

Put
\[
 B=xb,\qquad H=xh,\qquad L=x(b-h)=B-H,\qquad
 a^2=k^2/x,\quad t=L/k^2.
\]
The original gaps give
\[
 0<H<1,\qquad0<L<k^2(1+xr),\qquad B=H+L.
\]
Thus $B,H,L$ are bounded and $x/B=1/b$ is infinitesimal.
Set $A_c=1+xr$ and $S=A_c+L$, both positive bounded units. Also
\[
 g/b=1+(r-1)/b+H/b\longrightarrow1,
 \qquad j=1+xr-L/k^2>0,\quad j\le1+x.
\]
The following exact rearrangements of the original short polynomials
control every event scale:
\begin{align*}
 x^2N={}&(1+x)(B+xr)(H+k^2)-x(1+k^2)(1-r)L,\\
 \Delta/b={}&(1+x)B+2x(1+x)r+(1+x)(1+k^2)\\
 &+(x^2-k^2x)r^2/b+2x(1+k^2)r/b+(1+k^2)/b.
\end{align*}
Consequently $xN/b=x^2N/B$ and $\Delta/b$ are positive bounded
units, with respective residues
\[
 H_0+k_0^2,\qquad B_0+1+k_0^2.
\]
Here $H,L\le B$ controls the divided corrections, and $x/B\to0$.

The exact physical condition $n\ge0$ is
\[
 r(\Delta/b)\le(1+r/b)A_cS.
\]
It implies
\[
 r_0\le\frac{1+L_0}{1+k_0^2+B_0},\qquad
 1-r_0\ge\frac{k_0^2+H_0}{1+k_0^2+B_0}>0.
\]
Thus $1-r$ is a positive bounded unit without any lower bound on $r$.
The original formulas now give the exact comparabilities
\[
 m\asymp b^{-2},\qquad C\asymp L^{-1},\qquad A\asymp H^{-1},
 \qquad D\asymp B/j,\qquad\alpha\asymp B/(1-H).
 \tag{SR1}
\]
Every omitted factor in $\asymp$ is a positive bounded unit with
bounded reciprocal. For instance,
\[
 CL=\frac{A_c(xN/b)}{2\Delta/b},\qquad
 AH=\frac{k^2A_c(1+r/b)S}{2\Delta/b}
\]
are positive units; the other identities follow using $g/b\to1$
and $(b-k^2r)/b\to1$.

There are exactly two unbounded-$b$ scale leaves:
\begin{enumerate}
\item If $B=xb$ is infinitesimal, so are $H,L$, while $j\to1$
and $1-H\to1$. Thus $C,A$ are unbounded and $m,D,\alpha$ are
weak, with $m\asymp b^{-2}$ and $D\asymp\alpha\asymp xb$.
This is the one-weak-complementary-event coefficient domain, with weak
$\alpha$ and strong $C$ in this actual-angle orientation.
\item If $B$ is a positive bounded unit, then $C,A,D,\alpha$
are all strong and $m$ is weak. The new strong-complement pitchfork
applies to this entire leaf at all upper event scales. In particular
it includes $C$ unbounded when $L$ is infinitesimal and $\alpha$
unbounded when $1-H$ is infinitesimal.
\end{enumerate}
Unbounded $B$ is impossible by the original chart inequalities.
The homogeneous all-tests theorem of Section~\ref{pab6222b899ef:sec:c2-unbounded-b-small-B}
closes the remaining infinitesimal-$B$ leaf without assuming a bounded
actual state ratio. Thus both unbounded-$b$ leaves, as well as all
bounded-$b$ openings, are now covered. Literal $n=0$ is preserved.

\subsection{Homogeneous completion of the scalar graph}
\label{component:pbb29c73cc124}

\label{pbb29c73cc124:sec:unbounded-b-homogeneous-scalar}

This lemma is a scalar lifting statement.  It neither supplies the full
state normals nor proves positivity of its coefficient Gram by itself.
Its purpose is to avoid any bound on the contact state ratio.

Let $B,\sigma>0$ be pure coefficient infinitesimals, let $k>0$ be a
pure coefficient bounded unit, and put $x=\sigma B$, $\xi=u^2$.
Suppose
\[
 f=a(\xi,v)W^2+2u d(\xi,v)WZ+c(\xi,v)Z^2,
 \quad
 a=\frac\xi B Q+\sigma^2 a_0,\qquad
 c=B e+\xi b,                                      \tag{HL1}
\]
where $Q,a_0,d,e,b$ have bounded polynomial coefficients.  Bounded
pure testing unit denominators may be cleared beforehand.
Assume the value and all first derivatives of $f$ vanish at
\[
 u=\sqrt{x},\quad v=k,\quad
 W=\frac1k\sqrt{B/\sigma},\quad Z=1.                \tag{HL2}
\]
No bound on $B/\sigma$ is assumed.

Define the homogeneous roots
\[
 R=\frac{uW-(B/k)Z}{\sqrt B},\qquad
 S=uZ-\sigma kW,\qquad
 U=\sqrt B(v-k)Z,\qquad T=u(v-k)Z.                 \tag{HL3}
\]
They satisfy the exact conductors
\[
 (\xi-x)W=u\sqrt B R+(B/k)S,
 \qquad
 (\xi-x)Z=uS+\sigma k\sqrt B R.                    \tag{HL4}
\]
Then $f$ has an exact symmetric Gram on $(R,S,U,T)$ whose entries
are bounded polynomials.  In particular no inverse of $\sigma$,
$B$, a testing infinitesimal, or a mixed expression is required in
these Gram entries.  The roots themselves need not be bounded.

\paragraph{First set $v=k$.}
In this paragraph every scalar coefficient is evaluated at $v=k$.
State stationarity and $u$ stationarity give
\begin{gather*}
 a_c=-\sigma k d_c,\qquad c_c=-B d_c/k,\\
 B(a_\xi)_c+kd_c+2kx(d_\xi)_c
                      +\sigma k^2(c_\xi)_c=0,     \tag{HL5}
\end{gather*}
where $c$ on a subscript means evaluation at $\xi=x$.
These are identities in the pure coefficient field; deriving them
does not insert an inverse into a tensor Gram.

Set, by monic polynomial division in $\xi$,
\[
 s=\frac{c+Bd/k}{\xi-x},\qquad
 q=-kd-\sigma k^2s,\qquad
 h_2=\frac{Ba+k\xi d+\sigma k^2c+\sigma Bkd}
                    {(\xi-x)^2}.                 \tag{HL6}
\]
All three quotients have bounded polynomial coefficients.  The first
numerator vanishes at $x$ by (HL5), and the last numerator and its
$\xi$ derivative vanish there by the same three identities.
Moreover $Ba=\xi Q+B\sigma^2a_0$ is bounded.  Division by a monic
polynomial preserves bounded coefficients and is not evaluation of
an inverse of $\xi-x$.

Direct multiplication gives the exact identity
\[
 f(\xi,k)=qR^2+sS^2
                  +h_2\left(uR+\frac{\sqrt B}{k}S\right)^2.
                                                               \tag{HL7}
\]
For later residue computations, if $E_0=e+d/k$, then
\[
 s=b+B\frac{E_0(\xi)-E_0(x)}{\xi-x}
       +x\frac{b(\xi)-b(x)}{\xi-x}.                \tag{HL8}
\]
Thus $s_0=b_0$, $q_0=-k_0d_0$, and the $h_2$ correction has zero
residue in every Gram entry.  Here subscripts $0$ denote residues,
not evaluation at the contact polynomial variable.

\paragraph{The complete angular coquotient.}
Write $\tau=v-k$ and, for each of $Q,a_0,d,e,b$, form its bounded
monic difference in $v$, denoted respectively by
$Q_1,a_{01},d_1,e_1,b_1$:
\[
 Q(\xi,v)-Q(\xi,k)=\tau Q_1(\xi,v),
\]
and similarly for the other four functions.  Put
\[
 H_1=Q_1/k^2+2d_1/k+e_1,\qquad
 L_1=a_{01}/k^2+b_1,\qquad J_1=H_1+\sigma L_1.
\]
Angular stationarity at (HL2) gives exactly $J_1(x,k)=0$.
Consequently the successive monic differences
\[
 J_1=(\xi-x)J_\xi+\tau J_v,
 \quad
 J_\xi=\frac{J_1(\xi,v)-J_1(x,v)}{\xi-x},\qquad
 J_v=\frac{J_1(x,v)}{v-k}                            \tag{HL9}
\]
are bounded polynomials.  Let $D_1=L_1+BJ_\xi$.  The full exact
angular correction is
\begin{align*}
 f(\xi,v)-f(\xi,k)={}&
 \tau Q_1R^2+\frac{\tau a_{01}}{k^2}S^2\\
 &+\left\{2(Q_1/k+d_1)+\sigma kD_1\right\}RU\\
 &+\left\{D_1-2a_{01}/k^2\right\}ST+J_vU^2.       \tag{HL10}
\end{align*}
To check this without any division, use
$uW/\sqrt B=R+\sqrt B Z/k$ and
$\sigma W=(uZ-S)/k$.  The only term not already among the displayed
root products becomes
\[
 \tau Z^2\{BH_1+\xi L_1\}
 =D_1\tau(\xi-x)Z^2+BJ_v\tau^2Z^2
 =D_1(ST+\sigma kRU)+J_vU^2,
\]
by (HL4).  Equations (HL7) and (HL10) prove the asserted homogeneous
Gram lift for every state $W,Z$ simultaneously.  There is no state
normalization and no discarded endpoint $Z=0$.

The usual exact zero, now with the new source scale, is
\[
 T^2-\sigma U^2-ST(v-k)-\sigma kRU(v-k)=0.          \tag{HL11}
\]
This is the appropriate identity for redistributing the angular
diagonal.  Its $U^2$ coefficient is $\sigma$, which is infinitesimal;
replacing it by $1$ would erase the smaller radial scale.  Any
positivity argument must use the actual residue coefficients of
(HL7)--(HL10), rather than importing the bounded-state Schur test
without this scale correction.

\paragraph{Strict scalar corollary for the small-$B$ source.}
Suppose additionally that the coefficientwise leading functions on
the finite $u=0$ axis are
\[
 Q_0(v)=\frac{2K^2(1+v^2)}{1+K^2},\qquad
 d_0(v)=-2v,\qquad
 e_0(v)=\frac{2(v^2/K^2+K^2)}{1+K^2},
 \qquad b_0(K)>0,
\]
where $K=k_0$.  Then the $(R,U)$ residue block produced above is
\[
 2K^2R^2+\frac{4(K^2-1)}{1+K^2}RU+\frac2{K^2}U^2,
\]
with Schur margin $8/(1+K^2)^2>0$.  The $S^2$ coefficient has
residue $b_0(K)>0$, the $RS$ residue is zero, and the only other
possible residue entry is a bounded $ST$ coefficient, say $\ell$.
Add a positive real multiple $\delta$ of (HL11), with
$\delta>\ell^2/(4b_0(K))$.  Its correction to the $U^2$ entry has
zero residue because $\sigma$ is infinitesimal, whereas it supplies
the positive $T^2$ diagonal $\delta$.  Both residue blocks are then
strictly positive.  This proves the scalar strict bounded Gram for
all homogeneous states in the matched-angle chart.  Full normals
and the other testing charts remain separate requirements.

\subsection{The case $b$ unbounded and $xb$ infinitesimal}
\label{component:pab6222b899ef}

\label{pab6222b899ef:sec:c2-unbounded-b-small-B}

\begin{theorem}[The remaining unbounded-state rational C2 source]
Use the exact original rational C2 source.  Suppose $x>0$ is
infinitesimal, $k$ is a positive bounded unit, $b$ is positive
unbounded, $B=xb$ is infinitesimal, $0<r<1$, $0<h<b$, and all
original positive-chart inequalities hold.  Assume $n\geq0$ and
physicality of the source.  Then the full original parent has tensor
SOS for every independent real or complex projective state and
every independent projective pair of measurement angles.
\end{theorem}

The conclusion includes $n=0$, arbitrarily large $n$, arbitrary
complementary openings, and every value of the pure coefficient
ratio $B/\sigma$, where $\sigma=1/b$.  In particular, the actual
contact ratio $W/Z=k^{-1}\sqrt{B/\sigma}$ need not be bounded.
This proof uses the homogeneous scalar lemma in
Section~\ref{pbb29c73cc124:sec:unbounded-b-homogeneous-scalar}; it restores both
full state normals and exhausts the other testing charts.

\paragraph{Exact source scales.}
Put
\[
 \sigma=1/b,\quad x=\sigma B,\quad \eta=h/b,\quad
 y=B\eta,\quad z=B(1-\eta),\quad
 A_c=1+\sigma Br,\quad S_c=A_c+z.
\]
The positive chart gives $0<\eta<1$; both $y,z$ are infinitesimal.
Let $\widehat\Delta=\sigma\Delta$,
$\widehat N=\sigma^2 B N_{\rm src}$,
$\widehat g=\sigma g=1+\sigma(r-1+B\eta)$ and
$j=A_c-z/k^2$.  The cancelled exact formulas are
\begin{align*}
 nr&=\frac{(1+\sigma r)A_cS_c-r\widehat\Delta}
                       {2\widehat\Delta},\\
 m&=\frac{\sigma^2(1+k^2)(1-r)A_cS_c}
                       {2\widehat g\widehat\Delta},\\
 C&=\frac{A_c\widehat N}{2z\widehat\Delta},\qquad
 A=\frac{k^2A_c(1+\sigma r)S_c}{2y\widehat\Delta},\\
 D&=\frac{B(1+\sigma B)(1+\sigma r)(1-\sigma k^2r)S_c}
                       {2k^2j\widehat\Delta},\\
 \alpha&=\frac{B(1+\sigma r)\widehat N}
                       {2(1-y)\widehat\Delta}.
                                                               \tag{UB1}
\end{align*}
Here $\widehat\Delta,\widehat N,\widehat g,j,1-y$ are bounded
positive units.  Write $K=k_0>0$ and $D_K=1+K^2$.  Their relevant
residues imply
\begin{gather*}
 (nr)_0=\tfrac12(D_K^{-1}-r_0),\quad
 (m/\sigma^2)_0=(1-r_0)/2,\\
 (Cz)_0=(Ay)_0=K^2/(2D_K),\quad
 (D/B)_0=1/(2K^2D_K),\quad
 (\alpha/B)_0=K^2/(2D_K).                            \tag{UB2}
\end{gather*}
Since $n\geq0$, $r_0\leq D_K^{-1}<1$, so $1-r$ is a unit.
Set $N=4n+2\geq2$.  Exactly,
\[
 Nr=\frac{2(1+\sigma r)A_cS_c}{\widehat\Delta},      \tag{UB3}
\]
which is a bounded positive unit.  Thus $nr,nr^2$ are bounded even
when $n$ is unbounded, while $1/N$ and $1/\sqrt N$ are bounded.

All opening cancellations needed below are uniform:
\[
 BC(1-\eta)=Cz,\quad BA\eta=Ay,\quad
 BC(1-\eta)^2=Cz(1-\eta),\quad BA\eta^2=Ay\eta.
                                                               \tag{UB4}
\]
The quantity $H_B=B(C+A)$ is bounded below by a positive real;
its reciprocal and reciprocal square root are bounded.  No
assumption that $\eta$ or $1-\eta$ is a unit is made.

\paragraph{The whole finite $u=0$ axis.}
Use the original parent and the exact-contact graph
\[
 X_g=-ruZ,\qquad Y_g=\eta uW,
 \qquad E_X=\sqrt N(X+ruZ),\quad
 E_Y=\sqrt{C+A}(Y-\eta uW).                          \tag{UB5}
\]
The actual contact is
\[
 u=\sqrt{\sigma B},\quad v=k,\quad
 W=k^{-1}\sqrt{B/\sigma},\quad Z=1.
\]
Let $\xi=u^2$ and $f=F|_g$.  Exact polynomial coquotients give
\[
 f=a(\xi,v)W^2+2u d(\xi,v)WZ+c(\xi,v)Z^2,
 \quad a=\frac\xi B Q+\sigma^2a_0,\quad c=B e+\xi b.
                                                               \tag{UB6}
\]
Every coefficient of $Q,a_0,d,e,b$ is bounded by (UB3)--(UB4).
For example $a_0=4(m/\sigma^2)v^2$,
$e=4(D/B)v^2+4\alpha/B$, and
$Q=B(a-a|_{u=0})/u^2$.  These are monic polynomial divisions.
The full leading functions at $u=0$ are
\begin{align*}
 q(v)&=2K^2(1+v^2)/D_K,& d_0(v)&=-2v,\\
 e_0(v)&=2(v^2/K^2+K^2)/D_K,&
 (a_0)_0(v)&=2(1-r_0)v^2,\\
 b_0(v)&=2r_0(1+v^2)/D_K-4r_0+2.                  \tag{UB7}
\end{align*}

For remote tests use
\[
 P=uW/\sqrt B,\quad B_u=uZ,\quad
 V=\sigma vW,\quad E=\sqrt B Z.
\]
Then (UB6) has the exact bounded Gram
\[
 QP^2+b B_u^2+(a_0/v^2)V^2+eE^2+2dPE,             \tag{UB8}
\]
where $a_0/v^2$ is a polynomial coquotient.  The $(P,E)$ residue
block is strictly positive whenever $v_0^2\ne K^2$, since
\[
 q(v)e_0(v)-4v^2=4(v^2-K^2)^2/D_K^2.               \tag{UB9}
\]
This is a global identity for every finite residue $v$, not a
local-contact positivity inference.

\paragraph{Both complete state normals.}
The exact quadratic normal expansion is
\[
 F=f+E_X L_X+E_Y L_Y+H_{XX}E_X^2
                       +2H_{XY}E_XE_Y+H_{YY}E_Y^2.
                                                               \tag{UB10}
\]
Its entire hard block is bounded with residue
$\operatorname{diag}(D_{v,0},4D_{v,0})$.
Indeed the $X$ pivot is $ND_uD_v+O(1)$ before division by $N$,
with the remaining terms either infinitesimal in $u$ or of order
$m$; the $Y$ pivot from the strong positive square terms is exactly
$4D_v(1+[A/(C+A)]u^2)$.  All mixed entries have zero residue and
contain no $C,A$ contribution.

The complete unnormalized $X$ derivative has the form
\[
 F_X|_g=uZ A_X+W(\xi B_1+\sigma^2 B_0),             \tag{UB11}
\]
with bounded polynomial coefficients, where
\begin{align*}
 B_0&=-8(m/\sigma^2)v(1+\eta\xi),\\
 B_1&=v(8D+4)(1-\eta),\\
 A_X&=-8nrD_uD_v+8m(1-r)
       -(8D+4)v^2(1+r\xi)+4D_v(1-r).
\end{align*}
Thus in the remote roots its $P,V$ coefficients have zero residue;
its only possible nonzero residue coupling is to $B_u$.
The exact residue Schur expression for the $(B_u,E_X)$ block is
\[
 (N r^2D_v-4r+2)
       -\frac{(4-2NrD_v)^2}{4ND_v}
                 =2-\frac4{ND_v}.                  \tag{UB12}
\]
It is strictly positive for every finite nonzero residue $v$,
because $N\geq2$.  At $v_0=0$ it is strict if $n$ is unbounded
or has positive residue.  The sole zero case $n_0=v_0=0$ is
treated below.

For the other complete first normal,
\[
 L_Y=\frac1{\sqrt{H_B}}
 \{P(B_2+B e_w)+B_u(\sqrt B\,u e_z)+E e_*\},       \tag{UB13}
\]
where $e_w,e_z$ are bounded,
\[
 B_2=8D_v\{(BA)\eta D_u-(BC)(1-\eta)\},\qquad
 e_*=8v(D-\alpha).
\]
All coefficients have zero residue, by (UB4), $B\to0$, and
$u\to0$.  This verifies both full normals and proves the full
parent at every remote finite-axis test covered by (UB12).

\paragraph{Matched residues without any state-ratio restriction.}
For $v_0=K$, apply the homogeneous scalar lemma to (UB6).  Its
contact and full stationarity hypotheses hold exactly because the
graph (UB5) passes through the original stationary contact.
Use the homogeneous roots
\[
 R=\frac{uW-(B/k)Z}{\sqrt B},\quad
 S=uZ-\sigma kW,\quad
 U=\sqrt B(v-k)Z,\quad T=u(v-k)Z.                  \tag{UB14}
\]
They satisfy
\[
 (\xi-x)W=u\sqrt B R+(B/k)S,\qquad
 (\xi-x)Z=uS+\sigma k\sqrt B R.                    \tag{UB15}
\]
The lemma constructs an exact bounded scalar Gram, valid for
every $W,Z$ simultaneously.  Its $(R,U)$ residue is
\[
 2K^2R^2+\frac{4(K^2-1)}{D_K}RU+\frac2{K^2}U^2,
\]
whose Schur margin is $8/D_K^2>0$.  Its radial $S^2$ residue is
$s_0=2(1-r_0)>0$; its only other possible residue coupling is
$\ell ST$, with a bounded real residue $\ell$.

To lift the full $X$ normal, set
\[
 L=A_X+(B/k)B_1+(\sigma/k)B_0.
\]
Equations (UB11) and (UB15) give exactly
\[
 F_X|_g=u\sqrt B B_1R-(\sigma/k)B_0S+uZL.
\]
At the source contact $L(x,k)=0$, because $F_X=0$ there.
Write its bounded monic differences as
$L=(\xi-x)L_\xi+(v-k)L_v$.  Then
\[
 F_X|_g=
 u\sqrt B(B_1+\sigma kL_\xi)R
 +\{\xi L_\xi-(\sigma/k)B_0\}S+L_vT.             \tag{UB16}
\]
After division by $\sqrt N$, all coefficients are bounded.
The $R,S$ coefficients have zero residue; the $T$ coefficient
may have a finite residue $\gamma$ (it is
$-8K/(D_K\sqrt{N_0})$ when $N$ is bounded).
This nonzero full-normal coupling is retained, not silently
discarded.

For (UB13), put $A_Y=B_2+B e_w$ and
$L_Y'=A_Y/k+\xi e_z+e_*$.  Its numerator is exactly
$RA_Y+\sqrt B ZL_Y'$.  Exact normal contact gives
$L_Y'(x,k)=0$.  Monic differences and (UB15) yield bounded
coefficients on $R,S,U$, all of zero residue: the $R,S$ corrections
carry factors $\sigma B$ or $u\sqrt B$, and the $U$ coefficient
is the angular coquotient of an identically zero leading function.
Division by $\sqrt{H_B}$ preserves these bounds.

Finally add $\delta$ times the exact zero
\[
       T^2-\sigma U^2-ST(v-k)-\sigma kRU(v-k)=0.    \tag{UB17}
\]
Choose a real $\delta>\ell^2/(4s_0)+\gamma^2/(4D_K)$.
This supplies a $T^2$ residue $\delta$ at zero residue cost to
the $U^2$ coefficient, since $\sigma$ is infinitesimal.  The block
on $(S,E_X,T)$ has strictly positive Schur complement; the $(R,U)$
block and $E_Y$ diagonal are already strict.  Thus the complete
six-root Gram on $(R,S,U,T,E_X,E_Y)$ is strictly positive in
residue.  The calculation is homogeneous in all four state
coordinates and needs no bound on $W/Z$, no comparison of $B$ with
$\sigma$, and no state endpoint exclusion.  The exact simultaneous
sign copy treats $v_0=-K$.

\paragraph{The double-zero test, including literal $n=0$.}
Suppose now $n_0=v_0=0$.  Use the pure coefficient quantities
\[
 H_c=4n+4m+2,\quad \zeta=(4m+2)/H_c,\quad
 \xi_c=4m/H_c,\quad \eta_c=C/(C+A),
\]
and graph $X_g=\xi_cvW-\zeta uZ$, $Y_g=\eta_cuW$.
Normalize its $X$ normal by $\sqrt{H_c}$ and its $Y$ normal by
$\sqrt{C+A}$.  Put
\[
 \epsilon=4\alpha+4n(4m+2)/H_c>0.
\]
Then $\epsilon$ is infinitesimal, $n/\epsilon$ and $B/\epsilon$
are bounded, and the harmonic product $BCA/(C+A)$ is a unit.
Use the five scalar roots
\[
 P=uW/\sqrt B,\quad V=\sigma vW,\quad E=\sqrt B Z,
 \quad B_\epsilon=\sqrt\epsilon\,uZ,\quad T_0=uvZ.
\]
The exact inherited identity, with its original coefficients,
is
\[
 f_{zz}=4\alpha+\epsilon u^2+4n\zeta^2u^4
                       +4Dv^2+u^2v^2 B_b,
 \qquad (B_b)_0=2.
\]
The $W^2$ coefficient at $u=0$ is exactly
\[
 \frac{4m(2n+1)}{2n+2m+1}v^2
              +(4n+2)\xi_c^2v^4.
\]
Therefore its coefficient on $V^2$ has positive residue
$4(m/\sigma^2)_0=2(1-r_0)$.  The scalar residue on the five roots
is strictly diagonal, with entries
\[
 \left(\frac{2K^2}{D_K},\ 2(1-r_0),\
        \frac{2K^2}{D_K},\ 1,\ 2\right).
\]
The cross term is allocated using $uvWZ=vPE$, so it has zero residue.

Both full normals are restored, including when $n=0$ exactly.
At $v=0$ the $X$ derivative is $-8n\zeta u^3Z$, whose coefficient
on $B_\epsilon$ is infinitesimal since $n/\epsilon$ is bounded.
Its remaining $Z$ coefficient is $uv^2$ times a bounded function,
so it has zero residue on $T_0$.  Its $W$ coefficient is a sum of
$u^2vW$ times a bounded function and $mvW$ times a bounded
function divisible by $v^2$ at $u=0$.  Allocate the former to
$P$ with coefficient $u\sqrt B\,v$, and the latter to $V$ with
coefficient $m/\sigma$; both have zero residue.  In particular one
must not divide the former by $\sigma$.
The $Y$ derivative is the inherited harmonic-opening identity:
$B_2=8D_v(BA)\eta_cu^2$, with bounded residual coquotients.
Its coefficients on $P,T_0,E$ all have zero residue.  The hard
residue is $\operatorname{diag}(1,4)$.  This proves the full-parent
corner without deleting a state coordinate.

\paragraph{All remaining projective tests.}
For the other commuting charts use the four original outer-core
identities with $\rho$ replaced by $\sqrt B$, and cap $C,A$ at
$s/B$, where $s=\min(1,BC,BA)$ is a positive bounded unit.
Retain both original nonnegative excess positive square terms and the entire
original $n$ positive square term.  No positive lower bound on $m$ is needed:
the first two outer pivots are $4m_0+2=2$, while the other two
charts retain the whole $m$ positive square term and do not use it for strictness.
The reciprocal-$v$ chart uses $(D/B)_0>0$ exactly as before.
Thus the finite unit $u$ axis, both reciprocal axes, all double
corners, and literal projective endpoints have strict bounded
Grams.  The inherited physical noncommuting-residue theorem covers
all remaining testing pairs.

The finite-axis identities are homogeneous in the state and include
unbounded actual state ratios. Every mixed quotient used above is a
monic polynomial quotient; coefficient inverses and radicals remain in
one labelled field. Restore the retained squares and apply the strict
bounded-Gram lemma to each real and imaginary state part.

\section{One small actual angle: strong complementary events}
\subsection{Two unbounded retained events}
\label{component:p0c7a15f800d9}

\label{p0c7a15f800d9:sec:c2-two-large}

\begin{theorem}[The $n,D$-unbounded C2 pitchfork]
\label{p0c7a15f800d9:thm:c2-two-large}
Let the physical canonical C2 source over a real closed coefficient
field have $\lambda=\beta=0$, with $d,C,\alpha$ positive bounded
units, $A\geq0$ bounded, and $n,m,D\geq0$ arbitrary. Suppose there is
an actual finite noncommuting stationary zero with $u_c=k>0$
infinitesimal and $v_c=a>0$ a bounded unit. Then the entire original
parent is tensor SOS at every independent projective measurement
pair and every real or complex projective state.
There is no bound on $n,D$, or on their ratio. In fact physicality and
contact force $A$ to be a positive unit, $D$ to be bounded below by
a positive real, and $m/k^4$ to be a positive bounded unit.
\end{theorem}

This extends the preceding pitchfork by removing the upper bound on
$D$. A bound on $m$ need not be assumed: it follows from physicality.
The original $D$ positive square term supplies a state normal, so an unbounded
weight is not put into a bounded Gram coefficient.

Write the unchanged parent as
\[
\begin{split}
F={}&dH+4nD_uD_vX^2+4mR_2^2+4CD_vr^2
       +4AD_uD_vY^2+4DJ^2+4\alpha D_up^2,\\
H={}&2\{J^2+D_up^2+D_vB^2-D_uD_vZ^2\},\\
R_2={}&vW-X+uvY-uZ,\quad r=uW-Y,\quad
J=uW+uvX-Y-vZ,\\
p={}&vY-Z,\quad B=X+uZ,\quad D_u=1+u^2,\quad D_v=1+v^2.
\end{split}                                                    \tag{TD1}
\]

\paragraph{Entry and bounded contact coordinates.}
At $u=0,v=a$, the state block on $(X,Y,Z)$ of
$dH+4CD_vr^2+4\alpha D_up^2$ has positive definite residue.
Its $W$ column is $O(k)$ and its $W^2$ entry is $O(k^2)$.
Completing this ordinary bounded hard block gives, for positive
real constants $b,B_0$, the coefficient-field inequality
\[
 dH+4CD_vr^2+4\alpha D_up^2
 \geq b(X^2+Y^2+Z^2)-B_0 k^2W^2
\]
at the actual angles. All omitted positive square terms are nonnegative.
Thus an actual zero has $W\ne0$ and, on setting $W=1$,
$X,Y,Z=O(k)$. Since $R_2=a+O(k)$ is a unit, the same inequality
and its retained $m$ positive square term give $m=O(k^2)$.

The state equation $F_Y=0$ gives $DJ=O(k)$: every other term is
$O(k)$, using bounded $A,C,\alpha,d$, $m=O(k^2)$ and the bounds
just obtained. The equation $F_X=0$ then gives
$(d+2n)X=O(k^2)$, since the $J$ contribution is $O(kDJ)$ and the
$m$ contribution is $O(m)$, while $Z=O(k)$. Consequently, with
\[
 N=d+2n,\qquad x=X_c/k^2,\quad z=Z_c/k,\quad j=J_c/k,
\]
we have
\[
 Nx,\ z,\ j,\ Dj\quad\hbox{bounded},\qquad N^{-1}\ \hbox{bounded}.
                                                               \tag{TD2}
\]
The bound on $j$ follows directly from $X,Y,Z=O(k)$, even before a
lower bound for $D$ has been obtained. No inverse of $D$ has yet
been used.

For the scale entry, put $U=k^2,V=a^2$, $L=d+2D$,
$\theta=n/N$ and use
\[
\begin{gathered}
t_2=d+2(\alpha+m+A),\quad t_3=d+2(C+D),\quad
 t_4=d+2(\alpha+m),\\
t_6=d+2(\alpha+C+A),\quad t_8=d+2(C+A+D).
\end{gathered}
\]
The original Laurent-determinant stationary equations, divided only
by $N$, are
\[
\begin{split}
U^2\{4\theta A t_3t_4+4CA(1+2D/N)t_4V\}
 &=mV\{4\alpha t_8+4Dt_6V\},\\
4\alpha A(1+2m/N)t_3+4\theta A t_3t_4U
 &=\{4mDt_6/U+4CDt_2\}V^2.
\end{split}                       \tag{TD3}
\]
Here $t_2,t_4,t_6,t_3/L,t_8/L$ are positive bounded units.
After division by $L$, the braces in the first equation are
$A$ times a positive unit on the left and a positive unit on the
right. To check the lower bound uniformly when both $n,D$ diverge,
observe
\[
 \theta+\frac{1+2D/N}{L}
 =\frac{n+1}{N}+\frac{2n}{NL},                    \tag{TD4}
\]
which is a positive bounded unit; all its neighboring factors in
(TD3) have positive upper and lower real bounds. Thus
$m/U^2$ is $A$ times a positive unit. The second equation, again
divided by $L$, has the form
\[
 A\cdot(\hbox{positive unit})+O(AU)
    =(D/L)\cdot(\hbox{positive unit})+O(AU).       \tag{TD5}
\]
If $A$ were infinitesimal, then $D/L$ and hence $D$ would be
infinitesimal, as would $m$. This is incompatible with physicality:
on the test state $(W,X,Y,Z)=(1,0,u,zu)$ the $n$ positive square term vanishes
exactly, and at the residue source $m=A=D=0$ choose
\[
 z=\frac{v(2\bar\alpha+\bar d)}{2\bar\alpha}.
\]
The coefficient of $u^2$ in the full polynomial is then
$-\bar d v^2(2\bar\alpha+\bar d)/\bar\alpha<0$ for real $v\ne0$.
The bounded residue polynomial cannot be nonnegative for all real
small $u$. Therefore $A$ is a unit. Equations (TD3)--(TD5) now prove
that $D/L$ and $m/k^4$ are positive bounded units. In particular
$D>0$ and $D^{-1}$ is bounded. Set $\gamma=4ma^2/k^4$.

\paragraph{A graph that keeps both large positive square terms bounded.}
For arbitrary independent testing angles use the exact graph
\[
 X_g=xu^2W,\qquad Z_g=zuW,\qquad J_g=juW,
\]
which fixes
\[
 Y_g=u(1-j-vz+xu^2v)W.                            \tag{TD6}
\]
It matches the entire actual contact, including its $Y$ coordinate.
The three exact state normals are
\[
 E_X=\sqrt N(X-xu^2W),\qquad
 E_J=\sqrt D(J-juW),\qquad E_Z=Z-zuW.             \tag{TD7}
\]
The inverse uses only pure coefficient inverses:
\[
 X=xu^2W+E_X/\sqrt N,\quad Z=zuW+E_Z,\quad
 Y=u(1-j)W+uvX-vZ-E_J/\sqrt D.                   \tag{TD8}
\]
The $n$ positive square term on the graph is
$4nx^2u^4D_uD_vW^2$, and the $D$ positive square term is
$4Dj^2u^2W^2$. Both have bounded coefficients by (TD2).
Their linear coefficients are bounded too, since $nx$ and
$\sqrt D\,j=(Dj)/\sqrt D$ are bounded. The exact quadratic
normal block is bounded: its only large positive square terms become
$4(n/N)D_uD_vE_X^2$ and $4E_J^2$.

Set $W=1$, $K=k^2$, and
\[
 R=u^2-(K/a)v,\qquad T=u(v-a).
\]
The entire graph scalar has the exact shape
\[
 F|_g=4mv^2+u^2H_0(u^2,v),\qquad H_0\in B_c[u^2,v].             \tag{TD9}
\]
The three contact scalar jets give exactly the same monic division
as for the preceding pitchfork:
\[
 F|_g=\gamma R^2+u^2\{(u^2-K)^2H_{20}
                   +(u^2-K)(v-a)H_{11}+(v-a)^2H_{02}\},         \tag{TD10}
\]
with bounded polynomial coquotients. After
$u^2-K=R+(K/a)(v-a)$ this is a bounded Gram on $(R,T)$.
Its $R$ residue pivot is $\bar\gamma>0$.

The full first normals, before the indicated scale divisions, have
exact shapes
\[
 (F_X+uvF_Y)|_g=-8mv+u^2L_X(u^2,v),\quad
 (-F_Y)|_g=uL_J(u^2,v),\quad
 (F_Z-vF_Y)|_g=uL_Z(u^2,v).                       \tag{TD11}
\]
The coefficient of $E_X$ is the first expression divided by
$\sqrt N$; that of $E_J$ is the second divided by $\sqrt D$.
These scaled polynomials have bounded coefficients, as shown above.
Contact makes $L_X(K,a)=8ma/K$, and makes the scaled $J,Z$ functions
vanish at $(K,a)$. Monic division by $u^2-K,v-a$ therefore puts
every first normal in $(R,T)$ with bounded coquotients.
Every coefficient coupling $E_X$ to $R,T$ is infinitesimal; the
$R$ coefficients in the $E_J,E_Z$ rows are infinitesimal.
For example, the $X$ row is exactly
\[
\frac1{\sqrt N}\left[
 (8ma/K+u^2L_{X1})R
 +u\{(K/a)L_{X1}+L_{X2}\}T\right],              \tag{TD12}
\]
and the other two rows have form
$uL_{i1}R+\{(K/a)L_{i1}+L_{i2}\}T$ after their scale division.
State degree is two, so there are no omitted higher-normal terms.
Equations (TD7)--(TD12) give an exact bounded full-parent Gram on
$(E_X,E_J,E_Z,R,T)$.

\paragraph{The residue Gram is strict.}
The $X$ pivot has residue $2(1+\bar a^2)$ and the $R$ pivot has
residue $\gamma$. For the remaining hard pair the unscaled state
form is
\[
 B(v)[Y,Z]=\{4(C+A)+2d\}(1+v^2)Y^2
           +4D(Y+vZ)^2+4\alpha(vY-Z)^2.           \tag{TD13}
\]
In the scaled coordinates $Y=-vE_Z-E_J/\sqrt D$, $Z=E_Z$ its
coefficients are bounded and its residue is positive definite.
If $D$ is bounded, this follows from the invertible residue change
and the positive form (TD13). If $D$ is unbounded, its residue is
\[
 4E_J^2+\bigl(\{4(\bar C+\bar A)+2\bar d\}(1+v^2)v^2
                  +4\bar\alpha(1+v^2)^2\bigr)E_Z^2\succ0.
\]
The exact minimized boundary coefficient of $u^2$, before residue,
is
\[
 f(v)=\frac{8(1+v^2)}{\det B(v)/L}
       \left(\frac{Q_{10}}L+\frac{Q_0}L v^2+\frac{Q_{12}}L v^4\right),
                                                               \tag{TD14}
\]
where
\[
 Q_{10}=4\alpha A[d+2(C+D)],\qquad
 Q_{12}=4CD[d+2(\alpha+A)].
\]
The middle coefficient $Q_0$ is defined by the exact Schur numerator;
its explicit value is unnecessary. All quantities in (TD14) are
bounded, $\det B/L$ is a positive unit for bounded real $v$, and
$Q_{10}/L,Q_{12}/L$ are positive units. One may verify boundedness
of $Q_0/L$ directly from the determinant expansion, which is affine
in $D$. This is only a residue calculation of a real hard Schur
complement, not inversion of a mixed tensor element.

The full contact normal and scalar equations imply that the residue
of (TD6) at $v=\bar a$ minimizes the hard form, and that
$\bar f(\bar a)=\bar f'(\bar a)=0$. Thus the real quartic bracket
in (TD14) equals
\[
 \overline{Q_{12}/L}\,(v^2-\bar a^2)^2.
\]
Its second derivative at $\bar a$ is strictly positive. This is
exactly the Schur complement on $T$ of the remaining residue block
$(E_J,E_Z,T)$. Hence the whole five-root residue Gram is positive
definite. The bounded strict-residue Gram lemma proves the local
certificate at $\bar u=0,\bar v=\bar a$. The Bob sign copy covers
$\bar v=-\bar a$.

\paragraph{All other independent tests.}
For $u$ infinitesimal and $v$ bounded with
$\bar v\ne\pm\bar a$, retain the original $m,n$ positive square terms. Set
$K_J=\sqrt D J$ and substitute
\[
 Y=uW+uvX-vZ-K_J/\sqrt D.
\]
The exact remainder is a bounded quadratic Gram on
$(X,Z,K_J,uW)$. Its hard residue is the block just considered with
$X$ pivot $2\bar d(1+\bar v^2)$; its soft Schur complement is
$\bar f(\bar v)>0$. This includes $v=0$ and all infinitesimal $v$.
For $u,q=1/v$ infinitesimal use the homogeneous parent $q^2F(u,1/q)$.
Here $J_h=quW+uX-qY-Z$. Retain $m,n$ and substitute
$Z=quW+uX-qY-K_J/\sqrt D$ with $K_J=\sqrt D J_h$.
The remainder is a bounded Gram on $(X,Y,K_J,uW)$. Its soft residue
is $\lim_{q\to0}q^2\bar f(1/q)>0$, since
$\overline{Q_{12}/L}>0$ and
$\overline{\det B_\infty/L}>0$.
These are polynomial homogeneous expressions, so literal projective
endpoints are included.

At every other commuting test the capping proof from the preceding
all-tests corollary applies verbatim: choose
$\mu=\min(1,C/d,A/d,D/d,\alpha/d)$, a positive unit, retain every
excess and the whole $m,n$ positive square terms, and use its displayed open-axis
state forms and two corner quadratic forms. That proof used lower
bounds, not upper bounds, on these four source weights. Noncommuting
residue tests use the retained arbitrary-physical-source theorem.
Finally homogenize (TD6)--(TD12) in $W$, replacing $R,T$ by $RW,TW$;
restore pure testing projective factors and split complex states into
real and imaginary parts. This proves the theorem for all states,
including $W=0$, without comparing $k,n,D$ with independent testing
infinitesimals.

\paragraph{Exact scope.}
This is an all-tests theorem for the stated actual-contact domain.
It does not bound or remove the residual regimes where $C/d$ or
$\alpha/d$ is unbounded, $A/d$ is unbounded, or the actual contact
has a different angle-residue type. The C2 reflection
$(m,C,D,\alpha;u_c,v_c)\mapsto(D,\alpha,m,C;1/v_c,1/u_c)$ supplies
its corresponding reflected domain. No exhaustive C2 join follows
from these two domains alone.

\subsection{Three arbitrary event weights}
\label{component:p6c62a53b1454}

\label{p6c62a53b1454:sec:c2-three-large}

\begin{theorem}[All scales of the three retained events]
\label{p6c62a53b1454:thm:c2-three-large}
Let $K_c,K_t$ be independent real closed extensions of $\mathbb R$.
Let a physical canonical C2 source over $K_c$ have $\lambda=\beta=0$,
with $d,C,\alpha$ positive bounded units and
$n,m,A,D\geq0$ arbitrary. Suppose an actual finite noncommuting
stationary zero has $u_c=k>0$ infinitesimal and $v_c=a>0$ a bounded
unit. Then the entire original parent is tensor SOS at every
independent projective measurement pair and every real or complex
projective state. No ratios among $n,A,D$ are restricted.
\end{theorem}

The bounded-$A$ result is a special case. The mechanism is to keep
all three possibly large event amplitudes as state coordinates.
The exact parent and its amplitudes are (TD1); in particular
$J=uW+uvX-Y-vZ$.

\paragraph{Source entry with three large weights.}
The initial coercivity estimate in the preceding proof uses only
$d,C,\alpha$, not $A$. Thus after $W_c=1$,
\[
 X_c,Y_c,Z_c=O(k),\qquad m=O(k^2).
\]
Now use $F_Z=0$ first. The $A$ and $n$ positive square terms do not enter this
state equation. Every term other than $-8DvJ$ is $O(k)$, because
$v_c=a$ is a unit and the just displayed bounds hold. Therefore
$DJ_c=O(k)$. The $Y$ equation then gives $AY_c=O(k)$, and the $X$
equation gives $(d+2n)X_c=O(k^2)$. Write
\[
 N=d+2n,\quad T_A=d+2A,\quad L=d+2D,\qquad
 x=X_c/k^2,\quad y=Y_c/k,\quad j=J_c/k.
\]
Consequently
\[
 x,y,j,Nx,Ay,Dj\ \hbox{are bounded},\qquad N^{-1}\ \hbox{bounded}.
                                                               \tag{TA1}
\]
These estimates do not yet invert $A$ or $D$.

Use the exact stationary equations (TD3) and divide both by
$T_A L$. Now $t_2/T_A,t_6/T_A,t_3/L,t_4$ are positive bounded
units. In the first equation the coefficient multiplying $U^2$ is
$(A/T_A)$ times a positive bounded unit: its remaining bracket is
\[
 4t_4\left\{\theta\frac{t_3}{L}
            +CV\frac{1+2D/N}{L}\right\},
\]
whose lower bound follows from (TD4). The coefficient multiplying
$m$ is also a positive bounded unit:
\[
 4V\left\{\alpha\frac{t_8}{T_A L}
              +\frac{D}{L}\frac{t_6}{T_A}V\right\}.             \tag{TA2}
\]
Indeed, if $D/L$ is bounded below, the second summand is a unit.
Otherwise $D$ is infinitesimal and $L$ is a unit, while
$t_8/T_A=1+2(C+D)/T_A$ is a unit. All displayed factors have
bounded upper bounds in either case. Hence
\[
 \frac{m}{U^2}=\frac{A}{T_A}\cdot(\hbox{positive bounded unit}). \tag{TA3}
\]
The second stationary equation becomes
\[
 \frac{A}{T_A}\cdot(\hbox{positive unit})
       +O\left(\frac{A}{T_A}U\right)
 =\frac{D}{L}\cdot(\hbox{positive unit})
       +O\left(\frac{A}{T_A}U\right).                          \tag{TA4}
\]
Thus $A/T_A$ and $D/L$ are comparable. If they were infinitesimal,
$A,D,m$ would all be infinitesimal, contradicting exactly the
physicality test printed after (TD5). Therefore
\[
 A/T_A,\ D/L,\ m/k^4\quad\hbox{are positive bounded units}.       \tag{TA5}
\]
In particular $A,D$ have bounded reciprocals, but may be arbitrarily
large. Set $\gamma=4ma^2/k^4$, also a positive bounded unit.

\paragraph{The three-positive square term graph.}
On the local testing domain $\bar v=\bar a\ne0$, put
$b=1-j-y$ and use
\[
 X_g=xu^2W,\qquad Y_g=yuW,\qquad J_g=juW,\qquad
 Z_g=u\left(\frac b v+xu^2\right)W.                            \tag{TA6}
\]
This graph matches the whole actual contact: the definition of $J_c$
gives $Z_c/k=b/a+xk^2$. Its three hard normals are
\[
 E_X=\sqrt N(X-xu^2W),\quad
 E_Y=\sqrt A(Y-yuW),\quad E_J=\sqrt D(J-juW).                    \tag{TA7}
\]
The exact inverse is
\[
 X=xu^2W+E_X/\sqrt N,\quad Y=yuW+E_Y/\sqrt A,
\quad Z=\frac{uW+uvX-Y-juW-E_J/\sqrt D}{v}.                    \tag{TA8}
\]
Every denominator is either pure coefficient or the pure testing
unit $v$. All graph coefficients and all scaled normal coefficients
are bounded. For the three large positive square terms this follows explicitly
from
\[
 n x^2,\ nx,\ n/N,\quad A y^2,\ \sqrt A\,y,\quad
 D j^2,\ \sqrt D\,j\quad\hbox{bounded},                       \tag{TA9}
\]
using (TA1),(TA5). Their quadratic normal contributions are exactly
$4(n/N)D_uD_v E_X^2$, $4D_uD_vE_Y^2$, and $4E_J^2$.
Every other original source coefficient is bounded, including $m$.

For clarity, the finite Taylor division used below works in the
bounded Laurent polynomial ring in $v$: for positive integer $r$,
\[
 v^{-r}-a^{-r}
 =-(v-a)\frac{\sum_{i=0}^{r-1}v^i a^{r-1-i}}{a^rv^r}.
\]
Thus a Laurent polynomial vanishing at $(U,v)=(K,a)$ belongs to
$(U-K,v-a)$ with bounded Laurent coquotients; repeating the same
identity gives the second-order division. There is no localization
at a difference of the two fields.

Set $K=k^2$, $R=u^2-(K/a)v$ and $T=u(v-a)$. The entire graph
scalar is again $4mv^2+u^2H_0(u^2,v)$, now with bounded Laurent
coefficients in $v$. Its three scalar contact jets give (TD10),
and hence an exact bounded Gram on $(R,T)$ with $R$ residue pivot
$\bar\gamma$. The full first normals before their pure source
scale divisions are
\[
\begin{split}
 (F_X+uF_Z)|_g&=-8mv+u^2L_X(u^2,v),\\
 (F_Y-F_Z/v)|_g&=uL_Y(u^2,v),\\
 (-F_Z/v)|_g&=uL_J(u^2,v).                                     \end{split} \tag{TA10}
\]
Divide the rows by $\sqrt N,\sqrt A,\sqrt D$, respectively.
The scaled coefficients are bounded by (TA9). At contact they
vanish. Laurent Taylor division gives exactly the root membership
and infinitesimal $R$ couplings proved in (TD11)--(TD12).
The $X$ row has infinitesimal coefficients on both $R,T$ because
$m/K=O(K)$; the other rows have infinitesimal coefficients on $R$.
The exact quadratic state remainder is bounded by (TA8)--(TA9),
so the entire original parent has a bounded five-root Gram on
$(E_X,E_Y,E_J,R,T)$.

\paragraph{Uniform strictness without a largest-event choice.}
The $X$ and $R$ residue pivots are $2(1+\bar a^2)$ and
$\bar\gamma$. The remaining hard state form is still (TD13).
At $u=0$, (TA8) gives
\[
 Y=E_Y/\sqrt A,\qquad
 Z=-\frac{E_Y}{v\sqrt A}-\frac{E_J}{v\sqrt D}.
\]
In these coordinates (TD13) is the sum of the two explicit pivots
\[
 4(1+v^2)E_Y^2+4E_J^2
\]
and the nonnegative bounded form
$(4C+2d)(1+v^2)Y^2+4\alpha(vY-Z)^2$.
Its residue is therefore positive definite for every combination of
bounded or unbounded $A,D$.

The boundary Schur function is the exact expression
\[
 f(v)=\frac{8(1+v^2)}{\det B(v)/(T_A L)}
 \left(\frac{Q_{10}}{T_A L}+
       \frac{Q_0}{T_A L}v^2+\frac{Q_{12}}{T_A L}v^4\right),     \tag{TA11}
\]
with $Q_{10},Q_{12}$ as in (TD14). For $v$ a bounded unit all its
coefficients and denominator are bounded, and
\[
 \frac{Q_{10}}{T_A L}=4\alpha\frac A{T_A}\frac{t_3}{L},\qquad
 \frac{Q_{12}}{T_A L}=4C\frac D L\frac{t_2-2m}{T_A}
\]
are positive units. The determinant has the positive term
$16AD(1+v^2)v^2$, so its quotient by $T_A L$ has a positive lower
bound. Its other terms are bounded after that division.
The exact middle coefficient is polynomial of degree at most one
in each of $A,D$, so $Q_0/(T_A L)$ is bounded too.

The contact normal equations and scalar jets give
$\bar f(\bar a)=\bar f'(\bar a)=0$ as before. Therefore the real
quartic bracket equals
$\overline{Q_{12}/(T_A L)}(v^2-\bar a^2)^2$.
Its strict second derivative gives the positive soft Schur pivot on
$T$. The whole five-root residue Gram is strict. This proves the
local all-state certificate; the Bob sign identity gives its negative
sign copy. Homogenization in $W$ includes $W=0$.

\paragraph{Outer tests when both $A$ and $D$ are large.}
For $u$ infinitesimal and $v$ a bounded unit with
$\bar v\ne\pm\bar a$, retain the entire $m,n$ positive square terms and use the
four roots
\[
 X,\quad K_A=\sqrt A Y,\quad K_D=\sqrt D J,\quad uW.
\]
Substitute $Y=K_A/\sqrt A$ and
$Z=(uW+uvX-K_A/\sqrt A-K_D/\sqrt D)/v$.
The exact remainder is a bounded quadratic Gram. Its hard residue
has the positive $X$ pivot $2\bar d(1+\bar v^2)$ and the two
explicit pivots just displayed; its soft Schur complement is
$\bar f(\bar v)>0$ by (TA11).

At $u,v$ both infinitesimal, and at $u,1/v$ both infinitesimal,
one need not divide by $v$. Cap $C,A,D,\alpha$ at the positive
unit $\mu=\min(1,C/d,A/d,D/d,\alpha/d)$ after division by $d$,
and retain all excesses and the full $m,n$ positive square terms. The bounded
selected parent has an exact Gram on $(X,Y,Z,uW)$ in the finite
chart and on the same four roots in the homogeneous reciprocal
chart. At $(u,v)=(0,0)$ its soft Schur complement is
$8Q_{10,\mu}/\det B_\mu(0)>0$; at $(u,1/v)=(0,0)$ it is
$8Q_{12,\mu}/\det B_{\mu,\infty}>0$. The hard blocks are positive
by (TD13) and its homogeneous limit. Hence both entire endpoint
neighborhoods have strict-residue Grams, including literal zeros
and infinities. These expressions use only the bounded selected
parent; no residue of an unbounded $A$ or $D$ is taken.

For every other commuting test, the open-axis and corner forms of the
preceding proof apply with the four lower bounds (TA5). Noncommuting
residue tests use the physical-source theorem. Restore the positive pure
factors and split complex states into real and imaginary parts.

\paragraph{Coverage gained and remaining restriction.}
The theorem removes every relative-scale restriction among $n,A,D$
from the finite pitchfork. Its reflected form covers the corresponding
ordinary--infinite actual contacts. It still assumes $C/d$ and
$\alpha/d$ are bounded units and this actual-angle residue type.
Unbounded $C/d$ or $\alpha/d$, other actual-angle types, and the
exhaustive C2 join remain outside this theorem.

\subsection{Arbitrary complementary event scales}
\label{component:pf7943b34a7dc}

\label{pf7943b34a7dc:sec:strong-complement-pitchfork}

\begin{theorem}[No upper bound on the two complementary events]
\label{pf7943b34a7dc:thm:strong-complement-pitchfork}
Let $K_c,K_t\supset\mathbb R$ be independent real closed fields.
Let the original canonical C2 parent be physical over $K_c$, with
$\lambda=\beta=0$, $d$ a positive bounded unit, $C,\alpha>0$
having bounded reciprocals, and $n,m,A,D\geq0$ arbitrary.
Suppose it has an actual finite noncommuting stationary zero with
$u_c=k>0$ infinitesimal and $v_c=a>0$ a bounded unit.
Then the entire original parent is tensor SOS at every independent
projective measurement pair and every real or complex projective state.
There is no upper bound or relative-scale condition on any event weight.
\end{theorem}

The case of bounded $C,\alpha$ is the existing three-large-event
theorem. We prove both remaining cases, without modifying the parent
or identifying the source and testing fields. We use
\[
\begin{gathered}
 D_u=1+u^2,\quad D_v=1+v^2,\quad r=uW-Y,\quad p=vY-Z,\\
 J=uW+uvX-Y-vZ,\quad R_2=vW-X+up,\quad B=X+uZ,\\
 H=2\{J^2+D_up^2+D_vB^2-D_uD_vZ^2\},\\
 F=dH+4nD_uD_vX^2+4mR_2^2+4CD_vr^2
       +4AD_uD_vY^2+4DJ^2+4\alpha D_up^2.              
\end{gathered}
\tag{SC1}
\]
Every square in this formula retains its original coefficient.

\paragraph{Necessary reductions from the full parent.}
Cap only $C,\alpha$ at positive bounded units in the initial hard
state estimate. At the actual angles the bounded selected form is
bounded below by
$b(X^2+Y^2+Z^2)-B_0k^2W^2$, with positive real $b,B_0$.
Every retained positive square term is nonnegative. Hence the contact has $W_c\ne0$;
on setting $W_c=1$,
\[
 X_c,Y_c,Z_c=O(k),\qquad m=O(k^2).
\]
Each positive event positive square term is $O(k^2)$ as well. Thus, writing
$y=Y_c/k$, $z=Z_c/k$, $x_1=X_c/k$, $j=J_c/k$,
\[
 Ay^2,\ Dj^2,\ C(1-y)^2,\ \alpha(ay-z)^2=O(1),\qquad
 j=1-y-az+aX_c.                                      \tag{SC2}
\]
If $C$ is unbounded then $\bar y=1$, so $A$ is bounded. If $D$
were also unbounded, $j,z$ would be infinitesimal, and the residue of
\[
 H_c/(2k^2)=j^2+(1+k^2)(ay-z)^2
            +(1+a^2)(x_1+Z_c)^2-(1+k^2)(1+a^2)z^2
\]
would be $\bar a^2+(1+\bar a^2)\bar x_1^2>0$.
This contradicts $H_c\leq0$ at the zero of (SC1). Therefore
$C$ unbounded forces both $A,D$ bounded. If $\alpha$ is unbounded
and $A$ were unbounded, (SC2) would make $y,z$ infinitesimal and
$\bar j=1$; the same residue would instead be
$1+(1+\bar a^2)\bar x_1^2>0$. Thus $\alpha$ unbounded forces $A$
bounded. It remains to prove precisely
\[
 \text{(I) } C\text{ unbounded},\ A,D\text{ bounded};\qquad
 \text{(II) } C\text{ bounded},\ \alpha\text{ unbounded},\ A\text{ bounded}.
                                                               \tag{SC3}
\]

\paragraph{Uniform stationary scale entry.}
Set $U=k^2$, $V=a^2$, $N=d+2n$, $\theta=n/N$, and
\[
\begin{gathered}
 t_2=d+2(\alpha+m+A),\quad T=t_3=d+2(C+D),\quad
 t_4=d+2(\alpha+m),\\
 t_6=d+2(\alpha+C+A),\quad t_8=d+2(C+A+D).
\end{gathered}
\]
The exact original Laurent stationary equations are
\[
\begin{split}
 U^2 A t_4\{\theta T+C(1+2D/N)V\}
     &=mV\{\alpha t_8+Dt_6V\},\\
 \alpha A(1+2m/N)T+\theta A Tt_4U
     &=\{mDt_6/U+CDt_2\}V^2.                         
\end{split}\tag{SC4}
\]
Define only pure coefficient quantities
\[
 \delta=CD/T,\quad e=(D/T)(t_6/\alpha),\quad b=t_2/\alpha,
 \quad \mathcal B=\theta+(C/T)(1+2D/N)V,
 \quad \mathcal G=t_8/T+eV.
\]
Because one of $C,D$ is bounded in either case (SC3),
$\delta\leq\min(C,D)/2$ is bounded. The identities
\[
 e=(D/T)\{2+(d+2A)/\alpha\}+2\delta/\alpha,
 \quad \mathcal G=1+2A/T+eV,
 \quad \mathcal B=\theta+(C/T)V+(2\delta/N)V
\]
show boundedness of all these quantities. The quantities
$b,\mathcal G,t_4/\alpha$ are positive bounded units. So is
$\mathcal B$: since $C,N$ have bounded reciprocals and $d$ is
bounded,
\[
 \frac{T}{C(N+2D)}\leq\frac d{CN}+\frac2N+\frac1C=O(1),
\]
and consequently
\[
 \theta+(C/T)(1+2D/N)
 =\frac{n+C(N+2D)/T}{d+2n}
 \geq\frac{n+c_0}{d+2n}
\]
for a positive real $c_0$. The bounded-unit factor $V$ preserves
this lower bound. Dividing (SC4) by $\alpha T$ gives exactly
\[
 \frac m{U^2}=A\rho,\qquad
 A\mathcal H=\delta bV^2,\qquad
 \rho=\frac{(t_4/\alpha)\mathcal B}{V\mathcal G},\quad
 \mathcal H=1+2m/N+\theta(t_4/\alpha)U-U\rho eV^2.      \tag{SC5}
\]
Here $\rho$ is a positive bounded unit and $\mathcal H=1+O(U)$.
These identities were obtained before inverting $A$ or $D$.
They show $A\asymp\delta$. Since $C$ has a bounded reciprocal,
$\delta$ infinitesimal implies $D$ infinitesimal.
If $A$ were infinitesimal, therefore, so would $D,m$ be.
But the exact physical test
\[
 u=v=1,\qquad (W,X,Y,Z)=(1,0,1,1)
\]
has $X=r=p=0$, $R_2=1$, $J=-1$, $H=-2$, and hence
$F=-2d+4m+16A+4D$. Its residue would be negative. We conclude
\[
 A\asymp1,\qquad D^{-1}=O(1),\qquad
 m/k^4\asymp1,\qquad \gamma=4ma^2/k^4\asymp1.        \tag{SC6}
\]
In case (I), $D$ is a bounded unit; in case (II), it may be
arbitrarily large. In both cases $\delta$ is a positive bounded unit.

\paragraph{Full state equations and bounded graph data.}
For later use, the three entire state derivatives of (SC1) are
\[
\begin{split}
 F_X/4={}&D_v[d+2nD_u]X+(d+2D)uvJ+d uD_vZ-2mR_2,\\
 F_Y/4={}&-(d+2D)J+(d+2\alpha)vD_up-2CD_vr
                  +2AD_uD_vY+2muvR_2,\\
 F_Z/4={}&-(d+2D)vJ-(d+2\alpha)D_up+d uD_vB
                  -dD_uD_vZ-2muR_2.                
\end{split}\tag{SC7}
\]
In case (I), the $Z$ equation first yields $\alpha p_c=O(k)$,
because $D$ is bounded and every core coordinate is $O(k)$.
Then the $Y$ equation yields $Cr_c=O(k)$. In case (II), $C,A$
are bounded, so the $Y,Z$ equations are a system with bounded right
side $O(k)$ and unit coefficient matrix
\[
 \begin{pmatrix}-1&a\\-a&-1\end{pmatrix}
\]
on $((d+2D)J_c,(d+2\alpha)(1+k^2)p_c)$.
They yield $DJ_c,\alpha p_c=O(k)$. Finally the $X$ equation in
either case yields $NX_c=O(k^2)$. Thus put $x=X_c/k^2$ and
\[
\begin{array}{ll}
 \text{(I)}&r_0=r_c/k,\ p_0=p_c/k,
       \qquad Nx,Cr_0,\alpha p_0=O(1),\\
 \text{(II)}&j_0=J_c/k,\ p_0=p_c/k,
       \qquad Nx,Dj_0,\alpha p_0=O(1).
\end{array}                                                    \tag{SC8}
\]
All unweighted graph constants are bounded as well. The reciprocals
of $N,C,D,\alpha$ used below are bounded pure coefficient elements.

\paragraph{Two exact graphs retaining the large positive square terms.}
For independent finite testing angles, case (I) uses
\[
 X_g=xu^2W,\quad r_g=r_0uW,\quad p_g=p_0uW,
 \quad Y_g=(1-r_0)uW,\quad Z_g=\{v(1-r_0)-p_0\}uW.    \tag{SC9}
\]
Its exact normals and inverse are
\[
\begin{gathered}
 E_X=\sqrt N(X-xu^2W),\quad E_r=\sqrt C(r-r_0uW),
 \quad E_p=\sqrt\alpha(p-p_0uW),\\
 X=xu^2W+E_X/\sqrt N,\quad
 Y=(1-r_0)uW-E_r/\sqrt C,\quad
 Z=vY-p_0uW-E_p/\sqrt\alpha.                         
\end{gathered}
\tag{SC10}
\]
Case (II) uses
\[
 X_g=xu^2W,\quad J_g=j_0uW,\quad p_g=p_0uW,
 \quad Y_g=\frac{u(1-j_0+vp_0+xu^2v)}{D_v}W,
 \quad Z_g=vY_g-p_0uW.                              \tag{SC11}
\]
Here the exact normals and inverse are
\[
\begin{gathered}
 E_X=\sqrt N(X-xu^2W),\quad E_J=\sqrt D(J-j_0uW),
 \quad E_p=\sqrt\alpha(p-p_0uW),\\
 X=xu^2W+E_X/\sqrt N,\quad J=j_0uW+E_J/\sqrt D,
 \quad p=p_0uW+E_p/\sqrt\alpha,\\
 Y=(uW+uvX+vp-J)/D_v,\qquad Z=vY-p.                 
\end{gathered}
\tag{SC12}
\]
Both graphs match the entire actual contact. All denominators in
these formulas are pure coefficient factors or the pure testing
unit $D_v$. No inverse of a mixed tensor element occurs.
The possibly unbounded positive square terms become respectively
\[
\begin{array}{ll}
 \text{(I)}&4(n/N)D_uD_vE_X^2,quad4D_vE_r^2,quad4D_uE_p^2,\\
 \text{(II)}&4(n/N)D_uD_vE_X^2,quad4E_J^2,quad4D_uE_p^2
\end{array}
\]
in the quadratic normal remainder. Their graph and first-normal
coefficients are bounded by (SC8), for example
$nx,nx^2,Cr_0^2,\sqrt C r_0,\alpha p_0^2,\sqrt\alpha p_0$
in case (I), and $Dj_0^2,\sqrt D j_0$ in case (II).
Every remaining source coefficient is bounded.

\paragraph{Entire scalar and entire first-normal division.}
Set $W=1$, $K=k^2$, and
$R=u^2-(K/a)v$, $T_s=u(v-a)$.
Since $R_2|_g=v+(p_0-x)u^2$ for both graphs, the entire scalar has
the exact shape
\[
 \Phi=F|_g=4mv^2+u^2H_0(u^2,v),                    \tag{SC13}
\]
with bounded polynomial coefficients in case (I), and bounded
coefficients in $B_c[v,D_v^{-1}]$ in case (II).
The three scalar contact jets give
$H_0(K,a)=-\gamma K$, $H_{0,U}(K,a)=\gamma$, and
$H_{0,v}(K,a)=-2\gamma K/a$. In particular the explicit expression
\[
 \Psi=H_0-\gamma u^2+2\gamma(K/a)v
\]
satisfies $\Phi-\gamma R^2=u^2\Psi$, and its value and both first
derivatives vanish at $(u^2,v)=(K,a)$. Finite Taylor division of
$\Psi$ therefore gives
\[
 \Phi=\gamma R^2+u^2\{(u^2-K)^2H_{20}
       +(u^2-K)(v-a)H_{11}+(v-a)^2H_{02}\},           \tag{SC14}
\]
with bounded coquotients. This is an exact bounded Gram on
$(R,T_s)$, with separate residue pivot $\bar\gamma$ on $R$.
For the rational case, Taylor division is legitimate because
\[
 D_v^{-r}-D_a^{-r}
 =-(v-a)(v+a)
   \frac{\sum_{i=0}^{r-1}D_v^iD_a^{r-1-i}}{D_a^rD_v^r}.
\]
Iterating this identity gives the second-order division; only pure
factors have been inverted.

Here are the full first normals, rather than initial forms. In
case (I) they are
\[
 \mathcal N_X=F_X|_g,\quad
 \mathcal N_r=-(F_Y+vF_Z)|_g,\quad
 \mathcal N_p=-F_Z|_g,                              \tag{SC15}
\]
and in case (II) they are
\[
 \mathcal N_X=\{F_X+(uv/D_v)(F_Y+vF_Z)\}|_g,\quad
 \mathcal N_J=-(F_Y+vF_Z)|_g/D_v,\quad
 \mathcal N_p=(vF_Y-F_Z)|_g/D_v.                     \tag{SC16}
\]
The derivatives in (SC7) specify every term. In particular the
following exact positive square term expressions expose the cancellations:
\[
\begin{split}
 -(F_Y+vF_Z)/(4D_v)
   &=(d+2D)J+2Cr-2AD_uY-duvB+dvD_uZ,\\
 (vF_Y-F_Z)/(4D_v)
   &=(d+2\alpha)D_up-2Cvr+2AvD_uY+2muR_2-duB+dD_uZ.
\end{split}\tag{SC17}
\]
The actual coefficients of the normal roots are obtained by dividing
these rows by $\sqrt N,\sqrt C,\sqrt\alpha$ in case (I), and by
$\sqrt N,\sqrt D,\sqrt\alpha$ in case (II).
Using (SC8), their exact forms after those divisions are
\[
 \mathcal N_X/\sqrt N=-8mv/\sqrt N+u^2L_X(u^2,v),
 \qquad \mathcal N_i/\sqrt{w_i}=uL_i(u^2,v),          \tag{SC18}
\]
where all $L_i$ have bounded coefficients in the indicated ring.
The contact equations give
$L_X(K,a)=8ma/(K\sqrt N)$ and $L_i(K,a)=0$.
Writing $L_X=L_X(K,a)+(u^2-K)L_{X1}+(v-a)L_{X2}$, and similarly
for $L_i$, gives the full exact row identities
\[
\begin{split}
 \mathcal N_X/\sqrt N
  &=(8ma/(K\sqrt N)+u^2L_{X1})R
       +u\{(K/a)L_{X1}+L_{X2}\}T_s,\\
 \mathcal N_i/\sqrt{w_i}
  &=uL_{i1}R+\{(K/a)L_{i1}+L_{i2}\}T_s.             
\end{split}\tag{SC19}
\]
Thus the $X$ row has infinitesimal couplings to both soft roots,
and the other rows have infinitesimal couplings to $R$.
The exact state degree is two, so the bounded quadratic remainder
already accounts for every higher-normal term. Equations
(SC10)--(SC19) give an exact full-parent bounded five-root Gram.

\paragraph{A common exact real Schur certificate.}
The $X$ and $R$ residue pivots are $2(1+\bar a^2)$ and
$\bar\gamma$. The unscaled boundary hard form on $Y,Z$ is
\[
 \mathsf B(v)[Y,Z]
 =sD_vY^2+4D(Y+vZ)^2+4\alpha(vY-Z)^2,
 \qquad s=4(C+A)+2d.                               \tag{SC20}
\]
In case (I), $Y=-E_r/\sqrt C$,
$Z=-vE_r/\sqrt C-E_p/\sqrt\alpha$ at $u=0$.
The form has the explicit pivots $4D_vE_r^2+4E_p^2$ plus a
bounded nonnegative form. In case (II),
$Y=(vE_p/\sqrt\alpha-E_J/\sqrt D)/D_v$ and
$Z=(-E_p/\sqrt\alpha-vE_J/\sqrt D)/D_v$.
The explicit pivots are $4E_J^2+4E_p^2$, again with a bounded
nonnegative remainder. The hard residue is strict in both cases.

The determinant and the exact boundary Schur numerator are
\[
\begin{gathered}
 \det\mathsf B(v)=4(1+v^2)\{s(D\,v^2+\alpha)+4D\alpha(1+v^2)\},\\
 Q_{10}=4\alpha A[d+2(C+D)],\qquad
 Q_{12}=4CD[d+2(\alpha+A)],\\
 Q_0=8ACD+8AC\alpha-4ADd-2Ad^2-2Cd^2-4Cd\alpha
       -2Dd^2-4Dd\alpha-d^3-2d^2\alpha.             
\end{gathered}
\tag{SC21}
\]
Put $S=\alpha T$. Every coefficient of $\det\mathsf B/S$ and of
$Q_{10}/S,Q_0/S,Q_{12}/S$ is bounded. This follows directly from
(SC21), bounded $A$, bounded $C/T,D/T,\delta$, and bounded
$1/\alpha$. Moreover
\[
 Q_{10}/S=4A\asymp1,\qquad
 Q_{12}/S=4\delta[d+2(\alpha+A)]/\alpha\asymp1.       \tag{SC22}
\]
The normalized determinant has a positive lower bound for every
bounded unit real $v$. Indeed it equals $C/T$ times the determinant
of the scaled hard form in case (I), and equals
$(D/T)D_v^2$ times that determinant in case (II). The factors
$C/T$ and $D/T$, respectively, are positive units in their cases.

We now take residues of these bounded coefficients and use an
ordinary real variable $v$. The exact real hard Schur function is
\[
 \bar f(v)=
 \frac{8(1+v^2)}{\overline{\det\mathsf B/S}(v)}
 \left\{\overline{Q_{10}/S}+\overline{Q_0/S}v^2
                      +\overline{Q_{12}/S}v^4\right\}.          \tag{SC23}
\]
This formula is an ordinary real Schur calculation, never an
inverse of $\det\mathsf B(v)$ in $K_c\otimes K_t$.
The contact first-normal identities identify the graph at $\bar a$
with the unique minimizer of the scaled real hard form. The scalar
contact jets give $\bar f(\bar a)=\bar f'(\bar a)=0$.
Because $\bar a\ne0$, the quartic bracket in (SC23) is exactly
\[
 \overline{Q_{12}/S}(v^2-\bar a^2)^2.
\]
Its strict second derivative supplies the soft Schur pivot
\[
 \tfrac12\bar f''(\bar a)=
 \frac{32(1+\bar a^2)\bar a^2\overline{Q_{12}/S}}
      {\overline{\det\mathsf B/S}(\bar a)}>0.
\]
The entire five-root Gram therefore has strictly positive residue
when $u$ is infinitesimal or zero and $\bar v=\bar a$.
The strict-residue Gram lemma proves tensor SOS. The exact identity
\[
 F(u,-v;W,-X,Y,-Z)=F(u,v;W,X,Y,Z)
\]
supplies $\bar v=-\bar a$.

\paragraph{All other independent tests, including literal endpoints.}
If $u$ is infinitesimal and $v$ is a bounded unit with
$\bar v\ne\pm\bar a$, retain the full original $m,n$ positive square terms.
In case (I) use the roots
\[
 X,\quad K_C=\sqrt C r,\quad K_\alpha=\sqrt\alpha p,\quad uW
\]
and the exact inverse
$Y=uW-K_C/\sqrt C$, $Z=vY-K_\alpha/\sqrt\alpha$.
In case (II) use
\[
 X,\quad K_D=\sqrt D J,\quad K_\alpha=\sqrt\alpha p,\quad uW
\]
and the inverse in (SC12), with unshifted $J,p$.
The entire remaining parent is a bounded quadratic Gram on these
four roots. Its separate $X$ residue pivot is
$2\bar d(1+\bar v^2)$; its hard pair is the strict form (SC20)
in the scaled coordinates just displayed, and its soft Schur
complement is $\bar f(\bar v)>0$ by (SC23).

For the two endpoint neighborhoods $u,v\sim0$ and
$u,q=1/v\sim0$, divide by $d$, choose
\[
 \mu=\min(1,C/d,A/d,D/d,\alpha/d)>0,
\]
a bounded unit by (SC6), and retain all excesses together with
the full $m,n$ positive square terms. The selected bounded parent has signed
weight one, $C=A=D=\alpha=\mu$, and $m=n=0$.
In the finite chart it is exactly a Gram on $(X,Y,Z,uW)$.
At $(u,v)=(0,0)$ its hard block is positive and its soft Schur
complement is $8Q_{10,\mu}/\det\mathsf B_\mu(0)>0$.
For the reciprocal chart use the polynomial homogeneous parent
$q^2F_\mu(u,1/q)$, on the same four roots. At $(u,q)=(0,0)$
its hard block is
$\operatorname{diag}(2,4(C+A+\alpha)+2,4D)$ in the
$(X,Y,Z)$ variables, evaluated at the capped weights, and its
soft Schur complement is
$8Q_{12,\mu}/\det\mathsf B_{\mu,\infty}>0$.
These are bounded polynomial Grams, so they include the full
infinitesimal neighborhoods and literal $v=0$ or $q=0$.

For completeness, the other open commuting axes use the same cap;
here the other angle has a nonzero finite residue.
At $v=0$ the core is $2\{(uW-Y)^2+(X+uZ)^2\}$ and the
$A,\alpha$ positive square terms force $Y=Z=0$; finite nonzero $u$ then forces
$W=X=0$. At $u=\infty$ the homogeneous core is
$2\{(W+vX)^2+(vY-Z)^2\}$; the $C,A$ positive square terms force $W=Y=0$,
and finite nonzero $v$ forces $X=Z=0$. At $v=\infty$ the core
is $2(1+u^2)(X^2+Y^2)$; the $C,D$ positive square terms force $W=Z=0$
for finite nonzero $u$. Thus these complete state matrices are strict.

Only the two corners with $u=\infty$ remain. Put $p_u=1/u$.
At $(p_u,v)=(0,0)$ the selected state form is
$(2+8\mu)W^2+4\mu Y^2+(2+4\mu)Z^2$.
Its sole zero ray is $X$. In the chart $X=1$, the five small
variables $(p_u,v,W,Y,Z)$ have exact quadratic part
\[
 (2+4\mu)(W+v)^2+2(Z+p_u)^2+4\mu(W^2+Y^2+Z^2)\succ0.
\]
At $(p_u,q)=(0,0)$ the selected state form is
$4\mu W^2+(2+4\mu)X^2+(2+8\mu)Y^2$.
In its sole zero-ray chart $Z=1$, the variables
$(p_u,q,W,X,Y)$ have exact quadratic part
\[
 2X^2+(2+4\mu)(Y-q)^2+4\mu\{W^2+Y^2+(X-p_u)^2\}\succ0.
\]
There are no lower-order terms. Allocating every higher monomial
to two coordinate roots leaves an infinitesimal factor, giving a
strict-residue Gram for the entire selected polynomial at each
corner. Retain the original excess positive square terms throughout.
After pure testing projective normalization, states outside each
displayed zero ray have positive residue value and are covered by
the scalar strict-residue lemma.

Noncommuting-residue tests use the inherited arbitrary-physical-source
theorem. This exhausts all independent projective tests. Finally
homogenize the local graph identities in $W$, replacing $R,T_s$
by $RW,T_sW$; they then include $W=0$. Restore all pure positive
coefficient and testing factors, and split complex states into real
and imaginary parts. No source infinitesimal was compared with an
independent testing infinitesimal anywhere in the argument.

\paragraph{Domain.}
This proves the finite pitchfork when $C/d,\alpha/d$ have positive real
lower bounds, with arbitrary upper event scales, and its C2 reflection.
Other actual-angle types and vanishing complementary lower bounds remain
separate cases.

\subsection{Source bounds from an infinitesimal first angle}
\label{component:pf5726218aec8}

\label{pf5726218aec8:sec:strong-angle-entry}
Let $z=u_c^2$ be infinitesimal, $y=v_c^2>0$ arbitrary, and use the exact
positive C2 source. Set
\[
 \begin{gathered}
 H=zh\in(0,1),\quad L=yt>0,\quad B=zb=H+L,\\
 A_c=1+zr,\quad S=A_c+L,\quad j=A_c-L/y>0.
 \end{gathered}
\]
Here $0<r<1$. Define
\begin{align*}
 \delta={}&z+(1+z)B+(1+z)B^2+2z(1+z)Br+z^3r^2+2z^2r\\
 &+y\{(1+z)B+z+z^2r(2-r)\},\\
 G={}&B+z(H+r-1),\\
 M={}&(1+z)(B+zr)(H+y)-z(1+y)(1-r)L\\
 ={}&(1+z)(B+zr)H+yG+yzj+zrL(1+y)>0.
\end{align*}
The unchanged original coefficient formulas are
\begin{align*}
 n&=\frac{(B+zr)A_cS-r\delta}{2r\delta},&
 m&=\frac{z^2(1+y)(1-r)A_cS}{2G\delta},\\
 C&=\frac{A_cM}{2L\delta},&
 A&=\frac{yA_c(B+zr)S}{2H\delta},\\
 D&=\frac{(1+z)(B+zr)(B-zyr)S}{2yj\delta},&
 \alpha&=\frac{(B+zr)M}{2(1-H)\delta}.
\tag{GE1}
\end{align*}
Assume $C$ and $\alpha$ are strong, meaning each has a positive real
lower bound in the coefficient field. Then $B$ is a bounded positive
unit. Indeed
\[
 \delta\ge(1+y)(B+z),\qquad
 M\le(1+z)(B+z)(1+y)
\]
give $CL\le A_c(1+z)/2$, so $L$ and $B$ are bounded.
If $B$ were infinitesimal then so would $H$, while
$\alpha(1-H)\le(B+zr)(1+z)/2$ would force $\alpha$ infinitesimal,
a contradiction.

Next, $n\ge0$ gives $r\delta\le(B+zr)A_cS$ and therefore
$yr\le A_cS$, since $\delta\ge y(B+zr)$.
Thus $G$, $B-zyr$ and $\delta/(1+y)$ are positive units.
Since $L<yA_c$, the ratio $(H+y)/(1+y)$ is a positive unit: if $y$
is infinitesimal then $H=B-L$ is a unit, and otherwise $y/(1+y)$
is a unit. The positive second expression for $M$ now shows that
$M/(1+y)$ is a positive unit.

Finally $1-r$ is a unit. To see this, let
$E(r)=(B+zr)A_cS-r\delta$, holding $B,H,y,z$ fixed.
Its coefficients divided by $1+y$ are bounded, while direct expansion
at $r=1$ gives
\[
 E(1)=-(H+y)(B+z)(1+z).
\]
The right side divided by $1+y$ is negative with nonzero residue.
If $1-r$ were infinitesimal, $E(r)/(1+y)$ would have the same negative
residue, contradicting $n\ge0$.

Consequently the following exact scaled quantities are positive units:
\[
 \frac m{z^2},\quad CL,\quad\alpha(1-H),\quad
 \frac{AH(1+y)}y,\quad Dyj(1+y).
\tag{GE2}
\]
For the IH source subcase $y$ is unbounded, hence $r=O(1/y)$ is
infinitesimal, $j$ is a unit, and $Dy^2$ is a positive unit.
Put $\rho=B/y$, $s=r/\rho$, and $c_s=(1+L)/B$.
Then $s$ is bounded and, with $N=4n+2$ and $\nu=Nr$,
\[
 \frac\nu\rho\longrightarrow2c_s,\quad
 \frac{nr}\rho\longrightarrow\frac{c_s-s}{2},\quad
 \frac D{\rho^2}\longrightarrow\frac{c_s}{2},\quad
 CL\longrightarrow\frac12,\quad
 \frac{Bm}{z^2}\longrightarrow\frac{c_s}{2},
\]
\[
 AH\longrightarrow\frac{1+L}{2},\qquad
 \alpha(1-H)\longrightarrow\frac B2.
\tag{GE3}
\]
Here arrows mean equality of residues of bounded expressions, with
$B,H,L,s$ replaced on the right by their residues. The expressions
remain valid at $H_0=0$, $H_0=1$, or $L_0=0$ whenever allowed by $B_0>0$.
The reason is that the scaled rational functions have denominators
bounded away from zero after the displayed opening factors have
been absorbed; no unscaled $1/H$, $1/(1-H)$ or $1/L$ is asserted bounded.

\label{component:pfb3a83f3849a}
\label{pfb3a83f3849a:sec:common-ii}
\paragraph{The II source corollary.}
\label{sec:ii-source-bounds-refined}
For the original positive C2 source of GE1, suppose $z,y$ are
infinitesimal, without an initial strength assumption on $C$ or $\alpha$. Then $\alpha$ is strong if and only if $B=H+L$
is a positive unit.  On this face $C$ is unbounded, $H$ and
$1-r$ are positive units, and
\[
 \frac{m}{z^2},\quad\frac{A}{y},\quad CL,\quad Dyj,
 \quad\alpha(1-H)\asymp1,\qquad nr=O(1). \tag{DI7}
\]
Here every displayed comparison is in the coefficient field.
Indeed $0<L<yA_c$ makes $L$ infinitesimal and $B$ bounded.
If $B$ is infinitesimal, then $H$ is infinitesimal and GE1 gives
\[
 0<\alpha\le
 \frac{(B+zr)(1+z)(H+y)}{2(1-H)}\in\operatorname{Inf}_c,
\]
using $\delta\ge B+z$ and
$M\le(1+z)(B+zr)(H+y)$.  If $B$ is a unit, then
$H=B-L$ is a unit.  The positive expressions in GE1 show that
$G,\delta,M$ are units and hence that $CL$ and
$\alpha(1-H)$ are positive units.  Thus $C$ is unbounded
and $\alpha$ is strong, so the hypotheses used to prove GE2
hold.  That proof also gives $1-r$ a positive unit, from the
strictly negative value of $E(1)$ and $E(r)\ge0$.
Now GE2 yields the five bounds in DI7 because $H$ and $1+y$
are units.  Finally the exact GE1 identity
\[
 nr=\frac{(B+zr)A_cS-r\delta}{2\delta}
\]
is bounded.  No lower bound on $r$, $j$, or $1-H$ was imposed,
and the argument includes $n=0$.

\section{One small and one unbounded actual angle}
\subsection{Source bounds for weak $C$}
\label{component:p93b05ea6ad1e}

\label{p93b05ea6ad1e:sec:c2-weak-ih-frontier}

Write $z=u_c^2$, $y=v_c^2$, and use the exact contact quantities
\[
 H=v_cY_c=zh,\quad L=v_c(u_cW_c-Y_c)=yt,\quad
 B=H+L=u_cv_cW_c=zb,\quad R=-u_cX_c=zr.
\]
Thus $0<H<1$, $0<R<z$, $L>0$, and
\[
 L<y(1+R),\quad B>yR,\quad
 G=B+R-z(1-H)>0.
\]
Set $S=1+R+L$ and
\begin{align*}
 \widetilde\Delta={}&(1+z)B^2+2(1+z)BR+(1+z)(1+y)B\\
 &+(z-y)R^2+2z(1+y)R+z(1+y),\\
 \widehat N={}&(1+z)(B+R)(H+y)-(1+y)(z-R)L.
\end{align*}
These are exactly $z\Delta,z^2N_{\rm src}$. In particular,
\[
 \widehat N=(1+z)(B+R)H+yG+z\{y(1+R)-L\}+RL(1+y)>0.
\]
The original coefficients, without an approximation, are
\begin{align*}
 C&=\frac{(1+R)\widehat N}{2L\widetilde\Delta},&
 \alpha&=\frac{(B+R)\widehat N}{2(1-H)\widetilde\Delta},\\
 A&=\frac{y(1+R)(B+R)S}{2H\widetilde\Delta},&
 D&=\frac{(1+z)(B+R)(B-yR)S}
 {2\{y(1+R)-L\}\widetilde\Delta},\\
 m&=\frac{z(1+y)(z-R)(1+R)S}{2G\widetilde\Delta},&
 n+\frac12&=\frac{z(B+R)(1+R)S}{2R\widetilde\Delta}.
                                                        \tag{WI1}
\end{align*}

\paragraph{Uniform source reduction for IH.}
Suppose $z$ is positive infinitesimal, $y$ is positive unbounded,
and $n\ge0$. Then $C$ is weak if and only if $B$ is unbounded.
To prove the nontrivial direction, suppose $B$ is bounded. Since
$R<z$ and $G>0$, one has $z\le4(B+R)$: use $B\ge1/2$
or, otherwise, $B+R>z(1-H)>z/2$. Direct division in (WI1)
gives bounded $\widetilde\Delta/[y(B+R)]$. Moreover,
\[
 \frac{\widehat N}{y(B+R)}
 =(1+z)(1+H/y)
 -(1+1/y)\frac{(z-R)L}{B+R}
\]
has residue $1$, since $L\le B$ and $z$ is infinitesimal.
Consequently $C\ge c/L$ for a positive real $c$, which is strong
because $L$ is bounded.

Now let $B$ be unbounded. Put
\[
 \rho=B/y,\qquad e=1+R-L/y>0.
\]
Here $0<\rho<1+R+H/y$, so $\rho$ is bounded and its residue
lies in $[0,1]$. The following exact normalized quantities have
residue $1$:
\[
 \frac{\widetilde\Delta}{B(B+y)},\qquad
 \frac{\widehat N}{yB},\qquad
 \frac{L}{B},\qquad \frac{S}{B},\qquad\frac{G}{B}.
                                                        \tag{WI2}
\]
They give
\[
 n+\tfrac12=\frac{\rho}{2r(1+\rho)}(1+o(1)).
\]
The inequality $n\ge0$ therefore gives $r/\rho$ bounded and
$r\le\rho/[1+\rho](1+o(1))\le1/2+o(1)$.
In particular $1-r$ is a positive unit, and
$R/\rho=z(r/\rho)$ is infinitesimal. The normalized leading
relations, uniform at arbitrary openings $H,1-H,e$, are
\begin{align*}
 C&=\frac{1+o(1)}{2B(1+\rho)},&
 A&=\frac{B(1+o(1))}{2H(1+\rho)},\\
 \alpha&=\frac{B(1+o(1))}{2(1-H)(1+\rho)},&
 m&=\frac{z^2(1-r)(1+o(1))}{2B(1+\rho)},\\
 D&=\frac{\rho^2(1+o(1))}{2e(1+\rho)}.
                                                        \tag{WI3}
\end{align*}
Here each $1+o(1)$ means a positive coefficient unit of residue one;
it never denotes a comparison with a testing infinitesimal. In
particular $A,\alpha$ are unbounded, $m,C$ are infinitesimal, and
$Cm$ is infinitesimal for every weak-$C$ IH source. Thus the four
reciprocal-product exit can never admit this entire directed leaf.
If $\rho$ is infinitesimal then $e$ is a unit and $D$ is weak;
if $\rho$ is a unit then $D$ is strong. Literal $n=0$ remains
allowed throughout the reduction. No lower bound on $r$ was used.

\subsection{Interior double corner}
\label{component:pc00b84929b12}

\label{pc00b84929b12:sec:c2-weak-ih-interior-corner}

This proof covers the stated local testing chart; the remote commuting
charts are treated below.

\begin{theorem}[Full-parent IH corner, interior shape]
Use the exact positive C2 source with $d=1$, the omitted original
$\lambda,\beta$ coefficients zero, and $n\ge0$. Suppose its actual
angles have type IH, and $C$ is weak. Use the source notation of
\S\ref{p93b05ea6ad1e:sec:c2-weak-ih-frontier}; thus $z$ is infinitesimal and
$B$ unbounded. Write $\rho=B/y$ (a shape parameter, not an event
coefficient). Assume
\[
 \rho,\quad1-\rho
 \quad\hbox{are positive bounded units}.                 \tag{IC1}
\]
At every independent testing pair for which $u$ and the homogeneous
reciprocal coordinate $w=1/v$ are infinitesimal or zero, the entire
original parent has a separated-field tensor-SOS certificate, for
all real or complex state coordinates. The source openings $H,1-H$ and the positive radial parameter $r$
are arbitrary; $n$ may be unbounded. Literal
testing $u=0$ and $v=\infty$, either testing sign, and literal source
$n=0$ are included.
No state coordinate or certificate root is required to be bounded.
\end{theorem}

\begin{proof}
Put $N=4n+2$ and $\nu=Nr$. Source reduction (WI1)--(WI3)
shows that $\nu$ and $1-r$ are bounded positive units, even when
$N$ is unbounded. In residue,
\[
 \nu_0=\frac{2\rho_0}{1+\rho_0},\qquad 2r_0\le\nu_0,
 \qquad D_0=\frac{\rho_0^2}{2(1-\rho_0^2)}.             \tag{IC2}
\]
All assumptions and divisions in this sentence concern the coefficient
field alone.
The uniform coefficient estimates needed below are
\[
 BC\asymp1,\quad \frac{Bm}{z^2}\asymp1,\quad
 \frac{AH}B\asymp1,\quad\frac{\alpha(1-H)}B\asymp1,
 \quad D\asymp1,\quad nr\le\nu/4.
\]
In particular $C,m$ are infinitesimal, $nr,nr^2$ are bounded,
$mB/z\asymp z$, and $m\sqrt B/z\asymp z/\sqrt B$.
These estimates are uniform in $r,H,1-H$ and do not assert that
$A/B$ or $\alpha/B$ is bounded at an opening.

The homogeneous testing parent is exactly
\begin{align*}
 \widehat F={}&4nD_uD_wX^2+4mM^2+4CD_w(uW-Y)^2
     +4AD_uD_wY^2+4DJ^2+4\alpha D_u(Y-wZ)^2\\
 &+2\{J^2+D_u(Y-wZ)^2+D_w(X+uZ)^2-D_uD_wZ^2\},\\
 M={}&W-wX+uY-uwZ,\qquad
 J=uwW+uX-wY-Z,
 \qquad D_u=1+u^2,\quad D_w=1+w^2.
                                                               \tag{IC3}
\end{align*}
For $w\ne0$ it is $w^2F(u,1/w)$; (IC3) defines the literal
projective endpoint as well. Let $G$ be (IC3) restricted to the graph
\[
                         X=-ruZ,\qquad Y=HwZ.
\tag{IC4}
\]
Use the following five scalar roots, all defined using only pure
coefficient radicals and inverses:
\begin{align*}
 R_0&=(uW-BwZ)/\sqrt B,&
 R_1&=(zW-BuwZ)/\sqrt B,\\
 R_2&=uwW-\rho Z,&
 T&=(u^2-z)Z,&
 E&=\sqrt{BH}\,(w^2-1/y)Z.                               \tag{IC5}
\end{align*}
They vanish at every signed exact source contact. Define the pure
coefficient quantities
\begin{align*}
 K={}&4AH^2+8DHr+4\alpha(1-H)^2
      +4m(H+r-1)^2+2(H+r-1)^2+4nr^2,\\
 k={}&K-4B^2m/z^2+8znr^2+8AH^2/y.                      \tag{IC6}
\end{align*}
The exact source formulas and state/angular stationarity give the
following finite polynomial Gram identity:
\begin{align*}
G={}&\left(4BC-\frac{kz}{B}\right)R_0^2
     +\frac{4Bm}{z^2}R_1^2+2(2C+2D+1)R_2^2\\
 &+r^2\{4nw^2+2(2D+2n+1)\}T^2
  +\frac H B\{4Au^2+2(2A+2C+2D+1)\}E^2\\
 &+\frac{ku}{B}R_0R_1
  +\left\{-4r(2D+1)+\frac{k}{B}\right\}R_2T
  -\frac{4\sqrt H(2C+2D+1)}{\sqrt B}R_2E.                   \tag{IC7}
\end{align*}
For clarity, the only non-obvious algebraic step is the exact toric
conductor
\[
 B(u^2-z)(w^2-1/y)Z^2
 =(u^2-z)ZR_2-zR_0^2+uR_0R_1.                         \tag{IC8}
\]
Before (IC8), the last mixed radial term is
$k(u^2-z)(w^2-1/y)Z^2$. Expanding (IC7), its difference from
$G$ has only the four monomials $uwWZ,u^2Z^2,w^2Z^2,Z^2$;
the unchanged original source sets all four coefficients exactly to
zero. 

All coefficients of (IC7) are bounded: $BC$, $Bm/z^2$, $AH/B$,
$\alpha(1-H)/B$, $nr,D$ and $k/B$ are bounded. No inverse of $r$
is used, and $1/\sqrt N$ is bounded because $N\ge2$. Moreover,
\[
 (k/B)_0=\frac{2r_0}{1+\rho_0}.
\]
Add to the right side of (IC7) the following exact zero, obtained
by one more toric conductor:
\[
 \rho T^2+(u^2-z)TR_2
       -u(zR_0-uR_1)(uR_0-R_1)=0.                    \tag{IC8a}
\]
The new $T^2$ term has positive unit residue $\rho_0$. Every other
coefficient introduced by (IC8a) is infinitesimal. This step is what
makes the proof uniform when $r$ is infinitesimal and $n$ unbounded.

Suppress residue bars temporarily. The $R_0,R_1,E$ diagonal
residues are respectively
\[
 \frac2{1+\rho},\qquad
 \frac{2(1-r)}{1+\rho},\qquad
 \frac2{1+\rho},
\]
which are positive. The only surviving off-diagonal block is
\[
 (R_2,T):\qquad
 \begin{pmatrix}
 \dfrac2{1-\rho^2}&-\dfrac r{1-\rho}\\[3pt]
 -\dfrac r{1-\rho}&\rho+r\nu+4Dr^2
 \end{pmatrix}.                                      \tag{IC9}
\]
Using (IC2), its determinant is exactly
\[
 \frac{2\rho}{1-\rho^2}
  +\frac{r\{4\rho-r(3\rho+1)\}}
       {(1-\rho)(1+\rho)^2}>0.
\]
Indeed $0\le r\le\rho/(1+\rho)$, and consequently
$4\rho-r(3\rho+1)\ge\rho(3+\rho)/(1+\rho)>0$.
The first determinant summand is already strictly positive, including
$r_0=0$. Thus (IC7) after the exact zero (IC8a) has strictly positive
definite residue, including literal $n=0$ and unbounded $n$.

It remains to restore the full state normals; the graph certificate
alone is not the claimed conclusion. Put
\[
 E_X=\sqrt N(X+ruZ),\qquad E_Y=\sqrt{B/[H(1-H)]}(Y-HwZ).
\]
We give exact factorizations of both linear normal terms. Let
$U=u^2,V=w^2,s=1/y$ and define $L_X,L_Y$ by
$4L_X=\partial_X\widehat F|_{\rm graph}$ and
$4L_Y=\partial_Y\widehat F|_{\rm graph}$. Direct expansion gives
\begin{align*}
 L_X={}&(2D+1)uR_2-\frac{2m\sqrt B}{z}wR_1+uZP_X(U,V),\\
 P_X={}&(2D+1)\rho-\frac{2mB}{z}V
    -2nr(1+U)(1+V)-r(2D+1)U-2D-r\\
 &+V\{1-H-r-2DH-2m(H+r-1)\}.                         \tag{IC10}
\end{align*}
Every coefficient of $P_X$ is bounded, and source stationarity gives
$P_X(z,s)=0$. Write its exact monic division as
\[
 P_X=(U-z)f_X(V)+(V-s)g_X.
\]
Both $f_X,g_X$ have bounded coefficients. Explicitly,
\begin{align*}
 f_X(V)&=-2nr(1+V)-r(2D+1),\\
 g_X&=-2mB/z-2nr(1+z)+1-H-r-2DH-2m(H+r-1).
\end{align*}
Thus no inverse of $r$ or of either opening is introduced. We obtain
\[
 L_X=\left(2D+1+\frac{g_X}B\right)uR_2
       -\frac{2m\sqrt B}{z}wR_1
       +u f_X(V)T
       -\frac{uwg_X}{\sqrt B}R_0.                       \tag{IC11}
\]
Here we used the exact identity
$B(V-s)Z=R_2-\sqrt B\,wR_0$; this avoids dividing by $\sqrt H$
when $H$ is infinitesimal. Every displayed multiplier is infinitesimal
in the bounded tensor ring; in particular $m\sqrt B/z\asymp z/\sqrt B$.

For the other normal, put
\begin{align*}
 f(V)={}&2m-2C(1+V)-(2D+1)V,\\
 P_Y(U,V)={}&Bf(V)+2AH(1+U)(1+V)+2CH(1+V)
       +2DHV+2DrU+2D\\
 &-2\alpha(1-H)(1+U)+2m(H+r-1)U
       +(H+r-1)U+H(V+1).
\end{align*}
Then
\[
 L_Y=\sqrt B f(V)R_0+wZP_Y(U,V),\qquad P_Y(z,s)=0.
\]
The coefficients of $P_Y/B$ are bounded. Write
$P_Y=(U-z)f_Y(V)+(V-s)g_Y$. The two quotient coefficients are explicitly
\begin{align*}
 f_Y(V)&=2AH(1+V)+2Dr-2\alpha(1-H)
                   +(2m+1)(H+r-1),\\
 g_Y&=-B(2C+2D+1)+2AH(1+z)+2CH+2DH+H.
\end{align*}
The displayed source estimates show directly that $f_Y/B$ and
$g_Y/B$ have bounded coefficients, including both opening limits.
The elementary exact identity
\[
          (U-z)\sqrt B\,wZ=zR_0-uR_1
\]
therefore gives
\[
 \frac{L_Y}{\sqrt{B/[H(1-H)]}}
 =\sqrt{H(1-H)}\left\{f(V)+\frac{z f_Y(V)}B\right\}R_0
       -\frac{\sqrt{H(1-H)}u f_Y(V)}B R_1
       +\frac{\sqrt{1-H}w g_Y}B E.          \tag{IC12}
\]
Again every displayed multiplier is infinitesimal. This argument
retains the full $X,Y$ linear terms; none is estimated on the graph
and then discarded.

Finally the normal quadratic block, in $(E_X,E_Y)$, is bounded.
Its residue is diagonal with entries
\[
             1,\qquad \frac2{1+\rho_0}>0.
\]
Indeed the $X^2$ coefficient before scaling is
$4nD_uD_w+4mw^2+(4D+2)u^2+2D_w$; the $Y^2$ coefficient is
$4CD_w+4AD_uD_w+(4D+2)w^2+(4\alpha+2)D_u+4mu^2$.
Their mixed coefficient is $-4uw(2m+2D+1)$.
For the $Y$ normalization, if $L=B/[H(1-H)]$, then $A/L$ and
$\alpha/L$ are bounded and
\[
 \frac{4(A+\alpha)}L
 =4\left\{\frac{AH}B(1-H)
               +\frac{\alpha(1-H)}B H\right\}
 \longrightarrow\frac2{1+\rho_0}.
\]
Thus the stated normal residue remains strict even when $H_0=0$
or $H_0=1$; no denominator involving a vanishing opening was dropped.
Equations (IC11)--(IC12) show that all cross terms from these normals
to the five scalar roots have zero residue. Thus the exact seven-root
Gram for the entire (IC3) is bounded and has positive definite
residue. The bounded positive-residue Gram lemma yields tensor SOS.
Its roots may be unbounded; it uses no mixed inverse or radical.
The homogeneous chart and real/imaginary splitting give all endpoints
and states stated in the theorem.
\end{proof}

This calculation proves the interior $0<\rho_0<1$ testing corner. Section~\ref{component:p18dbedfff317} treats the source endpoints; Sections~\ref{component:p0fb32cd2f1bc} and~\ref{component:pf95b1c5fc32f} prove the other testing charts.

\subsection{The two source endpoints}
\label{component:p18dbedfff317}

\label{p18dbedfff317:sec:weak-ih-endpoints}
We use the exact parent, graph, roots, coefficients $K,k$, and both
normal factorizations (IC3)--(IC12) in the preserved interior-corner
proof.  Those identities are exact for every positive canonical C2
source; their validity does not require the interior-shape assumption.
Here $z$ and testing $u,w$ are infinitesimal (zero testing coordinates
are allowed), $B$ is unbounded, $H\in(0,1)$, $\rho=B/y$, and
$N=4n+2$, $\nu=Nr$. All source and testing rates are independent.

\paragraph{The endpoint $\rho_0=0$.}
Put $s=r/\rho$ and $\tau=1/B$. Nonnegativity of $n$ and the exact
source formulas imply that $s$ is bounded. The following normalized
source quantities have the indicated residues:
\[
 \nu/\rho\to2,\quad nr/\rho\to(1-s)/2,\quad
 D/\rho^2\to1/2,\quad BC,Bm/z^2,AH/B,\alpha(1-H)/B\to1/2,
\]
\[
 k/(B\rho)\to2s,\qquad 0\le s_0\le1.
\tag{WIE1}
\]
These are uniform opening identities: after substituting
$B=1/\tau$, $y=1/(\tau\rho)$, $r=s\rho$, their normalized rational
functions have a nonzero denominator at $z=\tau=\rho=0$, including
$H=0,1$. 
Thus they concern simultaneous, not iterated, infinitesimal limits.

Replace $T$ by $T_\rho=\rho T$ in (IC7). Its coefficients are bounded;
for example $r^2n/\rho^2=s(nr/\rho)$ is bounded. Add the exact zero
\[
 T_\rho^2+(u^2-z)T_\rho R_2
 -\rho u(zR_0-uR_1)(uR_0-R_1)=0.
\tag{WIE2}
\]
This is (IC8a) multiplied by $\rho$, and uses no mixed inverse.
The residue Gram has $R_0,R_1,E$ diagonal entries $2,2,2$, and
\[
 (R_2,T_\rho):\quad
 \begin{pmatrix}2&-s_0\\-s_0&1+2s_0\end{pmatrix}.
\]
Its determinant $2+4s_0-s_0^2\ge2$ is strictly positive on $[0,1]$.

The entire normal restoration survives this rescaling. In (IC11),
replace $f_XT$ by $(f_X/\rho)T_\rho$, where
\[
 f_X/\rho=-2(nr/\rho)(1+w^2)-s(2D+1)
\]
is bounded. Every multiplier still contains testing $u$ or $w$, and
$1/\sqrt N$ is bounded. The remaining coefficients in (IC11) are
unchanged and infinitesimal. Formula (IC12) uses the root $R_1$ rather
than $T$ and remains bounded and infinitesimal unchanged. The normal
quadratic residue is $\operatorname{diag}(1,2)$. Hence the full
seven-root Gram is strict. This includes arbitrary positive $r,H,1-H$
and literal $n=0$.

\paragraph{The endpoint $\rho_0=1$.}
Set $e=1+zr-\rho+H/y=j>0$, an infinitesimal, and write
\[
 \rho=\frac{1+zr-e}{1-H\tau},\qquad \tau=1/B.
\]
The exact source formulas give the simultaneous normalized residues
\[
 De\to1/4,\quad\nu\to1,\quad nr\to(1-2r)/4,\quad
 BC\to1/4,\quad Bm/z^2\to(1-r)/4,
\]
\[
 AH/B,\alpha(1-H)/B\to1/4,\quad
 0\le r_0\le1/2.
\tag{WIE3}
\]
Put $\kappa=k-8DHr$; then $\kappa/B\to r$.
The scaled rational functions have regular nonzero denominators at
$z=\tau=e=0$ after the displayed powers of $e$ are absorbed.

The large coefficient $D$ must be retained in a square, not treated
as a bounded Gram coefficient. Introduce the full-state root
\[
 F_D=\sqrt D\,(J+eZ),\qquad a=\sqrt{H/B},\qquad
 Q=R_2-rT-aE.
\]
On the graph, $Q=J+eZ$. Subtracting its square from (IC7) gives the
exact identity
\begin{align*}
 G-4DQ^2={}&(4BC-\kappa z/B)R_0^2+4BmR_1^2/z^2
       +2(2C+1)R_2^2\\
 &+r^2(4nw^2+4n+2)T^2
 +\frac HB\{4A(1+u^2)+4C+2\}E^2\\
 &+\kappa uR_0R_1/B+(-4r+\kappa/B)R_2T
       -4a(2C+1)R_2E.
\tag{WIE4}
\end{align*}
Add the exact zero (IC8a). Restore the full normals
\[
 E_X=\sqrt N(X+ruZ),\qquad
 E_Y=\sqrt{B/[H(1-H)]}(Y-HwZ).
\]
Use the exact relation
\[
 R_2=F_D/\sqrt D+rT+aE-uE_X/\sqrt N
                         +wE_Y/\sqrt{B/[H(1-H)]}.
\tag{WIE5}
\]
The inverse $1/\sqrt D$ is purely in the source field and infinitesimal.
All coefficients after this substitution are bounded.

For completeness, the first-normal quotients of
$\widehat F-4F_D^2$ are precisely the D-removed versions of
(IC11)--(IC12). Explicitly,
\begin{align*}
 \bar f_X&=-2nr(1+V)-r,\\
 \bar g_X&=-2mB/z-2nr(1+z)+1-H-r-2m(H+r-1),\\
 \bar f&=2m-2C(1+V)-V,\\
 \bar f_Y&=2AH(1+V)-2\alpha(1-H)+(2m+1)(H+r-1),\\
 \bar g_Y&=-B(2C+1)+2AH(1+z)+2CH+H.
\end{align*}
The constant term in the $Y$ monic polynomial acquires $2De$, but
its monic quotient coefficients are unchanged; the $X$ polynomial
has the constant correction $-2De$. Exact source stationarity
still sets the constant remainders to zero because $J+eZ$ vanishes
at the contact. Thus the displayed quotient formulas reconstruct
both linear normals. Their normalized coefficients are infinitesimal
by (WIE3); in particular $De$ is bounded and no bare $D$ survives.
The normal quadratic block of the D-removed expression has residue
$\operatorname{diag}(1,1)$.

In the order $(R_0,R_1,T,E,F_D,E_X,E_Y)$, the complete residue Gram
is diagonal:
\[
 \operatorname{diag}\{1,1-r_0,1+r_0-r_0^2,1,4,1,1\}.
\tag{WIE6}
\]
For instance the $T^2$ entry comes from substituting $R_2=r_0T$
in (WIE4): it is $r_0\nu_0-r_0^2$, plus the unit contribution from
(IC8a). Every entry is positive for $0\le r_0\le1/2$.
This proves full-parent tensor SOS at the corner also when
$1-\rho$ is infinitesimal. Along with the preserved interior proof
and the $\rho_0=0$ argument, every allowed source shape is covered
at the independent testing corner $u,w$ infinitesimal.

\subsection{The open testing axis}
\label{component:pf95b1c5fc32f}

\label{pf95b1c5fc32f:sec:weak-ih-open-axis}
\paragraph{Domain and source input.}
The source and testing real closed fields are independent.  Let $B_c,B_t$
be their rings of bounded elements and use $B=B_c\otimes_{\mathbb R}B_t$
for Gram coefficients.  Source square roots and testing square roots are
always taken in their respective fields.  In the weak-$C$ IH source chart,
write $z=u_c^2$ infinitesimal, $y=v_c^2$ unbounded,
$H=zh\in(0,1)$, $B_*=zb$ unbounded and $\rho=B_*/y$ bounded.
The source entry calculation gives
\[
 0\leq\rho_0\leq1,\qquad
 (B_*C)_0=\left(\frac{AH}{B_*}\right)_0
 =\left(\frac{\alpha(1-H)}{B_*}\right)_0
 =c:=\frac1{2(1+\rho_0)}.
\tag{OA1}
\]
In particular $C$ is infinitesimal and $A,\alpha$ are unbounded.
For $0<\rho_0<1$,
\[
 D_0=\frac{\rho_0^2}{2(1-\rho_0^2)};
\tag{OA2}
\]
for $\rho_0=0$, $D$ is positive infinitesimal; for $\rho_0=1$, $D$ is
unbounded.  These statements follow from
$D\sim\rho^2/[2e(1+\rho)]$ and $e_0=1-\rho_0$.
No lower bound on $n\geq0$ is imposed.

\paragraph{Statement.}
For every testing $u$ with nonzero finite residue and every infinitesimal
$w=1/v$, including $w=0$, the complete reciprocal-$v$ canonical C2 parent
is a tensor sum of squares for every state.  The same assertion holds for
complex states by applying the real quadratic form to real and imaginary
parts separately.

\paragraph{Exact parent and retained squares.}
Put $D_u=1+u^2$, $D_w=1+w^2$ and
\[
 M=W-wX+uY-uwZ,\qquad J=uwW+uX-wY-Z.
\]
The parent is
\begin{align*}
\widehat F={}&4nD_uD_wX^2+4mM^2+4CD_w(uW-Y)^2
 +4AD_uD_wY^2+4DJ^2+4\alpha D_u(Y-wZ)^2\\
&+2\{J^2+D_u(Y-wZ)^2+D_w(X+uZ)^2-D_uD_wZ^2\}.
\end{align*}
Retain the entire $n,m$ positive square terms. When $\rho_0=0$, retain the entire
$D$ positive square term as well and set $d=0$. When $0<\rho_0<1$, set $d=D$.
When $\rho_0=1$, choose a sufficiently large positive real $d$ below $D$
and retain $4(D-d)J^2$.  Denote the remainder by $F_*$.  All retained
terms are tensor squares or sums of tensor squares.

Use the graph $X=0$, $Y=HwZ$, and define
\[
 P=\frac{uW}{\sqrt{B_*}},\quad E=\sqrt{B_*}\,wZ,
 \quad V=uwW,\quad
 E_Y=\sqrt{\frac{B_*}{H(1-H)}}(Y-HwZ).
\tag{OA3}
\]
Crucially, the exact root relation is
\[
 PE=VZ.\tag{OA4}
\]
On the graph, the $C$, $A,\alpha$, $d$ positive square terms are respectively
\begin{align*}
 &4B_*CD_w(P-HE/B_*)^2,\\
 &\frac{4D_u}{B_*}\{AH^2D_w+\alpha(1-H)^2\}E^2,\\
 &4d\{V-(1+Hw^2)Z\}^2.
\end{align*}
The remaining core is exactly
\[
 2V^2-4(1+Hw^2)PE
 +\frac{2}{B_*}\{H^2+u^2(1-H)^2+H^2w^2\}E^2.
\tag{OA5}
\]
Thus every graph coefficient is bounded.  The last $E^2$ coefficient is
infinitesimal, uniformly for all $0<H<1$.

\paragraph{Restoring the $Y$ normal.}
Let $L=B_*/[H(1-H)]$ and $Y_*=HwZ$.
The first normal, divided by four, is
\begin{align*}
 L_Y:=\frac14\partial_YF_*\big|_{X=0,Y=Y_*}
={}&2CD_w(Y_*-uW)+2D_u\{AH D_w-\alpha(1-H)\}wZ\\
&-2dw\{V-(1+Hw^2)Z\}
 -w\{V-(1+Hw^2)Z\}+D_u(H-1)wZ.
\tag{OA6}
\end{align*}
Every coefficient of $L_Y/\sqrt L$ in $(P,E,V,Z)$ is infinitesimal
when $d$ is bounded, including $d=0$.  Indeed the only potentially finite contribution is
\[
 2D_u\frac{AH-\alpha(1-H)}{B_*}\sqrt{H(1-H)}E,
\]
and it has zero residue by (OA1); the additional $w^2$ contribution is
infinitesimal.  The other terms contain $C$, $w/\sqrt L$ or $1/B_*$.
The quadratic $E_Y$ coefficient has residue $4cD_{u,0}>0$, and the
$X E_Y$ coefficient has zero residue.  These statements remain valid
when $H$ or $1-H$ is infinitesimal.  No inverse of $d$, $H$ or $1-H$ enters a Gram coefficient
unaccompanied by its compensating factor.

\paragraph{The strict residue when $0<\rho_0<1$.}
Write $s=\sqrt{1+u_0^2}>1$ and $r=\rho_0$.
Using (OA4), split the core cross term as
$-4PE=-4aPE-4(1-a)VZ$, where $a$ is any real number satisfying
\[
 \frac{s}{s+r}<a<\min\left\{\frac{s}{1+r},\frac{s}{s-r}\right\}.
\tag{OA7}
\]
The interval is nonempty: both upper endpoints exceed $s/(s+r)$.
The residue Gram splits into the positive $E_Y$ pivot, a $(P,E)$ block
\[
 G_{PE}=\begin{pmatrix}2/(1+r)&-2a\\-2a&2s^2/(1+r)\end{pmatrix},
\]
and a $(V,Z,X)$ block
\[
 G_{VZX}=\begin{pmatrix}
 2+4d_0&-4d_0-2(1-a)&2u_0(1+2d_0)\\
 -4d_0-2(1-a)&4d_0&-4d_0u_0\\
 2u_0(1+2d_0)&-4d_0u_0&2s^2+4d_0u_0^2
 \end{pmatrix}.
\tag{OA8}
\]
The $(V,X)$ principal block is positive definite.  Its Schur complement
is positive exactly when
\[
 \frac{s}{s+r}<a<\frac{s}{s-r},
\]
as direct substitution of $d_0=r^2/[2(1-r^2)]$ shows.  Hence (OA7)
makes both blocks strictly positive.  This checks the entire $X$ normal,
not merely the scalar graph.

\paragraph{The endpoint $r=1$.}
For finite $d>0$, positivity of the $(V,Z,X)$ block is equivalent to
$a_-<a<a_+$, where
\[
 a_\pm(d)=
 \frac{(2d+1)s^2\pm s\sqrt{2d(2d+1)}}{(2d+1)u_0^2+1}.
\]
As $d\to\infty$, the lower endpoint tends to $s/(s+1)$ and the upper
endpoint to $s/(s-1)$.  Choose first a real $a$ strictly between
$s/(s+1)$ and $\min\{s/2,s/(s-1)\}$, then a sufficiently large finite
$d$.  The interval is nonempty because $s>1$.  Since source $D$ is
unbounded this cap is admissible, and the full residue is strict.

\paragraph{The endpoint $r=0$.}
Take $a=1$, $d=0$, having retained the whole $D$ positive square term, and omit
$Z$ from the root list. The residue is the direct sum of
\[
 2(P-E)^2+2u_0^2E^2,
 \qquad 2(V+u_0X)^2+2X^2,
 \qquad 2D_{u,0}E_Y^2.
\tag{OA9}
\]
It is strictly positive. The exact rewriting has bounded coefficients:
(OA5) uses only $P,E,V$, and the first $X$ normal from the core is
$uV+u(1-H)w^2Z=uV+u(1-H)wE/\sqrt{B_*}$.
Thus no unscaled $Z$ or inverse of $D$ remains.

\paragraph{Conclusion.}
In all three cases the original remainder has a bounded Gram matrix with
strictly positive real residue.  The bounded-Gram lifting lemma therefore
makes it a tensor sum of squares.  Adding the retained source-positive
positive square terms proves the assertion.  The argument is valid for arbitrary $n\geq0$,
including $n=0$, and literal projective $w=0$; no comparison between source
and testing infinitesimals was used.

\subsection{Complementary tests for weak $C$}
\label{component:p0fb32cd2f1bc}

\label{p0fb32cd2f1bc:sec:weak-ih-remote}
Assume the weak-$C$ IH source entry, with no restriction on
$\rho_0\in[0,1]$, on positive $r,H,1-H$, or on $n\ge0$.
The uniform source inputs are $B$ unbounded, $C,m$ infinitesimal,
$BC,AH/B,\alpha(1-H)/B$ positive units with common residue
$c=1/[2(1+\rho_0)]$, and $A,\alpha$ unbounded.
Every square retained below has nonnegative pure-source coefficient;
products $1+u^2$ and $1+v^2$ are themselves testing sums of squares.

\paragraph{Finite $v$, infinitesimal $u$.}
Write $Y_A=\sqrt A Y$ and $Z_\alpha=\sqrt\alpha(vY-Z)$, so
$Y=Y_A/\sqrt A$ and $Z=vY_A/\sqrt A-Z_\alpha/\sqrt\alpha$.
Retain the complete $n,C,D$ positive square terms. Use the five roots
\[
 U=uW,\quad X,\quad V=\sqrt m\,vW,\quad Y_A,\quad Z_\alpha.
\]
The $m$ positive square term is exactly
$4(V-\sqrt m X+u\sqrt{m/\alpha}Z_\alpha)^2$.
All coefficients in the remaining core are bounded after substituting
$Y,Z$, and all coefficients involving $Y$ or $Z$ are infinitesimal.
The full residue Gram is diagonal:
\[
 2U^2+2(1+v_0^2)X^2+4V^2
       +4(1+v_0^2)Y_A^2+4Z_\alpha^2.
\tag{IR1}
\]
This includes literal $u=0$ or $v=0$; redundant zero roots cause no
problem for a positive Gram representation.

\paragraph{Finite $v$, unit $u$.}
Retain $n,m,C,D$. Use $(W,X,Y_A,Z_\alpha)$.
The residue of the core and selected $A,\alpha$ positive square terms is
\[
 2\{u_0^2(W+v_0X)^2+(1+v_0^2)X^2\}
       +4(1+u_0^2)(1+v_0^2)Y_A^2
       +4(1+u_0^2)Z_\alpha^2,
\tag{IR2}
\]
which is strictly positive because $u_0\ne0$.

\paragraph{Finite $v$, unbounded $u$.}
Use the homogeneous coordinate $p=1/u$, infinitesimal or zero.
Retain $n,m,D$ and use
\[
 P=pX,\quad Q=W+vX,\quad S=\sqrt C W,\quad Y_A,\quad Z_\alpha.
\]
The $C$ positive square term becomes
$4(1+v^2)(S-p\sqrt{C/A}Y_A)^2$.
The entire selected expression has bounded coefficients, and residue
\[
 2Q^2+2(1+v_0^2)P^2+4(1+v_0^2)S^2
       +4(1+v_0^2)Y_A^2+4Z_\alpha^2.
\tag{IR3}
\]
Thus every finite-$v$ chart is complete, independently of $n$.

\paragraph{Both testing angles unbounded.}
Let $p=1/u$, $w=1/v$ be infinitesimal or zero and put
$D_p=1+p^2$, $D_w=1+w^2$.
Retain the entire $n,m$ positive square terms. Set $d=\min(1,D)>0$ in the source
field and retain $4(D-d)J^2$. The remainder is exactly
\begin{align*}
 F_*={}&4CD_w(W-pY)^2+4AD_pD_wY^2+4dJ^2
       +4\alpha D_p(Y-wZ)^2\\
 &+2\{J^2+D_p(Y-wZ)^2+D_w(pX+Z)^2-D_pD_wZ^2\},\\
 J={}&wW+X-pwY-pZ.
\end{align*}
Use the graph $X=-wW$, $Y=HwZ$, the normals
$E_X=X+wW$, $E_Y=\sqrt L(Y-HwZ)$ with
$L=B/[H(1-H)]$, and scalar roots
\[
 S=\sqrt C W,\quad E=\sqrt B\,wZ,\quad
 U=pwW,\quad P=\sqrt d\,pZ.
\tag{IR4}
\]
On the graph the exact expression is
\begin{align*}
 F_*|_{\rm graph}={}&4D_w(S-pH\sqrt{C/B}E)^2
 +\frac{4D_p}{B}\{AH^2D_w+\alpha(1-H)^2\}E^2\\
 &+4(1+Hw^2)^2P^2+2D_wU^2
 -\frac{4pD_w}{\sqrt{BC}}SE
 +\frac2B\{(1-H)^2+H^2p^2D_w\}E^2.
\tag{IR5}
\end{align*}
All coefficients are bounded and its residue is
$4S^2+4cE^2+2U^2+4P^2$.

The first $X$ normal, divided by four, is exactly
\[
 L_X=-pD_wU+\frac{p(1-H)w}{\sqrt B}E
               -2\sqrt d(1+Hw^2)P.
\tag{IR6}
\]
The first $Y$ normal is
\begin{align*}
 L_Y={}&-2p\sqrt C D_w S+\frac{E}{\sqrt B}
 \{2Cp^2HD_w+2D_p(AHD_w-\alpha(1-H))\\
 &\hspace{38mm}+(2d+1)p^2(1+Hw^2)-D_p(1-H)\}.
\tag{IR7}
\end{align*}
After dividing by $\sqrt L$, all its coefficients are infinitesimal;
the only possible nonzero term is proportional to
$(AH-\alpha(1-H))/B$, whose residue is zero.
The quadratic normal coefficients have residues
$2+4d_0$ and $4c$, and their mixed coefficient has zero residue.
The only nonzero scalar-normal coupling is the $(P,E_X)$ block
\[
 \begin{pmatrix}4&-4\sqrt{d_0}\\-4\sqrt{d_0}&2+4d_0\end{pmatrix},
\tag{IR8}
\]
whose determinant is exactly $8$. Thus the full Gram is strict even
when $D$ is infinitesimal and even when $n=0$. The normalization uses
$\sqrt d$ but never $1/\sqrt d$ in a Gram coefficient.

Each displayed Gram is a representation of the complete selected
state quadratic form, not just its graph restriction. The bounded
positive-residue lemma and the retained squares prove tensor SOS.
These charts, the open-axis theorem, and the three local source-shape
proofs cover every independent testing pair.

\subsection{Local completion for strong $C$}
\label{component:pf3e0ca8bf054}

\label{pf3e0ca8bf054:sec:strong-ih-local}
Assume the strong IH source entry of \S\ref{pf5726218aec8:sec:strong-angle-entry}:
$z$ is infinitesimal, $y$ unbounded, $B=H+L$ a bounded positive unit,
$0<H<1$, $L>0$, and $C,\alpha$ strong. The source parameter
$\rho=B/y$ is infinitesimal, $s=r/\rho$ bounded, and $c_s=(1+L)/B$
is a positive unit. Testing $u,w=1/v$ are infinitesimal or zero.
We use the exact graph, five roots and both full normal identities
(IC3)--(IC12) of the preserved corner calculation; none of those
algebraic identities assumes $C$ weak.

The source quantities obey (GE3), and direct rational simplification
gives additionally
\[
 \frac{k}{B\rho}\longrightarrow
 \frac{2(B^2s+BHs+BH-Bs-H^2+H)}{B^2}.
\tag{SL1}
\]
No opening denominator occurs in this normalized expression.

Set $T_\rho=\rho T$ and
$\eta=\sqrt{H(1-H)/B}$, so $Y=HwZ+\eta E_Y$.
The full Gram's entire explicit dependence on $C$ is exactly
\[
 4C(U_1^2+U_2^2),\quad
 U_1=uW-Y-LwZ=\sqrt B R_0-\eta E_Y,
\]
\[
 U_2=w(uW-Y)-(L/y)Z
       =R_2-\sqrt{H/B}E-w\eta E_Y.
\tag{SL2}
\]
The scalar and normal identities verifying this assertion use
\[
 \sqrt H R_2-\sqrt{BH}\,wR_0-\sqrt B E=0.
\tag{SL3}
\]
Choose a positive real cap $c_*\le C$ and retain
$4(C-c_*)(U_1^2+U_2^2)$. In the remaining exact Gram replace explicit
$C$ by $c_*$, leaving $k$ and all other original source quantities
unchanged. The distinction matters: this is an exact square
subtraction from the complete Gram, not a perturbation of the source.
Add $\lambda$ times the exact zero (WIE2), where $\lambda>0$ is a
real number to be chosen. All resulting coefficients are bounded,
using (GE3), (SL1), and the opening-normalized roots.

We display the entire residue. The roots $R_1$ and $E_X$ have isolated
positive pivots $2c_s$ and $1$. The $(R_0,E_Y)$ block is
\[
 \begin{pmatrix}
 4Bc_*&-4c_*\sqrt{H(1-H)}\\
 -4c_*\sqrt{H(1-H)}&
 (4c_*+4A+4\alpha+2)H(1-H)/B
 \end{pmatrix}.
\tag{SL4}
\]
Here and below all expressions in a residue matrix mean their bounded
residues; products such as $AH$ are evaluated by (GE3), even when
$A$ itself is unbounded. Its determinant is exactly $8c_*(1+L)>0$.
Put $a_*=2(2c_*+1)$. The $(R_2,E)$ block is
\[
 \begin{pmatrix}
 a_*&-a_*\sqrt{H/B}\\
 -a_*\sqrt{H/B}&H(4A+a_*)/B
 \end{pmatrix},
\tag{SL5}
\]
with determinant $2a_*c_s>0$. The only remaining coupling is between
$R_2$ and $T_\rho$, whose Gram entry is
$c_s(H/B-s)$. The $T_\rho$ diagonal is $\lambda+2c_s s$.
Choose $\lambda$ larger than the finite Schur-complement threshold
of the positive block (SL5). This produces a strictly positive
residue matrix. Since this choice is made after taking a real
residue, it is a real scalar and introduces no mixed inverse.

We spell out why no normal term is hidden in this block description.
In (IC11), $f_X/\rho$ is bounded by (GE3) and every scalar-root
multiplier contains $u$ or $w$. In (IC12), the only nonzero limiting
term is the explicit $C$ contribution to $R_0E_Y$, already included
in (SL2)--(SL4). After the cap, all other normal multipliers contain
$z$, $u$, or $w$. The unbounded factors $A$ and $\alpha$ always occur
as $AH$ and $\alpha(1-H)$ after normalization, hence remain bounded
also at $H_0=0,1$. The normal $X$ coefficient divided by $N=4n+2$
has residue one, and its coupling to $E_Y$ vanishes. The remaining
$Y$ coefficient is exactly the lower-right entry of (SL4).
These facts follow directly from the displayed exact normal
factorizations, not from scalar-graph positivity.

The bounded positive-residue Gram lemma now gives tensor SOS for the
complete original parent, after restoring the retained positive
$C-c_*$ squares. This holds for all states and all independent local
tests, for both opening endpoints, for arbitrary $n\ge0$, and at
literal $u=0$ or $w=0$.

\subsection{Complementary tests for strong $C$}
\label{component:p7405dbf89341}

\label{p7405dbf89341:sec:strong-ih-remote}
Use the source entry of \S\ref{pf5726218aec8:sec:strong-angle-entry}. In particular
$B=H+L$ is a bounded positive unit, $m,D$ are infinitesimal,
$C,A,\alpha$ are strong, and the bounded products satisfy
\[
 CL\to\tfrac12,\qquad AH\to(1+L)/2,
 \qquad\alpha(1-H)\to B/2.
\tag{SR1}
\]
Whenever an event coefficient is unbounded, a positive finite real cap
may be used and its excess retained in its original positive square.
When a coefficient is bounded it may be used unchanged.

\paragraph{Finite-penalty principle.}
Suppose a real symmetric quadratic form is strictly positive on the
common kernel of finitely many real linear forms. Adding sufficiently
large positive multiples of their squares makes the complete form
strictly positive. To verify this without a limiting assertion, choose
an orthogonal splitting into their common kernel and its complement.
The upper block has a positive lower eigenvalue $c$, and the penalty
matrix on the complementary block has a positive lower eigenvalue
$e$. If the remaining lower block has norm at most $b$ and the cross
block norm is at most $a$, any penalty $M$ with $Me>b+a^2/c$ makes the
Schur complement strictly positive. All choices are real scalars.
This principle makes the opening cases below uniform; no unbounded
coefficient is inserted into a bounded Gram.

\paragraph{Testing $u$ infinitesimal, $v$ finite.}
Retain the $n,D$ positive square terms and excess caps of $C,A,\alpha$.
Write their finite selected coefficients as $c,a,b>0$. Use roots
$U=uW,Y,Z,X,V=\sqrt m\,vW$. The $m$ positive square term is exactly
\[
 4(V-\sqrt m X+\sqrt m\,uvY-\sqrt m\,uZ)^2.
\]
Together with the selected positive square terms and the original core this is an
exact bounded Gram expression. Its residue has isolated $X,V$ pivots
$2(1+v^2),4$ and the following $(U,Y,Z)$ block (residue notation suppressed):
\[
 \begin{pmatrix}
 2+4cD_v&-2-4cD_v&-2v\\
 -2-4cD_v&2D_v+4cD_v+4aD_v+4bv^2&-4bv\\
 -2v&-4bv&4b
 \end{pmatrix}.
\tag{SR2}
\]
In the interior, use $c=1/(2L)$, $a=(1+L)/(2H)$ and
$b=B/[2(1-H)]$. Its leading principal minors are
\[
 \frac{2(L+D_v)}L,\quad
 \frac{4(L+1)(D_v-H)(L+D_v)}{HL(1-H)},\quad
 \frac{8B(L+1)^2D_v}{HL(1-H)},
\tag{SR3}
\]
all positive. At an opening, apply the finite-penalty principle:
\begin{itemize}
\item $H_0=0$: the large $A$ positive square term forces $Y=0$; the restricted
      two-dimensional determinant is $4(L+1)>0$.
\item $H_0=1$: the large $\alpha$ positive square term forces $Z=vY$; the restricted
      determinant is $4(L+1)^2D_v/L>0$.
\item $L_0=0$: the large $C$ positive square term forces $U=Y$; the restricted
      determinant is $4/(1-H)>0$.
\item $L_0=0$, $H_0=1$: both constraints leave the strictly positive
      form $2Y^2$.
\end{itemize}
The case $L_0=H_0=0$ is excluded by $B_0>0$. Thus (SR2) can always be
made strictly positive with bounded selected coefficients.

\paragraph{Testing $u$ a unit, $w=1/v$ infinitesimal.}
Retain $n,m$ and the excess caps of $C,A,\alpha$; retain the full
$D$ positive square term using its natural root. Put
$U=uW$, $E=wZ$, $P_Z=\sqrt D Z$. In roots $(U,Y,E,X,P_Z)$ the exact
selected expression is
\begin{align*}
 &4cD_w(U-Y)^2+4aD_uD_wY^2+4bD_u(Y-E)^2\\
 &\quad+4\{\sqrt D(wU+uX-wY)-P_Z\}^2\\
 &\quad+2\{(wU+uX-wY)^2+D_u(Y-E)^2+D_wX^2-E^2
                   -2UE+2YE+2uwXE\}.
\tag{SR4}
\end{align*}
Its isolated $X,P_Z$ residue pivots are $2D_u,4$. The $(U,Y,E)$ block is
\[
 \begin{pmatrix}
 4c&-4c&-2\\
 -4c&4c+(4a+4b+2)D_u&-4bD_u-2u^2\\
 -2&-4bD_u-2u^2&4bD_u+2u^2
 \end{pmatrix}.
\tag{SR5}
\]
For interior openings the leading minors are
\[
 2/L,\qquad4(1+L)D_u/[HL(1-H)],\qquad
 8u^2(1+L)^2D_u/[HL(1-H)],
\tag{SR6}
\]
strict because $u_0\ne0$. The opening constraints and restricted
positive determinants are respectively
\[
 \begin{array}{c|c|c}
 H_0=0&Y=0&4u^2(1+1/L)\\
 H_0=1&Y=E&4u^2(1+L)/L\\
 L_0=0&U=Y&4u^2D_u/[H(1-H)]
 \end{array}
\]
and at $L_0=0,H_0=1$ the remaining form is $2u^2Y^2$.
Finite penalties therefore give a strict full Gram at every opening.

\paragraph{Both testing angles unbounded.}
Use (IR4)--(IR7), replacing $C$ by a finite positive cap $c_*$ and
setting $d=D$ (which is infinitesimal in the strong IH source).
Retain the excess $C$ positive square term and the entire $n,m$ positive square terms.
The roots are
\[
 \begin{gathered}
 S=\sqrt{c_*}W,\quad E=\sqrt B\,wZ,\quad U=pwW,\quad P=\sqrt D\,pZ,\\
 E_X=X+wW,\qquad E_Y=\sqrt{B/[H(1-H)]}(Y-HwZ).
 \end{gathered}
\]
The exact graph identity (IR5) remains valid. The $E^2$ coefficient
has residue $2c_s$, $c_s=(1+L)/B$, because its core contribution
$2(1-H)^2/B$ must be retained when $B$ is a unit.
In (IR7), the potentially nonzero first-normal coefficient is now
\[
 2\{AH-\alpha(1-H)\}-(1-H),
\]
whose residue is zero by (SR1). The normal $E_Y$ pivot is also $2c_s$.
All scalar-normal cross coefficients are infinitesimal, since
$\sqrt D$ is infinitesimal. Thus the full diagonal residue is
\[
 \operatorname{diag}(4,2c_s,2,4,2,2c_s)>0.
\tag{SR7}
\]

\paragraph{The remaining remote charts.}
When $u$ is unbounded and $v$ is a unit, the homogeneous core has
residue $2\{(W+vX)^2+(vY-Z)^2\}$. Positive finite caps on the $C$
and $A$ positive square terms supply $W^2$ and $Y^2$, respectively, so the residue
is strict on $(W,X,Y,Z)$. When $u$ is unbounded and $v$ infinitesimal,
put $p=1/u$ and use the five roots
$(P,Q,W,Y,Z)=(pX,W+vX,W,Y,Z)$. Positive caps $c,a,b$ give the residue
\[
 2Q^2+2(P+Z)^2+4cW^2+4aY^2+4bZ^2>0.
\tag{SR8}
\]
When $u$ is a unit and $v$ infinitesimal, the core has residue
$2\{(uW-Y)^2+(X+uZ)^2\}$; the positive $A$ and $\alpha$ caps give
$Y^2$ and $Z^2$, making the residue strict. In these expressions all
excess event coefficients are retained in their positive squares.
When both testing angles are units, the inherited noncommuting-test
theorem applies independently of source scales.

\paragraph{Exhaustive join.}
Classify each testing angle in the projective line as infinitesimal,
a bounded unit, or unbounded, with literal zero and infinity included
in the corresponding homogeneous chart. The infinitesimal/unbounded
cell is the local theorem of \S\ref{pf3e0ca8bf054:sec:strong-ih-local}.
The infinitesimal/finite cells follow from (SR2); the unit/unbounded
cell from (SR4); the unbounded/unbounded cell from (SR7); the
unbounded/finite and unit/infinitesimal cells from (SR8) and its
preceding argument; and the unit/unit cell from the inherited theorem.
These nine cells are exhaustive. Both exact sign gauges
$(u,W,X)\mapsto(-u,-W,-X)$ and
$(v,W,Y)\mapsto(-v,-W,-Y)$ preserve the full parent, so every
signed testing chart is included. Every certificate retains the full
state normals and the omitted original $\lambda,\beta$ remain zero.
Therefore the entire strong IH source leaf has all-tests tensor-SOS
case, including $n=0$ and every projective endpoint.

\section{One ordinary and one small actual angle}
\subsection{Source reduction at an ordinary--small contact}
\label{component:pc3daa35e9a9b}

\label{pc3daa35e9a9b:sec:ui-full-parent-reduction}

This section records the original source, every parent positive square term, and the
bounded products used in the following reciprocal and local certificates.
It is an exact algebraic reduction, not an all-tests case. The
coefficient and testing fields remain independent real closed extensions
of $\mathbb R$.

\paragraph{Source definitions and physical domain.}
Take the actual source angles $u_c=a$ and $v_c=k$, where $a>0$ is a
bounded unit and $k>0$ is infinitesimal. Set $z=a^2$ and
$\varepsilon=k^2/z$. In the original canonical C2 chart put
\[
 h=b-\varepsilon t,\quad A_c=1+zr,\quad j=A_c-t,\quad
 g=(1+z)b-1+r-\varepsilon zt,\quad S=A_c+\varepsilon zt.
 \tag{UR1}\label{pc3daa35e9a9b:eq:ui-source-gaps}
\]
The native positive gaps are $r,1-r,t,h,g,j,1-zh$; assume all original
event signs, including $n\ge0$, and assume $C,\alpha$ are bounded below
by positive real numbers. No positive lower bound on a native opening
is imposed. Define the two complete source polynomials
\begin{align*}
 \Delta={}&1+(1+z)b+z(1+z)b^2+2z(1+z)br+z^2r^2+2zr\\
 &+\varepsilon z\{(1+z)b+1+zr(2-r)\},\\
 N={}&(1+z)(b+r)\{b+\varepsilon(1-t)\}
          -\varepsilon(1+\varepsilon z)(1-r)t.
 \tag{UR2}\label{pc3daa35e9a9b:eq:ui-source-polynomials}
\end{align*}
With $d=1$ and $\lambda=\beta=0$, the unchanged original coefficients
are
\[
 \begin{aligned}
 n&=\frac{(b+r)A_cS-r\Delta}{2r\Delta},&
 m&=\frac{(1+k^2)(1-r)A_cS}{2g\Delta},\\
 C&=\frac{A_cN}{2\varepsilon t\Delta},&
 A&=\frac{\varepsilon A_c(b+r)S}{2h\Delta},\\
 D&=\frac{(1+z)(b+r)(b-k^2r)S}{2\varepsilon j\Delta},&
 \alpha&=\frac{z^2(b+r)N}{2(1-zh)\Delta}.
 \end{aligned}                                                   \tag{UR3}
 \label{pc3daa35e9a9b:eq:ui-original-source}
\]
These are the original rational C2 source formulas, with its angle-square
parameter replaced by $z$ and its squared angle-ratio parameter by
$\varepsilon$; all quantities formerly supplied by a source-file
reference are displayed in (UR1)--(UR3).

The source bounds used below follow directly from these definitions.
First $t<A_c$, $h<1/z$, and $b=h+\varepsilon t$ give bounded $t,h,b$.
The inequality $(1+z)b+r>1$ gives a positive lower bound on $b+r$;
(UR2) makes $\Delta$ a positive bounded unit. The exact positive identity
\[
 N=(1+z)(b+r)h+
       \varepsilon\{g+j+rt(1+\varepsilon z)\}                 \tag{UR4}
\]
shows that infinitesimal $h$ would make $b,N$ infinitesimal and
$1-zh$ a unit, contradicting strong $\alpha$. Thus $h,b,N$ are units.
The numerator of $2rn\Delta$ evaluated at $r=1$ is exactly
\[
 -z(\varepsilon+h)(b+1)(z+1).
\]
It has negative unit residue, so $n\ge0$ forces $1-r$ to be a unit.
Consequently
\[
 mg,\quad A/\varepsilon,\quad \varepsilon Ct,\quad
 \varepsilon Dj,\quad \alpha(1-zh)
 \quad\hbox{are positive bounded units},\qquad nr\text{ is bounded}.
 \tag{UR5}\label{pc3daa35e9a9b:eq:ui-source-bounds}
\]
In particular $m$ is strong, $A$ is infinitesimal, and $C,D$ may have
arbitrarily large positive sizes. Literal $n=0$ is retained.

\paragraph{The entire parent and exact contact.}
For independent testing angles and a homogeneous state define
\[
 \begin{gathered}
 D_u=1+u^2,\quad D_v=1+v^2,\quad
 r_{\rm amp}=uW-Y,\quad Q=uX-Z,\\
 J=r_{\rm amp}+vQ=uW+uvX-Y-vZ,\\
 M=vW-X+uvY-uZ,\quad p=vY-Z,\quad B=X+uZ.
 \end{gathered}
\]
The signed core and the full original parent are
\[
 \mathcal K=2\{J^2+D_up^2+D_vB^2-D_uD_vZ^2\},
\]
\[
 F=4nD_uD_vX^2+4mM^2+4CD_vr_{\rm amp}^2
       +4AD_uD_vY^2+4DJ^2+4\alpha D_up^2+\mathcal K.
 \tag{UR6}\label{pc3daa35e9a9b:eq:ui-full-parent}
\]
Normalize the actual contact by $W_c=1$. Introduce the bounded pure-source
quantities
\[
 x=-r/b,\quad \rho=t/(ab),\quad \sigma=-j/(ab),\quad
 s=\sigma-\rho=-A_c/(ab),\quad
 \mu=g/b,\quad \pi=(zh-1)/(ab).                       \tag{UR7}
 \label{pc3daa35e9a9b:eq:ui-contact-amplitudes}
\]
Here $s$ is a nonzero bounded unit. The exact stationary contact is
\[
 (u_c,v_c;W_c,X_c,Y_c,Z_c)
 =\left(a,k;1,kx,a-\rho k^2,\frac{k}{ab}\right),       \tag{UR8}
\]
and its amplitudes are
\[
 r_{{\rm amp},c}=\rho k^2,\quad J_c=\sigma k^2,\quad
 M_c=\mu k,\quad p_c=\pi k,\quad Q_c=sk.
\]
Substitution of (UR3) gives $F=0$ and all four state and both angle
first derivatives zero at (UR8). These are the unchanged native
contact identities; no equality of independent testing rates is used.

\paragraph{Products with all possible opening poles cancelled.}
Define two positive bounded units
\[
 c_1=\frac{zA_cN}{2\Delta},\qquad
 c_2=\frac{z(1+z)(b+r)(b-k^2r)S}{2\Delta}.
\]
The exact products required by the local cap argument are
\[
 \begin{aligned}
 nx&=-\frac{(b+r)A_cS-r\Delta}{2b\Delta},&
 m\mu&=\frac{(1+k^2)(1-r)A_cS}{2b\Delta},\\
 \alpha\pi&=-\frac{z^2(b+r)N}{2ab\Delta},&
 C\rho k^2&=\frac{c_1}{ab},\qquad
 D\sigma k^2=-\frac{c_2}{ab}.
 \end{aligned}                                                   \tag{UR9}
 \label{pc3daa35e9a9b:eq:ui-cancelled-products}
\]
Their products with the bounded factors $x,\mu,\pi$ also show
$nx^2,m\mu^2,\alpha\pi^2$ bounded. Crucially, cancellation before
estimating gives the further exact identity
\[
 Ct-Dj=
 \frac{(1+z)(b+r)\{A_cj-zt(b-\varepsilon zr)\}
       -A_c(1+\varepsilon z)(1-r)t}{2\Delta}.          \tag{UR10}
 \label{pc3daa35e9a9b:eq:ui-cancelled-CD-difference}
\]
Thus $C\rho+D\sigma=(Ct-Dj)/(ab)$ is bounded for every positive rate
of $t,j$, including infinitesimal openings.

Set
\[
 L=C+D,\qquad q=D/L,\qquad H=CD/L,\qquad
 \tau=(C\rho+D\sigma)/L,\qquad \gamma=k^2H.
\]
Here $q,1/L,\tau,L\tau$ are bounded, and
\[
 \gamma=\frac{c_1c_2}{c_1j+c_2t}\asymp1,
 \qquad t+j=A_c\asymp1.                              \tag{UR11}
 \label{pc3daa35e9a9b:eq:ui-harmonic-unit}
\]
The denominator in (UR11) is a positive unit even if either opening
is infinitesimal. At finite source residues it gives precisely $H_0$
in (RC3). Equations (UR9)--(UR11) also show that the bounded finite
part extends continuously to the allowed boundary residues: no
inverse of $r,g,1-zh,t,j$ remains in these products.

\begin{proposition}[Exact bounded UI remainder]
\label{pc3daa35e9a9b:prop:ui-bounded-remainder}
For bounded source caps $0\le n_*\le n$ and positive bounded-unit
caps $m_*\le m$, $\alpha_*\le\alpha$, the parent (UR6) is the sum
of explicit positive centered positive square terms and a bounded polynomial
remainder $G$. Every centered positive square term vanishes with its full first jet
at (UR8), so $G$ has exact zero value and all first derivatives there.
\end{proposition}
\begin{proof}
The exact original $C,D$ combination is
\[
 CD_vr_{\rm amp}^2+DJ^2
   =L(r_{\rm amp}+qvQ)^2+Hv^2Q^2+Cv^2r_{\rm amp}^2.
\]
Retain the three nonnegative positive square terms
\[
 \begin{aligned}
 P_L&=4L(r_{\rm amp}+qvQ-\tau k^2W)^2,\\
 P_H&=4H(vQ-sk^2W)^2,\\
 P_C&=4Cv^2(r_{\rm amp}-\rho k^2W)^2.
 \end{aligned}                                                   \tag{UR12}
\]
Their subtraction leaves exactly
\begin{align*}
 G_{CD}={}&8L\tau k^2W(r_{\rm amp}+qvQ)-4L\tau^2k^4W^2\\
 &+8Hs k^2vWQ-4Hs^2k^4W^2\\
 &+8C\rho k^2v^2Wr_{\rm amp}-4C\rho^2k^4v^2W^2.
 \tag{UR13}
\end{align*}
All its coefficients are bounded by (UR9)--(UR11), including
$L\tau^2=(L\tau)^2/L$. Centering at $k^2W$ is essential to this
boundedness assertion.

Also retain
\[
 \begin{aligned}
 P_n&=4(n-n_*)D_uD_v(X-xvW)^2,\\
 P_m&=4(m-m_*)(M-\mu vW)^2,\\
 P_\alpha&=4(\alpha-\alpha_*)D_u(p-\pi vW)^2.
 \end{aligned}                                                   \tag{UR14}
\]
The exact corresponding remainders are
\[
 \begin{aligned}
 G_n={}&4n_*D_uD_vX^2+8(n-n_*)xvWD_uD_vX
                    -4(n-n_*)x^2v^2W^2D_uD_v,\\
 G_m={}&4m_*M^2+8(m-m_*)\mu vWM
                    -4(m-m_*)\mu^2v^2W^2,\\
 G_\alpha={}&4\alpha_*D_up^2+8(\alpha-\alpha_*)\pi vWD_up
                    -4(\alpha-\alpha_*)\pi^2v^2W^2D_u.
 \end{aligned}                                                   \tag{UR15}
\]
These coefficients are bounded by (UR9) and the bounded caps. With
\[
 G=G_{CD}+G_n+G_m+G_\alpha+4AD_uD_vY^2+\mathcal K,
\]
the entire exact identity is
\[
 F=P_L+P_H+P_C+P_n+P_m+P_\alpha+G.                 \tag{UR16}
 \label{pc3daa35e9a9b:eq:ui-entire-reduction}
\]
It includes the signed core and every original state-normal term.
All six retained roots vanish at (UR8), proving the contact-jet
assertion. This proposition does not claim positivity of $G$.
\end{proof}

For the local theorem retain $P_L,P_C,P_n,P_m,P_\alpha$ outside the
bounded Gram and combine $G$ with
\[
 P_H=4\gamma R_h^2,\qquad R_h=(vQ-sk^2W)/k.
 \tag{UR17}
\]
The factor $1/k$ is pure coefficient. The exact graph normals are
\[
 E_X=X-xvW,\quad E_Z=Z-(ux-s)vW,\quad
 E_r=r_{\rm amp}-\rho k^2W.
\]
Their inverse is polynomial:
\[
 X=E_X+xvW,\quad Z=E_Z+(ux-s)vW,\quad
 Y=uW-\rho k^2W-E_r,\quad Q=svW+uE_X-E_Z.
\]
Thus no state is discarded, including $W=0$. The following reciprocal
lemma identifies the leading relative-angle completion; the local
certificate supplies the complete normal lift, cap choice, and strict
Gram. The outer-angle and other commuting-axis join remains a separate
obligation, so this reduction adds no all-tests source case.

\subsection{Reciprocal completion of a singular angle}
\label{component:pa142a027a282}

\label{pa142a027a282:sec:ui-reciprocal-contact}

The following algebraic identity applies at the indicated singular angle.

\begin{lemma}[Reciprocal contact completion]
Let $A_0,A_2$ be symmetric matrices over a real closed field, let
$q_0$ be a column vector, and consider
\[
 f(q,V,W)=q^T(A_0+V^2A_2)q+2VW b^Tq+(c_0+c_2V^2)W^2.
\]
Suppose $(q,V,W)=(q_0,1,1)$ is a zero with all first derivatives
zero. Then
\[
 b=-(A_0+A_2)q_0,\qquad
 c_0=q_0^TA_2q_0,\qquad c_2=q_0^TA_0q_0,
\]
and the entire polynomial has the exact decomposition
\[
 \boxed{f=(q-Vq_0W)^TA_0(q-Vq_0W)
             +(Vq-q_0W)^TA_2(Vq-q_0W).}             \tag{RC1}
\]
Consequently $A_0,A_2\succeq0$ imply a sum of squares. If
$A_0\succ0$ and $q_0^TA_2q_0>0$, its nonzero zero states have
$V=\pm1$ and $q=Vq_0W$.
\end{lemma}
\begin{proof}
The $q$ equation gives the displayed value of $b$. The $V$
equation then gives $c_2=q_0^TA_0q_0$, and the $W$ equation gives
$c_0=q_0^TA_2q_0$. Expanding the right side of (RC1) proves the
identity. The zero statement follows first from the strictly positive
$A_0$ term and then from $(V^2-1)^2W^2q_0^TA_2q_0$.
\end{proof}

\paragraph{Homogeneous form without a mixed inverse.}
For a positive pure-coefficient $k$, substitute $q=Q/k$ and
$V=v/k$, and multiply (RC1) by $k^2$. It becomes
\[
 (Q-vq_0W)^TA_0(Q-vq_0W)
 +\frac1{k^2}(vQ-k^2q_0W)^TA_2(vQ-k^2q_0W).       \tag{RC2}
\]
Thus its roots use only the pure-coefficient factor $1/k$.
There is no inverse or radical of a mixed tensor element and no
comparison between a testing rate and $k$.

\paragraph{The bounded-source C2 UI leading model.}
Here $u_c=a$ is a positive bounded unit and $v_c=k$ is
infinitesimal, with $d=1$. To state a residue model, first restrict
to the subdomain where the original $n,m,\alpha$ are bounded and
the source parameters $r,g,1-a^2b$ have nonzero residues. In the
following formulas $a,b,r$ denote those residues, $z=a^2$, and
\[
 A_c=1+zr,\quad g=(1+z)b-1+r,\quad
 \Delta=1+(1+z)b+z(1+z)b^2+2z(1+z)br+z^2r^2+2zr.
\]
The source formulas give
\[
\begin{gathered}
 n=\frac{(b+r)A_c^2-r\Delta}{2r\Delta},\qquad
 m=\frac{(1-r)A_c^2}{2g\Delta},\qquad
 \alpha=\frac{z^2(1+z)b(b+r)^2}{2(1-zb)\Delta},\\
 H_0=\overline{k^2CD/(C+D)}
       =\frac{z(1+z)(b+r)b}{2\Delta},\qquad
 A_0^{\rm weak}=\overline{A/k^2}
       =\frac{A_c^2(b+r)}{2zb\Delta}.
\end{gathered}                                                   \tag{RC3}
\]
Normalize the actual contact by $W=1$. With $v=kV$, $X=k\xi$,
$Z=k\zeta$ and the leading constraint $Y=aW$, the leading
polynomial of $F/k^2$ is
\begin{align*}
 f={}&4n(1+a^2)\xi^2
 +4m\{V(1+a^2)W-\xi-a\zeta\}^2
 +4H_0V^2(a\xi-\zeta)^2\\
 &+4A_0^{\rm weak}(1+a^2)a^2W^2
 +(4\alpha+2)(1+a^2)(VaW-\zeta)^2\\
 &+2(\xi+a\zeta)^2-2(1+a^2)\zeta^2.               \tag{RC4}
\end{align*}
It has the form of the lemma, with $q=(\xi,\zeta)^T$,
\[
 q_0=(-r/b,1/(ab))^T,
\]
\[
 q^TA_0q=4n(1+a^2)\xi^2+(4m+2)(\xi+a\zeta)^2
                      +4\alpha(1+a^2)\zeta^2,
 \qquad q^TA_2q=4H_0(a\xi-\zeta)^2.
\]
Here $A_0\succ0$ and $H_0>0$. Put
$s_0=aq_{0,1}-q_{0,2}=-(1+a^2r)/(ab)\ne0$.
The exact source identities (RC3) verify all three coefficient
relations in the lemma. Hence the apparently problematic relative
normal and scale difference combine into the single root
\[
             \frac{v(aX-Z)-s_0k^2W}{k}.             \tag{RC5}
\]
The ordinary centered normals $(X,Z)-vq_0W$ supply the remaining
positive directions. Requiring separate positive pivots for the
relative normal and the scale difference would obscure this
completion.

This proves an exact leading-model identity on its stated bounded
subdomain. It does not discard the additional large source scales:
those are handled algebraically by the centered-positive square term reduction in
Section~\ref{pc3daa35e9a9b:sec:ui-full-parent-reduction}. The complete local lifted
Gram is proved in Section~\ref{p76545700d90b:sec:local-ui-certificate}; the outer
projective join remains a separate obligation.

\subsection{Local contact certificate}
\label{component:p76545700d90b}

\label{p76545700d90b:sec:local-ui-certificate}

This section uses the exact full-parent reduction and source notation
in Section~\ref{pc3daa35e9a9b:sec:ui-full-parent-reduction}. It proves a local independent-test
statement. The other commuting axis and the outer projective join
are not included.

\begin{theorem}[Local ordinary--small C2 certificate]
Use the unchanged physical C2 source with $d=1$, actual contact
$u_c=a$ a positive bounded unit and $v_c=k$ a positive infinitesimal,
and $C,\alpha$ bounded below by positive real numbers. Then the full
parent is tensor SOS for every independent test with $v$ infinitesimal
and $u-a$ infinitesimal, and every projective state. The original
$n=0$ boundary and all positive source opening rates are included.
No upper bound on $n,m,C,D,\alpha$ is imposed.
\end{theorem}

Here $u-a$ is the difference of the corresponding tensor factors.
The residue condition means that their real residues agree; it does
not identify the source and testing fields or compare their rates.
The exact sign gauges give the corresponding negative-angle charts.

\paragraph{Bounded full remainder and its exact graph.}
The positive square term reduction leaves a bounded polynomial remainder $G$ and
the positive relative square
\[
 4\gamma R_h^2,\qquad
 \gamma=k^2CD/(C+D)\asymp1,\qquad
 R_h=\frac{v(uX-Z)-s k^2W}{k}.
\]
Every other retained positive square term is an exact positive square vanishing
at the full actual contact. For positive bounded caps $m_*,\alpha_*$
and $0\le n_*\le n$, set
\[
 E_X=X-xvW,\quad E_Z=Z-(ux-s)vW,\quad
 E_r=uW-Y-\rho k^2W,
 \qquad Q_N=uE_X-E_Z.
\]
The full inverse is polynomial. All terms of $G$, including its
complete first-normal rows and quadratic normal block, are retained.
There are no higher state-normal terms.

Put $w=u-a$, $\xi=v^2$, and use the additional separated roots
\[
 S=vwW,\qquad T=kwW,\qquad R=(v^2-k^2)W/k.
\]
The exact identities
\[
 kR=\frac{kR_h-vQ_N}{s},\qquad
 RT=\frac{R_hT-SQ_N}{s},\qquad
 S^2-T^2=\frac w s(R_hT-SQ_N)                      \tag{UI1}
\]
use only the bounded pure-source inverse $1/s$.

\paragraph{Complete scalar and normal lifting.}
On the graph $E_X=E_Z=E_r=0$, write $G=\phi(u,\xi)W^2$.
Its bounded polynomial coefficients and the original parity imply
$\phi\in(k^2,\xi)$. The exact contact identities give
$\phi=\phi_u=\phi_\xi=0$ at $(u,\xi)=(a,k^2)$.
Successive monic Taylor division gives
\[
 \phi=w^2(k^2A_0+\xi A_1)
             +w(\xi-k^2)B+(\xi-k^2)^2C_0,          \tag{UI2}
\]
with bounded coefficients. To see the asserted first factor directly,
divide first in $u-a$; the second $u$ coquotient stays in
$(k^2,\xi)$. The constant and first $u$ terms are then divided in
$\xi-k^2$. No source infinitesimal is inverted by these operations.
Equations (UI1) allocate the three terms of (UI2) on
$(S,T,R_h,E_X,E_Z)$.

The full first $E_X,E_Z$ coefficients are $vL_X(u,\xi)$ and
$vL_Z(u,\xi)$. Exact contact gives $L_i(a,k^2)=0$, so they are
\[
 v\{wL_{i0}+(\xi-k^2)L_{i1}\}.
\]
The first summand is allocated on $S$, and the second uses
$vkR=v(kR_h-vQ_N)/s$. The entire $E_r$ coefficient belongs to
$(k^2,\xi)$ and vanishes at contact. The same division writes it as
\[
 w(k^2L_{r0}+\xi L_{r1})+(\xi-k^2)L_{r2}.
\]
It consequently uses $kT,vS,kR$. All these $E_r$ couplings have
zero residue. Thus (UI1)--(UI2) produce an exact bounded Gram on
\[
             (E_r,E_X,E_Z,R_h,S,T).                 \tag{UI3}
\]
This is a full-parent lift, not a restriction to a state kernel.

\paragraph{The strict matched Hessian.}
First take finite positive residue source weights. The leading model
before setting the ordinary test angle equal to $a$ is (RC4), with
every testing occurrence of $a$ replaced by $u$; its source
coefficients remain those of (RC3). Let $\mathsf G$ denote one half
of its angular Hessian after eliminating its two state directions,
in variables $(V,u)$ at the exact contact. Direct differentiation and
two-by-two Schur elimination give
\[
 \mathsf G_{uu}=
 \frac{8(1-r)(1+zr)^2}{zb\Delta}>0,                 \tag{UI4}
\]
\[
 \det\mathsf G=
 \frac{16(b+r)^2(1+zr)^3
                \{3z+4-(z+2)^2r\}}{b^2\Delta^2}.  \tag{UI5}
\]
The other two entries, making boundedness explicit, are
\begin{align*}
 \mathsf G_{Vu}&=
 \frac{4(1+zr)^2\{zb+zr+2r-2\}}{ab\Delta},\\
 \mathsf G_{VV}&=
 \frac{8(1+zr)\{z^2b^2+z^2br+zb^2+2zbr-zb-r+1\}}
      {b\Delta}.
\end{align*}
For a short check of strictness, the exact leading $n$ numerator is
$bQ_n$, where
\[
 Q_n=1+r(z-1)-z(z+2)r^2-z(z+1)br
     =(1-(z+1)r)(1+zr)-zr g.
\]
Since $Q_n\ge0$ and $g\ge0$, one has $r\le1/(z+1)$.
Consequently
\[
 3z+4-(z+2)^2r\ \ge\ \frac{z(2z+3)}{z+1}>0.       \tag{UI6}
\]
All denominators in (UI4)--(UI6) are positive bounded units.
These formulas therefore give a strict, uniformly bounded-away-from-zero
angular Schur form on the source residue domain, including $n=0$.
They contain no inverse of $r$, $g$, or $1-zb$.

\paragraph{Large weights and finite real caps.}
We use the elementary finite-dimensional penalty lemma: if a real
symmetric form $B$ is strictly positive on the common kernel of the
rows of a matrix $V$, then $B+M V^TV$ is positive definite for some
finite real $M>0$. A proof is to split into $\ker V$ and its
orthogonal complement. The first diagonal block is strict; taking
$M$ sufficiently large makes its Schur complement strict too.

The bounded finite part obtained by the centered positive square terms is
continuous at all source boundary residues. This assertion is
algebraic: its only large-weight products are
\[
 nx,\ nx^2,\ m\mu,\ m\mu^2,\ \alpha\pi,\ \alpha\pi^2,
 \quad C\rho k^2,\ D\sigma k^2,\ k^2H,
\]
and their bounded cap analogues. Substitution of the original source
formulas cancels the possible $r,g,1-zh,t,j$ poles in those products;
their denominators are the displayed positive units. For example,
\[
 m\mu=\frac{(1+k^2)(1-r)A_cS}{2b\Delta},\qquad
 \alpha\pi=-\frac{z^2(b+r)N}{2ab\Delta}.
\]
The other cancellations are
\[
 nx=-\frac{(b+r)A_cS-r\Delta}{2b\Delta},\quad
 nx^2=x(nx),\quad m\mu^2=\mu(m\mu),\quad
 \alpha\pi^2=\pi(\alpha\pi),
\]
\[
 C\rho k^2=\frac{c_1}{ab},\qquad
 D\sigma k^2=-\frac{c_2}{ab},\qquad
 c_1=\frac{zA_cN}{2\Delta},\quad
 c_2=\frac{z(1+z)(b+r)(b-k^2r)S}{2\Delta}.
\]
Finally $C\rho+D\sigma=(Ct-Dj)/(ab)$, using the exact canceled
formula in the full-parent reduction. Thus even this difference has
no hidden small-opening denominator.
For $k^2H$, write $k^2C=c_1/t$, $k^2D=c_2/j$, with
$c_1,c_2$ positive bounded units. Then
$k^2H=c_1c_2/(c_1j+c_2t)$, whose denominator is a unit because
$t+j=A_c$ is a unit. Its residue is exactly $H_0$ in (RC3).

At most two of $n,m,\alpha$ can diverge. Indeed $m$ can diverge
only at $g=0$, while $\alpha$ can diverge only at $1-zb=0$;
these two equalities would give $g=1/z+r>0$. The hard-state parts
of the three possible diverging centered rows are, up to nonzero
unit factors,
\[
                   E_X,\qquad E_X+aE_Z,\qquad E_Z.
\]
Every possible pair is independent with a unit determinant.

Here is why their constrained limiting finite form is strict.
Approach the boundary by the positive leading-source domain. The
state block is bounded below by a positive real form because $m$ and
$\alpha$ are strong. Its minimizing hard vector for a fixed angular
vector is bounded: solve the at most two diverging positive square term equations
in the hard variables, using the unit determinants just displayed;
this supplies a bounded competitor, and all remaining coefficients
stay bounded. If $h_0$ is that competitor and $h_{\min}$ the state
minimizer for angular vector $q$, the exact identity
\[
 (h_0-h_{\min})^TP(h_0-h_{\min})
       =F(h_0,q)-q^T\mathsf Gq
\]
and $P\succeq cI$ make boundedness of $h_{\min}$ explicit.
Together with (UI4)--(UI6), the state Schur factorization
then bounds the full matched Hessian below by a positive real multiple
of the Euclidean norm. For a vector in the common limiting positive square term
kernel, adjust only its hard coordinates to lie in the exact nearby
positive square term kernels. The adjustment tends to zero and annihilates the
large weights exactly. Taking the limit proves strict positivity of
the bounded finite part on that common kernel. The penalty lemma now
chooses finite real caps for every diverging weight. Bounded source
weights can be retained in full. All caps are below their original
weights, so the remaining centered excesses are positive squares.

This limiting argument uses only real finite-dimensional forms to
choose caps. No inverse of a tensor matrix and no mixed root are
used in the full-parent identity. The boundary residues are limits
of the positive leading-source domain: $g$ and $1-zb$ specify the
interval for $b$, and decreasing a positive $r$ if necessary makes
$Q_n\ge0$ strict. Thus the argument also covers a literal original
$n=0$ and all combinations of the allowed boundary residues.

\paragraph{Separate angular roots and full strictness.}
The residue of the hard block in (UI3) is
\[
 2E_r^2+4n_*(1+a^2)E_X^2
 +(4m_*+2)(E_X+aE_Z)^2+4\alpha_*(1+a^2)E_Z^2
 +4\gamma R_h^2,
\]
which is strict. In the residue Gram, $S$ couples only to the two
ordinary hard normals, and $T$ only to $R_h$; the $E_r$ couplings
vanish. Setting $S=T$ gives the matched leading Hessian under an
invertible bounded change of coordinates. Indeed at $V=1$ the
ordinary normals are its two state deviations, $R_h$ supplies the
relative-angle deviation with nonzero coefficient $2s$, and
$S=T$ supplies the ordinary-angle deviation. The chosen caps therefore
make the sum of the two angular Schur diagonals strictly positive.

Add an appropriate real multiple of the last zero identity in (UI1).
It redistributes mass between $S^2$ and $T^2$; all accompanying
coefficients contain $w$ and have zero residue. Since their Schur
diagonal sum is positive, choose this multiple so that both Schur
diagonals are positive. The entire six-root residue Gram is now
positive definite. The strict-residue Gram lemma proves the exact
bounded remainder plus relative positive square term is tensor SOS.

Restore the centered excesses and the other C/D squares. The construction
is homogeneous in $W$ with a polynomial state inverse, so it includes
$W=0$ and every state. Real and imaginary parts give complex states.
The complementary angle charts are proved next.

\subsection{Complementary angle charts}

\label{p6b22e944bd7c:sec:ui-all-tests-recovered}

\begin{theorem}[Strong C2 UI, all independent tests]
Use the original C2 source (UR1)--(UR3), with $d=1$, actual source
angles $u_c=a>0$ a bounded unit and $v_c=k>0$ infinitesimal, and
$C,\alpha$ strong. Then the entire parent (UR6) is a sum of squares
in the tensor product of the independent source and testing fields,
for every testing angle, including projective endpoints, and every
homogeneous state. The statement retains the literal boundary $n=0$
and arbitrary positive rates of every source opening.
\end{theorem}

Section~\ref{p76545700d90b:sec:local-ui-certificate} covers $v$ infinitesimal
and $u$ with the same residue as $a$. Its complement follows from a
rank-one allocation and positive frames at the testing endpoints.

\paragraph{A rank-one allocation lemma.}
Let $P$ be a positive definite real symmetric matrix, let $\ell,b$ be
real columns, and let $H,c_0>0$. Set
\[
 d=\ell^TP^{-1}\ell,\qquad e=\ell^TP^{-1}b,\qquad
 h=c_2-b^TP^{-1}b.
\]
Suppose
\[
 p(T)=c_0+(h+4Hdc_0)T+4H(dh+e^2)T^2>0\quad(T\ge0),
 \qquad dh+e^2>0.                                      \tag{UA1}
\]
There exists a real $\lambda$ such that
\[
 |\lambda|<\sqrt{4Hc_0},\qquad
 h+2e\lambda-d\lambda^2>0.                             \tag{UA2}
\]
Indeed put $s=\sqrt{4Hc_0}$. If $d>0$ and $|e|<ds$, take
$\lambda=e/d$. If $|e|\ge ds$, the value at the appropriate endpoint
is $h+2|e|s-ds^2$. Were this nonpositive, writing
$x=|e|-ds\ge0$ would give
$h+ds^2\le-2sx$ and $dh+e^2\le x^2$; the quadratic in (UA1)
would then have a nonpositive value at a positive argument. Thus a
nearby interior $\lambda$ works. This includes equality $|e|=ds$.
The scalar interval argument also holds in a limiting case $d=0$;
then $e\ne0$ follows from $dh+e^2>0$. (For a finite positive $P$
itself, $d=0$ would force $\ell=e=0$, so this case occurs only
when using the continuous scalar extension at a boundary.)

Consequently the two real matrices
\[
 \begin{pmatrix}P&b-\lambda\ell\\(b-\lambda\ell)^T&c_2\end{pmatrix},
 \qquad
 \begin{pmatrix}4H&\lambda\\\lambda&c_0\end{pmatrix}       \tag{UA3}
\]
are positive definite. This is an elementary real lemma. Its inverse
and square root are never taken in the tensor ring.

\paragraph{The finite ordinary angle away from contact.}
Assume $v$ infinitesimal and $u$ a bounded unit whose real residue has
square different from $a_0^2$. Keep the centered positive square terms
$P_L,P_C,P_n,P_m,P_\alpha$ of (UR12), (UR14) outside the Gram, and
combine $P_H$ with the exact remainder. Use the six separated roots
\[
 r_0=uW-Y,\quad X,\quad Z,\quad V_0=vW,\quad K_0=kW,
 \quad R_0=\frac{v(uX-Z)}{k}.                           \tag{UA4}
\]
They satisfy the exact zero relation
\[
 R_0K_0-(uX-Z)V_0=0.                                  \tag{UA5}
\]
All divisions in (UA4) are pure source divisions.

For completeness the boundedness of the $C,D$ remainder in this frame
is explicit. Write $\ell_0=L\tau$, $c_R=C\rho k^2$, and
$Q=uX-Z$. Then (UR13) gives exactly
\begin{align*}
 G_{CD}+P_H={}&4\gamma R_0^2
 +8\ell_0kK_0r_0+8\ell_0qk^2V_0Q-4\ell_0\tau k^2K_0^2\\
 &+8c_RvV_0r_0-4c_R\rho k^2V_0^2.                  \tag{UA6}
\end{align*}
All displayed coefficients are bounded by (UR9)--(UR11); all terms
except the first have zero residue. The $A$ positive square term is exactly
$4(A/k^2)D_uD_v(uK_0-kr_0)^2$. Substituting
$Y=uW-r_0$ in (UR15) and the signed core gives a bounded Gram on
(UA4), with no discarded normal term. Its residue has an isolated
$2r_0^2$ term. First suppose the leading source weights are finite.
The remaining residue is
\[
 q_s^TPq_s+2V_0b_*^Tq_s+c_2V_0^2+4H R_0^2+c_0K_0^2,
 \quad q_s=(X,Z)^T,                                  \tag{UA7}
\]
where $D_u=1+u^2$, $H=H_0$, $A_w=A_0^{\rm weak}$ are (RC3), and
\[
\begin{gathered}
 P=\begin{pmatrix}
 4nD_u+4m+2&(4m+2)u\\
 (4m+2)u&(4m+2)u^2+4\alpha D_u
 \end{pmatrix},\quad \ell=(u,-1)^T,\\
 b_*=\binom{-4mD_u}{-D_uu(4m+4\alpha+2)},\quad
 c_0=4A_wD_uu^2,\quad
 c_2=4mD_u^2+(4\alpha+2)D_uu^2.                      
\end{gathered}
\tag{UA8}
\]
The matrix $P$ is strictly positive because $m,\alpha$ are strong.
The last sign-changing part of (UA7) is handled by adding
$2\lambda$ times (UA5). This gives precisely the two blocks (UA3).

Here is a direct proof of (UA1) for every ordinary angle away from
contact. Put $N_*=2n+1$, $T_*=2m+2\alpha+1$, $U=u^2$, and
$T=V^2$. The leading state polynomial (RC4), with its testing $a$
replaced by $u$ but its source unchanged, has determinant
\[
 \det f=8D_u^2u^2V^2\,\mathcal L(U,T),                \tag{UA9}
\]
where
\begin{align*}
\mathcal L={}&\frac{4\alpha mN_*}{U}
 +4A_wHT_*U+\frac{4A_w\alpha(2m+N_*)}{T}
 +\frac{4A_wnT_*U}{T}\\
 &+2HN_*T_*T+\frac{4HmN_*T}{U}
 +4A_wH(2m+N_*)-N_*T_*.
\end{align*}
After division by $N_*T_*$ all possible opening poles cancel. Define
\[
 i_N=N_*^{-1},\quad t_m=m/T_*,\quad t_{am}=\alpha m/T_*,
 \quad q_\alpha=\frac{\alpha(2m+N_*)}{N_*T_*},\quad i_\alpha=1/\alpha.
\]
Writing $z=a^2$, $A_c=1+zr$, and using the leading $\Delta$ of (RC3),
these quantities are exactly
\[
\begin{aligned}
 i_N&=\frac{r\Delta}{(b+r)A_c^2},&
 t_m&=\frac{(1-r)(1-zb)}{2b(b+r)(z+1)},\\
 t_{am}&=\frac{z^2(b+r)(1-r)}{4\Delta},&
 q_\alpha&=\frac{bz^2\{(z+1)b+(z+2)r-1\}}{2A_c^2},\\
 i_\alpha&=\frac{2(1-zb)\Delta}{z^2(z+1)b(b+r)^2}.
\end{aligned}\tag{UA10}
\]
All denominators here are units. The normalized Laurent polynomial is
\[
 \frac{\mathcal L}{N_*T_*}
 =\frac{4t_{am}}U+4A_wHi_NU+\frac{4A_wq_\alpha}T
  +2A_w(1-i_N)\frac UT+2HT+4Ht_m\frac TU
  +4A_wHq_\alpha i_\alpha-1.                         \tag{UA11}
\]
Its value and its two first derivatives vanish at $(U,T)=(z,1)$,
as direct substitution verifies. In logarithmic coordinates the
nonconstant terms have nonnegative coefficients. The coefficients of
$U^{-1}$ and $T$ are positive units. If $i_N>0$, the coefficients of
$U$ and $T^{-1}$ are positive; if $i_N=0$, that of $U/T$ is positive.
The function is therefore strictly convex in logarithmic coordinates
and has its unique minimum at $(z,1)$. This proves strict positivity
at fixed $U\ne z$ for every $T>0$. The constant and leading
coefficients of the equivalent quadratic (UA1) are strictly positive:
$c_0>0$, and the positive $T$ coefficient in (UA11) gives the second
assertion. Thus (UA2)--(UA3) prove the required strict residue Gram.

\paragraph{Opening boundaries and divergent source weights.}
The preceding proof extends to all allowed opening residues, rather
than assuming $n,m,\alpha$ bounded. The following formulas make the
needed uniformity checkable. Put
\[
 J_0=q_\alpha+\frac{u^2(1-i_N)}2.
\]
Then $J_0$ is a positive unit, and direct two-by-two inversion gives
\[
\begin{aligned}
 \frac{\det P}{N_*T_*}&=8D_uJ_0,\\
 d&=\frac{u^2i_N+q_\alpha i_\alpha}{4J_0},&
 e&=\frac{uD_u}{2J_0},&
 h&=\frac{D_u(4t_{am}-u^2)}{J_0},\\
 (P^{-1}b_*)_1&=\frac{i_N(u^2-4t_{am})}{2J_0},&
 (P^{-1}b_*)_2&=-\frac{u\{4t_{am}i_N+(1-i_N)u^2+1\}}{2J_0}.
\end{aligned}\tag{UA12}
\]
To check the unit assertion, if $i_N=0$ then $J_0=u^2/2>0$.
If $i_N>0$, the weight $N_*$ is finite and the strong lower bounds
on $m,\alpha$ give $q_\alpha>0$: this follows directly from
$\alpha(2m+N_*)/[N_*(2m+2\alpha+1)]$, whether both weights are
finite or one diverges. Thus (UA12) has only unit denominators on
the entire source residue domain.

The strictly positive leading coefficient is also explicit:
\[
 dh+e^2=\frac{D_u(u^2+2t_m)}{2J_0}>0.
\]
Equations (UA10)--(UA12) prove continuity of $d,e,h$ and of the
minimizing hard vector. The convexity proof for (UA11) remains strict
at every allowed boundary. Choose $\lambda$ from the strict interval
(UA2) at that boundary; the same strict inequalities hold nearby.
Moreover $P$ is bounded below by a positive real matrix, while
$P^{-1}b_*$ and $P^{-1}\ell$ are bounded. Its Schur factorization
therefore bounds the first block in (UA3) below by a positive real
multiple of the identity, uniformly in the nearby source. The second
block is a strictly positive real matrix as well.

The finite leading sources used in this limiting argument are valid
real interior sources. This point must be established separately:
the actual source has $b=h+(k^2/z)t$, and it is not legitimate to
substitute that actual $b$ into a leading formula with denominator
$1-zb$ when this denominator has an infinitesimal sign change.
For the limiting real parameters, put
\[
 f(b,r)=1+r(z-1)-brz(z+1)-r^2z(z+2).
\]
Their domain is $b>0$, $0\le r<1$, $g=(1+z)b-1+r\ge0$,
$1-zb\ge0$, and $f(b,r)\ge0$. For $0<s<1$, let
\[
 b_s=sb+(1-s)/(1+z),\qquad r_s=sr.
\]
Direct expansion gives
\[
 g_s=sg,\qquad 1-zb_s>0,\qquad
 f(b_s,r_s)=(1-s)(1+s-sr)+s^2f(b,r)>0.
\]
Increase $b_s$ by a sufficiently small $\epsilon>0$. This makes
$g_s$ strict, decreases the last two strict margins by respectively
$z\epsilon$ and $r_sz(z+1)\epsilon$, and therefore preserves them.
If $r=0$, then increase $r_s$ by a sufficiently small positive
amount; $g$ increases and $f$ remains positive by continuity.
Letting these perturbations tend to zero as $s\to1$ supplies valid
finite real leading sources approaching every allowed residue.
Their $t$ may be chosen as $(1+zr)/2$, since the leading harmonic
weight is independent of $t$. The normalized leading identity
$2rn\Delta=bf(b,r)$ verifies the nonnegative-$n$ requirement.
The same construction supplies the real-interior approximation
used in the finite-cap step of the local UI certificate.

At most two of $n,m,\alpha$ diverge: $m$ can diverge only at $g=0$,
and $\alpha$ only at $1-zb=0$, which cannot occur together. Their
centered excess rows have hard-state parts proportional to
\[
 X,\qquad X+uZ,\qquad Z.
\]
Any possible pair is independent because $u$ is a unit. For a vector
in their common limiting kernel, change only $(X,Z)$ by a vanishing
amount to annihilate the corresponding nearby exact rows. Their
large penalties then vanish exactly. Taking the limit in the uniform
strict lower bound proves that the bounded finite part is strictly
positive on that limiting kernel. The elementary finite-dimensional
penalty lemma used in the local theorem now supplies finite real caps
$n_*,m_*,\alpha_*$. Their centered excesses are precisely the retained
positive square terms (UR14). The exact bounded remainder with those caps, after
adding (UA5), has a strictly positive residue Gram on (UA4).
This proves the finite ordinary-angle complement for all source
openings and literal $n=0$.

\paragraph{Endpoint frames.}
Choose positive real $\mu,m_*,\alpha_*$ below, respectively,
$k^2C,k^2D$, $m$, and $\alpha$. The products $k^2C,k^2D$ are strong
by (UR5), and $A_w=A/k^2$ is a positive unit. Retain $n$ and all
positive excesses in full. The following formulas concern the parent
with $C=D=\mu/k^2$, $m=m_*$, $\alpha=\alpha_*$ and $A=A_wk^2$.
They therefore suffice. The root coordinates need not themselves be
bounded; only their displayed Gram coefficients must be bounded.

At $u,v$ infinitesimal use
\[
 E=(uW-Y)/k,\quad R=v(uX-Z)/k,\quad V=vW,\quad K=kuW,
 \quad X,Z.
\]
The exact amplitudes are
\[
 J=k(E+R),\quad M=D_uV-X-uZ-uvkE,\quad
 vY-Z=uV-Z-vkE,\quad kY=K-k^2E.
\]
Their full residue form is
\[
 4\mu\{E^2+(E+R)^2\}+4m_*(V-X)^2+4\alpha_*Z^2
       +4\overline{A_w}K^2+2X^2,                    \tag{UA13}
\]
which is strict. Literal zero testing angles are included without
any division by a testing angle.

For bounded $u$ and finite unit $v$, put
$E=(uW-Y)/k$, $F=J/k$. Then
\[
 Y=uW-kE,\quad Z=uX+k(E-F)/v.
\]
The roots $(E,F,W,X)$ have strict residue
\[
 4\mu D_vE^2+4\mu F^2+
 \{4m_*D_u^2+(4\alpha_*+2)D_uu^2\}(vW-X)^2
 +2D_vD_uX^2.                                      \tag{UA14}
\]
For bounded $u$ and $v=\infty$ write $q=1/v$ and homogenize by $q^2$.
The same roots are defined using the homogeneous $J$, and
\[
 Y=uW-kE,\qquad Z=uX+kqE-kF.
\]
The strict residue is (UA14) with $D_v=1$ and $(vW-X)^2$ replaced
by $W^2$. This includes the literal endpoint $q=0$.

For $u=\infty$, put $p=1/u$, $D_p=1+p^2$, and homogenize by $p^2$.
Write $r_p=W-pY$, $Q_p=X-pZ$, and $P=vY-Z$. Then
\[
 J=r_p+vQ_p,\qquad M=D_pP-pQ_p+pv r_p.
\]
For $v$ infinitesimal use the six roots
\[
 E=r_p/k,\quad F=vQ_p/k,\quad P,\quad R=pQ_p,
 \quad K=kY,\quad S=pZ.
\]
The signed core is exactly
\[
 2k^2(E+F)^2+2D_pP^2+2D_vR^2-4D_pD_vRP
    +4D_pD_vpFK+2D_pD_vS^2.
\]
Its cross-term reduction uses $RvY=pFK$, an exact identity.
Adding every capped positive square term gives the strict residue
\[
 4\mu\{E^2+(E+F)^2\}+(4m_*+2)(P-R)^2
 +4\alpha_*P^2+4\overline{A_w}K^2+2S^2.             \tag{UA15}
\]
For finite unit $v$, instead use $E=r_p/k$, $F=J/k$, $P,K,S$.
Since $Q_p=k(F-E)/v$, the strict residue is
\[
 4\mu D_vE^2+4\mu F^2+(4m_*+4\alpha_*+2)P^2
 +4\overline{A_w}D_vK^2+2D_vS^2.                   \tag{UA16}
\]
Finally, at both infinite endpoints put $q=1/v$ and homogenize by
$p^2q^2$. Use
\[
 E=(W-pY)/k,\quad F=(qW+X-pqY-pZ)/k,\quad
 P=Y-qZ,\quad K=kY,\quad S=pZ.
\]
Here $M=D_pP+pkD_qE-pqkF$ and the entire signed core is
\[
 2k^2F^2+2D_pP^2+2D_qp^2k^2(F-qE)^2
 +4D_qD_pk(F-qE)S+2D_qD_pS^2.
\]
The strict residue is (UA16) with $D_v=1$.

\paragraph{Exhaustion and restoration.}
A testing angle is infinitesimal, a bounded unit, or projectively
infinite. Equations (UA13)--(UA16) cover every pair except ordinary
$u$ with infinitesimal $v$. The local theorem covers the equal-residue
contact, and (UA4)--(UA12) cover its complement. Exact sign gauges are
\[
 (u,W,X)\mapsto(-u,-W,-X),\qquad
 (v,W,Y)\mapsto(-v,-W,-Y);
\]
all unspecified variables are fixed. Both preserve the full parent,
so negative ordinary contact residues are included. All endpoint
frames have been written after homogenization, including literal
projective endpoints. Each proof has an exact bounded Gram with
strict real residue; the strict-residue Gram lemma gives tensor SOS.
Restoring all retained positive square terms completes the original parent.
Real and imaginary parts treat complex states. This closes the
entire strong C2 UI case and its exact symmetry image HU.

\section{An unbounded first angle and a small second angle}
\subsection{Source entry at the HI corner}
\label{component:p54b8f20ad2dd}

Use the original positive rational C2 source with $d=1$, actual
$u_c$ unbounded and $v_c$ infinitesimal. Put
\[
 q=u_c^{-2},\quad y=v_c^2,\quad H=u_c^2h\in(0,1),
 \quad R=r/q,\quad B=b/q,\quad L=yt=B-H,\quad T=B-yR.
\]
The already proved source identities (WE1)--(WE3) show that $R$ is
bounded using only $n\geq0$, $q$ infinitesimal and $y$ bounded.
Here is their short contradiction argument, to make this use
explicit. The positive source gaps imply
\[
 0<T<H+y,\quad g=H+(1+y)r-1+qT>0,
\]
and $n=-qP/(2rQ)$ with $Q>0$ and
\[
 P\ge r^2(1+y)(H+yr)-Tqr(y+2)-q^2(T^2+T).
\]
If $q/r$ were infinitesimal, then $T/(H+yr)$ and
$(T^2+T)/(H+yr)$ would be bounded. For unit $r$ this follows
from $T<H+y\le(H+yr)/r$; for infinitesimal $r$, the inequality
$g>0$ forces $H$ to have residue one. Division by $r^2(H+yr)$
would give $P>0$, contrary to $n\geq0$. Thus $R$ is bounded.
(The exact positive remainder in this lower bound is printed in
(WE3), so this argument uses no asymptotic sign substitution.)

Now $r=qR$ is infinitesimal and $g>0$ gives $H_0=1$.
Define the positive opening
\[
 E=(1-H)/q.
\]
The same gap implies
$E<(1+y)R+T$, so $E$ is bounded. Since $j=1+R-t>0$,
$t$ is bounded; consequently
\[
 H=1-qE,\quad B=H+yt\sim1,\quad
 j=1+R-t,\quad K=B+R-E=g/q>0.
\]
All of $E,t,j,K$ are positive bounded, while
$t+j=1+R$ and $E+K=B+R$ are positive units.

Let $\Delta$ and $N$ denote the unchanged original source
polynomials, and set $\mathcal N=N/q$. Their exact expressions are
\begin{align*}
\Delta={}&(1+q)\{B^2+2BR+(1+y)B\}
 +(1-yq)R^2+2(1+y)R+1+y,\\
\mathcal N={}&(1+q)(B+R)(H+y)-(1+y)(1-qR)yt.
\end{align*}
Thus $\Delta$ is a positive bounded unit and
\[
 \bar\Delta=(1+\bar R)(3+\bar R),\qquad
 \overline{\mathcal N}=1+\bar R>0.
\]
Set $S=1+R+yt$, also a positive unit. Direct substitution into
the original coefficients gives the following exact equalities:
\[
\begin{aligned}
 mqK&=\frac{(1+y)(1-qR)(1+R)S}{2\Delta},&
 Cyt&=\frac{(1+R)\mathcal N}{2\Delta},\\
 \frac{A}{qy}&=\frac{(1+R)(B+R)S}{2H\Delta},&
 Dyj&=\frac{(1+q)(B+R)(B-yR)S}{2\Delta},\\
 \alpha qE&=\frac{(B+R)\mathcal N}{2\Delta},&
 n&=\frac{(B+R)(1+R)S-R\Delta}{2R\Delta}.
\end{aligned}                                                   \tag{HI-entry}
\]
Every one of the first five right-hand sides is a positive bounded
unit; in particular $B-yR$ has residue one. The last numerator has
residue $(1+\bar R)(1-\bar R)$, so $n\geq0$ implies
$0\le\bar R\le1$. This does not imply $R\le1$ as a source
inequality, and no such stronger claim is used.

Since $q,y$ are infinitesimal and $E,t,j,K$ bounded, the positive
numbers $qK,yt,yj,qE$ are infinitesimal. Dividing their unit
products in (HI-entry) proves
\[
 m,C,D,\alpha\quad\hbox{are all unbounded},\qquad
 A\asymp qy.
\]
This conclusion follows from the original source formulas and
$n\geq0$; no additional strong-weight assumption is needed for
an actual HI source.

\subsection{The HI contact for every $n\ge0$}
\label{component:p888d9f79ba3b}
\label{component:p4363feef60dd}
\label{p888d9f79ba3b:sec:r7-hi-n0-retained}
\label{p4363feef60dd:sec:r7-hi-full-retained}
Use the HI source estimates of Section~\ref{component:p54b8f20ad2dd}.
Let $q=u_c^{-2}$ and $y=v_c^2$ be positive infinitesimals, and put
$p=1/u$ for the testing angle. The full target here is
$\widehat F=p^2F(1/p,v)$, with its polynomial value at $p=0$.
Here the independent testing $p,v$ are infinitesimal or zero.
The source numbers $R=r/q$, $H$, and the opening quantities $t,j,K,E$
are those of the source calculation; $R$ is bounded, $H\sim1$, and
\[
 S_c=1+R+yt=j+t(1+y)=K+E(1+q).
\]
Set $d_p=1+p^2$, $d_v=1+v^2$, $\delta_p=p^2-q$,
$\delta_v=v^2-y$, and use the bounded positive source products
\[
 c=Cyt,\quad d=Dyj,\quad \mu=mqK,\quad a_*=\alpha qE.
\]
The original $n$ is included from the start. Define
\[
 \nu=2nR={(B_{\rm src}+R)(1+R)S_c-R\Delta\over\Delta},
 \qquad \zeta=R+\nu,\quad h_*=H-\nu.
\]
The exact source identities are
\[
 {2AH\over qy}=\zeta,\qquad
 c-d={y\over2}\{\zeta(1+q)-t\},\qquad
 \mu-a_*={q\over2}\{E-\zeta(1+y)\}.
\]
They cancel the apparent opening denominators before any normal is
estimated. At $R_0=r_0\in[0,1]$ they give
\[
 \nu_0={1-r_0\over3+r_0},\qquad
 c_0=d_0=\mu_0=(a_*)_0={1+r_0\over2(3+r_0)}.
\]
Thus these products remain units at $r_0=0$, even when $n$ is
unbounded. Literal $n=0$ is retained in every formula.

Define the state coordinates
\[
 A_0=W-pY,\quad B_0=vY-Z,\quad
 S=X-vW-pZ+pvY,\quad s=-S/S_c,\quad Q=\sqrt{qy}\,Y/H,
\]
\[
 P=p/\sqrt q,\quad V=v/\sqrt y,\quad T=PV,
 \quad N_c=(A_0-tv s)/\sqrt y,\quad
 N_m=(B_0+Ep s)/\sqrt q,
\]
\[
 N_D=d_vN_c+t\delta_vVs,\qquad
 N_M=d_pN_m-E\delta_pPs.
\]
The exact source has $P=V=1$, $s=Q=\sqrt q$ and all four $N$'s zero.
Retain the positive squares
\[
 \mathcal E={4cj d_v\over tS_c}N_c^2
 +{4dt(1+y)\over jS_c}N_D^2
 +{4\mu E(1+q)\over KS_c}N_M^2
 +{4a_*K d_p\over ES_c}N_m^2
\]
and $\mathcal N=4nd_pd_vx^2$, where
\[
 x=X+(R/H)pvY=-g s+(H+R)TQ+v\sqrt yN_c-p\sqrt qN_m,
 \qquad g=S_c-tv^2-Ep^2.
\]
The complete difference of the original $n$ term and this square is
\[
 \begin{split}
 4nd_pd_v(X^2-x^2)={}&4\nu d_pd_vg\,TsQ
 -2\nu(2H+R)d_pd_vT^2Q^2\\
 &-4\nu d_pd_vTQ(v\sqrt yN_c-p\sqrt qN_m).
 \end{split}
\]
All its coefficients are bounded, regardless of the size of $n$.

The exact quadratic-state expansion is
\[
 \begin{split}
 \widehat F-\mathcal E-\mathcal N={}&h_cN_c^2+h_mN_m^2
 -4\sqrt{qy}\,pv d_pd_vN_cN_m\\
 &+N_c(C_sVs+\beta_cTQ)+N_m(M_sPs+\beta_mTQ)+\mathcal G,
 \end{split}
\]
where
\[
 h_c={4d_v\over S_c}\{c(1+y)+dd_v\}+2yd_v(1+v^2+p^2v^2),
\quad
 h_m={4d_p\over S_c}\{\mu d_p+a_*(1+q)\}+2qd_p(1+p^2+p^2v^2),
\]
\[
 C_s=4d_v[y\{\zeta(1+q)-d_pg\}+2dt\delta_v/S_c],
 \quad M_s=4d_p[q\{d_vg-\zeta(1+y)\}-2\mu E\delta_p/S_c],
\]
\[
 \beta_c=4h_*\sqrt y\,v d_pd_v,
 \qquad\beta_m=-4h_*\sqrt q\,p d_pd_v,
\]
and
\[
 \mathcal G=f_AP^2s^2+f_BV^2s^2+f_CT^2Q^2+f_DQ^2-eTsQ,
\]
\[
 \begin{split}
 f_A={}&4[\mu K+a_*Ed_p-2\mu E\delta_p
                   +\mu E^2\delta_p^2/S_c]
       +2qd_v(E\delta_p-K+tv^2)^2-2E^2qv^2d_p,\\
 f_B={}&4[ctd_v+dj-2dt\delta_v+dt^2\delta_v^2/S_c]
                   +2y(t\delta_v-j)^2,\\
 f_C={}&2\{H^2-\nu(2H+R)\}d_pd_v,\quad
 f_D=2H(R+\nu)d_pd_v,\quad e=4h_*d_pd_vg.
 \end{split}
\]
This identity follows by the exact quadratic-state extraction rule:
substitute the displayed inverse coordinates, subtract the five
retained squares, collect $N_c,N_m$, and use the three source identities
above. It specifies the full scalar, both complete linear normals and
the entire hard block. Its coefficients are polynomials in $p,v$ with
bounded source coefficients. In particular $C_s,M_s,\beta_c,\beta_m$
are infinitesimal. The hard diagonal residues are $4/(3+r_0)>0$ and
\[
 (f_A)_0=(f_B)_0=(f_C)_0=(f_D)_0
 =\lambda(r_0)={2(1+r_0)^2\over3+r_0},\qquad e_0=4\lambda(r_0).
\]

For completeness the finite radical division is given here once. A
star means $p=\sqrt q$, $v=\sqrt y$ and $\sigma=\sum f_i-e$.
Because every subtracted square vanishes at the actual contact, the
original state and angular contact equations give exactly
\[
 f_A^*+f_B^*=f_C^*+f_D^*=e^*/2,
 \quad q^{1/4}\sigma_{\sqrt p}^*=-4(f_A^*-f_D^*),
 \quad y^{1/4}\sigma_{\sqrt v}^*=-4(f_B^*-f_D^*),
\]
\[
 C_s^*+\beta_c^*=M_s^*+\beta_m^*=0.
\]
If $sQ>0$, a pure-testing sign choice makes $s,Q>0$.
Put $(A_1,A_2,A_3,A_4)=(Ps,Vs,TQ,Q)$,
$G=(A_1A_2A_3A_4)^{1/4}$, $E_i=A_i-G$. Each radical separates into
pure-source and pure-testing radicals. The exact circuit is
\[
 L=\sum A_i^2-4G^2=\sum E_i^2+2G\sum E_i.
\]
For every pairing $ij|k\ell$,
\[
 G\sum E_i=G(\sqrt{A_i}-\sqrt{A_j})^2
 +G(\sqrt{A_k}-\sqrt{A_\ell})^2
 +2G\{(A_iA_j)^{1/4}-(A_kA_\ell)^{1/4}\}^2.
\]
Also $L=(A_1-A_3)^2+(A_2-A_4)^2
+2(\sqrt{A_1A_3}-\sqrt{A_2A_4})^2$, and for $rr'=G^2$,
$(r-r')^2=(r-G)^2+(r'-G)^2+2G(\sqrt r-\sqrt{r'})^2$.
A finite average supplies positive pivots simultaneously on the $E_i$
and on $F_p=\sqrt{A_1A_3}-G$, $F_v=\sqrt{A_2A_3}-G$.
The exact relations
\[
 G(\sqrt p-q^{1/4})=q^{1/4}F_p,
 \qquad G(\sqrt v-y^{1/4})=y^{1/4}F_v
\]
turn finite monic Taylor division in $\sqrt p-q^{1/4}$ and
$\sqrt v-y^{1/4}$ into bounded infinitesimal corrections to these
positive pivots. More explicitly put $z=\sqrt p$, $w=\sqrt v$,
$z_0=q^{1/4}$, $w_0=y^{1/4}$ and $k=(f_C^*+f_D^*)/2$.
Each monomial has the finite expansion
\[
 z^iw^j=\sum_{a=0}^i\sum_{b=0}^j
 {i\choose a}{j\choose b}z_0^{i-a}w_0^{j-b}
 (z-z_0)^a(w-w_0)^b.
\]
All its coefficients are bounded. In $\mathcal G-kL$, terms with two
differences multiplying $G^2$ become bounded infinitesimal multiples
of $F_p^2,F_pF_v,F_v^2$, because each replacement contributes a
source factor $z_0$ or $w_0$. Terms with one difference multiplying
$GE_i$ become infinitesimal multiples of $F_pE_i,F_vE_i$.
Terms multiplying $E_i^2$ have infinitesimal coefficients since
$f_i-k$ has zero residue. The only possible first-jet expression is
\[
 G\{2\sum_i(f_i^*-k)E_i+z_0\sigma_z^*F_p+w_0\sigma_w^*F_v\}.
\]
Using the displayed contact jets, it equals exactly
\[
 \begin{split}
 2(f_D^*-k)G\sum_iE_i
 &-{z_0\sigma_z^*\over2}
       \{\sqrt G(\sqrt{A_1}-\sqrt{A_3})\}^2\\
 &-{w_0\sigma_w^*\over2}
       \{\sqrt G(\sqrt{A_2}-\sqrt{A_3})\}^2.
 \end{split}
\]
All three source coefficients are infinitesimal, and the finite circuit
average already supplies strictly positive pivots on their roots.
For each normal, its two infinitesimal coefficients have zero sum at
contact. Write their two raw amplitudes as $G+E_i$; finite monic
division of that sum and the same two relations express the whole
normal coupling by infinitesimal entries on the already positive
roots. Thus the hard block and circuit block have strict residue.
For $sQ<0$ the raw diagonal and $\sqrt{T|sQ|}$ already have positive
residue coefficients $\lambda,\lambda,\lambda,\lambda,4\lambda$;
for $sQ=0$ omit the last root. The strict-Gram lemma proves all three
cases. Restore the five positive squares. Section~\ref{component:pcee1372511ed}
then proves the other measurements, including the literal axes and
infinities, for every $n\ge0$.

\subsection{Tests outside the HI critical corner}
\label{component:pcee1372511ed}

\label{pcee1372511ed:sec:hi-remote}
\begin{theorem}[All tests outside the HI critical corner]
Let $K_c,K_t$ be independent real closed extensions of $\mathbb R$.
Consider the original C2 parent with $d=1$, $\lambda=\beta=0$,
$n,A\geq0$, and $m,C,D,\alpha$ positive unbounded source numbers.
Then at every projective testing pair except possibly
$1/u$ and $v$ both nonzero infinitesimal, and every real or complex testing
state, the
original parent is a sum of squares in $K_c\otimes_{\mathbb R}K_t$.
No comparison of source and testing rates is made.
\end{theorem}

\begin{proof}
For any positive real $\mu$, retain the four nonnegative event
positive square terms with coefficients $m-\mu,C-\mu,D-\mu,\alpha-\mu$, and
retain the $n$ and $A$ positive square terms in their entirety. It suffices to
prove the capped parent obtained by setting
$m=C=D=\alpha=\mu$ and $n=A=0$ is tensor SOS at the given test.
The real constant $\mu$ may be chosen after the real testing
residue is fixed, since each of the original four source
coefficients is larger than every real number.

Write $D_u=1+u^2$, $D_v=1+v^2$.
First suppose $v$ is infinitesimal, including zero. Use the roots
\[
 S=vW,\qquad E=Y-uW,\qquad X,\quad Z.
\]
All four capped event amplitudes are bounded linear forms in
these roots:
\[
 M=D_uS-X-uZ+uvE,\quad uW-Y=-E,\quad
 J=-E+uvX-vZ,\quad vY-Z=uS+vE-Z.
\]
The remaining core depends only on the same four roots, so the
exact capped parent has a bounded four-root Gram.
Let $a=\bar u$ and $d_a=1+a^2$. Its residue is
\[
 (8\mu+2)E^2+
 4\mu(d_aS-X-aZ)^2+(4\mu+2)d_a(aS-Z)^2
       +2(X+aZ)^2-2d_aZ^2.
\]
Put $T=X+aZ$ and $R=aS-Z$; this is an invertible real change of
coordinates. The last three-variable form is
\[
 4\mu(d_aS-T)^2+4\mu d_aR^2+2T^2
       -2d_a a^2S^2+4d_a aSR.
\]
The $T$ and $R$ pivots are $4\mu+2$ and $4\mu d_a$.
The remaining Schur pivot is exactly
\[
 \frac{d_a[4\mu^2-a^2(4\mu+1)]}{\mu(2\mu+1)}.
\]
Taking $\mu=1+a^2$ makes its numerator positive, because
$4\mu^2-a^2(4\mu+1)=3a^2+4$.
Thus this bounded Gram has strictly positive real residue and
is tensor SOS by the inherited strict-residue Gram lemma.

Next suppose $v$ is a bounded unit. At the real residues
$(a,b)=(\bar u,\bar v)$ the four event forms have common kernel
\[
 (W,X,Y,Z)=W(1,b,a,ab).
\]
The core restricted to this line is
$2(1+a^2)(1+b^2)b^2W^2$, which is strictly positive away from zero.
Consequently a sufficiently large positive real $\mu$ makes the
whole capped residue form positive definite. For completeness,
this elementary assertion has a direct quantitative proof.
Let $E_0\succeq0$ be the event Gram, split the state space into
$\ker E_0$ and its orthogonal complement, and let $\sigma>0$ be
its smallest eigenvalue on that complement. Let $\kappa>0$ be
the core's eigenvalue on the one-dimensional kernel and bound
the norms of the two remaining core blocks by $M$.
Then $4\mu\sigma>M+M^2/\kappa$ gives a positive Schur
complement. Every constant here belongs to $\mathbb R$.
The exact state Gram is bounded and has this positive residue.

Finally suppose $v$ is unbounded, including the projective
endpoint. Set $w=1/v$, homogenizing before setting $w=0$.
The four event forms become
\[
 W-wX+uY-uwZ,\quad uW-Y,\quad
 uwW+uX-wY-Z,\quad Y-wZ,
\]
with bounded positive factors $D_w$ or $D_u$ where appropriate.
At $w=0$ their common kernel is
$(W,X,Y,Z)=(0,X,0,aX)$, on which the homogenized core is
$2(1+a^2)X^2$. The same real Schur argument gives a positive
residue Gram for a sufficiently large $\mu$.
For finite unbounded $v$, undo homogenization by the pure
positive testing square $v^2$; for the projective endpoint the
homogeneous expression is already the required parent.
These three cases exhaust every projective Bob test when $u$ is bounded.

It remains to consider $p=1/u$ infinitesimal, after homogenizing
by $p^2$. If $v$ is a bounded unit, use the four roots
\[
 S=pY,\qquad T=vY-Z,\qquad W,\quad X.
\]
The exact event amplitudes are
\[
 W-S,\quad J=W+vX-D_vS+pvT,\quad
 M=pvW-pX+T,\quad T.
\]
The signed core also has bounded coefficients in these roots:
\[
 2J^2+2D_pT^2+2D_v
 [p^2X^2+2vSX-2pXT-v^2S^2+2pvST-p^2T^2].
\]
At $p=0$, set $b=\bar v\ne0$, $A_0=W-S$, and
$B_0=X-bS$. The full capped residue is
\[
 (8\mu+2)T^2+4\mu(1+b^2)A_0^2
 +(4\mu+2)(A_0+bB_0)^2
 +4(1+b^2)bSB_0+2(1+b^2)b^2S^2.
\]
As $\mu$ grows, the first three positive penalties constrain
$T=A_0=B_0=0$, and the residual kernel has positive coefficient
$2(1+b^2)b^2S^2$. The quantitative real Schur argument above
therefore gives a finite $\mu$ with a strictly positive residue.

If $u,v$ are both projectively unbounded, put $w=1/v$ and
homogenize by $p^2w^2$. Use $P=Y-wZ$, $S=pZ$, $W,X$.
Now the exact event amplitudes are
\[
 M=pW-pwX+P,\quad W-pP-wS,\quad
 J=wW+X-pwP-D_wS,\quad P.
\]
The exact core is
\[
 2J^2+2D_pP^2+2D_w[p^2X^2+2SX-S^2].
\]
At $p=w=0$ the capped residue is exactly
\[
 4\mu W^2+(8\mu+2)P^2+4\mu(X-S)^2+2X^2,
\]
which is strict for every $\mu>0$. The literal projective
endpoints cause no division. The literal critical axes also have direct exact certificates:
\begin{align*}
 F|_{p=0}={}&4nD_vX^2+(4m+4\alpha+2)(vY-Z)^2
 +4CD_vW^2+4AD_vY^2+(4D+2)(W+vX)^2,\\
 F|_{v=0}={}&4nD_pX^2+(4m+2)(pX+Z)^2
 +(4C+4D+2)(W-pY)^2+4AD_pY^2+4\alpha D_pZ^2.
\end{align*}
Thus only $p,v$ both nonzero infinitesimal remain untreated.
Applying the real-state result separately to real and imaginary
state components proves the complex-state claim.
\end{proof}

For a strong C2 HI source, the exact source-entry lemma below proves
that all four coefficients in the theorem are unbounded. Hence the
theorem applies without any additional source-shape assumption.
Its precise uncovered complement is the critical testing corner
$1/u$ and $v$ both nonzero infinitesimal. Exact sign gauges may
make both positive.
The original HI source leaf is therefore narrowed but not closed.

\subsection{Further source reductions}
\label{component:pe56fb394b522}

\label{pe56fb394b522:sec:c2-angle-reductions}

All angles in this section are actual source contacts. They are strictly
positive finite field elements; ``infinite'' below means unbounded in
the coefficient field, not a literal projective endpoint. Write
$z=c^2=u_c^2$, $y=a^2z=v_c^2$, and
$A_c=1+zr$, $S=A_c+a^2zt$. The exact source identities used are
\[
 n\ge0\ \Longleftrightarrow\ (b+r)A_cS-r\Delta\ge0,
 \qquad
 N=(1+z)(b+r)h+a^2\{g+j+rt(1+y)\}>0.
\]

\begin{theorem}[The ordinary--ordinary source orbit is already closed]
If both actual C2 half-angles are positive bounded units, then the
source belongs either to the two-weak-complementary-event all-tests
exit or to the five-consecutive-strong-event all-tests exit. No upper
bound on an original event coefficient is required.
\end{theorem}
\begin{proof}
Here $a,z$ are units. The positive gaps give $0<h<1/z$ and
$0<t<1+zr$, so $h,t,b=h+a^2t$ are bounded. From
$g=b-1+r+zh>0$ we obtain $1<b+r+zh\le(1+z)b+r$;
thus $b+r$ is a positive unit. The positive expansion of $\Delta$
makes it a bounded unit too. The original $A$ formula has unit
numerator and bounded denominator $h$, so $A$ is strong.

The exact numerator for $n$, after substituting $t=(b-h)/a^2$,
satisfies the polynomial specialization
\[
 \left[(b+r)A_cS-r\Delta\right]_{r=1}
     =-z(a^2+h)(b+1)(z+1).
\]
If $1-r$ were infinitesimal, taking residues would contradict
$n\ge0$: the right side is strictly negative. Hence $1-r$ is a
unit. The original $m$ numerator is now a unit and $g\Delta$ is
bounded, so $m$ is strong. Apply the exact C2 reflection; it preserves
ordinary actual angles and exchanges $m,D$. Thus $D$ is strong.

If either $C$ or $\alpha$ is weak, its original rational formula
implies $N$ infinitesimal: all the multiplying denominator factors
are bounded and $z,a,b+r,\Delta$ have bounded inverses. The positive
identity for $N$ then forces $h,j$ infinitesimal. Consequently
$t=A_c-j$ and $1-zh$ are units. Both original complementary formulas
therefore give $C,\alpha$ infinitesimal. The two-weak theorem applies.
Otherwise $C,\alpha$ are both strong; with $m,A,D$ strong, the
five-consecutive-event theorem applies. These cases exhaust the source.
\end{proof}

\paragraph{Some directed one-weak angle types are empty.}
Orient an exactly-one-weak source so that $C$ is weak and $\alpha$
strong. In addition to the ordinary--ordinary exclusion above, its
actual angle types cannot be $(\mathrm{unit},\mathrm{Inf})$,
$(\mathrm{Inf},\mathrm{unit})$ or $(\mathrm{Inf},\mathrm{Inf})$.
Here $\mathrm{Inf}$ denotes positive infinitesimal.

For $(\mathrm{unit},\mathrm{Inf})$, $z$ is a unit and $a^2$ is
infinitesimal. Again $t,h,b,\Delta$ are bounded, $\Delta$ is a unit,
and $b+r$ is a unit. Weak $C$ implies $N/a^2$ infinitesimal.
The positive identity forces $j$ and $rt$ infinitesimal. Since
$t=A_c-j$ is a unit, $r$ is infinitesimal. Then $g>0$ makes
$b+zh$ a unit, and $b=h+a^2t$ makes $h$ a unit. The first positive
summand of $N$ is consequently a unit, contradicting $N/a^2$
infinitesimal.

For $(\mathrm{Inf},\mathrm{unit})$, bounded $b$ is a unit by the
inherited source-entry inequality. Explicitly,
\[
 xC(b-h)=\frac{A_c(xN)}{2\Delta},\qquad
 xN=(1+x)(b+r)(xh+k^2)-x(1+k^2)(1-r)(b-h).
\]
The last numerator has residue $k_0^2(b_0+r_0)>0$, independently
of any lower bound on $g$. Thus $xC(b-h)$ is a unit and $C$ is strong. For unbounded
$b$, (SR1) gives $C\asymp1/L$ with bounded $L$, again strong.
This exhausts $b$ and excludes weak $C$ in this angle orientation.

For $(\mathrm{Inf},\mathrm{Inf})$, both $z,y$ are infinitesimal
and $t<A_c$ is bounded. Put $H=zh$, $L=z(b-h)=yt$ and $B=zb=H+L$.
Here $0<H<1$, $L$ is infinitesimal, and $B$ is bounded. If $H$ is
a unit then $b$ is unbounded; direct division of the original
$\Delta$ by $b$ gives a positive unit of residue $1+B_0$.
The positive first summand of $N$ consequently gives
$C\ge c_0H/L$ for some positive real $c_0$, contradicting weak $C$.
Thus weak $C$ forces $H,B$ infinitesimal. The exact ratio is then
\[
 \frac\alpha C=
 \frac{z^2(b+r)(b-h)}{A_c(1-H)}
 =\frac{(B+zr)L}{A_c(1-H)},
\]
which is infinitesimal. Therefore $\alpha$ is weak as well, contrary
to the selected one-weak orientation.

After these exclusions, an exactly-one-weak source oriented with weak
$C$ has only the following candidate actual types:
\begin{gather*}
 (\mathrm{Inf},\mathrm{unbounded}),\quad
 (\mathrm{unit},\mathrm{unbounded}),\quad
 (\mathrm{unbounded},\mathrm{Inf}),\\
 (\mathrm{unbounded},\mathrm{unit}),\quad
 (\mathrm{unbounded},\mathrm{unbounded}).
\end{gather*}
The second type is the reflection of the single-angle domain and
is now fully covered by the homogeneous unbounded-$b$ theorem and the
other single-angle cases. The four other listed directed types
remain conditional source leaves; no nonemptiness is inferred from
their listing.

\paragraph{Other strong-event angle types are not all empty.}
For $x>0$ infinitesimal, the exact positive chart
\[
 c=\sqrt x,\quad a=1,\quad b=1/x,\quad r=1/4,\quad t=1/2
\]
has $h=1/x-1/2$, $g=1/x+1/4-x/2$, $j=1/2+x/4$ and all
remaining chart gaps positive. Its original coefficients are
\begin{align*}
 n&=\frac{8-5x^2}{16(1+x)^2},&
 m&=\frac{3x^2(3x+4)}{16(1+x)(4+x-2x^2)},\\
 C=\alpha&=\frac{4-x}{8x},&
 A&=\frac{x(x+4)(3x+4)}{32(2-x)(1+x)^2},\\
 D&=\frac{(2-x)(3x+4)}{16x(1+x)}.
\end{align*}
Every coefficient is positive. The unchanged exact contact is
\[
 (u_c,v_c;W,X,Y,Z)
 =(\sqrt x,\sqrt x;1/x,-\sqrt x/4,1/\sqrt x-\sqrt x/2,1).
\]
Original source stationarity, positive transverse state inertia and
Laurent convexity prove physicality, including projective limits.
Here $C,\alpha,D$ are unbounded, $m\sim3x^2/16$, $A\sim x/4$,
and $n\to1/2$. Thus a genuine strong-complementary source can have
two infinitesimal actual angles and weak $m,A$. This is an obstruction
to deleting that angle orbit, not a counterexample to tensor SOS.

\subsection{A strong $C$ event from an unbounded first angle}
\label{component:pb7189b9c1f6e}

\label{pb7189b9c1f6e:sec:c2-weak-hi-hu-empty}

\begin{theorem}[The directed weak-$C$ types HI and HU are empty]
Use the exact positive rational C2 source, with $d=1$ and $n\geq0$.
If the first actual half-angle $u_c$ is unbounded and the second
actual half-angle $v_c$ is bounded, then $1/C$ is bounded.
Consequently no exactly-one-weak source oriented with $C$ weak can
have directed actual type HI or HU.  Literal $n=0$ is included.
\end{theorem}

This is an original-source impossibility theorem.  It makes no
claim that a scalar contact identity by itself is a tensor-SOS
certificate.  All independent testing variables and state coordinates
remain unrestricted; the two source leaves are empty before any
testing restriction is imposed.

\begin{proof}
Write
\[
 q=u_c^{-2}>0,\qquad y=v_c^2>0,
 \qquad H=u_c^2h=v_cY_c\in(0,1),
\]
and retain the original $r\in(0,1)$.  Here $q$ is infinitesimal
and $y$ is bounded; $y$ may itself be infinitesimal.  Put
\[
 R=r/q=-u_cX_c,\quad B=u_cv_cW_c,\quad
 L=v_c(u_cW_c-Y_c),\quad T=B-yR.
\]
The exact positive gaps give
\[
 B=H+L=yR+T,\qquad 0<T<H+y,\qquad L>0.
\]
In particular $H,T,y,r$ are bounded.  The original source coordinates
become exactly
\[
 z=1/q,\quad a^2=yq,\quad b=yr+qT,\quad h=qH,
 \quad t=r/q+(T-H)/y,
\]
and the two gaps satisfy
\[
             g=H+(1+y)r-1+qT>0,
 \qquad j=(H+y-T)/y>0.                             \tag{WE1}
\]
The displayed expressions are pure coefficient identities; no
inverse of a mixed source/testing element is used.

We first show that $R$ is bounded.  Set $Q=q^2\Delta>0$.
Direct expansion of the unchanged original formula for $n$ gives
\[
                         n=-\frac{qP}{2rQ},        \tag{WE2}
\]
where the following exact decomposition of $P$ is sufficient:
\begin{align*}
P={}&r^2(1+y)(H+yr)-Tqr(y+2)-q^2(T^2+T)+P_+,\\
P_+={}&HTq^2+HTqr+Hqr(y+1)+T^2q^2r\\
     &+Tq^2r(y+1)+2Tqr^2(y+1)+qr^2y(y+1)>0.        \tag{WE3}
\end{align*}
Suppose for contradiction that $R$ is unbounded, and put
$\kappa=q/r=1/R$, which is infinitesimal.  Let $S=H+yr>0$.
The ratio $T/S$ is bounded.  Indeed, if $r$ is a unit, then
\[
                T<H+y\leq S/r.
\]
If $r$ is infinitesimal, (WE1) and boundedness of $y,T$ imply
$H_0=1$, so $S$ is a unit and $T/S$ is again bounded.
Since $T$ is bounded, $(T^2+T)/S$ is bounded as well.  Dividing
the lower bound in (WE3) by the positive $r^2S$ gives
\[
 \frac{P}{r^2S}\geq (1+y)
       -\kappa\frac{T(y+2)}S
       -\kappa^2\frac{T^2+T}S.
\]
The last two terms are infinitesimal, while $1+y\geq1$.
Thus $P>0$.  Equation (WE2) contradicts $n\geq0$, including
the equality case $n=0$.  Therefore $R$ is bounded.

Now $r=qR$ is infinitesimal.  Equation (WE1) becomes
\[
             g=H-1+q\{(1+y)R+T\}>0,
\]
so $H_0=1$.  Both $B=yR+T$ and $L=B-H$ are bounded.
The original denominator has the exact form
\begin{align*}
 \Delta={}&(1+q)\{B^2+2BR+(1+y)B\}
           +(1-yq)R^2+2(1+y)R+1+y.                \tag{WE4}
\end{align*}
It is a positive bounded unit: all quantities are bounded, and
$1-yq$ is a positive unit, with constant contribution $1+y\geq1$.
The unchanged numerator $N_{\rm src}$ satisfies
\[
 \frac{N_{\rm src}}q
  =(1+q)(B+R)(H+y)-(1+y)(1-qR)L.                  \tag{WE5}
\]
Since $H_0=1$ and $L=B-H$, its residue is
\[
       (N_{\rm src}/q)_0=(1+y_0)(1+R_0)>0.
\]
Thus $N_{\rm src}/q$ is a positive bounded unit.  Finally,
$A_c=1+R$ and $a^2t=qL$, so the original coefficient is exactly
\[
          C=\frac{(1+R)(N_{\rm src}/q)}{2L\Delta}.
                                                               \tag{WE6}
\]
The numerator has bounded positive inverse, and $L\Delta$ is
bounded.  Therefore $1/C$ is bounded, as asserted.  No lower bound
on $L$ was required; if $L$ is infinitesimal the conclusion is
stronger.
\end{proof}

The exact C2 reflection preserves directed type HI and exchanges
$C,\alpha$.  Applying the theorem in both orientations shows that
every physical HI source has both complementary positive square terms strong.
For type HU the reflected type is UI; the reflected consequence is
that $\alpha$ is automatically strong at every UI actual contact.
These observations respect the directed weak-$C$ orientation; they
do not identify IH with HI.

\section{Sources with five lower-bounded events}
\subsection{Five consecutive lower-bounded events}
\label{component:p80253fe9bc1b}

\label{p80253fe9bc1b:sec:five-retained-events}

The following argument permits every retained event coefficient to be
unbounded. It requires no canonical omission, actual-contact coordinates,
or source-shape entry. Five consecutive events in the eight-cycle suffice.

\begin{theorem}[Five consecutive lower bounds]
\label{p80253fe9bc1b:thm:five-retained-events}
Let the original canonical source be physical, with $d>0$ and all eight
event coefficients nonnegative. Suppose
\[
 d/m,\ d/c,\ d/A,\ d/D,\ d/\alpha
 \quad\hbox{are positive bounded coefficient elements}.
\]
Then the entire original parent is tensor SOS at every independent
projective measurement pair and every real or complex projective state.
There is no upper bound on any event coefficient divided by $d$.
By exact CHSH relabeling the same conclusion holds if any five
consecutive event coefficients in the cycle
$(\lambda,D,c,A,\alpha,m,\beta,n)$ have these lower bounds.
\end{theorem}

\begin{proof}
Divide by the positive pure coefficient $d$, so $d=1$.
Let $\mu=\min(1,m,c,A,D,\alpha)$. It is a positive bounded unit.
Select the parent $F_\mu$ with $m=c=A=D=\alpha=\mu$ and
$\lambda=n=\beta=0$. Keep every nonnegative excess and all three
unselected coefficients as their original event squares. This uses only
pure coefficient square roots. It makes no assertion that $F_\mu$ is
globally physical.

At noncommuting residue tests apply
the inherited noncommuting-residue theorem to the full physical source. It
performs its own positive coefficient normalization if the event
weights are unbounded.
It remains to certify $F_\mu$ near every commuting residue test. Its
coefficients are bounded, so a positive-definite real residue Gram is
sufficient by the bounded positive-residue Gram lemma.

Write $\mathcal H$ for the signed core with $d=1$ and
$D_u=1+u^2,D_v=1+v^2$. On the four open commuting axes its exact
homogeneous restrictions are
\[
\begin{aligned}
\mathcal H|_{u=0}&=2D_v(X^2+Y^2),\\
\mathcal H|_{v=0}&=2\{(uW-Y)^2+(X+uZ)^2\},\\
\mathcal H|_{u=\infty}&=2\{(W+vX)^2+(vY-Z)^2\},\\
\mathcal H|_{v=\infty}&=2D_u(X^2+Y^2).
\end{aligned}                                                   \tag{NW1}
\]
The remaining finite angle has nonzero residue. At $u=0$, the retained
$m,D$ positive square terms force $W=Z=0$ on the core kernel. At $v=0$, the retained
$A,\alpha$ positive square terms force $Y=Z=0$, and then $W=X=0$. At $u=\infty$,
the retained $c,A$ positive square terms force $W=Y=0$, and then $X=Z=0$. At
$v=\infty$, the retained $c,D$ positive square terms force $W=Z=0$ after the core
has forced $X=Y=0$. Thus the
selected residue state matrix is positive definite on every open axis,
for either sign of the other finite angle.

Normalize a nonzero testing state by a coordinate of largest absolute
value. Its coordinates are then bounded and one is a unit. At the four
double corners, write $F_{ab}$ for the homogeneous restriction of
$F_\mu$ to $(u,v)=(a,b)$. Directly,
\[
\begin{aligned}
F_{00}&=(2+4\mu)X^2+(2+12\mu)Y^2+4\mu Z^2,\\
F_{\infty0}&=(2+8\mu)(W^2+Z^2)+4\mu Y^2,\\
F_{0\infty}&=4\mu(W^2+Z^2)+2X^2+(2+12\mu)Y^2,\\
F_{\infty\infty}&=4\mu W^2+(2+4\mu)X^2+(2+12\mu)Y^2.
\end{aligned}                                                    \tag{NW2}
\]
Here infinity means the polynomial homogeneous chart, obtained by
dividing the finite parent by the appropriate pure-testing squared
angle. Every form is strict except its explicitly missing state
coordinate. Consequently the only nontrivial residue states are the
$W$ ray at $(0,0)$, the $X$ ray at $(\infty,0)$,
and the $Z$ ray at $(\infty,\infty)$.

At $(0,0)$ normalize $W=1$. The five local coordinates
$(u,v,X,Y,Z)$ are infinitesimal. The complete polynomial $F_\mu$
has no constant or linear term, and its degree-two part is
\[
 (2+8\mu)(u-Y)^2+2X^2
       +4\mu\{(v-X)^2+Y^2+Z^2\}.                    \tag{NW3}
\]
Its residue is positive definite: vanishing successively gives
$X=Y=Z=v=u=0$.

At the corner $(\infty,0)$ normalize $X=1$ and put $p=1/u$.
The five infinitesimal coordinates are $(p,v,W,Y,Z)$. The complete
homogeneous polynomial again has no terms below degree two, and its
quadratic part is
\[
 (2+4\mu)(W+v)^2+(2+4\mu)(p+Z)^2
                         +4\mu(W^2+Y^2+Z^2).        \tag{NW4}
\]
This is positive definite. At the corner $(\infty,\infty)$,
normalize $Z=1$ and put $p=1/u,q=1/v$. On $(p,q,W,X,Y)$ the exact
quadratic part is
\[
 (2+8\mu)(q-Y)^2+2X^2
                         +4\mu\{(p-X)^2+Y^2+W^2\}. \tag{NW5}
\]
It is positive definite as well. One can obtain (NW5) directly from
(NW3) by the exact symmetry
$(u,v;W,X,Y,Z)\mapsto(1/v,1/u;Z,X,Y,W)$, which permutes the five
equal retained weights.

For completeness, these quadratic calculations control the entire
polynomial. In each local chart, allocate every higher monomial to a
product of two of the five coordinate roots. The remaining monomial
has positive degree, hence is infinitesimal; its coefficient is bounded
because it is a real polynomial in the bounded unit $\mu$. Symmetrizing
gives an exact finite five-root Gram whose residue is (NW3), (NW4), or
(NW5). No inversion, root, or comparison mixing the two fields occurs.
The bounded positive-residue Gram lemma therefore certifies the full
selected polynomial, including all higher terms.

Restore the excess event squares, the positive pure-testing homogeneous
and projective state factors, and the positive coefficient $d$. The
zero state is immediate. Splitting a complex state into real and
imaginary parts proves the complex assertion.

Finally the exact full-parent Alice reflection acts on the displayed
cycle indices as $i\mapsto7-i$, and the swap/reflection
$(u,v;W,X,Y,Z)\mapsto(1/v,1/u;Z,X,Y,W)$ acts as $i\mapsto6-i$.
Their composition gives a one-step cycle rotation. Their full-parent
identities have only positive pure-testing homogeneous factors, so they
transport the five-positive square term conclusion to every cyclic translate.
\end{proof}

This theorem is a sufficient source domain, not an exhaustive zero-$d$
source cover. It requires five selected events to be positive and
noninfinitesimal relative to $d$. Its all-tests conclusion is independent
of actual contact angles and of all original event hierarchies. In
particular, it contains both C1 and C2 sources with all six retained
events bounded below relative to $d$.

\section{Two small actual angles: common identities}
\subsection{Source coefficients and normal arrays}
\label{component:pd709027c7805}

The following are the exact native scalar and state-normal coefficient arrays.
Put $r=1-e$, $z=aB$, $y=tB$, and
\[
 c_a=1+aBr,\quad S_c=1+B(1-h+ar),\quad
 g=1+a(hB+r-1),\quad j=c_a-(1-h)/t,
\]
\begin{align*}
\Delta={}&a+(1+aB)+(1+aB)B+2aB(1+aB)r+a^3B^2r^2+2a^2Br\\
 &+tB\{1+a+aB+a^2Br(2-r)\},\\
M={}&(1+aB)(1+ar)(h+t)-a(1+tB)e(1-h).
\end{align*}
The physical stationary event coefficients are
\[
\begin{aligned}
n&=\frac{(1+ar)c_aS_c-r\Delta}{2r\Delta},&
m&=\frac{a^2(1+tB)e c_aS_c}{2g\Delta},\\
C&=\frac{c_aM}{2(1-h)\Delta},&
Ah&=\frac{t c_a(1+ar)S_c}{2\Delta},\\
D&=\frac{(1+aB)(1+ar)(1-atBr)S_c}{2tj\Delta},&
\alpha&=\frac{B^2(1+ar)M}{2(1-hB)\Delta}.
\end{aligned}
\]
The following scalar and normal arrays are exact rational functions of
the native stationary source. Each case certificate proves the source
bounds it needs at its own independent openings; the arrays themselves
have no bounded-$a$ or old R9-domain restriction.

\paragraph{Scalar arrays.}
The notation $Ah$ in these formulas means the product of $A$ and $h$.
\[
\begin{aligned}
q_s&=nr^2+(m+\tfrac12)e^2,\qquad q=q_s+\alpha,\\
a_w&=(C+D+\tfrac12)(1-h)^2+Ah\,h,\\
b_2&=2mh+C(1-h)^2+Ah\,h+(\alpha+\tfrac12)h^2,\\
b_4&=h^2(m+\alpha+\tfrac12)+Ah\,h,\\
d_2&=(n+\tfrac12)r^2+2Dr,\qquad
d_4=(n+D+\tfrac12)r^2,\\
e_0&=me+D(1-h)+\alpha h+\tfrac12,\\
e_1&=mhe+r(D+\tfrac12)(1-h)+(\alpha+\tfrac12)h,\\
c_F&=(b_2-a^2nr^2)/t^2,\qquad
\lambda=(m-a^2q)/(ta^2B),\\
f_{A0}&=Ah\,h/t^2+d_2-2e_1/t+2anr^2/t,\\
f_{C0}&=(\alpha-c_Ft^2B^2)/(taB^2),\\
f_G'&=-\{a_w/t^2+D-2e_0/t+2aq/t+2c_FtB\}/(aB).
\end{aligned}
\]
The potentially singular $a,B$ factors cancel in the combinations
used in RC1. One direct proof is to clear the denominators in the displayed source
coefficients and the printed monomials, and then expand; all three resulting numerator polynomials are
zero. Equivalently, the original value and two first derivatives at
contact give precisely RC1. Thus RC1 is an exact identity, including
arbitrarily small $a$ and literal $n=0$.

\paragraph{All first-normal arrays.}
For the angular weight $q=q_s+\alpha$ define
\[
\begin{aligned}
K_X'&=-nr+(m+\tfrac12)e-aq,\\
K_E'&=a_w/t^2-e_0/t+aq/t,\\
K_Y&=\{-(C+D+\tfrac12)(1-h)+Ah\}/t+D-\alpha,\\
x_3&=2\{(D+\tfrac12)[(1-h)/t-1]+e/2-nr(1+a/t)-mh/t\},\\
x_5&=-2r(n+D+\tfrac12),\\
a_E&=2Ah\,h/t^2-2e_1/t,\qquad b_E=2b_2/t^2,\qquad c_E=2b_4/t^2,\\
a_Y&=2\{Ah/t-me+r(D+\tfrac12)-\alpha-\tfrac12\},\\
b_Y&=2\{-C(1-h)+Ah+m+(\alpha+\tfrac12)h\}/t,\\
c_Y&=2\{Ah+h(m+\alpha+\tfrac12)\}/t.
\end{aligned}
\]
\paragraph{The six exact contact balances.}
With $x=aB$, $y=tB$, the scalar balances are
\[
 f_{A0}=\lambda-2d_4x-b_4y/t^2,\qquad
 f_{C0}=\lambda+b_4y/t^2,\qquad
 f_G'=3\lambda-d_4x+b_4y/t^2.\tag{RC1}
\]
The complete first-normal balances are
\[
 2K_E'+a_Ex+b_Ey+c_Exy+2\lambda x=0,
\]
\[
 2K_X'+x_3y+x_5xy-2t\lambda x=0,\qquad
 2K_Y+a_Yx+b_Yy+c_Yxy=0.\tag{RC6}
\]
Substitution of the native weights and multiplication by the displayed
source denominators gives six zero polynomial numerators. Equivalently
these are the original contact value and its five first derivatives.
They underlie the generic whole-parent identity BA7 (also printed as
R10.9). BA6 explicitly gives its entire hard quadratic. Those identities
hold before any source size assumption; the bounded-$a$, R8 and R10
certificates separately establish their coefficient bounds and pivots.

\subsection{The five-root identity}
\label{component:p9ec3d6bbd717}
Set $\kappa=1/t$, $\ell=1/a$. The letters $e_0,e_1$ are the native
scalar coefficients above. In this subsection $a_g$ denotes the scalar
coefficient called $a$ in the five-root table, and not the source ratio.
\paragraph{Five homogeneous contact roots.}
Put $U=u^2$, $V=v^2$, and define
\[
 \xi=X+ruZ,\qquad q=\sqrt A\,(Y-huW),\qquad
 R_1=uW-\kappa vZ,\quad R_2=vW-\ell uZ,\quad
 T=uvW-BZ.                                             \tag{WA3}
\]
Every root vanishes at the actual source contact. Inverses and
radicals in their coefficients are pure source. The large $A$
positive square term is retained in $q$; no positive square term is subtracted from
the parent and no physicality hypothesis is transferred to a
modified remainder.

Use the following source coefficients for the exact scalar graph:
\begin{align*}
 a_g&=(C+D+\tfrac12)(1-h)^2+Ah^2,&a_4&=Ah^2,\\
 b_2&=2mh+C(1-h)^2+(A+\alpha+\tfrac12)h^2,
 &b_4&=h^2(m+A+\alpha+\tfrac12),\\
 c&=\alpha,&c_2&=nr^2+(m+\tfrac12)(1-r)^2+\alpha,\\
 c_4&=nr^2,&d_2&=(n+\tfrac12)r^2+2Dr,\\
 d_4&=(n+D+\tfrac12)r^2.
\end{align*}
Thus the graph on $\xi=q=0$ is
\[
 \begin{split}
 F_g/4={}&[a_gU+a_4U^2+(m+b_2U+b_4U^2)V]W^2\\
 &+[c+c_2U+c_4U^2+(D+d_2U+d_4U^2)V]Z^2
       -2uv(e_0+e_1U)WZ,
 \end{split}
\]
where $e_0,e_1$ are exactly the original coefficients in (WS3).
Set
\[
 \begin{gathered}
 \beta_s=\frac{Bd_4+\kappa c_4}{\kappa\ell^2},\qquad
 \gamma_s=b_4+\frac{d_4}{\kappa\ell},\\
 j_0=\frac{Bd_4-a_4\kappa^2\ell+d_2\ell}
                   {2\kappa\ell^2},\qquad
 j_1=\frac{d_4}{2\kappa\ell},\qquad J_s(U)=j_0+j_1U,\\
 f_s=\frac12\left(b_2-\frac{c_4}{\ell^2}
                -\frac c{B^2}-\frac{Bd_4}{\kappa\ell^2}\right).
 \end{gathered}                                        \tag{WA4}
\]
The four original scalar contact equations (value and the three
logarithmic derivatives in $u,v,W/Z$) give the exact identity
\[
 F_g/4=(R_1,R_2,T)\,\mathsf S\,(R_1,R_2,T)^T,\qquad
 \mathsf S=
 \begin{pmatrix}
 a_g+a_4U&uvJ_s&-vJ_s\\
 uvJ_s&m+\beta_sU&u(f_s-j_1U)\\
 -vJ_s&u(f_s-j_1U)&c/B^2+\gamma_sU
 \end{pmatrix}.                                       \tag{WA5}
\]
This is a different scalar identity from the excluded polynomial
two-block splitter (JG7). In particular, the last diagonal retains
the full radial mass $\alpha/B^2$ and the Gram depends on the
independent testing variables. The source bounds for this exact identity are established in each coefficient domain below.

\paragraph{Both full first normals and the complete hard block.}
For the $\xi$ normal define
\[
 \begin{gathered}
 p_0=\frac{D+r(n+1/2)}\kappa,\qquad
 p_1=\frac{r(n+D+1/2)}\kappa,\qquad
 q_1=\frac{nr+Bp_1}\ell,\\
 f_x=-mh+(D+\tfrac12)(1-h)-p_0-q_1,\\
 \mathsf L_x=\bigl(uv(p_0+p_1U),\ -m+q_1U,\
                                u(f_x-p_1U)\bigr).
 \end{gathered}                                        \tag{WA6}
\]
For the unnormalized $\eta=Y-huW$ normal set
\[
 \begin{gathered}
 p_y=-(C+D+\tfrac12)(1-h)+Ah,\qquad p_y'=Ah,\\
 k_v=-C(1-h)+Ah+m+(\alpha+\tfrac12)h,\qquad
 k_{uv}=h(A+m+\alpha+\tfrac12),\\
 k_z=-m(1-r)+r(D+\tfrac12)-\alpha-\tfrac12,\\
 q_y=\frac{-k_z-\kappa p_y'-Bk_{uv}}\ell,\qquad
 f_y=k_v-q_y,\\
 \mathsf L_y=\bigl(p_y+p_y'U,\ uvq_y,\
                                      v(f_y+k_{uv}U)\bigr).
 \end{gathered}                                        \tag{WA7}
\]
The two remaining original source stationarity equations, in $X,Y$,
give these exact representations of the complete first normals.
No first-normal term is discarded. Define also
\[
 \begin{split}
 h_{xx}&=n(1+U)(1+V)+m+(D+\tfrac12)UV+\tfrac12(1+V),\\
 h_{xy}&=-uv(m+D+\tfrac12),\\
 h_{yy}&=A(1+U)(1+V)+C(1+V)+D+\tfrac12
                   +mUV+(\alpha+\tfrac12)(1+U)V.
 \end{split}
\]
Then the \emph{entire original four-state parent} is exactly
\[
 F/4=\mathcal R^T\mathsf G(u,v)\mathcal R,\qquad
 \mathcal R=(\xi,q,R_1,R_2,T)^T,
\]
\[
 \mathsf G=
 \begin{pmatrix}
 h_{xx}&h_{xy}/\sqrt A&\mathsf L_x\\
 h_{xy}/\sqrt A&h_{yy}/A&\mathsf L_y/\sqrt A\\
 \mathsf L_x^T&\mathsf L_y^T/\sqrt A&\mathsf S
 \end{pmatrix}.                                       \tag{WA8}
\]
Block notation in (WA8) has sizes $1,1,3$. Every coefficient of this finite polynomial matrix is pure source.
 It then verifies (WA5)--(WA7) from the corresponding
four plus two contact equations. Thus (WA8) is not merely a graph
or Hessian certificate.

\subsection{Division at the source-ratio zero}
\label{component:p296e824cb97a}
Use the native coefficient dictionary and put $\mu=m/a^2$,
$\alpha_*=\alpha/B^2$.

The only new potentially unbounded local coefficient is
$(b_2-\alpha_*)/a$, where
\[
 b_2=2a^2\mu h+C(1-h)^2+(Ah)h+(B^2\alpha_*+1/2)h^2.
\]
At $a=0$, with $d_*=1+B(1+t)$, the exact values are
\[
 C=\frac{h+t}{2(1-h)d_*},\quad
 Ah=\frac{t[1+B(1-h)]}{2d_*},\quad
 \alpha_* =\frac{h+t}{2(1-Bh)d_*}.
\]
Substitution gives $b_2=\alpha_*$ identically for arbitrary
$B,h,t$. Clearing the unit denominators in (WE1), the numerator
of $b_2-\alpha_*$ is therefore divisible by $a$ in the polynomial
ring. Its quotient has a unit denominator at $B=a=h=0$.
Hence
\[
 Q_*=(b_2-\alpha_*)/a\quad\hbox{is bounded}.       \tag{WE3}
\]
This is a joint regularity argument, so it permits arbitrary
$h/a$, including an unbounded ratio. 

\paragraph{The complete local five-root Gram survives.}
Keep $\kappa=B/y$ and every coefficient and normal in
(WA3)--(WA8), but set $\ell=1/a$ and replace $R_2$ by
\[
 \widetilde R_2=aR_2=avW-uZ.
\]
The entire full-parent matrix is the exact congruence
\[
 \widetilde G=EGE,\qquad E=\operatorname{diag}(1,1,1,a^{-1},1),
 \qquad \widetilde{\mathcal R}=(\xi,q,R_1,\widetilde R_2,T).
                                                               \tag{WE4}
\]
Its newly rescaled coefficients are
\[
\begin{gathered}
 m/a^2=\mu,\qquad
 \beta_s/a^2=(Bd_4+\kappa c_4)/\kappa,\\
 J_s(U)/a=\frac{aBd_4-a_4\kappa^2+d_2+d_4U}{2\kappa},\qquad
 f_s/a=\tfrac12[Q_*-ac_4-aBd_4/\kappa],\\
 (-m+q_1U)/a=-a\mu+(nr+Bp_1)U,\\
 q_y/a=-k_z-\kappa p_y'-Bk_{uv}.                  \end{gathered}\tag{WE5}
\]

\subsection{The angular circuit}
\label{sec:c2-short-angular-identity}
For $\lambda=(m-a^2q)/(ta^2B)>0$ and $q=nr^2+(m+1/2)e^2+\alpha$,
the following identity holds before any source-size assumption.
Set $U=u^2$, $V=v^2$, $k=\sqrt{x}$, and $c=\sqrt{a/t}$.
Use the state-linear forms
\[
 E=tuW-vZ,\quad \eta=Y-huW=\sqrt h\,N_y,\quad
 \xi=X+ruZ,\quad F_v=(V-y)Z,\quad R_D=v(U-x)Z,
\]
\[
 J=(uv-x/c)Z,\quad Q=\sqrt q(avW-uZ-a\xi),\quad
 N=\sqrt n\{r[(a/t)V-U]Z+u\xi\},
\]
\[
 R=\sqrt\lambda\,tk(cv-u)W,\qquad
 T=\sqrt{2\lambda/c}\,k\sqrt{uv}(tcW-Z).              \tag{BA4}
\]
The square root in $T$ is pure testing. Initially take $u,v>0$;
literal axes are addressed below.

For a compact coefficient table put
\[
 \begin{aligned}
 b_2&=2mh+Cl^2+(A+\alpha+\tfrac12)h^2,&
 b_4&=h^2(m+A+\alpha+\tfrac12),\\
 d_4&=(n+D+\tfrac12)r^2,&
 e_1&=mhe+r(D+\tfrac12)l+(\alpha+\tfrac12)h,\\
 c_F&=(b_2-a^2nr^2)/t^2,&
 a_E&=2Ah^2/t^2-2e_1/t,\\
 b_E&=2b_2/t^2,&c_E&=2b_4/t^2,\\
 a_Y&=2[Ah/t-me+r(D+\tfrac12)-\alpha-\tfrac12],\\
 b_Y&=2[-Cl+Ah+m+(\alpha+\tfrac12)h]/t,\\
 c_Y&=2[Ah+h(m+\alpha+\tfrac12)]/t,\\
 x_3&=2[(D+\tfrac12)(l/t-1)+e/2-nr(1+a/t)-mh/t],&
 x_5&=-2r(n+D+\tfrac12).
 \end{aligned}                                                     \tag{BA5}
\]
Let $Y_*=hE/t+\eta$ and define the entire hard quadratic explicitly by
\[
 \begin{split}
 H_0={}&CD_v(lE/t-\eta)^2+AD_uD_vY_*^2
 +(D+\tfrac12)(lE/t-\eta+uv\xi)^2\\
 &+(\alpha+\tfrac12)D_uVY_*^2
 +[n(1+V+UV)+D_v/2]\xi^2\\
 &+2m(VE/t-uv\xi)Y_*+mUVY_*^2.
 \end{split}                                                       \tag{BA6}
\]
In particular BA6 has no inverse testing angle and no $a$ denominator.
Upon replacing $\eta$ by $\sqrt hN_y$, every occurrence of $A$
is through $Ah$, $\sqrt hAh$, or $hAh$.

With $\sigma=t\lambda x/\sqrt q$, the exact identity is
\[
 \begin{split}
 P={}&Q^2+N^2+T^2+R^2+\lambda J^2+d_4R_D^2
 +(c_F+b_4x/t^2)F_v^2+(b_4/t^2)vR_DF_v\\
 &+H_0-\lambda xE^2-\mu\xi^2-2\sigma Q\xi\\
 &+E[(a_E+c_EV)R_D+(b_E+c_Ex)vF_v]\\
 &+\sqrt hN_y[(a_Y+c_YV)R_D+(b_Y+c_Yx)vF_v]\\
 &+\xi[x_5uvR_D+u(x_3+x_5x)F_v].
 \end{split}\tag{BA7}
\]
This is the generic RC/R10 identity, not a scalar restriction.
One pen-and-paper verification is to substitute BA4 into BA7 and use
\[
 R^2+T^2+\lambda J^2
 =\lambda\{t^2x[U+(a/t)V]W^2-4txuvWZ+(UV+xy)Z^2\}.
\]
The hard quadratic is then exactly BA6. The remaining six scalar
and first-normal coefficient identities follow by substituting the native source fractions in Section~\ref{component:pd709027c7805};
a common denominator leaves zero polynomial numerators. Equivalently
these six identities are the original value and all five first-contact
conditions. Both the full identity and those six numerator equalities
have also been independently expanded during verification.

\subsection{Homogeneous identities for unit $t>1$}
The following polynomial identities hold before any source-size
specialization. Their strict residues in this subsection use
$t_0=\tau>1$, $C\sim c_0=\tau/2$, $D\sim d_0=1/[2(\tau-1)]$,
$A^{-1/2}$ infinitesimal and $n,m,\alpha$ infinitesimal. Each source
calculation establishes its own coefficient bounds; in particular R10.6
below supplies the exact modification when $m$ is merely bounded.
Put $\widehat Y=\sqrt A Y$, so $Y=A^{-1/2}\widehat Y$.

First suppose $u$ is a bounded unit or unbounded and $v$ is
infinitesimal. Put $p=1/u$, $D_p=1+p^2$ and use
\[
 N_\infty=\sqrt n X,\quad V_X=vX,\quad K=pX+Z,\quad
 R_\infty=\sqrt\alpha pX,\quad W,\quad\widehat Y.
\]
Let $J_\infty=W+D_pV_X-pY-pvK$. The entire homogeneous parent is
\[
 \begin{split}
 p^2P(1/p,v)={}&D_pD_vN_\infty^2+m(pvW+vY-K)^2
 +CD_v(W-pY)^2+D_pD_v\widehat Y^2\\
 &+DJ_\infty^2+D_p[R_\infty+\sqrt\alpha(vY-K)]^2\\
 &+\tfrac12\{J_\infty^2+(1-p^2v^2)K^2\\
 &\hspace{9mm}+D_p[v^2Y^2-2vYK+2pYV_X+2pvKV_X-p^2V_X^2]\}.
 \end{split}                                                       \tag{BA11}
\]
It follows directly by substituting $Z=K-pX$ in the quarter-parent $P=F/4$ defined by \eqref{eq:c2-master-parent}. Its bounded
six-root residue, with $d_p=1+p_0^2$, is
\[
 d_p(N_\infty^2+R_\infty^2+\widehat Y^2)+K^2/2
 +c_0W^2+(d_0+\tfrac12)(W+d_pV_X)^2-d_pp_0^2V_X^2/2.
\]
The $W,V_X$ determinant is $d_pc_0(d_0+1/2)>0$, using
$c_0(d_0+1/2)=(c_0+d_0+1/2)/2$. All other pivots are positive.
No inverse of $n$ or $\alpha$ occurs, and either root may vanish.
This covers literal $u=\infty$ by $p=0$ as well.

Next suppose $u$ is infinitesimal and $v$ is a bounded unit.
Use the five roots
\[
 (S,R_u,X,\widehat Y,Z)=(\sqrt mW,uW,X,\sqrt A Y,Z).
\]
The exact full polynomial is obtained from the quarter-parent $P=F/4$ defined by \eqref{eq:c2-master-parent} by the two replacements
\[
 m(vW-X+uvY-uZ)^2=(vS-\sqrt mX+\sqrt m\,uvY-\sqrt m\,uZ)^2,
 \qquad uW=R_u
\]
in all other terms. This operation uses no inverse of $m$ or $u$.
Its separate residue pivots on $S,X,\widehat Y$ are
$V,(1+V)/2,1+V$, where $V=v_0^2>0$. The remaining $R_u,Z$ block is
\[
 \begin{pmatrix}
 c_0(1+V)+d_0+1/2&-(d_0+1/2)v_0\\
 -(d_0+1/2)v_0&d_0V
 \end{pmatrix},\qquad \det=c_0d_0V^2>0.             \tag{BA12}
\]
Here $c_0d_0-d_0/2-1/4=0$. If instead $v$ is unbounded, use
$q_v=1/v$ and the same five roots. Multiplying the same exact
polynomial by $q_v^2$ gives, at $u=q_v=0$, the diagonal residue
\[
 (S,R_u,X,\widehat Y,Z):\quad (1,c_0,1/2,1,d_0)>0.   \tag{BA13}
\]
This homogeneous substitution is polynomial and also applies at
literal $q_v=0$.

If both finite testing residues are nonzero, the hard
$\widehat Y$ pivot is $(1+U)(1+V)$. The complete remaining
three-state residue on $(W,X,Z)$ is
\[
 c_0(1+V)UW^2+(d_0+\tfrac12)(u_0W+u_0v_0X-v_0Z)^2
 +\tfrac12(1+V)(X+u_0Z)^2-\tfrac12(1+U)VZ^2,
\]
where $U=u_0^2,V=v_0^2$. Its determinant is
\[
 (1+U)(1+V)UV\left\{\frac{c_0(2d_0+1)}4U
                              +\frac{c_0d_0}2V\right\}>0.        \tag{BA14}
\]
The $Z=0$ restriction is strictly positive, so the full matrix is
positive definite. Every coefficient is bounded after the hard-state
normalization.

The one-infinity cases not already covered by BA11 or BA13 are
strict directly. At $u=\infty$ and a bounded unit $v$, the residue is
\[
 c_0(1+V)W^2+(d_0+\tfrac12)(W+v_0X)^2+Z^2/2
 +(1+V)\widehat Y^2\succ0.
\]
At $v=\infty$ and a bounded unit $u$, it is
\[
 c_0UW^2+(1+U)X^2/2+d_0(u_0X-Z)^2
 +(1+U)\widehat Y^2\succ0.
\]
These are the degree-two homogeneous values of the quarter-parent $P=F/4$ defined by \eqref{eq:c2-master-parent} in the respective
angles; nearby reciprocal coordinates are pure-testing infinitesimals.

For two infinite angles put $p=1/u$, $q_v=1/v$ and introduce
$P_\infty=pZ$, $Q_\infty=q_vZ$. The entire exact polynomial is
\[
 \begin{split}
 p^2q_v^2P(1/p,1/q_v)={}&n(1+p^2)(1+q_v^2)X^2
 +m(pW-pq_vX+Y-Q_\infty)^2\\
 &+C(1+q_v^2)(W-pY)^2
 +(1+p^2)(1+q_v^2)\widehat Y^2\\
 &+D(q_vW+X-pq_vY-P_\infty)^2
 +\alpha(1+p^2)(Y-Q_\infty)^2\\
 &+\tfrac12\{(q_vW+X-pq_vY-P_\infty)^2
 +(1+p^2)(Y-Q_\infty)^2\\
 &\hspace{9mm}+(1+q_v^2)[p^2X^2+2XP_\infty-P_\infty^2]\}.
 \end{split}                                                       \tag{BA15}
\]
Its residue on $(W,X,\widehat Y,P_\infty,Q_\infty)$ is
\[
 c_0W^2+(d_0+\tfrac12)X^2-2d_0XP_\infty
 +d_0P_\infty^2+\widehat Y^2+Q_\infty^2/2\succ0.
\]
The nontrivial determinant is $d_0/2>0$. At $p=q_v=0$ the two
extra roots vanish, respecting the original null state rather than
inventing a positive $Z$ pivot.

\subsection{Bounded source ratio}
\label{sec:c2-bounded-a-new}

Let $K_c,K_t$ be independently labelled real closed extensions of
$\mathbb R$. All source letters below belong to $K_c$, and all testing
letters belong to $K_t$. No algebraic-disjointness hypothesis is imposed.
Write $\mathrm I$ for positive infinitesimals and call a positive bounded
element a unit when its reciprocal is bounded.

We use the actual stationary C2 source, normalized by $d=1$, with
$\lambda_{\rm event}=\beta=0$. Put
\[
 x=aB,\quad y=tB,\quad l=1-h,\quad e=1-r,\quad
 g=1-ae+haB,\quad j=1+aBr-l/t.
\]
The source assumptions are
\begin{equation}
 \begin{gathered}
 B,x,y\in\mathrm I,\qquad a>0\text{ bounded},\qquad
 t\text{ a positive bounded unit},\\
 0<h,r<1,\qquad g,j,1-txr>0,\qquad n\ge0,\qquad
 C\ge c_*>0\quad(c_*\in\mathbb R).
 \end{gathered}
 \label{eq:c2-ba-domain}
\end{equation}
The last strict gap is exactly the original $B-xyr>0$ divided by $B$.
The source formulas, which also fix our normalization, are
\[
 c_a=1+aBr,\quad S=1+B(l+ar),
\]
\[
 \begin{split}
 \Delta={}&a+(1+aB)+(1+aB)B+2aB(1+aB)r+a^3B^2r^2+2a^2Br\\
 &+tB\{1+a+aB+a^2Br(2-r)\},\\
 M={}&(1+aB)(1+ar)(h+t)-a(1+tB)el,
 \end{split}
\]
\begin{equation}
 \begin{aligned}
 n&={ (1+ar)c_aS-r\Delta\over2r\Delta},&
 m&={a^2(1+tB)e c_aS\over2g\Delta},\\
 C&={c_aM\over2l\Delta},&
 Ah&={t c_a(1+ar)S\over2\Delta},\\
 D&={(1+aB)(1+ar)(1-taBr)S\over2t j\Delta},&
 \alpha&={B^2(1+ar)M\over2(1-hB)\Delta}.
 \end{aligned}
 \label{eq:c2-ba-source}
\end{equation}
The original source is physical; no physicality assertion about a
polynomial obtained by a formal coefficient substitution will be used.

\begin{proposition}
\label{prop:c2-bounded-a}
Under \eqref{eq:c2-ba-domain}, the entire C2 parent is a finite sum of
squares in $K_c\otimes_{\mathbb R}K_t$ for every independent projective
testing pair and every state. The conclusion includes $n=0$, every
relative rate of $g,j,h,l,r,e$, and $a$ infinitesimal in residue.
\end{proposition}

For finite tests set $D_u=1+u^2$, $D_v=1+v^2$ and define the quarter-parent
\begin{equation}
 \begin{split}
 P={}&m(vW-X+uvY-uZ)^2+nD_uD_vX^2+CD_v(uW-Y)^2+AD_uD_vY^2\\
 &+(D+\tfrac12)(uW+uvX-Y-vZ)^2
 +(\alpha+\tfrac12)D_u(vY-Z)^2\\
 &+\tfrac12D_v(X+uZ)^2-\tfrac12D_uD_vZ^2.
 \end{split}
 \label{eq:c2-ba-parent}
\end{equation}
All identities below are identities of this four-state polynomial.
\paragraph{Gram completion and real caps.}
Apply Lemma~\ref{p25af2c58da7e:lem:gram} to the bounded remainder after
retaining the displayed nonnegative squares. The roots may be unbounded;
only its matrix and positive-definite real residue enter this lemma.

For real caps, let a real quadratic form be positive definite on the
common kernel of finitely many squared rows. A sufficiently large real
multiple of those squares makes the form positive definite: decompose
into the kernel and its orthogonal complement and take the Schur
complement. Each unbounded source weight admits that cap; a finite weight
admits a sufficiently close smaller cap when the finite-residue form is
strict. Retain every nonnegative excess.

\paragraph{Source bounds when $e$ is a unit.}
The denominator has residue $1+a_0>0$. The physical numerator has the
exact identity
\begin{equation}
 2r\Delta n=eQ_s-BK_s\ge0,
 \quad
 \begin{aligned}
 Q_s={}&1+B(1+t+ah+2ar+at)\\
       &+B^2\{a(h+t)+a^2[h(r+1)+t(1+r-r^2)]\},\\
 K_s={}&(1+a)(h+t)(1+aB).
 \end{aligned}
 \label{eq:c2-ba-physical}
\end{equation}
Consequently $(nr)_0=e_0/[2(1+a_0)]>0$. In particular $n$ has a
positive real lower bound and $n=0$ is impossible on this branch.
The displayed source formulas give positive bounded units
\begin{equation}
 nr,\qquad mg/a^2,\qquad Ah,\qquad Cl,\qquad Dj,
 \qquad \alpha_*:=\alpha/B^2.
 \label{eq:c2-ba-unit-products}
\end{equation}
For $Cl$ and $\alpha_*$ one needs the actual strong-$C$ assumption:
if $l$ is a unit it makes $M$ a unit; if $l$ is infinitesimal then
$h_0=1$ and $M_0=(1+a_0r_0)(1+t_0)>0$. This explicitly excludes the
weak-$C$ corner where $M_0=0$. Thus $A,C,D,n$ have positive real
lower bounds. If $g$ is a unit, $m/a^2$ is a positive unit; if $g$ is
infinitesimal, $a$ is a unit and $m$ is unbounded.

\paragraph{One exact five-root matrix for the unit-$e$ local chart.}
Use the complete WA matrix, including both normal rows, in the roots
\[
 z=(\xi,\eta,R_1,R_2',T)^T,
 \quad \xi=X+ruZ,\quad\eta=Y-huW,
\]
\[
 R_1=uW-vZ/t,\qquad R_2'=avW-uZ,\qquad T=uvW-BZ,
 \qquad P=z^TGz.
\]
The precise generic input here is WA3--WA8: replace its normalized
$q=\sqrt A\eta$ by $\eta$, and its $R_2$ by $R_2'/a$.
It is an exact full-parent identity derived from the six original
contact equations, not an assertion about a Hessian.

The full divergent $D$ channel has the exact identity
\begin{equation}
 P=(D-d_0)L_D^2+
 z^T\{G|_{D=d_0}+(D-d_0)jH_D\}z,
 \label{eq:c2-ba-D-five}
\end{equation}
where $0<d_0<D$ is real,
\[
 L_D=uv\xi-\eta+lR_1+ruvR_2'-arvT,
\]
and the only nonzero symmetric entries of $H_D$ are
\[
 (H_D)_{13}=tuv,\quad (H_D)_{15}=-tu,\quad
 (H_D)_{34}=rtuv,\quad (H_D)_{35}=-artv.
\]
For completeness the coefficient identity proving this is
\[
 \partial_DG=dd^T+\tfrac12(sk^T+ks^T)+jH_D,
\]
\[
 \begin{aligned}
 d&=(uv,-1,l,ruv,-arv)^T,\\
 s&=(0,0,(aB-u^2)v,-(B-v^2/t)u,u^2-av^2/t)^T,\\
 k&=(-2rtu,0,0,-r^2tu,ar^2t)^T.
 \end{aligned}
\]
The root syzygy $s^Tz=0$ and $d^Tz=L_D$ prove
\eqref{eq:c2-ba-D-five} including every normal entry. Since $Dj$ is
bounded, the correction is bounded and vanishes in local residue.

If $g$ is infinitesimal, also retain $(m-K)R_{mc}^2$, where $K>0$ is
real and
\[
 R_{mc}=-\xi+uv\eta+{1-g\over a}R_2'+huT,
 \qquad \tau=(m-K)g.
\]
The remaining exact matrix is $G|_{m=K}+\tau H_m$, before the
independent $D$ operation. In unscaled $R_2=R_2'/a$ coordinates the
only nonzero symmetric entries are
\[
 (H_m)_{14}=-1,\quad (H_m)_{25}=v,\quad
 (H_m)_{44}=2-g,\quad(H_m)_{45}=hu.
\]
This is the full centered-$m$ identity R5.2--R5.5. It commutes with
\eqref{eq:c2-ba-D-five}; neither correction depends on the other
coefficient. Here $a$ is a unit and $\tau$ is bounded.

Put $N=n+1/2$, $L=A+C+d_0+1/2$ and normalize
$\widehat\xi=\sqrt N\xi$, $\widehat\eta=\sqrt L\eta$.
Every coefficient of the resulting five-root matrix is bounded.
Here are the complete bounds needed from the WA table. Its scalar
entries involve $nr^2$, $Cl^2$, $Ah^2$, $\alpha_*$, the capped $m$,
and bounded $D$; its first normals involve $nr$, $Cl$, $Ah$ and these
same bounded products; the hard block uses only $n/N$, $A/L$, $C/L$
and bounded entries divided by $\sqrt N$ or $\sqrt L$. All are bounded
by \eqref{eq:c2-ba-unit-products}. If $a$ is infinitesimal, then $g$
is a unit and the only apparently singular scaled entry is
\begin{equation}
 {b_2-\alpha_*\over a},\qquad
 b_2=2mh+Cl^2+(Ah)h+(B^2\alpha_*+1/2)h^2.
 \label{eq:c2-ba-a-divisibility}
\end{equation}
The WE3 numerator vanishes exactly at $a=0$, for arbitrary
$B,h,r,t$. Cancel that factor in the actual source formula. Its remaining
denominator uses only $g,\Delta,t,1-hB$, all units here; the $h,l,r,j$
denominators have already canceled in the displayed products. The WE5
rewrites cover both normal rows. Thus this is joint regularity, including
all simultaneous openings, rather than separate one-variable limits.

At $u=v=0$ the residue splits into two $2\times2$ blocks and a scalar.
The $(\widehat\eta,R_1)$ block has determinant
\begin{equation}
 {A(C+d_0+1/2)\over A+C+d_0+1/2},\qquad
 {1\over\det}={1\over A}+{1\over C+d_0+1/2}.
 \label{eq:c2-ba-eta-determinant}
\end{equation}
Its upper-left entry is $1$, its other entries are bounded, and the
determinant is a positive unit. If $g$ is a unit, the
$(\widehat\xi,R_2')$ block is
\[
 \begin{pmatrix}
 1+m/N&-m/(a\sqrt N)\\-m/(a\sqrt N)&m/a^2
 \end{pmatrix},\qquad \det=m/a^2>0.
\]
If $g$ is infinitesimal, its determinant is
\[
 {K+2\tau_0-\tau_0^2\nu\over a_0^2}>0,
\]
where $\nu=\operatorname{st}(N^{-1})$, also when $N$ is unbounded.
Choose $K>2\tau_0^2+1$, since $0\le\nu\le2$.
The remaining $T$ pivot is $\alpha_*>0$. Therefore the complete bounded
matrix has strict residue. This proves every local independent test on
the unit-$e$ branch without any contact conductor.

\paragraph{The small-$e$ branch: source limits and physical Schur bound.}
Now $e$ is infinitesimal and $r_0=1$. The actual source formulas imply
\begin{equation}
 n,m,\alpha\in\mathrm I\cup\{0\},\quad g_0=1,\quad
 (Ah)_0=t_0/2,\quad (Cl)_0=(h_0+t_0)/2,\quad
 (Dj)_0=1/(2t_0),\quad
 (\alpha_*)_0=(1+a_0)(h_0+t_0)/2.
 \label{eq:c2-ba-small-e-products}
\end{equation}
In particular $A,C,D$ have positive real lower bounds, even if the
strong-$C$ hypothesis had not been imposed. From $j>0$ we have $t_0\ge l_0$.
Define
\[
 q=nr^2+(m+1/2)e^2+\alpha,\quad
 \mu=m-a^2q,\quad \lambda={\mu\over ta^2B},\quad
 \sigma={t\lambda aB\over\sqrt q}.
\]
The $a^2B$ factor in $\mu$ cancels exactly in the native rational
formula. Its canceled denominator at $B=0,r=1$, with $h$ free, is
$2t(a+1)$; hence
\begin{equation}
 \lambda_0={h_0+t_0\over2t_0}>0,\qquad \mu\in\mathrm I.
 \label{eq:c2-ba-free-h-lambda}
\end{equation}
This includes $a_0=0$ and does not assume $h$ infinitesimal.

The exact inequality \eqref{eq:c2-ba-physical} gives
$e/B\ge K_s/Q_s$, where $Q_s\sim1$ and
$K_s\sim(1+a)(h+t)$. Since $q\ge e^2/2+\alpha$, it follows that
\begin{equation}
 \operatorname{st}(\sigma^2)
 \le {a_0^2(h_0+t_0)\over
 2(1+a_0)\{1+(1+a_0)(h_0+t_0)\}}.
 \label{eq:c2-ba-free-h-sigma}
\end{equation}
If $e/B$ is unbounded the left side is zero. In all cases
\begin{equation}
 \operatorname{st}(1/2-\mu-\sigma^2)
 \ge {1+a_0+(h_0+t_0)(1+2a_0)\over
 2(1+a_0)\{1+(1+a_0)(h_0+t_0)\}}>0.
 \label{eq:c2-ba-free-h-schur}
\end{equation}
No division by $n$ occurs, so $n=0$ is retained.

\paragraph{The exact small-$e$ full-parent decomposition.}
Use the generic BA4--BA7 identity and its complete displayed $H_0$.
In particular take
\[
 E=tuW-vZ,\quad\eta=Y-huW,\quad\xi=X+ruZ,
 \quad F_v=(v^2-y)Z,\quad R_D=v(u^2-x)Z,
\]
\[
 J=(uv-\sqrt{xy})Z,\quad
 Q=\sqrt q(avW-uZ-a\xi),\quad
 N=\sqrt n\{r[(a/t)v^2-u^2]Z+u\xi\}.
\]
With $k=\sqrt x$ and $c=\sqrt{a/t}$ retain the whole nonnegative positive square term
\[
 T^2={2\lambda\over c}k^2uv(tcW-Z)^2,
 \qquad R=\sqrt\lambda\,tk(cv-u)W.
\]
Initially $u,v>0$; the literal axes and the other signs are restored below.
The generic identity has the exact circuit
\[
 R^2+T^2+\lambda J^2
 =\lambda\{t^2x[u^2+(a/t)v^2]W^2-4txuvWZ
                     +(u^2v^2+xy)Z^2\}.
\]
Let $\mathscr P=P-N^2-T^2$. Its entire $D$ channel is
\begin{equation}
 \mathscr P=D L_D^2+\mathscr P|_{D=0}-2Dj\,uF_v\xi,
 \qquad L_D=lE/t-rR_D-\eta+uv\xi.
 \label{eq:c2-ba-D-circuit}
\end{equation}
This identity holds before any source specialization. Indeed
the $D$ terms are
\[
 D(lE/t-\eta+uv\xi)^2+Dr^2R_D^2
 -2Dr(lE/t-\eta+uv\xi)R_D
 +2D(l/t-1-rx)uF_v\xi,
\]
whose sum is \eqref{eq:c2-ba-D-circuit} by $j=1+xr-l/t$.

Choose positive real caps $A_c<A$, $C_c<C$, $d_0<D$. Subtract the
following exact nonnegative positive square terms:
\begin{equation}
 N^2+T^2+(D-d_0)L_D^2
 +(A-A_c)D_uD_v\eta^2+(C-C_c)D_v\eta^2.
 \label{eq:c2-ba-excesses}
\end{equation}
The remainder is the full BA quadratic with $D$ formally set to $d_0$,
minus $2(D-d_0)j\,uF_v\xi$, and with precisely the two displayed
centered-$\eta$ squares removed. Every coefficient is bounded. The
apparent $A,C$ poles are now only $Ah,Cl$; for example
\[
 \begin{split}
 AD_uD_v(hE/t+\eta)^2-(A-A_c)D_uD_v\eta^2
 &=D_uD_v\{A_c\eta^2+2(Ah)E\eta/t+(Ah)hE^2/t^2\},\\
 CD_v(lE/t-\eta)^2-(C-C_c)D_v\eta^2
 &=D_v\{C_c\eta^2-2(Cl)E\eta/t+(Cl)lE^2/t^2\}.
 \end{split}
\]
Every other entry of BA5--BA7 uses the bounded products in
\eqref{eq:c2-ba-small-e-products}, $a,t,t^{-1}$, or the bounded
$\lambda,\sigma$. The correction in \eqref{eq:c2-ba-D-circuit} is
bounded because $Dj$ is bounded. There is no $a$, $h$, $l$, or $j$
inverse left in the matrix coefficients. Pure-source inverse factors
inside the separately retained positive square term $T^2$ do not need to
be bounded.

At $u=v=0$ write $d_+=d_0+1/2$ and use residues of the bounded source
letters on the right. The entire nontrivial hard block is exactly
\begin{equation}
 d_+(R_D-lE/t+\eta)^2
 +{t+hl\over2t^2}E^2-{h\over t}E\eta+(A_c+C_c)\eta^2.
 \label{eq:c2-ba-hard-completion}
\end{equation}
The remaining blocks are
\[
 Q^2-2\sigma_0Q\xi+\tfrac12\xi^2,
 \quad {h_0+t_0\over2t_0^2}F_v^2,
 \quad \lambda_0J^2,\quad R^2.
\]
The first is strict by \eqref{eq:c2-ba-free-h-schur}. To make
\eqref{eq:c2-ba-hard-completion} strict it suffices to choose the caps
so that $2(A_c+C_c)(t_0+h_0l_0)>h_0^2$. If $h_0,l_0>0$, the actual
finite source residues satisfy
\[
 A_0={t_0\over2h_0},\quad C_0={h_0+t_0\over2l_0},\quad
 2(A_0+C_0)(t_0+h_0l_0)-h_0^2
 ={t_0(t_0+h_0)\over h_0l_0}>0.
\]
Thus sufficiently close smaller real caps work. If $h_0=0$ or
$l_0=0$, respectively $A$ or $C$ is unbounded, so a sufficiently large
real cap works. This proves strictness also when $j$ is infinitesimal
or opens simultaneously with either endpoint. The full bounded-matrix
completion principle proves the local small-$e$ claim.

\paragraph{A single critical family supplies every weak-$m$ outer join.}
When $e$ is infinitesimal, or when $e$ is a unit and $a$ infinitesimal,
$m$ is infinitesimal. The finite coefficients $A,C,D$ or their
unbounded limits have the common critical source residues
\begin{equation}
 A_0={t_0\over2h_0},\qquad C_0={h_0+t_0\over2l_0},\qquad
 D_0={1\over2(t_0-l_0)},\qquad h_0+l_0=1,\quad t_0\ge l_0.
 \label{eq:c2-ba-critical-family}
\end{equation}
Each quotient is used only when its denominator is positive. Its zero
denominator means the corresponding coefficient is unbounded. We retain
the whole $\alpha$ positive square term. We retain the $n$ positive square term too, or cap it at
a positive real lower bound on the unit-$e$ branch. In both cases the
remaining $X$ mass from the signed core is at least $D_v/2$.

For $u$ infinitesimal and $v$ a bounded unit put $S_m=\sqrt m W$ and
$R_u=uW$. Rewrite the entire $m$ positive square term as
\[
 (vS_m-\sqrt mX+\sqrt m\,uvY-\sqrt m\,uZ)^2
\]
and replace every other $uW$ by $R_u$. These operations introduce no
inverse. The $S_m$ pivot is $v_0^2>0$ and the $X$ pivot is at least
$D_v/2>0$. The remaining $(R_u,Y,Z)$ residue is
\begin{equation}
 \begin{split}
 F_v^{\rm crit}={}&C_0(1+V)(R_u-Y)^2+A_0(1+V)Y^2
 +D_0(R_u-Y-v_0Z)^2\\
 &+\tfrac12\{(R_u-Y)^2+VY^2-2v_0R_uZ\},\qquad V=v_0^2>0.
 \end{split}
 \label{eq:c2-ba-outer-v}
\end{equation}
In the finite interior its $Z=0$ restriction is strict and
\[
 \det F_v^{\rm crit}
 ={V^2(h_0+t_0)^2(1+V)\over8h_0l_0(t_0-l_0)}>0.
\]
At $h_0=0<t_0-1$, the unbounded $A$ row imposes $Y=0$ and the
remaining determinant is $C_0D_0V^2>0$. At $l_0=0$, the unbounded
$C$ row imposes $R_u=Y$ and the determinant is
$V^2(t_0+1)/(4t_0)>0$. At $t_0=l_0$, $h_0>0$, the unbounded $D$
row imposes $R_u=Y+v_0Z$ and the determinant is
$V^2(1+V)/(4h_0l_0)>0$. At the only simultaneous endpoint
$h_0=0$, $t_0=l_0=1$, both $A,D$ are unbounded; on their common
kernel $Y=0$, $R_u=v_0Z$ the remaining form is $V^2Z^2/2>0$.
These exhaust the boundary of \eqref{eq:c2-ba-critical-family} because
$t_0>0$. The cap principle now supplies a bounded strict full matrix
in every case, with all positive original excesses restored.

For $u$ a bounded unit, or unbounded, and $v$ infinitesimal set
$p=1/u$, $K=pX+Z$, $B_v=vX$, and $D_p=1+p^2$. The complete identity is
\begin{equation}
 \begin{split}
 p^2P(1/p,v)={}&nD_pD_vX^2+m(pvW+vY-K)^2
 +CD_v(W-pY)^2+AD_pD_vY^2\\
 &+D J_p^2+\alpha D_p(pX+vY-K)^2\\
 &+\tfrac12\{J_p^2+(1-p^2v^2)K^2\\
 &\hspace{10mm}+D_p[v^2Y^2-2vYK+2pYB_v+2pvKB_v-p^2B_v^2]\},\\
 J_p={}&W+D_pB_v-pY-pvK.
 \end{split}
 \label{eq:c2-ba-outer-u-full}
\end{equation}
The $n,m,\alpha$ positive square terms are retained exactly. On writing
$L=W-pY$, the remaining residue is $K^2/2$ plus
\begin{equation}
 C_0L^2+D_pA_0Y^2+(D_0+1/2)(L+D_pB_v)^2
 +D_ppYB_v-\tfrac12D_pp^2B_v^2.
 \label{eq:c2-ba-outer-u}
\end{equation}
In the finite interior its $L,Y$ diagonal block is strict and its
$B_v$ Schur complement is $D_p(h_0+t_0)/(2t_0)>0$, since
\[
 {C_0(D_0+1/2)\over C_0+D_0+1/2}
 =\tfrac12+{1\over4A_0}={h_0+t_0\over2t_0}.
\]
The same four boundary kernels give the following positive transverse
Schur coefficients: $D_p/2$ when $h_0=0<t_0-1$;
$D_p(D_0+1/2)$ when $l_0=0$;
$D_pC_0$ when $t_0=l_0$, $h_0>0$;
and $D_p/2$ at $h_0=0$, $t_0=l_0=1$.
Thus the cap principle proves every such chart, including literal $p=0$.
There is no inverse of $v$ or of a source opening in these roots.

If $e$ is a unit and $a$ is a unit, $m$ has a positive real lower
bound instead. For finite tests with only one nonzero residue, cap
$n,m,A,C,D$ at positive real lower bounds and restore their exact
excess positive square terms. At $v_0=0$, $u_0\ne0$ a strict selected residue is
\[
 n_cD_uX^2+A_cD_uY^2+(C_c+D_c+1/2)(u_0W-Y)^2
                         +\tfrac12(X+u_0Z)^2.
\]
At $u_0=0$, $v_0\ne0$ it is
\[
 (n_c+1/2)D_vX^2+m_c(v_0W-X)^2
 +(A_c+C_c+1/2)D_vY^2+D_c(Y+v_0Z)^2.
\]
For both finite testing residues nonzero, apply the common
all-noncommuting-residues theorem to the unchanged physical parent.
Its arbitrary-coefficient-scale version is the only remote physical
theorem imported here; the cap arguments above do not assert that
an arbitrary clipped source is physical.

\paragraph{All remaining projective charts.}
On the unit-$e$ branch, $n,A,C,D$ have positive real lower bounds.
After capping them and retaining all excesses and the whole $m,\alpha$
positive square terms, the three exact homogeneous templates R4E14--R4E16 have strict
residues
\[
 n_cD_vX^2+C_cD_vW^2+A_cD_vY^2
 +(D_c+1/2)(W+v_0X)^2+\tfrac12(v_0Y-Z)^2
 \quad(u=\infty),
\]
\[
 (n_c+1/2)D_uX^2+C_c(R_u-Y)^2+(A_c+1/2)D_uY^2
 +D_c(u_0X-Z)^2\quad(v=\infty),
\]
\[
 (n_c+1/2)X^2+C_cW^2+A_cY^2+D_c(X-P_\infty)^2
 +\tfrac12(Y-Q_\infty)^2\quad(u=v=\infty),
\]
where $R_u=uW$, $P_\infty=Z/u$, $Q_\infty=Z/v$ are legal roots;
the displayed matrices are strict on their respective root lists.
Their entire homogeneous polynomials have bounded coefficients.

On the small-$e$ branch, $m$ is infinitesimal and $n$ is retained whole.
The $u$-infinite, $v$-small chart is already
\eqref{eq:c2-ba-outer-u-full} at $p_0=0$. For $u$ infinite and $v$ a
bounded unit, the selected residue is
\[
 C_cD_vW^2+A_cD_vY^2+(D_c+1/2)(W+v_0X)^2
 +\tfrac12(v_0Y-Z)^2\succ0.
\]
For $v$ infinite and $u$ bounded use $S_m=\sqrt mW$ and $R_u=uW$.
The complete homogeneous polynomial obtained from
\eqref{eq:c2-ba-parent} has strict selected residue
\[
 S_m^2+\tfrac12D_uX^2+C_c(R_u-Y)^2
 +(A_c+1/2)D_uY^2+D_c(u_0X-Z)^2\succ0.
\]
For two infinite tests put $p=1/u$, $q_v=1/v$,
$P_\infty=pZ$, $Q_\infty=q_vZ$. The exact half-core is
\[
 \tfrac12\{(q_vW+X-pq_vY-P_\infty)^2
 +(1+p^2)(Y-Q_\infty)^2
 +(1+q_v^2)(p^2X^2+2XP_\infty-P_\infty^2)\}.
\]
Together with the capped $A,C,D$ positive square terms its residue is
\[
 C_cW^2+(D_c+1/2)X^2-2D_cXP_\infty+D_cP_\infty^2
 +A_cY^2+\tfrac12(Y-Q_\infty)^2\succ0;
\]
the $X,P_\infty$ determinant is $D_c/2>0$. All omitted original
positive square terms are exact nonnegative squares. Reciprocal deviations are pure
testing infinitesimals, so the same completion principle applies. The
literal endpoints are the homogeneous polynomial values, not a division
by a zero coordinate.

Finally, the literal finite axes obey the complete identities
\begin{equation}
 P|_{u=0}=(n+1/2)D_vX^2+m(vW-X)^2
 +(C+A+1/2)D_vY^2+D(Y+vZ)^2+\alpha(vY-Z)^2,
\label{eq:c2-literal-u-axis}
\end{equation}
\begin{equation}
 P|_{v=0}=nD_uX^2+(m+1/2)(X+uZ)^2
 +(C+D+1/2)(uW-Y)^2+AD_uY^2+\alpha D_uZ^2.
\label{eq:c2-literal-v-axis}
\end{equation}
The sign gauges $(u,Y,Z)\mapsto(-u,-Y,-Z)$ and
$(v,X,Z)\mapsto(-v,-X,-Z)$ preserve the parent. Applying the real
identity to the real and imaginary parts proves the complex-state
statement. The two local branches, the two finite one-unit-angle
branches, the noncommuting finite branch, and the displayed homogeneous
charts exhaust every independent source and test type. This proves the
proposition.

\subsection{State signs and projective endpoints}
\label{component:pf77de02d80c4}

\label{pf77de02d80c4:sec:rederived-state-join}

We reduce the complete weak-C2 parent to a positive state-ratio chart
and evaluate its projective measurement endpoints. The subsequent local
certificates treat the remaining positive chart.

Let $K_c,K_t\supset\mathbb R$ be real closed fields, occupying separately
labelled tensor factors.  All quantities $u,v,W,X,Y,Z$ below belong to $K_t$;
all six coefficients $n,m,A,C,D,\alpha$ belong to $K_c$ and are nonnegative.
Write $D_u=1+u^2$, $D_v=1+v^2$, and define the complete parent
\begin{align}
 F={}&nD_uD_vX^2+m\{v(W+uY)-X-uZ\}^2
       +AD_uD_vY^2+CD_v(uW-Y)^2\\
    &+D(uW+uvX-Y-vZ)^2+\alpha D_u(vY-Z)^2+K,\tag{SJ1}\\
 K={}&\tfrac12(uW+uvX-Y-vZ)^2
       +\tfrac12D_u(vY-Z)^2
       +\tfrac12D_v(X+uZ)^2-\tfrac12D_uD_vZ^2.\tag{SJ2}
\end{align}
Products involving source and testing quantities mean tensor products.
For example, $nD_uD_vX^2$ means
$n\otimes(D_uD_vX^2)$; a testing-only term means $1\otimes K$.
The positive square terms in (SJ1) are sums of squares in
$K_c\otimes_{\mathbb R}K_t$, since each coefficient and each testing
multiplier has a square root in its own real closed field.

\paragraph{Sign gauges preserve the entire parent.}
For $\varepsilon,\eta\in\{1,-1\}$, the substitution
\[
 (u,v,W,X,Y,Z)\longmapsto
 (\varepsilon u,\eta v,W,\eta X,\varepsilon Y,
  \varepsilon\eta Z)                                      \tag{SJ3}
\]
preserves every summand of (SJ1), including (SJ2).  Indeed the roots of the
$m,C,D,\alpha$ positive square terms acquire, respectively, the signs
$\eta,\varepsilon,\varepsilon,\varepsilon\eta$, and the other roots also
change only by a sign.  Thus $u,v$ may both be taken nonnegative without
altering the source coefficients or the full expression.

\paragraph{An exact positive hard block.}
Set
\[
 t=\frac{X+uZ}{D_u},\qquad r=\frac{Y-uW}{D_u},\qquad
 w=\frac{W+uY}{D_u},\qquad z=\frac{Z-uX}{D_u}.             \tag{SJ4}
\]
The inverse transformation is
$X=t-uz$, $Z=z+ut$, $Y=r+uw$, $W=w-ur$.
Let $H=1+u^2+v^2$ and $d=H^2-u^2v^2$.  Both
$H-uv=1+(u-v)^2+uv$ and $H+uv$ are positive for $u,v\geq0$;
in fact $H>|uv|$ for arbitrary $u,v$.  Consequently $d>0$.
Define
\[
 \mathsf A=\frac12\begin{pmatrix}H&-uv\\-uv&H\end{pmatrix},\qquad
 \mathsf B=\frac12\begin{pmatrix}-u^2v&-uv^2\\uv^2&u^2v\end{pmatrix}.
\]
The matrix $\mathsf A$ is positive definite.  Direct expansion and square
completion give the exact identity
\begin{align}
 \frac K{D_u}
 ={}&\left(\binom tr+\mathsf A^{-1}\mathsf B\binom wz\right)^{\!T}
       \mathsf A
       \left(\binom tr+\mathsf A^{-1}\mathsf B\binom wz\right)\\
    &+\frac{D_uD_v}{2d}
       \{u^2v^2(w^2+z^2)-2uvH wz\}.                    \tag{SJ5}
\end{align}
The rational matrix used in the shift is explicitly
\[
 \mathsf A^{-1}\mathsf B=
 \frac1d\begin{pmatrix}
 -u^2v(1+u^2)&-uv^2(1+v^2)\\
 uv^2(1+v^2)&u^2v(1+u^2)
 \end{pmatrix}.                                         \tag{SJ6}
\]
For any $a,b\in K_t$, the hard quadratic form is
\[
 \binom{a}{b}^{\!T}\mathsf A\binom{a}{b}
   =\frac{H-uv}{4}(a+b)^2+\frac{H+uv}{4}(a-b)^2.         \tag{SJ7}
\]
This also gives an explicit square decomposition of the first term of
(SJ5), with square roots taken only in $K_t$.

If $u,v\geq0$ and $wz\leq0$, the second term of (SJ5) is nonnegative.
It too is a square in $K_t$ (or can be split into its three nonnegative
terms).  Hence $K\geq0$, and (SJ1) is a tensor sum of squares immediately.
In particular the axes $w=0$ and $z=0$ are included, without division by
either state coordinate.  The same identity, without sign gauges, proves
the sharper implication $K<0\Longrightarrow uvwz>0$.

\paragraph{Reduction of the remaining states.}
The only remaining case after (SJ3) is $K<0$, $u,v>0$, $wz>0$.
Then $w\ne0$.  Divide all four state coordinates by $w$.  Equation (SJ4)
now has $w'=1$ and $z'=z/w>0$.  The ratio $z'$ is arbitrary: it may be
infinitesimal, bounded, or unbounded in $K_t$.  No estimate in this
reduction replaces $z'$ by a bounded coefficient.

The parent is homogeneous of degree two in the state coordinates, so
\[
 F(u,v;W,X,Y,Z)=w^2F(u,v;W/w,X/w,Y/w,Z/w).               \tag{SJ8}
\]
If the parent on the right has a certificate $\sum_i h_i^2$ in
$K_c\otimes_{\mathbb R}K_t$, then the parent on the left has certificate
$\sum_i((1\otimes w)h_i)^2$.  This proves that certificates for the full
parent with $u,v>0$, $w=1$, and arbitrary positive $z$ suffice for all
finite affine testing parameters and all real states.

\paragraph{Literal projective endpoints.}
For clarity, the projective endpoints can be checked without a limiting
argument.  Homogenize $u=U_1/U_0$, $v=V_1/V_0$, and put
$D_U=U_0^2+U_1^2$, $D_V=V_0^2+V_1^2$, with neither pair identically zero.
Multiplication by $U_0^2V_0^2$ turns (SJ2) into
\begin{align}
 K_h={}&\tfrac12(U_1V_0W+U_1V_1X-U_0V_0Y-U_0V_1Z)^2\\
       &+\tfrac12D_U(V_1Y-V_0Z)^2
        +\tfrac12D_V(U_0X+U_1Z)^2-\tfrac12D_UD_VZ^2.     \tag{SJ9}
\end{align}
At the four axes this is, exactly,
\begin{align}
 U_0=0:\quad &K_h=\tfrac12U_1^2\{(V_0W+V_1X)^2+(V_1Y-V_0Z)^2\},\\
 V_0=0:\quad &K_h=\tfrac12V_1^2D_U(X^2+Y^2),\\
 U_1=0:\quad &K_h=\tfrac12U_0^2D_V(X^2+Y^2),\\
 V_1=0:\quad &K_h=\tfrac12V_0^2\{(U_1W-U_0Y)^2+(U_0X+U_1Z)^2\}.
                                                               \tag{SJ10}
\end{align}
All other homogenized positive square terms in (SJ1) retain nonnegative coefficients
and are sums of squares.  Thus the endpoints are proved directly,
including the intersections of these axes.  Equations (SJ3)--(SJ10)
establish the stated full-state, sign, and endpoint reduction.  They do
not establish the still-required positive-ratio chart certificate.

\subsection{Source coordinates and the two-square decomposition}
\label{component:p667905169990}

\label{p667905169990:sec:r2-reconstructed-source-entry}

\paragraph{Actual source and conventions.}
The native source formulas give the following coordinate identities
and boundedness estimates. They provide the source parameters used
in the subsequent R2 contact and state-normal certificates.

Use the original weak-C2 native source $(a,B,t,r,h)$, with
$0<r,h<1$, $a,B,t>0$, the original strict source gaps, and $n\geq0$.
In R2, $t$ is positive unbounded while
\[
 B,\qquad z=aB,\qquad y=tB
\]
are positive infinitesimals.  There is no bound on $a$ and no lower
bound on $r$, $h$, $1-h$, or $e=1-r$.  Put $l=1-h$ and $S=1/t$.
Define the actual native arrays
\[
 c=1+zr,\quad s_c=1+B(l+ar),\quad
 R=1+hz,\quad g=R-ae>0,\quad j_D=c-l/t>0,
\]
\begin{align*}
 \Delta={}&a+(1+z)+(1+z)B+2z(1+z)r
          +a^3B^2r^2+2a^2Br\\
          &+tB\{1+a+z+a^2Br(2-r)\},\\
 M={}&(1+z)(1+ar)(h+t)-a(1+y)e\,l.
\end{align*}
The six weights used throughout this section are exactly
\[
\begin{aligned}
 n&=\frac{(1+ar)c s_c-r\Delta}{2r\Delta},&
 m&=\frac{a^2(1+y)e c s_c}{2g\Delta},\\
 C&=\frac{cM}{2l\Delta},&
 A&=\frac{tc(1+ar)s_c}{2h\Delta},\\
 D&=\frac{(1+z)(1+ar)(1-atBr)s_c}{2t j_D\Delta},&
 \alpha&=\frac{B^2(1+ar)M}{2(1-hB)\Delta}.
\end{aligned}                                                     \tag{RER1}
\]
All denominators are pure source.  The original positivity assumptions
are retained without adding a rate comparison.

\paragraph{Physicality controls the apparently unbounded source scale.}
The exact numerator identity is
\[
 (1+ar)(1+zr)\{1+B(l+ar)\}-r\Delta=eQ-BK,                \tag{RER2}
\]
where
\[
 K=(a+1)(h+t)(1+z),
\]
\[
 Q=1+B+y+z(h+t+2r)
       +z^2\{t(1+er)+h(1+r)\}+Bz(h+t).                 \tag{RER3}
\]
Every term of $Q$ is nonnegative.  Since $r\Delta>0$, original
physicality is exactly $eQ\geq BK$, including equality when $n=0$.
Put $\kappa=ay=tz$.  Before assuming that $\kappa$ is bounded, direct
termwise inspection of (RER3), using $B,y,z$ infinitesimal, gives
$Q\leq C_0(1+\kappa)$ for a fixed positive real $C_0$.
The strict gap $g>0$ gives $ae<R\sim1$, so multiplication of the
physicality inequality by $a$ yields
\[
 \kappa(a+1)(1+hS)(1+z)\leq aeQ<RQ
                                      \leq C_1(1+\kappa)           \tag{RER4}
\]
for a fixed positive real $C_1$.

If $a$ is bounded, $\kappa=ay$ is infinitesimal.  If $a$ is unbounded,
move $C_1\kappa$ to the left in (RER4).  Since
$a+1-C_1\geq(a+1)/2$, one obtains
$\kappa\leq2C_1/(a+1)$, again infinitesimal.  In either case,
(RER4) also bounds $(a+1)\kappa$.  Returning to (RER3) and then to
$eQ\geq BK$ proves
\[
 ay\in\mathrm I,\qquad a(a+1)y=O(1),\qquad Q\sim1,
 \qquad \frac{(a+1)y}{e}=O(1).                         \tag{RER5}
\]
No division by $n$ occurred; the argument includes the face $n=0$.

\paragraph{Uniform source normalizations.}
Put
\[
 k=\frac{1+ar}{1+a},\qquad p=\frac{ae}{R}\in(0,1),
 \qquad \mu=m(1-p)=\frac{mg}{R},\qquad\delta_n=nr.
\]
The quantity $k$ is a positive bounded unit.  Indeed $k\leq1$; if
$a\leq1$, then $k\geq1/(1+a)\geq1/2$.  If $a>1$, use
$ae<R=1+hz$ to get $k>(a-hz)/(1+a)$, which has a positive real lower
bound.  Direct termwise division of the displayed polynomial
$\Delta$ by $1+a$ gives $\Delta/(1+a)\sim1$.  Also
\[
 \frac{M}{t(1+ar)}\sim1:
\]
the first term has ratio $(1+z)(1+hS)\sim1$, while the absolute
ratio of the second is at most $R(1+y)S/(1+ar)$ and hence is
infinitesimal.  Since $c,s_c,j_D,1-hB\sim1$, (RER1) now gives
\[
 \frac{SC\,l}{k}\sim\frac12,\quad
 \frac{SA\,h}{k}\sim\frac12,\quad
 \frac{D}{Sk}\sim\frac12,\quad
 \frac{\alpha}{Sy^2(1+a)k^2}\sim\frac12.               \tag{RER6}
\]
For the $D$ relation use $atB=ay=\kappa\in\mathrm I$ from
(RER5).  Thus all original $h$ and $l$ openings are present.

Two exact bounded products remove the remaining nonuniform source
poles:
\[
 mg=\frac{a^2e(1+y)c s_c}{2\Delta},\qquad
 \delta_n=\frac{eQ-BK}{2\Delta}.                       \tag{RER7}
\]
The quantities $mg$, $mg/a$, $mg/a^2$, $\mu$, $\delta_n$, and
$nr^2$ are bounded.  For example
$a^2e/(1+a)\leq R a/(1+a)=O(1)$, while
$0\leq\delta_n\leq eQ/(2\Delta)=O(1)$.
These estimates do not imply that $m$ or $n$ themselves are bounded.

\paragraph{Exact extraction from the entire parent.}
Let testing $u,v,W,X,Y,Z$ be arbitrary, independently of the source;
put $D_u=1+u^2$, $D_v=1+v^2$.  The complete original quarter-parent is
\begin{align*}
 \mathcal P={}&m(vW-X+uvY-uZ)^2+nD_uD_vX^2
     +CD_v(uW-Y)^2+AD_uD_vY^2\\
 &+(D+\tfrac12)(uW+uvX-Y-vZ)^2
     +(\alpha+\tfrac12)D_u(vY-Z)^2\\
 &+\tfrac12D_v(X+uZ)^2-\tfrac12D_uD_vZ^2.
\end{align*}
Define the mixed rows
\[
 q_a=v(W+uY),\qquad N=X+uZ-pq_a,\qquad \xi=X+ruZ.
\]
Then the following identity retains both actual positive square terms,
without setting either normal to zero:
\[
 \mathcal P=mpN^2+nD_uD_v\xi^2+G,                       \tag{RER8}
\]
where
\begin{align*}
 G={}&\mu\{N^2-2q_aN+(1-p)q_a^2\}
        +\delta_nD_uD_v(-2uZX-ru^2Z^2)\\
 &+CD_v(uW-Y)^2+AD_uD_vY^2\\
 &+(D+\tfrac12)(uW+uvX-Y-vZ)^2
        +(\alpha+\tfrac12)D_u(vY-Z)^2\\
 &+\tfrac12D_v(X+uZ)^2-\tfrac12D_uD_vZ^2.
\end{align*}
It follows simply by expanding
$m\{(1-p)q_a-N\}^2-mpN^2$ and
$n(X^2-\xi^2)=nr(-2uZX-ru^2Z^2)$.
Both retained positive square terms are finite separated SOS: their source weights
are nonnegative, $D_uD_v$ is pure testing, and $N,\xi$ are allowed
mixed rows.  When $n=0$ the second positive square term vanishes literally.
The exact identity and all estimates above are the exact
structural conclusion; no positivity conclusion about $G$ is asserted.

\subsection{The C1--C2 boundary identity}
\label{component:pc024f5300388}

\label{pc024f5300388:sec:r13-source-overlap}

This section records an exact source-only identity. It does not by itself
admit any new independent testing domain. Its purpose is to identify the
same physical boundary in the two remaining source charts and to retain
the nonnegative slack when that boundary opens.

Use the original weak-C2 variables $B=H+L$, $z,y>0$ and $0<r<1$,
with the unchanged source gaps and nonnegative $n$. Write $e=1-r$,
$K=1+z$, $T_c=B+zr$, and let $\delta$ be the original positive
C2 denominator. Define the positive source polynomial
\[
 P_n=B+B^2+2Bzr+BHz+yB(1+z)+Hrz^2+yz^2re.
 \tag{OP1}
\]
Then the following identity holds on the entire original source:
\[
 \boxed{\quad eP_n=(1+z)(H+y)(B+zr)+2r\delta n.\quad}
 \tag{OP2}
\]
Every factor in the first summand is positive, and the second summand
is nonnegative, including literal $n=0$. In particular the pure-source
weights
\[
 q={ (1+z)(H+y)(B+zr)\over eP_n}>0,
 \qquad s_n={2r\delta n\over eP_n}\ge0
 \quad\hbox{satisfy}\quad q+s_n=1.                   \tag{OP3}
\]
Both square roots are legitimate source-field elements. Neither weight
is assumed to have nonzero residue.

\begin{proof}
The unchanged formula for $n$ is
\[
 2r\delta n=(B+zr)(1+zr)(1+zr+B-H)-r\delta.
\]
Substitution of the printed polynomial for $\delta$, followed by
collecting terms, gives OP2. This is a polynomial identity, so no
source boundary is obtained by continuity and no opening has been
inverted except the displayed positive pure-source factors.
\end{proof}

\paragraph{Exact correspondence of the source contacts.}
Use a subscript $1$ for native C1 notation in this paragraph. The map
\[
 x_1=z,\qquad \eta_1={y\over B},\qquad
 \zeta_1={B\over z},\qquad r_1={H\over B},\qquad
 \theta_1=e,
 \tag{OP4}
\]
with $k_1=\sqrt z$ and the native parameter
$\Delta_1=B/(z\sqrt y)$, gives
\[
 (u_c,v_c;W_c,X_c,Y_c,Z_c)
 =\left(\sqrt z,\sqrt y;{B\over\sqrt{zy}},
                       -r\sqrt z,{H\over\sqrt y},1\right)
\]
from the original native C1 contact formula. Thus it is the same
contact, including all state coordinates. Direct substitution in the
native positive polynomial $P_d$ gives $BP_d=P_n$.
The native C1 beta coefficient at this formal contact therefore obeys
\[
 \beta_1=-{r\delta n\over eP_n}=-{s_n\over2},
 \qquad \beta_1+\tfrac12={q\over2}>0.                 \tag{OP5}
\]
For $n=0$, the original C2 parent has
$\lambda=n=\beta=0$. It is literally an original C1 parent,
up to an optional positive normalization of $C$. The native source
entry and uniqueness at stationary contact identify its coefficients
with OP4. For $n>0$, OP5 makes the formal C1 beta negative. Hence
OP4 does \emph{not} map that source to the nonnegative C1 event cone,
and no physical C1 all-tests theorem may be imported there.

\section{Two small actual angles: the R1 sector}
\subsection{Scalar completion at unbounded source ratio}
\label{component:p9e725a065583}

\label{p9e725a065583:sec:reconstructed-r1-qo-source-scalar}

\paragraph{Scope of the dictionary and budgets.}
The identities below use the actual native weak-C2 stationary coefficients,
with $d=1$, $\lambda_{\rm event}=\beta=0$, all strict native gaps and
$n\ge0$. Source quantities lie in $K_c$ and testing quantities in a
separately labelled real closed extension $K_t$; all inverses and radicals
are pure source or pure testing. The rational identities have no source-size
restriction. Their R1 estimates assume $B,T,x,y$ infinitesimal; the later
RL and EN certificates establish their own replacement estimates.

For the native source write $a_s,B_s,t_s,r,h$, and put
$c_s=1+a_sB_sr$, $R_s=1+ha_sB_s$, $K_s=1+a_sB_s$.
The formerly used auxiliary ratio is simply
\[
 q_{\rm old}=R_s^2/(a_st_sK_s^2).\tag{RSR1}
\]
No size restriction on $q_{\rm old}$ or $q_z$ is imposed in this
refined dictionary. The actual full-R1 estimates, including arbitrary
positive $q_z=jT$, are justified by R1N1--R1N11.

\subsubsection*{The exact compact source and its uniform bounds}

Define pure-source quantities
\[
\begin{gathered}
 j=\frac{a_sc_s}{R_s},\quad B=\frac{B_sR_s}{c_s},\quad
 T=\frac{t_sc_s}{R_s},\\
 p=\frac{a_s(1-r)}{R_s},\quad H=\frac{1-h}{t_sc_s},\\
 x=jB,\quad y=TB,\quad q_z=jT .
\end{gathered}
 \tag{RSR2}
\]
The native gaps give $0<p<\min(1,j)$ and $0<H<1$.
In the R1 estimates $B,T,x,y$ are positive infinitesimals, while
$j$ and $q_z=jT$ are arbitrary positive. The exact dictionary itself
also applies at the different source scales used in RL and EN. Direct inversion gives
\[
 r=\frac{j-p}{j+px},\quad h=\frac{1-HT}{1+HTx},\quad
 a_s=\frac{j+px}{1+HTx},\quad
 B_s=\frac{B(1+HTx)}{1+pB},\quad
 t_s=\frac{T(1+pB)}{1+HTx}.
 \tag{RSR3}
\]
In particular the compact variables describe the actual source,
including its $r\to0$, $p\to1$, $H\to0$, and $H\to1$ openings.
Set
\[
 K=j+1-p,\qquad E=1+B(HT+p-1),\qquad
 W_0=(1-H)j+1-p.
 \tag{RSR4}
\]
Then $\tfrac12(j+1)<K<j+1$, $E\sim1$, and $W_0>0$.
The following polynomials contain the entire stationary dictionary:
\begin{align*}
\mathcal P={}&j+1+B\{(j+1)^2-p\}\\
 &+TB\{Hj+j+1+B[Hj(2j+2-p)+p(2j+1-p)]\}\\
 &+HjT^2B^2(1+Bp+BHj),\\
\mathcal L={}&K+T(W_0-H)+BT[pK+Hj(1-p)]+BHT^2(W_0-1),\\
\mathcal N={}&p+B(-j^2+3jp-j-p^2+p)\\
 &+BT(Hj^2+Hj-j^2+jp-j+p)\\
 &+B^2T(2Hj^2p-Hj^2-Hjp^2+2Hjp
             -2j^2p+3jp^2-jp-p^3+p^2)\\
 &+B^2HT^2j(Hj-j+p)(1+Bp).
 \tag{RSR5}
\end{align*}
With $\mu=m(1-p)$, the exact six original stationary weights are
\[
\begin{aligned}
 \mu&=\frac{jp(1+y)(1+x)(1+Hy)}{2\mathcal P},&
 n&=\frac{\mathcal N}{2(j-p)\mathcal P},\\
 C&=\frac{\mathcal L}{2HT\mathcal P},&
 A&=\frac{T(1+x)(1+Hy)(K+Hjy)}{2(1-HT)\mathcal P},\\
 D&=\frac{(1+Hy)(K+Hjy)\{1-y[(1-H)j-p]\}}
             {2T(1-H)\mathcal P},&
 \alpha&=\frac{B^2(K+Hjy)\mathcal L}{2E\mathcal P}.
\end{aligned}
 \tag{RSR6}
\]
For example, the original stationary denominator and its source mass
satisfy
\[
 \Delta=\frac{(1+x)\mathcal P}{(1+Bp)(1+HTx)},\qquad
 M=\frac{(1+x)\mathcal L}{(1+HTx)^2}.
\]
The two native strict gaps become
\[
 g=\frac{(1-p)(1+x)}{1+HTx},\qquad
 j_{\rm gap}=\frac{(1-H)(1+x)}{1+Bp}.
\]
These equations verify the source origin of RSR6 directly; no new
stationary solution has been substituted for the original one.

Every term after $j+1$ in $\mathcal P$, divided by $j+1$, is
infinitesimal: use $B(j+1)=B+x\in\mathrm I$, $T\in\mathrm I$,
$p<1$, and $H<1$. The same termwise estimate, using
$K>(j+1)/2$, applies to $\mathcal L-K$. Thus
\[
 \mathcal P\sim j+1,\qquad \mathcal L\sim K.
 \tag{RSR7}
\]
These are relative equivalences, valid also when $j$ is infinitesimal
or unbounded. They show, in particular, that $A/T$, $THC$, and
$T(1-H)D$ are positive bounded units, $\mu/p$ is bounded, and
$h_0=\mu+1/2$ is a positive bounded unit. Moreover $\alpha/B$ is
infinitesimal. Only the displayed source openings can make $m,n,C,D$
unbounded.

The original physicality condition gives a useful uniform restriction.
Regroup the exact polynomial $\mathcal N$ as
\[
\mathcal N=p\{1+B(3j+1-p)+BT(j+1)\}
 -Bj(j+1)\{1+(1-H)T\}+\mathcal R.
 \tag{RSR8}
\]
Here $\mathcal R$ is precisely the last two lines of $\mathcal N$
in RSR5. Since $0<p<\min(1,j)$, each coefficient in its $B^2T$
line is bounded in absolute value by a fixed multiple of $j(j+1)$.
The $B^2HT^2$ line has the same bound, multiplied by the bounded
factor $T(1+Bp)$. Consequently
\[
 |\mathcal R|\leq C_0BT\,Bj(j+1)
\]
for a fixed real $C_0$. All corrections multiplying $p$ and
$Bj(j+1)$ in RSR8 are infinitesimal. As $n\geq0$ is exactly
$\mathcal N\geq0$, it follows that
\[
 \eta_*:=\frac{Bj(j+1)}p\ \hbox{is bounded},\qquad
 \eta_*\in\mathrm I\ \Longrightarrow\
 n\sim\frac{p}{2(j-p)(j+1)}>0.
 \tag{RSR9}
\]
Indeed $\mathcal N=p(1+\varepsilon_1)-Bj(j+1)(1+\varepsilon_2)$
with $\varepsilon_i\in\mathrm I$. This proves both assertions without
dividing by $n$ and includes literal $n=0$ in the bounded-$\eta_*$
alternative. It also gives $nr$ bounded.

\subsubsection*{The complete graph and its actual contact circuit}

The quarter-parent, before any positive square term is removed, is
\begin{align*}
\mathscr P={}&m\{vW+uvY-X-uZ\}^2
 +CD_v(uW-Y)^2+AD_uD_vY^2\\
&+(D+\tfrac12)(uW+uvX-Y-vZ)^2
 +(\alpha+\tfrac12)D_u(vY-Z)^2+nD_uD_vX^2\\
&+\tfrac12D_v(X+uZ)^2-\tfrac12D_uD_vZ^2,
\qquad D_u=1+u^2,\ D_v=1+v^2.
 \tag{RSR10}
\end{align*}
The source-free rotated state coordinate is $w=(W+uY)/D_u$.
On its nonzero chart set $w=1$ and $z=(Z-uX)/D_u$. Define
the two complete state normals
\[
 N=X+uZ-pD_uv,\qquad M=uW-Y-HD_uvz.
\]
Their exact inverse is
\[
\begin{aligned}
 X&=N/D_u+pv-uz,& Z&=z+uN/D_u+puv,\\
 Y&=u-M/D_u-Hvz,& W&=1+uM/D_u+uHvz.
\end{aligned}
 \tag{RSR11}
\]
The scalar studied here is exactly
$F_{\rm full}=(\mathscr P/D_u)|_{N=M=0}$.
No $n$ positive square term, $m$ positive square term, or state-normal term is silently removed.
Substitution of RSR11 into RSR10 gives the complete full-parent
normal expression RFN7 in Section~\ref{component:p261f71af685d}, with
$N_0=N/D_u$, $M_0=M/D_u$ and the retained positive square term
$mpD_u^2N_0^2$ restored.

Use the separated rows
\[
 U=u/\sqrt x,\quad V=v/\sqrt y,\quad Z_*=z/\sqrt{q_z},\qquad
 Q_*=U-VZ_*,\quad B_*=V-UZ_*,\quad R_*=Z_*-UV,
\]
\[
 \Psi=U^2+V^2+Z_*^2+U^2V^2Z_*^2-4UVZ_*.
 \tag{RSR12}
\]
These are rows; $U,V,Z_*$ are not assumed bounded coefficients.
The actual contact is $(U,V,Z_*;N,M)=(1,1,1;0,0)$.
The exact native formulas make the complete parent value and all five
first derivatives vanish there. On the positive core,
\[
 \Psi=(U-V)^2+Z_*^2(1-UV)^2+2UV(1-Z_*)^2
 \tag{RSR13}
\]
is a legal separated sum of squares: the last root is
$\sqrt2\,(uv)^{1/2}(xy)^{-1/4}(1-Z_*)$, with its radicals in
the two separate real closed fields.

Let the ten positive monomial coefficients and the seven circuit
masses be exactly RB6--RB8. In particular
\[
 M_a=Ax,\quad M_b=\{\mu(1-p)+(n+\tfrac12)p^2\}y,\quad
 M_c=\alpha q_z,\quad M_{ac}=nj^2y,
\]
\[
 f_c=M_b-M_{ac},\quad q_c=M_a-f_c,\quad
 e_c=(n+\tfrac12)p^2y^2-g_c,\quad d_c=M_{b^2c}-k_c.
\]
The original contact balances give the exact scalar identity
\begin{align*}
 F_{\rm full}={}&q_cQ_*^2+M_{ac}B_*^2+r_cR_*^2
  +d_cV^2Q_*^2+e_cV^2B_*^2\\
 &+f_c\Psi+g_cV^2\Psi+k_cZ_*^2(V^2-1)^2.
 \tag{RSR14}
\end{align*}
The identities RB7 and the two odd coefficient balances are independently
checked on the actual source. Thus RSR14 is a finite polynomial identity,
not a Taylor approximation to the contact.

Here is an explicit source-only formula for its radial mass. Put
$\mathcal F=K+T\mathcal F_1+T^2\mathcal F_2$, where
\begin{align*}
\mathcal F_1={}&-2BHjp+BHj+2Bjp-2Bp^2+2Bp\\
 &-Hj-Hp-H+j+1-p^2/j+p/j,\\
\mathcal F_2={}&-B^2H^2jp+B^2Hjp-B^2Hp^2+B^2Hp\\
 &-BH^2j-BH^2p+BHj-BHp^2/j+BHp/j.
 \tag{RSR15}
\end{align*}
Every displayed term divided by $K$ is bounded, using $p/j<1$;
terms containing $Bj$ are even infinitesimal. Hence $\mathcal F\sim K$.
Direct source substitution yields
\[
\begin{aligned}
 f_c&=\frac{Tx^2\mathcal F}{2\mathcal P}>0,&
 q_c&=Ax-f_c\sim\frac{TxK}{2(j+1)}>0,\\
 r_c&=\frac{B^2jT(1-p)(1+x)(1+Hy)\mathcal L}
                {2E\mathcal P}>0,&
 g_c&=\frac{HT^2j^2B^3\mathcal L}{2\mathcal P}>0.
\end{aligned}
 \tag{RSR16}
\]
In particular
\[
 f_c\sim\frac{Tx^2K}{2(j+1)},\qquad
 \frac{r_c}{f_c}\sim\frac{1-p}{j},\qquad
 \frac{g_c}{f_c}=Hy\frac{\mathcal L}{\mathcal F}.
 \tag{RSR17}
\]

For the remaining signed mass define
\begin{align*}
\mathcal K={}&pW_0+
 B(1-p)\{j[H(p+1)-p]-(1-p)(p-H)\}\\
 &+TBH\{(p+1)(1-H)j+(1-p)(BHj+p+1-H)\}\\
 &+T^2B^2H^2j(1-H).
 \tag{RSR18}
\end{align*}
Then
\[
 k_c=-\frac{T^2B^2j(1+x)\mathcal K}{2E\mathcal P}.
\]
The identity
\begin{align*}
\mathcal K-pW_0+Bp(1-p)(j+1)
={}&B(1-p)\{Hj(p+1)+H(1-p)+p^2\}\\
 &+TBH\{(p+1)(1-H)j+(1-p)(BHj+p+1-H)\}\\
 &+T^2B^2H^2j(1-H)
 \tag{RSR19}
\end{align*}
has a nonnegative right side. Since $W_0\geq1-p$ and
$B(j+1)\in\mathrm I$, it proves $\mathcal K>0$ and $k_c<0$.
The displayed finite arrays also give
\[
 \mathcal K/K=O(p+B),\qquad
 |k_c|/(q_cy)=O(p+B).
 \tag{RSR20}
\]
All constants in these estimates are real and uniform over the source
openings; neither $j-p$ nor $1-H$ has been bounded below.

\subsubsection*{Paying the signed masses, with the full current reserved}

Put
\[
 h_0=\mu+\tfrac12,\qquad
 M_e=\frac{nh_0j^2}{n+h_0},\qquad J_c=M_ey,
 \qquad e_0=\tfrac12p^2y^2-g_c.
 \tag{RSR21}
\]
These are pure-source quantities and $0\leq J_c\leq M_{ac}$.
Writing $e_c=np^2y^2+e_0$ preserves the entire second correlated
current positive square term. The conclusion is the following precise scalar
statement: after subtracting the reserved positive square terms
\[
 (M_{ac}-J_c)B_*^2+np^2y^2V^2B_*^2,
 \tag{RSR22}
\]
the remainder is a separated SOS and retains, for example, one half
of $q_cQ_*^2+r_cR_*^2+f_c\Psi+J_cB_*^2$.
When $n=0$ the two reserved positive square terms and $J_c$ vanish literally.
This assertion supplies precisely the effective current mass needed
when the full $N$ normal is later completed against the original $n$
positive square term. It is not a completion of that normal itself.

A budget has infinitesimal cost when its Gram coefficients are
bounded and have zero residue. A finite number of such corrections
leaves any fixed fraction smaller than one of each protected square.
The remaining Gram matrix has positive-definite real residue, so
Lemma~\ref{p25af2c58da7e:lem:gram} gives its exact sum of squares.
Restoring the budget identities and the retained positive source
terms proves the required identity.

First $e_0^-\leq g_c$. The exact budget
\[
 (1+V^2)\Psi-V^2B_*^2
 =Q_*^2+R_*^2+V^2(1-UVZ_*)^2
 \tag{RSR23}
\]
pays its negative part from the circuit at cost
\[
 \frac{e_0^-}{f_c}+\frac{e_0^-}{f_cy}v^2
 \ \leq\ Hy\frac{\mathcal L}{\mathcal F}
       +H\frac{\mathcal L}{\mathcal F}v^2\in\mathrm I.
\]
Its positive part is retained as a square. Also $d_c>0$ because
$M_{b^2c}>0$ and $k_c<0$; the $d_c$ and $g_c$ terms cause no loss.

For $k_c$, use the exact identity
\[
 Z_*(V^2-1)=-VQ_*-R_*.
 \tag{RSR24}
\]
The $VQ_*$ part is paid from $q_cQ_*^2$ with coefficient
$2|k_c|v^2/(q_cy)$, which is infinitesimal by RSR20. The remaining
$R_*$ part is paid in the following exhaustive three cases. Fixed
real multipliers such as the factor two do not change any cost claim.

\paragraph{Case 1: $1-p$ has a positive real lower bound.}
The exact ratio is
\[
 \frac{|k_c|}{r_c}
 =\frac{T\mathcal K}{(1-p)(1+Hy)\mathcal L}\in\mathrm I.
 \tag{RSR25}
\]
This uses only $\mathcal K/\mathcal L=O(p+B)$ and $T\in\mathrm I$.
Neither current mass nor a bound on $j$ is required.

\paragraph{Case 2: $1-p\in\mathrm I$ and $\eta_*$ has a positive real
lower bound.}
Here $p\sim1$, so $Bj(j+1)$ is a bounded positive unit. Since $B$ is
infinitesimal, $j$ is unbounded and $j^2B$ is a positive bounded unit.
Equations RSR16--RSR20 give
\[
 \frac{|k_c|}{f_c}\in\mathrm I,\qquad
 \frac{|k_c|}{f_cx}
 =O\!\left(\frac{T(p+B)}{j^2B}\right)\in\mathrm I.
 \tag{RSR26}
\]
The exact circuit budget
\[
 (1+U^2)\Psi=R_*^2+B_*^2+U^2Q_*^2+U^2(1-UVZ_*)^2
 \tag{RSR27}
\]
therefore pays $|k_c|R_*^2$ at cost
$|k_c|/f_c+(|k_c|/(f_cx))u^2\in\mathrm I$.
This case includes any intersection with literal $n=0$.

\paragraph{Case 3: $1-p,\eta_*\in\mathrm I$.}
Now $p\sim1$, $j>p$, and RSR9 gives $n>0$ together with its exact
relative estimate. The source mass $M_e$ is a positive bounded unit.
For the lower bound use
\[
 \frac{p^2}{M_e}
 =\frac{p^2}{h_0j^2}+\frac{p^2}{nj^2}=O(1).
\]
For the upper bound use $M_e\leq h_0j^2$ when $j$ is bounded,
and $M_e\leq nj^2\sim p/2$ when $j$ is unbounded. In particular,
\[
 \frac{|k_c|}{f_c}\in\mathrm I,\qquad
 \frac{|k_c|}{q_cy}=O(1),\qquad
 \frac{|k_c|}{J_cx}
 =O\!\left(\frac{T(p+B)(1+1/n)}{j^2}\right)=O(1).
 \tag{RSR28}
\]
The last bound follows directly from
$1/n\sim2(j-p)(j+1)/p$, with $p\sim1$ and $j>p$.
There is an exact separated SOS budget
\[
 2\{\Psi+V^2Q_*^2+U^2B_*^2\}-R_*^2\quad\hbox{is SOS}.
 \tag{RSR29}
\]
For completeness, set $t_*=UV$, $\delta_*=U-V$,
$E_*=Z_*(1-t_*)$, $T_*=t_*(1-Z_*)$, and
$J_*=T_*-\delta_*^2Z_*$. Its SOS decomposition is
\[
 (E_*+T_*)^2+2\delta_*^2+4t_*(1-Z_*)^2
                       +2J_*^2+4t_*\delta_*^2Z_*^2.
\]
All roots are separated as in RSR13. Applying RSR29 costs a multiple
of $|k_c|/f_c$ on the circuit, $|k_c|v^2/(q_cy)$ on $Q_*$,
and $|k_c|u^2/(J_cx)$ on $B_*$. Each is infinitesimal by RSR28.

The three cases exhaust the source because $0<1-p<1$ and $\eta_*$
is bounded. A finite number of the displayed allocations therefore
proves the scalar assertion following RSR22 with a strict retained
fraction. All coefficient bounds depend only on the source and
$u,v$; no bound on the independent state coordinate $z$ was used.
The proof allows arbitrary relative rates of all source and testing
infinitesimals. Literal $n=0$ is covered by Cases 1 and 2; Case 3
excludes it by the original physicality equation itself.

\subsection{Completion of both state normals}
\label{component:p261f71af685d}

\label{p261f71af685d:sec:fn-reconstructed}

\paragraph{Use of the complete normal calculation.}
The parent identities below are exact for the compact stationary dictionary.
The source estimates and square budgets give the full local R1 certificate
under the full-R1 source hypotheses verified in R1N1--R1N11, with arbitrary
positive $jT$. The later RL and EN bodies reuse only the printed exact
identities and supply their own source estimates.

\paragraph{Fields, source scope, and notation.}
Let $K_c,K_t\supset\mathbb R$ be real closed fields, with bounded rings
$B_c,B_t$ and infinitesimal ideals $\mathfrak m_c,\mathfrak m_t$.
The ring of certificate expressions is
$\mathcal T=K_c\otimes_{\mathbb R}K_t$, and its bounded coefficient ring
is $\mathcal B=B_c\otimes_{\mathbb R}B_t$.  A bounded coefficient with
zero real residue belongs to
$\mathfrak n=\mathfrak m_c\otimes B_t+B_c\otimes\mathfrak m_t$.
Every radical or inverse used below is pure source, pure testing, or
an explicitly separated product of such factors.  No order on
$\mathcal T$ and no inverse of a mixed expression is used.

Use the positive compact source coordinates
\[
 x=jB,\qquad y=TB,\qquad q_z=jT,\qquad
 0<p<\min\{1,j\},\qquad 0<H<1,                         \tag{RFN1}
\]
where $B,T,x,y\in\mathfrak m_c$ and $q_z>0$ is arbitrary. These are
the compact coordinates of the complete original R1 sector; the absence
of any $q_z$ restriction is verified in R1N1--R1N11.  The original
strict source gaps remain in force and $n\geq0$, including $n=0$.
For reference, their native substitution is
\[
 r=\frac{j-p}{j+px},\quad
 h=\frac{1-HT}{1+HTx},\quad
 a_s=\frac{j+px}{1+HTx},\quad
 B_s=\frac{B(1+HTx)}{1+pB},\quad
 t_s=\frac{T(1+pB)}{1+HTx}.                              \tag{RFN2}
\]
The six weights $n,m,C,A,D,\alpha$ are always the original source
weights at (RFN2), rather than weights of an auxiliary stationary
source.  Put $\mu=m(1-p)$, $h_0=\mu+\tfrac12$,
$k=(j+1-p)/(j+1)$, and $\eta_*=Bj(j+1)/p$.
In this source sector RS proves
\begin{gather}
 \tfrac12<k<1,\quad \eta_*=O(1),\quad
 \mu=O(p),\quad h_0\asymp1,\quad A=O(T),\quad
 \alpha\asymp B^2(j+1),                                  \tag{RFN3a}\\
 TCH\sim k/2,\qquad T(D+\tfrac12)(1-H)\sim k/2,\qquad
 \frac{n(j-p)}j=O(1),                                   \tag{RFN3b}\\
 f_c\asymp Tx^2,\qquad q_c\asymp Tx,\qquad
 r_c\asymp TjB^2(1-p),                                  \tag{RFN3c}\\
 \eta_*\in\mathfrak m_c\quad\Longrightarrow\quad
 n\sim\frac{p}{2(j-p)(j+1)}>0.                           \tag{RFN3d}
\end{gather}
Here $\asymp$ means that the ratio is a positive bounded unit, and
$\sim$ means that the ratio has residue one.  All estimates are uniform
at the independent openings $H$, $1-H$, $j-p$, and $1-p$.
In particular $B(j+1)=B+x$ and $\alpha$ are infinitesimal.
No comparison between source and testing infinitesimals is allowed.

\subsubsection*{The complete parent and the correlated source mass}

Write $D_u=1+u^2$, $D_v=1+v^2$.  The original quarter-parent is
\begin{align}
 \mathcal P={}&m(vW-X+uvY-uZ)^2+nD_uD_vX^2
  +CD_v(uW-Y)^2+AD_uD_vY^2\\
 &+(D+\tfrac12)(uW+uvX-Y-vZ)^2
  +(\alpha+\tfrac12)D_u(vY-Z)^2\\
 &+\tfrac12D_v(X+uZ)^2-\tfrac12D_uD_vZ^2.               \tag{RFN4}
\end{align}
Use the two mixed normal rows
\[
 N_0=\frac{X+uZ-pD_uv}{D_u},\qquad
 M_0=\frac{uW-Y-HD_uvz}{D_u}.                            \tag{RFN5}
\]
Their exact inverse, with $w=1$, is
\[
 \begin{aligned}
 X&=N_0+pv-uz,&Z&=z+uN_0+puv,\\
 W&=1+u(M_0+Hvz),&Y&=u-M_0-Hvz.
 \end{aligned}                                           \tag{RFN6}
\]
Only the pure testing unit $D_u$ is inverted.  Retain the nonnegative
positive square term $mpD_u^2N_0^2$ and define
$G=(\mathcal P-mpD_u^2N_0^2)/D_u$.  Direct substitution gives the
\emph{entire} expression
\begin{align}
 G={}&\mu D_u\{N_0^2-2vN_0+(1-p)v^2\}
       +nD_v(N_0+pv-uz)^2\\
 &+CD_uD_v(M_0+Hvz)^2+A D_v(u-M_0-Hvz)^2\\
 &+(D+\tfrac12)D_u(M_0-(1-H)vz)^2\\
 &+(\alpha+\tfrac12)
   \{u(1-p)v-(1+Hv^2)z-uN_0-vM_0\}^2\\
 &+\tfrac12D_uD_v(N_0+pv)^2
       -\tfrac12D_v(z+uN_0+puv)^2.                      \tag{RFN7}
\end{align}
Thus $G$ is exactly quadratic in $(N_0,M_0)$:
\[
 G=F+2L_NN_0+2L_MM_0+H_NN_0^2+H_MM_0^2
                  +2(\alpha+\tfrac12)uvN_0M_0,          \tag{RFN8}
\]
where $F=G|_{N_0=M_0=0}$ and
\[
 \begin{split}
 H_N&=h_0D_u+\tfrac12v^2+\alpha u^2+nD_v,\\
 H_M&=D_u(CD_v+D+\tfrac12)+AD_v+(\alpha+\tfrac12)v^2.
 \end{split}                                           \tag{RFN9}
\]
There are no omitted higher state-normal terms.

\paragraph{Exact scalar interface, before spending the $n$-mass.}
Set
\[
 U=\frac u{\sqrt x},\quad V=\frac v{\sqrt y},\quad
 Z_* =\frac z{\sqrt{q_z}},\qquad
 Q_*=U-VZ_*,\quad B_*=V-UZ_*,\quad R_*=Z_*-UV,
\]
\[
 \Psi=U^2+V^2+Z_*^2+U^2V^2Z_*^2-4UVZ_*.                 \tag{RFN10}
\]
All three normalized testing variables are positive, but none is
assumed bounded.  A legal separated certificate is
\[
 \Psi=(U-V)^2+Z_*^2(1-UV)^2+2UV(1-Z_*)^2.               \tag{RFN11}
\]
For example, $\sqrt{UV}$ is the product of the source factor
$(xy)^{-1/4}$ and the testing factor $\sqrt{uv}$.

The exact RS/RB scalar identity for the actual source is
\begin{align}
 F={}&q_cQ_*^2+nj^2y B_*^2+r_cR_*^2+d_cV^2Q_*^2
   +(np^2y^2+e_0)V^2B_*^2\\
   &+f_c\Psi+g_cV^2\Psi+k_cZ_*^2(V^2-1)^2,
 \qquad e_0=\tfrac12p^2y^2-g_c.                        \tag{RFN12}
\end{align}
For an entirely algebraic definition of the masses, put
$K_0=CH^2+(D+\tfrac12)(1-H)^2$ and
\[
\begin{array}{ll}
 M_a=Ax,&M_b=\{\mu(1-p)+(n+\tfrac12)p^2\}y,\\
 M_c=\alpha q_z,&
 M_{ab}=\{A+\alpha(1-p)^2+\mu(1-p)+\tfrac12(1-p)^2\}xy,\\
 M_{ac}=nxq_z=nj^2y,&
 M_{bc}=\{K_0+AH^2+2H\alpha+H-\tfrac12\}yq_z,\\
 M_{b^2}=(n+\tfrac12)p^2y^2,&M_{abc}=(K_0+n)xyq_z,\\
 M_{b^2c}=H^2(A+C+\alpha+\tfrac12)y^2q_z,&
 M_{ab^2c}=CH^2xy^2q_z.
\end{array}
\]
Then
\[
\begin{split}
 k_c&=(M_a-M_b+M_{ac}+M_{ab^2c}-M_{bc})/2,\\
 r_c&=M_c-M_b+M_{ac}-k_c,\quad f_c=M_b-M_{ac},\quad q_c=M_a-f_c,\\
 d_c&=M_{b^2c}-k_c,\qquad g_c=M_{ab^2c}.
\end{split}
\]
The source contact identities make (RFN12) exact, including both odd
monomials of $F$.  It is checked coefficientwise against (RFN7), not
only at the contact point.

Define the unnormalized current $b=v-uz/j=\sqrt y B_*$ and the
pure source quantities
\[
 c_b=\frac{nj}{n+h_0},\qquad M_e=\frac{nh_0j^2}{n+h_0}.
                                                               \tag{RFN13}
\]
The whole scalar pieces $nj^2b^2+np^2v^2b^2$ in (RFN12) must be kept
until their correlations with $N_0$ have been completed.  The full
$N_0$-linear term below contains $njb+npv^2b$, so these pieces and the
corresponding hard terms give
\[
 n(N_0+jb)^2+n v^2(N_0+pb)^2.
\]
Complete only the first square, using a pure source inverse:
\[
 h_0N_0^2+n(N_0+jb)^2
   =(n+h_0)(N_0+c_bb)^2+M_e b^2.                       \tag{RFN14}
\]
Keep $nv^2(N_0+pb)^2$ whole.  Put
$\widehat N=\sqrt{n+h_0}(N_0+c_bb)$.
The other hard terms in (RFN9) are the nonnegative whole squares
\[
 h_0(uN_0)^2+\tfrac12(vN_0)^2+\alpha(uN_0)^2.             \tag{RFN15}
\]
They are still kept whole after the shift; no term $nu^2$ or $nv^2$
is treated as a bounded perturbation.

For $n>0$ one has the exact bounded-shift identity
\[
 \frac{c_b^2}{M_e}=\frac{n}{h_0(n+h_0)}=O(1).
                                                               \tag{RFN16}
\]
At $n=0$, set $c_b=M_e=0$ and omit the effective current row.  Formula
(RFN14) remains literal, and no division by $n$ or $\sqrt{M_e}$ is
made in that case.

RS pays the signed $e_0$ and $k_c$ terms in (RFN12) at infinitesimal
cost after (RFN14), with the effective current mass $M_e y$, and leaves
fixed positive real fractions of
\[
 q_cQ_*^2,\qquad M_e b^2,\qquad r_cR_*^2,\qquad f_c\Psi. \tag{RFN17}
\]
The second term is omitted when $n=0$.  This is the precise scalar
allocation needed here; an allocation that spends the raw $nj^2y$
mass before (RFN14) would not suffice.

\subsubsection*{Finite exact budgets for the full normal rows}

In the rest of the proof, a row is called free if any prescribed fixed
real multiple of its square can be purchased at infinitesimal cost
from (RFN17).  This has an exact meaning: a finite SOS identity writes
its square below a finite sum of those retained squares, with bounded
separated coefficients in $\mathfrak n$.  Subtracting any fixed real
multiple changes their strict Gram diagonals only infinitesimally.

\paragraph{Three elementary circuit budgets.}
Direct expansion gives
\begin{align}
 (1+U^2)\Psi-R_*^2
     &=B_*^2+U^2Q_*^2+U^2(1-UVZ_*)^2,                  \tag{RFN18a}\\
 (1+V^2)\Psi-V^2B_*^2
     &=Q_*^2+R_*^2+V^2(1-UVZ_*)^2.                   \tag{RFN18b}
\end{align}
There is also a useful opposite-current budget.  Put
$t_*=UV$, $\delta_*=U-V$, $E_*=Z_*(1-t_*)$,
$T_*=t_*(1-Z_*)$, and $J_*=T_*-\delta_*^2Z_*$.  Then
\begin{align}
 2(\Psi+V^2Q_*^2+U^2B_*^2)-R_*^2
 ={}&(E_*+T_*)^2+2\delta_*^2+4t_*(1-Z_*)^2\\
   &+2J_*^2+4t_*\delta_*^2Z_*^2.                      \tag{RFN18c}
\end{align}
Every term on the right is a legal separated square.  Thus a
nonnegative multiplier $a$ provides $aR_*^2$ at infinitesimal cost whenever
\[
 a/f_c\in\mathfrak n,\qquad a/(q_cy)=O(1),\qquad
 a/(M_e yx)=O(1),                                      \tag{RFN18d}
\]
since the last two diagonal costs have the additional testing factors
$v^2$ and $u^2$.  This criterion is used only when $M_e>0$.

\paragraph{An exact radial budget.}
The row $R_u=v(u^2-x)$ is free in all source subcases, including the
weak opening $1-p\in\mathfrak m_c$.  Here is an explicit finite
certificate, so that this assertion does not rest on pointwise
nonnegativity.  Write $D_*=V(U^2-1)$ and
\[
\begin{split}
 t_7={}&(Z_*-U^2V,\ V-U^2V,\ U-U^2V,\ UZ_*-U^2V,\\
       &\hspace{25mm}UV-U^2V,\ UVZ_*-U^2V,\ U^2-U^2V)^T,\\
 t_3={}&(1-UV,\ V-UV,\ U-UV)^T,
\end{split}
\]
\[
 G_7=\begin{pmatrix}
 4&-2&-3/2&0&-1/2&0&1/2\\
 -2&3&0&-3/2&0&-1&2\\
 -3/2&0&4&3/2&-2&-3&0\\
 0&-3/2&3/2&4&-3/2&0&-3/2\\
 -1/2&0&-2&-3/2&3&2&-1\\
 0&-1&-3&0&2&8&-4\\
 1/2&2&0&-3/2&-1&-4&4
 \end{pmatrix},\qquad
 G_3=\begin{pmatrix}3&-2&-2\\-2&2&2\\-2&2&3\end{pmatrix}.
\]
The exact identity is
\[
 4\Psi+4U^2Q_*^2+4U^2R_*^2-D_*^2
   =t_7^TG_7t_7+4VZ_*(1-U)^2+UZ_*t_3^TG_3t_3.         \tag{RFN19}
\]
The leading principal minors of $G_7$ are
$4,8,101/4,1049/16,907/16,3279/16,743/16$; those of $G_3$ are
$3,2,2$.  Both are positive definite real matrices.  Their real
Cholesky factors and the separated square roots of $VZ_*$ and $UZ_*$
therefore make the right side a finite legal SOS.  Denote it by
$\mathcal W(U,V,Z_*)$.

Since $R_u=x\sqrt y D_*$, the exhaustive source cases are as follows.

\emph{Case 1: $1-p$ is a positive unit.}
Multiply (RFN19) by $x^2y$.  The costs relative to
$f_c\Psi,q_cQ_*^2,r_cR_*^2$ are, respectively,
\[
 O(B),\qquad O(Bu^2),\qquad O(u^2/(1-p)).
\]
All belong to $\mathfrak n$.

\emph{Case 2: $p\sim1$ and $\eta_*$ is a positive unit.}
Then $j$ has a positive real lower bound, and $j^2B$ is a positive
unit.  Adding $4U^2$ times (RFN18a) to (RFN19) gives
\[
 D_*^2\ \preceq_{\rm sos}\
 4(1+U^2+U^4)\Psi+4U^2Q_*^2.
\]
After multiplication by $x^2y$, the costs are
\[
 O\!\left(B+u^2/j+u^4/(j^2B)\right),\qquad O(Bu^2),
\]
again infinitesimal.  The notation $a\preceq_{\rm sos}b$ means
$b-a\in\Sigma\mathcal T^2$ and is only shorthand for the displayed
exact identities.

\emph{Case 3: $p\sim1$ and $\eta_*\in\mathfrak m_c$.}
Formula (RFN3d) gives $n>0$ and $1/M_e=O(1)$; it also gives
$p^2/M_e=O(1)$.  Put $E^{\rm rad}=Z_*(U^2-1)$.
Permuting $V$ and $Z_*$ in (RFN19) yields
\[
 \mathcal W(U,Z_*,V)
   =4\Psi+4U^2Q_*^2+4U^2B_*^2-(E^{\rm rad})^2.
\]
The identity $D_*=UE^{\rm rad}+(U^2-1)B_*$ then gives, exactly,
\begin{align}
 &8U^2\Psi+8U^4Q_*^2+
   \{8U^4+2(U^2-1)^2\}B_*^2-D_*^2\\
 &\hspace{10mm}=2U^2\mathcal W(U,Z_*,V)
                 +\{UE^{\rm rad}-(U^2-1)B_*\}^2.        \tag{RFN20}
\end{align}
After multiplication by $x^2y$, the three costs are
\[
 O(u^2/j),\qquad O(u^4/j),\qquad
 \frac{8u^4+2(u^2-x)^2}{M_e}.
\]
They are infinitesimal.  These three cases prove the asserted free
radial row without assuming a rate for any testing variable.

\paragraph{The smaller current helpers.}
If $\eta_*\in\mathfrak m_c$, (RFN3d) implies
\[
 \frac{p^2}{M_e}
 =\frac{p^2}{h_0j^2}+\frac{p^2}{nj^2}=O(1),
\]
because $p/j<1$ and
$p(j-p)(j+1)/j^2\leq p+p/j$.
Thus $(pv^2b)^2$ costs only $(p^2/M_e)v^4$ of the effective current.
If $\eta_*$ is a positive unit instead, then
\[
 \frac{p^2y}{f_c}=O\!\left(\frac{p+p/j}{\eta_*}\right)=O(1).
\]
Multiply (RFN18b) by $p^2y^2v^2$.  This provides $(pv^2b)^2$ with circuit
cost
\[
 \frac{p^2y}{f_c}(yv^2+v^4)\in\mathfrak n.
\]
The same identity multiplied by $Tp^2u^2y^2$ shows that
$\sqrt T\,puvb$ is free: its cost is either
$(p^2/M_e)Tu^2v^2$ or
$(p^2y/f_c)Tu^2(y+v^2)$.

\subsubsection*{The entire two-normal conductor}

\paragraph{First normal: exact source cancellation and allocation.}
Put
\[
 q=vz-Tu,\qquad R=z-uv/B,\qquad R_u=v(u^2-x),
\]
\[
 K_c=\frac p2+j\{H(\alpha+\tfrac12)-\tfrac12\}+n(p-j),
 \qquad
 K_R=-\mu-(\alpha+\tfrac12)(1-p)+\frac TjK_c.
\]
Both $K_c/j$ and $K_R$ are bounded by (RFN3).  The complete first normal
is exactly
\[
 L_N=\{j(n-\alpha)+p(n+\tfrac12)v^2\}b
                  +K_RR_u+\frac{K_c}{j}uvq.            \tag{RFN21}
\]
In particular the sign of the $\mu$ term in $K_R$ is negative.
For a direct source check, the difference between this formula and
the derivative of (RFN7) is $vS_N$, where
\[
 S_N=np-\mu+\tfrac p2+
       \{-\mu-(\alpha+\tfrac12)(1-p)\}x
       +j(\alpha-n)+yK_c=0.
\]
The last equality is an exact rational identity of the actual native
source under (RFN2), not an asymptotic contact equation.

After the two correlated $n$-squares are formed, the remaining linear
coefficient is
\[
 L_N^{\rm rem}=-j\alpha b+\tfrac12pv^2b
                       +K_RR_u+\frac{K_c}{j}uvq.
\]
The $q$ row is strict because
$q_cQ_*^2=(q_c/(Tx))(q^2/T)$, with $q_c/(Tx)$ a positive unit.
Its coefficient in $L_N^{\rm rem}$ relative to $q/\sqrt T$ is
$(K_c/j)uv\sqrt T\in\mathfrak n$.
The $pv^2b$ helper is free by the preceding paragraph.
For the remaining small row, the exact identity
\[
 j\alpha b=-\frac\alpha B R_u-\alpha uR
\]
has $\alpha/B=O(B+x)\in\mathfrak m_c$.
When $1-p$ is a unit, the square of $\alpha uR$ costs
$O(B^2(j+1)^2u^2/(1-p))\in\mathfrak n$ against $r_cR_*^2$.
When $p\sim1$, use $B_*^2\preceq_{\rm sos}(1+U^2)\Psi$ from
(RFN18a).  The circuit cost of $(j\alpha b)^2$ is
\[
 O\!\left(B^3(j+1)^2+B^2(j+1)^2u^2/j\right)\in\mathfrak n,
\]
since $j$ is bounded below and $B(j+1)=B+x$ is infinitesimal.
Finally, $K_RR_u$ has bounded coefficient; provide an arbitrarily large
fixed real $R_u^2$ pivot using (RFN19)--(RFN20).

The shift $N_0=\widehat N/\sqrt{n+h_0}-c_bb$ introduces no new
unbounded Gram coefficient: $1/\sqrt{n+h_0}=O(1)$, and for $n>0$
$c_b/\sqrt{M_e}=O(1)$ by (RFN16).  The only bounded, potentially
non-infinitesimal couplings are therefore those of $R_u$ to
$\widehat N$ and $\sqrt{M_e}b$.  A sufficiently large fixed real radial
pivot absorbs both while preserving a strict positive real residue.
For $n=0$ the latter coupling is exactly absent.

\paragraph{Second normal: cancellation of both source openings.}
Define
\[
 K_d=CH-(D+\tfrac12)(1-H),\qquad
 C_q=K_d+AH+\tfrac12
       +v^2\{CH(1+u^2)+AH+(\alpha+\tfrac12)H\}.
\]
The complete second normal is
\[
 L_M=C_q q+\mu uvb+v(\alpha-BTCHu^2)R.                  \tag{RFN22}
\]
The actual source satisfies the three exact identities
\[
\begin{split}
 S_0&=T(K_d+AH+\tfrac12)-A=0,\\
 S_1&=jK_d+\mu+xTCH=0,\\
 S_2&=jK_d+xTCH+T\{CH+AH+(\alpha+\tfrac12)H\}
          -A-(\alpha+\tfrac12)(1-p)+\alpha/B=0.
\end{split}
\]
Before imposing these identities, the complete derivative minus
(RFN22) is exactly
$uS_0+u^2vzS_1/j+uv^2(S_2-S_1)$.
This records all state terms and rules out a hidden $j$-weighted
radial error.  In particular $S_1$ gives
$K_d=-\mu/j-BTCH=O(1)$ even when $C$ or $D$ is unbounded.
Consequently $TC_q=O(T+v^2)\in\mathfrak n$.

Normalize
\[
 M_0=\sqrt{TH(1-H)}\,\widehat M.
\]
The full diagonal (RFN9), not only its zero-angle value, obeys
\[
 TH(1-H)H_M\sim k/2>0.                                 \tag{RFN23}
\]
Indeed the $C$ and $D$ terms contribute
$(k/2)D_u\{(1-H)D_v+H\}$ up to infinitesimal bounded coefficients,
and all remaining terms are infinitesimal after this normalization.
This is uniform at both independent openings.
The normalized $C_qq$ coupling has coefficient
$T\sqrt{H(1-H)}C_q\in\mathfrak n$ on the strict row $q/\sqrt T$.
The normalized $\mu uvb$ row is free because $\mu=O(p)$ and
$\sqrt T\,puvb$ has already been purchased.

It remains to allocate the entire opposite row, including its
higher testing terms.  Put
\[
 J=B(j+1)+u^2\in\mathfrak n,\qquad
 a_R=Tq_zv^2(\alpha-BTCHu^2)^2.
\]
The factor $H(1-H)$ may be discarded from the cost bound, since it is
bounded.  Since $\alpha-BTCHu^2=O(BJ)$, the exact squared multiplier
$a_R$ satisfies
\[
 \frac{a_R}{f_c}=O((T/j)v^2J^2),\qquad
 \frac{a_R}{q_cy}=O(v^2J^2).                            \tag{RFN24a}
\]
If $1-p$ is a unit, its direct cost against $r_cR_*^2$ is
$O(Tv^2J^2/(1-p))\in\mathfrak n$.
If $p\sim1$ and $\eta_*$ is a unit, apply (RFN18a); the cost is
$(a_R/f_c)(1+u^2/x)\in\mathfrak n$, using that $j^2B$ is a unit.
If $p\sim1$ and $\eta_*\in\mathfrak m_c$, (RFN3d) gives
\[
 \frac{a_R}{M_e yx}
     =O\!\left(\frac{T(1+1/n)}{j^2}v^2J^2\right)=O(1),  \tag{RFN24b}
\]
because $(1+1/n)/j^2=O(1)$ in this source case.  The other two
conditions in (RFN18d) follow from (RFN24a).  The exact budget
(RFN18c) therefore provides this row at infinitesimal cost as well.
These estimates treat the whole square
$(\alpha-BTCHu^2)^2$; there is no truncation at a testing degree.

\paragraph{Full Gram, legal roots, and restoration.}
After the normalizations, the entire remaining mixed hard term is
\[
 2(\alpha+\tfrac12)uv\sqrt{TH(1-H)}
    \left(\frac{\widehat N}{\sqrt{n+h_0}}-c_bb\right)\widehat M.
\]
Its coefficients relative to $\widehat N$, $\widehat M$, and, when
present, $\sqrt{M_e}b$ are infinitesimal; (RFN16) controls the shift.
All other remaining normal couplings were allocated above.  Choose a
fixed sufficiently large radial pivot, and then purchase all the
finitely many free helpers.  Their scalar costs lie in $\mathfrak n$,
so fixed positive fractions of (RFN17) survive.  The complete bounded
Gram has strict positive real residue: the $\widehat N$ pivot is one,
the $\widehat M$ pivot has residue $\overline{k}/2$, the retained scalar pivots
are positive, and the only bounded radial off-diagonals have been
absorbed by the chosen real radial pivot.

The separated strict bounded-Gram lemma now gives a finite sum of
squares in $\mathcal T$.  Its row entries need not be bounded.  All
budget remainders are explicit finite separated sums of squares, all
source-only completions use positive pure source denominators, and
all state-coordinate inverses are pure testing.  Restore the whole
squares (RFN15), the retained $nv^2(N_0+pb)^2$ square, the scalar
budget remainders, the $mpD_u^2N_0^2$ positive square term, and the pure testing
factor $D_u$.  This certifies the complete original parent (RFN4).
Every bounded coefficient used in the argument depends only on the
source and on $u,v$.  The arbitrary state ratio $z$ remains inside
retained rows throughout.  Hence the proof has no bounded-state
restriction, and it includes literal $n=0$ without any limiting
argument.

\subsection{All source ratios and testing angles}
\label{component:pbaab204a8ab1}

\label{pbaab204a8ab1:sec:new-r1-all-tests}

\paragraph{Statement and exact source.}
Use the original common-II source, with $d=1$, $\lambda=\beta=0$,
the six native coefficients of the actual
Section~\ref{component:pd709027c7805}, every original strict gap,
and $n\geq0$.  In the original R1 sector,
\[
 t_s,\quad B_s,\quad a_sB_s,\quad t_sB_s>0
 \quad\hbox{are infinitesimal}.                         \tag{R1N1}
\]
There is no bound on $a_s$, on any original positive opening, or on
the source ratio previously called $q_{\rm old}$.  The conclusion is:
the entire original parent is a finite separated tensor sum of
squares at every independent real projective measurement pair and
every real or complex state.  This proof treats R1 only.

Source and testing quantities belong to separate real closed fields
$K_c,K_t\supset\mathbb R$.  All SOS assertions are in
$K_c\otimes_{\mathbb R}K_t$.  We use the actual actual strict
bounded-Gram lemma: a bounded finite Gram with positive definite
real residue gives tensor SOS, with no boundedness requirement on
its row entries.  No order or mixed inverse is imposed on the
tensor ring.

Put $c_s=1+a_sB_sr$, $R_s=1+ha_sB_s$, and define
\[
\begin{aligned}
 j&=\frac{a_sc_s}{R_s},\quad B=\frac{B_sR_s}{c_s},\quad
 T=\frac{t_sc_s}{R_s},\quad p=\frac{a_s(1-r)}{R_s},\\
 H&=\frac{1-h}{t_sc_s},\qquad x=jB,\quad y=TB,\quad q_z=jT.
\end{aligned}                                          \tag{R1N2}
\]
The native strict gaps give $0<p<\min(1,j)$ and $0<H<1$.
They also give $1-h<t_sc_s$, so $h\sim1$.
Consequently $B,T,x,y$ are infinitesimal, while
\emph{$q_z$ is an arbitrary positive source quantity}.
In fact $x=a_sB_s$ and $y=t_sB_s$ exactly.
The inverse source substitution is
\[
 r=\frac{j-p}{j+px},\quad h=\frac{1-HT}{1+HTx},\quad
 a_s=\frac{j+px}{1+HTx},\quad
 B_s=\frac{B(1+HTx)}{1+pB},\quad
 t_s=\frac{T(1+pB)}{1+HTx}.                             \tag{R1N3}
\]
All arguments below therefore concern the actual original source.

The source bounds are RSR7--RSR9: their proofs use only $B,T,Bj$ infinitesimal and $p,H<1$, and impose no condition on $jT$.

\subsubsection*{The complete small-angle theorem extends to all R1 sources}

The exact proof Section~\ref{component:p261f71af685d}, RFN4--RFN24,
is a full-parent proof, not just a restriction to a scalar graph.
We now justify its application under the larger source domain
(R1N1), rather than importing its former domain label.

First, all of its exact identities hold for the rational dictionary
(R1N3), without any size assumption on $q_z$.  The parent, its two
normals, their quadratic block, the scalar circuit, the radial Gram
and the three source cancellation identities were all proved as
identities of this dictionary.  Second, every uniform source bound
used in that proof is exactly (R1N4)--(R1N7) or the following scalar
bound:
\[
 f_c\asymp Tx^2,\quad q_c\asymp Tx,\quad
 r_c\asymp TjB^2(1-p),\quad
 g_c/f_c=O(Hy),\quad |k_c|/(q_cy)=O(p+B).                \tag{R1N9}
\]
To check the latter bounds without $jT\in\mathrm I$, use the explicit
finite polynomials RSR15 and RSR18: each term of
$\mathcal F_1/K,\mathcal F_2/K$ is bounded using $p/j<1$
and $Bj=x$, and $\mathcal K/K=O(p+B)$ by the same inequalities.
Thus $\mathcal F\sim K$ and all statements (R1N9) follow from
RSR16--RSR20.  No $jT$ factor occurs in these estimates.

All source and normal estimates in RFN18--RFN24 use precisely these
source bounds. In particular the apparently unbounded factor $q_z=jT$
cancels in the exact ratios RFN24a--RFN24b; their proofs use $Bj=x$,
$B,T,x,y$ infinitesimal, and no restriction on $jT$. The effective
$n$ square is completed before the signed scalar terms are allocated.
Consequently that full scalar-and-two-normal calculation proves every
R1 source in the local testing chart, including $n=0$.

\subsubsection*{Outer testing angles: two exact current identities}

We include the outer mechanism and its precise positivity condition.
Let
\[
 R=uW-Y,\quad C_t=v(uX-Z),\quad B_t=X+uZ,\quad
 N=R+h_*C_t,\quad D_u=1+u^2,\quad D_v=1+v^2,
\]
and let $K$ be the signed quarter-core
\[
 K=\tfrac12(R+C_t)^2+\tfrac12D_u(vY-Z)^2
             +\tfrac12D_v B_t^2-\tfrac12D_uD_vZ^2.       \tag{R1N12}
\]
All unselected original nonnegative events are kept as whole squares.

If $1/v$ is bounded, use
\[
 K_0=2\sqrt{g_*}C_t,\quad
 B_0=\sqrt{2D_v/D_u}\,B_t,\quad
 c_u=D_u^{-1/2},\quad t_u=uD_u^{-1/2},\quad
 \eta_v=\sqrt{D_v}/v .
\]
The exact full selected-parent identity is
\begin{align*}
 4CD_vR^2+4D(R+C_t)^2+4K
 ={}&4(C+D)N^2+4Cv^2R^2+2(R+C_t)^2\\
 &+2D_u(vY-Z)^2+(K_0,B_0)Q(K_0,B_0)^T,\\
 Q={}&
 \begin{pmatrix}
 1-c_u^2\eta_v^2/(2g_*)&t_u\eta_v/\sqrt{2g_*}\\
 t_u\eta_v/\sqrt{2g_*}&1
 \end{pmatrix}.
\tag{R1N13}
\end{align*}
Since $1/g_*$ is infinitesimal and all the testing factors are bounded,
$\overline Q=I$.  This treats every other angle and every state,
including Bob infinity by homogeneous clearing.

For $v$ infinitesimal put $\epsilon=1/g_*$, $q=\sqrt\epsilon$,
$\gamma=Ag_*$, and
\[
 K_1=2\sqrt{g_*}C_t+\sqrt\epsilon Y,\quad
 B_1=B_t-uvY,\quad V_0=vY,\quad Y_1=\sqrt{\epsilon D_u}\,Y,
\]
\[
 L_0=\tfrac12qc_u(K_1-c_uY_1)-t_uvB_1+c_uvV_0.
\]
Then $L_0=\sqrt{D_u}\,v(vY-Z)$ exactly, and
\begin{align*}
 &4CD_vR^2+4D(R+C_t)^2+4AD_uD_vY^2+4K\\
 &\quad=4(C+D)N^2+4Cv^2R^2+2(R+C_t)^2+\mathcal G,\\
 \mathcal G={}&K_1^2+2D_v(B_1^2+V_0^2)
 +(4\gamma D_v-c_u^2)Y_1^2\\
 &+2qv(K_1-c_uY_1)V_0-2L_0^2 .
\tag{R1N14}
\end{align*}
The complete four-root Gram has residue
\[
 \operatorname{diag}(1,2,2,\overline{k}^{\,2}
                                      -\overline{c_u}^{\,2}).    \tag{R1N15}
\]
Thus it is strict whenever the last entry is positive.  In particular
it covers every unbounded $u$, and covers every noninfinitesimal
$u$ whenever $\overline k=1$.  All coefficients in the identity
are bounded functions of the displayed normalized homogeneous Alice
coordinates, so it includes the literal Alice-infinite endpoint.

\subsubsection*{The remaining bounded nonzero Alice chart}

The only outer chart left by (R1N13)--(R1N15) has $v$ infinitesimal,
$u$ a bounded positive unit, and $\overline k<1$.  Since
$1-k=p/(j+1)$, this implies that $p,j$ are positive units and $j$
is bounded.  Therefore $\eta_*\in\mathrm I$ and (R1N6) gives the
exact relative critical limits
\[
 m\sim\frac{jp}{2(1-p)(j+1)},\qquad
 n\sim\frac{p}{2(j-p)(j+1)}.                              \tag{R1N16}
\]
These statements allow either or both coefficients to be unbounded.

Choose bounded positive real caps $m_c\le m$, $n_c\le n$ as specified
below and keep their nonnegative excesses.  Retain all of $\alpha$
as its original square.  Define six rows
\[
 \widehat N=\sqrt{C+D}(R+h_*C_t),\quad
 P=C_t/\sqrt T,\quad Y_0=\sqrt T\,Y,\quad
 X,\quad Z,\quad V=vY.                                  \tag{R1N17}
\]
They need not be independent or bounded.  Put
$\delta=(C+D)^{-1/2}=O(\sqrt T)$, so
$R=\delta\widehat N-h_*\sqrt T\,P$.
The selected \emph{entire} parent is exactly
\begin{align*}
 \mathcal P_{\rm sel}={}&
 \widehat N^2+Cv^2R^2+\tfrac12(R+C_t)^2
       +(Tg_*)P^2+(A/T)D_uD_vY_0^2\\
 &+m_c\{(v/u)R+(D_u/u)V-X-uZ\}^2+n_cD_uD_vX^2\\
 &+\tfrac12D_u(V-Z)^2+\tfrac12D_v(X+uZ)^2
                                      -\tfrac12D_uD_vZ^2 .
\tag{R1N18}
\end{align*}
This is obtained by the exact completion
$CD_vR^2+D(R+C_t)^2=(C+D)N^2+g_*C_t^2+Cv^2R^2$.
It keeps every state-normal term.

The row relation
\[
 PY_0=V(uX-Z)                                            \tag{R1N19}
\]
allows an arbitrary real zero shift
$2b(PY_0-V(uX-Z))$.  Keep the two explicitly nonnegative squares
$Cv^2R^2$ and $(R+C_t)^2/2$ whole.  The remaining complete Gram
on (R1N17) is bounded: $u,1/u$ are bounded, $Tg_*,A/T$ are
units, and the $R$ contribution in the $m_c$ square has coefficients
$\delta v/u$ and $h_*\sqrt T v/u$, both infinitesimal.
Its residue is the direct sum of a unit pivot and the two blocks
\[
 \begin{pmatrix}k_0/2&b\\b&k_0(1+u_0^2)/2\end{pmatrix},\qquad
 Q_3(m_c,n_c;b),                                        \tag{R1N20}
\]
where subscripts denote real residues and $Q_3$ is the quadratic
form in $(X,Z,V)$
\begin{align*}
 Q_3={}&m_c\{(D_u/u)V-X-uZ\}^2+n_cD_uX^2\\
       &+\tfrac12D_uV^2-D_uVZ+\tfrac12(X+uZ)^2
                                            -2bV(uX-Z).
\tag{R1N21}
\end{align*}
In (R1N21), $u=u_0>0$ is real.  The choices of $b,m_c,n_c$
that make both blocks strict are now explicit.

\paragraph{Interior openings.}
Suppose $p_0<1$ and $j_0>p_0$.  Use initially
$m_0=j_0p_0/[2(1-p_0)(j_0+1)]$ and
$n_0=p_0/[2(j_0-p_0)(j_0+1)]$, and take $b=k_0/2$.
The first leading principal minor of $Q_3$ is
$m_0+n_0D_u+1/2>0$; the other two are
\[
 \frac{p_0u_0^2D_uK_0}
      {4(j_0-p_0)(j_0+1)^2(1-p_0)},\qquad
 \frac{p_0u_0^2D_uK_0^2}{8(j_0-p_0)(j_0+1)^3},          \tag{R1N22}
\]
where $K_0=j_0+1-p_0$.  They are positive.
The first block in (R1N20) has determinant
$k_0^2u_0^2/4>0$.  Positive definiteness is open, so fixed
real caps slightly below $m_0,n_0$ preserve strictness and are
smaller than the actual source coefficients.

\paragraph{Only the \(n\) opening is infinite.}
If $j_0=p_0<1$, then $n$ is unbounded and
$m_0=j_0^2/[2(1-j_0)(j_0+1)]$ is finite positive.
At $b=k_0/2=1/[2(j_0+1)]$, restrict (R1N21) to the kernel
$X=0$ of the $n$ positive square term.  Its first diagonal on $(Z,V)$ is
$u_0^2(m_0+1/2)>0$, and its determinant is
$u_0^2k_0^2/4>0$.  Choose $m_c<m_0$ sufficiently close and
then a sufficiently large fixed real $n_c$.  The whole block is
strict, with both caps below the original weights.

\paragraph{Only the \(m\) opening is infinite.}
If $p_0=1<j_0$, then $m$ is unbounded and
$n_0=1/[2(j_0^2-1)]$ is finite positive.
Take a real $b$ just above $k_0/2=j_0/[2(j_0+1)]$, with
\[
 b<\min\{k_0\sqrt{D_u}/2,\ j_0/[2(j_0-1)]\}.
\]
This interval is nonempty.  The first block in (R1N20) is strict.
On the kernel $(D_u/u)V-X-uZ=0$ of the $m$ positive square term, use $(X,V)$
coordinates.  The first diagonal of (R1N21) is
$D_u/[2(j_0^2-1)]>0$, and its determinant is
\[
 -\frac{D_u^2[\,2b(j_0-1)-j_0\,][\,2b(j_0+1)-j_0\,]}
        {4u_0^2(j_0^2-1)}>0.                            \tag{R1N23}
\]
Choose $n_c<n_0$ sufficiently close and then a sufficiently large
real $m_c$.  This gives the required full strict block.

\paragraph{Both openings are infinite.}
If $p_0=j_0=1$, both $m,n$ are unbounded and $k_0=1/2$.
Choose
$1/4<b<\sqrt{D_u}/4$.
The common kernel of the two event positive square terms is
$X=0$, $Z=D_uV/u_0^2$.  On it the rest of (R1N21) is exactly
\[
 \frac{D_u(4b-1)}{2u_0^2}V^2>0.                        \tag{R1N24}
\]
A sufficiently large real cap on both positive square terms therefore gives
a strict block.

The elementary finite-cap fact used in these cases is as follows:
if a real positive semidefinite event Gram is positive definite on
the complement of its kernel, and the remaining real quadratic
form is positive definite on that kernel, a sufficiently large
real multiple of the event Gram makes their sum positive definite.
It follows by a block Schur complement, and works for the sum of
two event Grams as well.  Thus the choices above prove strictness
of the \emph{complete} bounded Gram (R1N18), without restricting
any normal or state coordinate.

\subsubsection*{Exhaustive testing and state join}

The sign gauges SJ3 preserve every original source positive square term and the
whole core, so affine $u,v$ may be taken nonnegative.
If $v$ is not infinitesimal, (R1N13) applies.
If $v$ is infinitesimal and $u$ is not infinitesimal, either
(R1N15) applies or the complete six-root argument
(R1N16)--(R1N24) applies.  These cases exhaust the outer angles.
If $u,v>0$ are both infinitesimal, the enlarged local FN theorem
proved above applies after SJ's full-state reduction to $w=1,z>0$;
SJ gives the other state signs and the axes directly.
Literal homogeneous measurement endpoints are given by SJ10 or
the homogeneous current identities, without a limiting argument.
The zero state is immediate.  For a complex state, the real
parent matrix splits its value into its real and imaginary
quadratic evaluations, to which the same proof applies.

Restore every capped excess, the original $n,m,\alpha$ positive square terms
not used in a particular chart, and all cleared positive pure
testing homogeneous factors.  All normalizations, radicals and
inverses are pure source or pure testing.  This proves the
complete original R1 source row at all independent tests and
states, without a source-ratio split or a comparison of source
and testing rates.

\section{Two small actual angles: the R2 sector}
\subsection{Bounded $a$ or an infinitesimal paired-event parameter}
\label{component:p393d3faa269a}

\label{p393d3faa269a:sec:r2-shifted-local}

\paragraph{Exact statement.}
Use the actual original weak-C2 source and all its original strict
gaps, with \(d=1\), \(\lambda=\beta=0\), and \(n\geq0\).
Suppose \(t_s\) is positive unbounded and
\(B_s,a_sB_s,t_sB_s\) are positive infinitesimals.
Put \(e=1-r\), \(R_s=1+ha_sB_s\).  On the source subset
\[
       a_s=O(1)\quad\hbox{or}\quad
       p=a_s e/R_s\in\mathrm I_c,                     \tag{RL1}
\]
the entire parent is separated tensor SOS for every independent
projective testing pair and every real or complex state.
Literal \(n=0\) and every opening of \(r,h,1-h,g\) are retained.
For the all-testing conclusion, the local proof below is joined
explicitly to R2O1--R2O24.

Use the field and bounded-Gram conventions of EN1--EN4.

\subsubsection*{Actual source entry and bounds}

Use the exact compact source dictionary RSR2--RSR8, with
\[
 j=a_sc_s/R_s,\quad B=B_sR_s/c_s,\quad T=t_sc_s/R_s,\quad
 p=a_se/R_s,\quad H=(1-h)/(t_sc_s),\quad c_s=1+a_sB_sr.
\]
Set
\[
 S=1/T,\quad \ell=HT,\quad \bar h=1-\ell,\quad
 x=jB,\quad y=TB,\quad B=Sy,\quad q_z=j/S .
                                                               \tag{RL2}
\]
Thus \(S,x,y\) are positive infinitesimals,
\(0<p<\min(1,j)\), \(0<\ell<1\), and \(j\asymp a_s\).
Both \(\ell\) and \(\bar h\) may be arbitrarily small.
The original physicality inequality RER2--RER7 gives \(jy\in\mathrm I_c\).
In terms of the actual finite polynomials \(\mathcal P,\mathcal L,
\mathcal N,\mathcal F,\mathcal K\) of RSR, put
\[
 K=j+1-p,\quad k=K/(j+1),\quad
 F_b=j+1+p(1-p)/j,\quad E=1+B(\ell+p-1).
                                                               \tag{RL3}
\]
Here \(1/2<k<1\) and \(1\le F_b/(j+1)<2\).  Direct substitution
\(B=Sy,T=1/S,H=S\ell\) in those finite polynomials gives
\[
 \mathcal P/(j+1)\sim1,\qquad
 S\mathcal L/K\sim1,\qquad S\mathcal F/F_b\sim1.         \tag{RL4}
\]
For clarity, the terms after division contain only bounded
\(p/j,j/(j+1),K/(j+1)\), multiplied by \(S,y\), or \(jy\).
No inverse of \(\ell\), \(\bar h\), or \(j-p\) is involved.

The same substitution in the exact numerator of \(n\) yields
\[
 \mathcal N=p(1+\epsilon_p)
                -yj(j+1)(1+\epsilon_y),\qquad
 \epsilon_p,\epsilon_y\in\mathrm I_c.                  \tag{RL5}
\]
For an explicit finite division certificate, the zero-\(p\) term is
\[
 \mathcal N|_{p=0}
 =-yj(j+1)\{1+S(1-\ell)\}
             \{1+S\ell yj/(j+1)\}.
\]
The quotient \((\mathcal N-\mathcal N|_{p=0})/p\) is a
finite polynomial.  This grouping can be read termwise from RSR5:
the first
coefficient is \(1+y(j+1)+O(Sy(j+1)+Sy^2(j+1))\);
the second coefficient is \(1+O(S+y+Syj)\).
Using \(\mathcal N\geq0\), therefore,
\[
 \eta_2=\frac{yj(j+1)}p=O(1),\qquad
 \eta_2\in\mathrm I_c\ \Longrightarrow\
 n\sim\frac{p}{2(j-p)(j+1)}>0 .                        \tag{RL6}
\]
The estimate includes the face \(\mathcal N=0\), which can occur
only in the positive-unit \(\eta_2\) case.

Let \(\mu=m(1-p)\) and \(h_0=\mu+1/2\).
The complete source formulas now imply
\[
\begin{gathered}
 \mu\sim\frac{jp}{2(j+1)},\quad h_0\asymp1,\quad
 CH\sim k/2,\quad SA\bar h\sim k/2,\quad D/S\sim k/2,\\
 \alpha\sim\frac{Sy^2K^2}{2(j+1)},\quad
 f_c\sim\frac{j^2y^2F_b}{2(j+1)},\quad
 r_c\sim\frac{jy^2(1-p)K}{2(j+1)} .
\end{gathered}                                                   \tag{RL7}
\]
These use the exact masses in RFN12 and the closed RSR formulas.
All displayed positive masses are strictly positive.
Write \(\kappa=-k_c>0\).  The exact positive-surplus formula
RSR19 proves this sign.  Its substituted finite polynomial gives
\[
 \frac{\kappa}{r_c}=O(p+Sy)\quad(p\in\mathrm I_c),\qquad
 \frac{\kappa}{jy^2}=O(1).                              \tag{RL8}
\]
For the remaining case in (RL1), \(p\) and \(j\) are positive
bounded units, and
\[
 \frac{\kappa}{f_c}\sim\frac{pK}{jF_b}\leq\frac12,
 \qquad jF_b-2pK=(j-p)K\geq0.                          \tag{RL9}
\]
The inequality is a bound on real residues when combined with
the preceding asymptotic equality.  On this case \(\eta_2\) is
infinitesimal, and the effective current defined below is a
positive bounded unit.

\subsubsection*{The entire parent and its harmonic shift}

Retain the complete original quarter-parent RFN4.  Its exact inverse
state change at \(w=(W+uY)/(1+u^2)=1\) is
\[
 X=N_0+pv-uz,\quad Z=z+uN_0+puv,\quad
 W=1+u(M_0+Hvz),\quad Y=u-M_0-Hvz.                    \tag{RL10}
\]
Here \(z>0\) is arbitrary.  Retain \(mpD_u^2N_0^2\) and divide
the remaining whole parent by \(D_u=1+u^2\).
The exact expression is RFN7, with both complete state normals,
hard entries RFN9, and scalar RFN12.

Use the additional exact change
\[
 q=vz-Tu,\qquad L=M_0+Hq,\qquad
 M_0=L-Hq,\qquad L=\sqrt{S\ell\bar h}\,\widehat L .
                                                               \tag{RL11}
\]
All divisions are pure source or pure testing.  In particular
this change does not invert the harmonic mixed unit
\(1+\rho u^2\).
The complete diagonal of \(\widehat L\) is
\[
 S\ell\bar h\{D_u(CD_v+D+\tfrac12)+AD_v+
                   (\alpha+\tfrac12)v^2\},
                                                               \tag{RL12}
\]
with positive real residue \(\overline k/2\) and bounded
coefficient array through both openings.

Use the circuit coordinates \(U,V,Z_*,Q_*,B_*,R_*,\Psi\)
of SH1--SH2, and \(b=v-uz/j=\sqrt y B_*\).
Before spending any scalar \(n\)-mass, complete the exact correlated
block, as in RFN13--RFN16:
\[
 h_0N_0^2+n(N_0+jb)^2
 =(n+h_0)(N_0+c_bb)^2+M_e b^2,\quad
 c_b=\frac{nj}{n+h_0},\quad M_e=\frac{nh_0j^2}{n+h_0}.  \tag{RL13}
\]
Keep \(nv^2(N_0+pb)^2\) and all three additional positive hard
squares RFN15 whole, including after the change
\(\widehat N=\sqrt{n+h_0}(N_0+c_bb)\).
At \(n=0\), put \(c_b=M_e=0\) and omit the current row.
For \(n>0\), \(c_b^2/M_e=n/[h_0(n+h_0)]\) is bounded.

\paragraph{The exact shifted scalar.}
Let \(F\) be RFN12.  The shifted scalar is
\(F-2HqL_M+H^2q^2H_M\).
Use the whole \(d_cV^2Q_*^2\) term in its \(q^2\) diagonal.
Using the exact actual-source identity \(S_0=0\) in RFN22,
the resulting coefficient on \((\sqrt S q)^2\) is exactly
\[
 D_q=S\{A\bar h^2+\ell^2(C+D+\tfrac12)D_u\}
       -\frac{f_c}{Tx}-TH^2C\,u^2v^2
       -\frac{k_cv^2}{yTx}.                            \tag{RL14}
\]
It has bounded array and positive real residue \(\overline k/2\).
Indeed \(f_c/(Tx)=O(jy)\) is infinitesimal, \(TH^2C=\ell CH\)
is bounded, and \(\kappa/(yTx)=\kappa/(jy^2)\) is bounded.
Its equivalent \(Q_*^2\) mass has order \(jy\).
The only extra scalar cross terms are
\[
 -2Hq\{\mu uvb+v(\alpha-yCHu^2)R\},\qquad
 R=z-uv/(Sy).                                        \tag{RL15}
\]
The budgets below allocate their entire expressions.

The remaining signed scalar terms are exactly
\(e_0V^2B_*^2+k_cZ_*^2(V^2-1)^2\), where
\(e_0=p^2y^2/2-g_c\).
The negative part of \(e_0\) is at most \(g_c\).
SH2 purchases it at relative circuit cost
\[
 \frac{g_c}{f_c}(1+V^2)
   =O\bigl(S(y+v^2)\bigr)\in\mathrm I.                \tag{RL16}
\]
Use \(Z_*(V^2-1)=-VQ_*-R_*\) for the other signed term.
If \(p\) is infinitesimal, (RL8) pays its opposite part
directly from \(r_cR_*^2\), and its \(VQ_*\) cost is
\(O(v^2)\).
Otherwise \(p,j,M_e\) are positive units.  Young's exact
square identity with parameter \(1/6\), followed by SH1 at
\(a=3/2\), gives the legal SOS domination
\[
 \kappa Z_*^2(V^2-1)^2
 \preceq_{\rm sos}
 \frac74\kappa\Psi+
 \frac{21}{2}\kappa V^2Q_*^2+
 \frac72\kappa V^2B_*^2 .                             \tag{RL17}
\]
The circuit cost is at most \(7/8+o(1)\), by (RL9).
The two other relative costs are \(O(v^2)\).
Thus fixed positive real fractions of \(D_q(\sqrt S q)^2\),
\(r_cR_*^2\), \(f_c\Psi\), and \(M_eb^2\) remain;
the last is absent only at \(n=0\).
The opposite mass \(r_c\) was not spent in the second case.

\subsubsection*{All helper rows and both complete normals}

As usual, a free row permits any fixed real square pivot at
infinitesimal cost from the preceding retained scalar forms.
SH2 proves that both \(E_u=u(uvz-jy)\) and \(E_v=v(uvz-jy)\)
are free.  The row \(R_u=v(u^2-x)\) is free as well.
For \(p\) infinitesimal, use the actual RFN19 rational Gram:
its three costs are \(O(S^2y)\), \(O(Syu^2)\),
and \(O(Su^2/(1-p))\).
When \(p,j,M_e\) are units, apply SH1 to the \(R_*^2\)
term in RFN19.  The resulting costs are
\[
 O(S^2y+Su^2/j),\quad O(Syu^2+Su^2v^2),\quad
 O(jSu^2v^2/M_e),                                    \tag{RL18}
\]
all infinitesimal.

The rows \(pv^2b\) and \(\sqrt S\,\mu uvb\) are free.
If \(\eta_2\) is infinitesimal, (RL6) and (RL13) imply
\[
 \frac{p^2}{M_e}
  =\frac{p^2}{h_0j^2}+\frac{p^2}{nj^2}=O(1).
                                                               \tag{RL19}
\]
They then cost \(O(v^4)\) and \(O(Su^2v^2)\).
If \(\eta_2\) is a positive unit, the actual source gives
\(p^2y/f_c=O(1)\) and \(\mu=O(p)\).
The second identity in SH2 provides the same rows with costs
\[
 O\bigl(v^2(y+v^2)\bigr),\qquad
 O\bigl(Su^2(y+v^2)\bigr),                             \tag{RL20}
\]
including \(n=0\).

The full opposite term \(\sqrt S\,v(\alpha-yCHu^2)R\)
is free by SH4--SH8.  Also \(j\alpha b\) is free:
\[
 j\alpha b=-(\alpha/B)R_u-\alpha uR,\qquad
 \alpha/B=O(y(j+1))\in\mathrm I_c.
                                                               \tag{RL21}
\]
For \(p\) infinitesimal, the direct \(r_c\) cost for
\(\alpha uR\) is \(O(Sy^2(j+1)^2u^2)\).
In the unit \(p,j,M_e\) case use SH1; the three costs
are \(O(Sy^2u^2)\), \(O(Sy^2u^2v^2)\), and
\(O(Sjy^2u^2v^2/M_e)\).  All are infinitesimal.
These statements account for the whole source \(\alpha\), not
only its leading coefficient.

\paragraph{The first normal.}
The exact RFN21 source cancellations give, after (RL11),
\[
 L_N'=[j(n-\alpha)+p(n+\tfrac12)v^2]b
             +K_sR_u+\kappa_NE_v,                     \tag{RL22}
\]
\[
 K_s=-\mu+(\alpha+\tfrac12)(\ell+p-1),\qquad
 \kappa_N=-(n+\tfrac12)(1-p/j).
\]
Both coefficients are bounded.  For \(\kappa_N\), use the
exact numerator of \(n\) and
\(n(j-p)/j=O(p/[j(j+1)])\); this cancels the \(r\) opening.
After the entire correlated completion (RL13), the remaining
normal is
\(-j\alpha b+\tfrac12pv^2b+K_sR_u+\kappa_NE_v\).
Every row is free by (RL18)--(RL21).
Division by \(\sqrt{n+h_0}\) is bounded.
The shift generates couplings to \(\sqrt{M_e}b\) with bounded
coefficients \(c_b/\sqrt{M_e}\); they are absent at \(n=0\).
Large fixed real helper pivots absorb these couplings by exact
Young identities.

\paragraph{The second normal and the entire hard block.}
The complete RFN22 normal after (RL11) is
\[
 L_L=(C_q-HH_M)q+\mu uvb+v(\alpha-yCHu^2)R,\qquad
 C_q-HH_M=-D-H(C+D+\tfrac12)u^2 .                      \tag{RL23}
\]
The coefficient of \(\sqrt S q\) after the normalization
of \(L\) is infinitesimal; \(D\) is infinitesimal and
\(H(C+D+1/2)\) is bounded.
The other two terms are free by (RL20) and SH6--SH8.
Exactly the same rows pay (RL15), whose first factor is
\(\sqrt S q\) with a bounded source multiplier.

The remaining hard off-diagonal is
\((\alpha+1/2)uv\sqrt{S\ell\bar h}/\sqrt{n+h_0}\),
which is infinitesimal.  Its \(c_b b\) correction has the
same bound after division by \(\sqrt{M_e}\).
The complete finite Gram therefore has a positive definite
real residue, with principal normal pivots
\[
                   1,\quad\overline k/2,\quad\overline k/2
                                                               \tag{RL24}
\]
on \(\widehat N,\widehat L,\sqrt S q\), together with the
retained strict scalar and purchased helper pivots.
All extra positive hard squares and both original retained
current positive square terms are restored whole.  The strict bounded-Gram
lemma proves tensor SOS for the entire parent.

\subsubsection*{State and outer-chart exhaustion}

The fresh full-state sign lemma SJ1--SJ10 permits positive
\(u,v\), and on the remaining negative-core chart gives
\(wz>0\).  Divide by the nonzero pure-testing \(w\) to obtain
the normalization used above with arbitrary \(z>0\).
No estimate in this proof uses its size.  Literal axes and
homogeneous endpoints have the direct original-core SOS exit.
Real and imaginary state parts separate.
The exact outer theorem R2O1--R2O24 covers every other testing
pair on all original R2 sources.  This proves the stated
all-tests conclusion on (RL1).

The complementary source \(a_s\) unbounded and \(p\) a positive
unit is not inferred from this proof.  It requires the separate
TU source comparison and additional-current conductor.

\subsection{Circuit budgets for the opposite normal}
\label{component:p09b2107cc05e}

Use the compact actual-source dictionary of RSR1--RSR8.  In the R2
source put \(S=1/T\), \(x=jSy\), \(B=Sy\), and \(H=S\ell\), where
\(0<\ell<1\).  Both testing angles \(u,v\) are positive
infinitesimals.  Normalize the rotated testing coordinate \(w=1\);
the other rotated coordinate \(z>0\) is arbitrary, including
unbounded values.  No coefficient estimate below involves \(z\).

Set
\[
 U=\frac{u}{\sqrt x},\quad V=\frac v{\sqrt y},\quad
 Z_*=\frac z{\sqrt{j/S}},\quad
 Q_*=U-VZ_*,\quad B_*=V-UZ_*,\quad R_*=Z_*-UV,
\]
\[
 \Psi=U^2+V^2+Z_*^2+U^2V^2Z_*^2-4UVZ_*.
\]
All monomials and radicals separate between the two fields.

\paragraph{A parameterized opposite-square identity.}
For every positive real \(a>1\), let \(K=a/(a-1)\),
\(d=U-V\), \(t=UV\), \(E=Z_*(1-t)\), \(T_*=t(Z_*-1)\).
Exact multiplication gives
\[
\begin{split}
 a\Psi+K V^2(Q_*^2+B_*^2)-R_*^2
 ={}&a d^2+(a-1)(E-T_*/(a-1))^2\\
 &+(2at+KV^4)(1-Z_*)^2+2KZ_*(t-V^2)^2 .
\end{split}                                                     \tag{SH1}
\]
The right side is a finite legal separated SOS.  For example
\(a=3/2\) gives \(K=3\).  This is a different budget from the
older one with \(U^2B_*^2\); neither budget dominates the other
uniformly at independent source and testing rates.

The two symmetric current budgets are
\[
\begin{split}
 (1+U^2)\Psi&=R_*^2+B_*^2+U^2Q_*^2
                         +U^2(1-UVZ_*)^2,\\
 (1+V^2)\Psi&=Q_*^2+R_*^2+V^2B_*^2
                         +V^2(1-UVZ_*)^2 .
\end{split}                                                     \tag{SH2}
\]
In particular the two actual rows
\[
 E_u=u(uvz-jy),\qquad E_v=v(uvz-jy)                       \tag{SH3}
\]
are free from a fixed positive fraction of \(f_c\Psi\) whenever
\(f_c\asymp j^2y^2\).  Their respective relative circuit costs
are \(O(x+u^2)\) and \(O(y+v^2)\), both infinitesimal.  Here
\emph{free} means that an arbitrarily large fixed real multiple of
the row square can be inserted at infinitesimal cost in the
strict circuit Gram.  It is not a pointwise bound without an
SOS witness: multiply (SH2) by \(j^2y^2x\) or \(j^2y^3\).

\paragraph{The full opposite row after the harmonic shift.}
Let
\[
 \begin{gathered}
 q=vz-u/S,\quad R=z-\frac{uv}{Sy},\quad R_v=u(v^2-y),\\
 d_s=\frac{BCH(1-p)(1+x)}{E},\qquad E=1+B(HT+p-1).
 \end{gathered}
\]
The exact actual source identity is
\[
                 \frac{\alpha}{y}-CHx=d_s.             \tag{SH4}
\]
Indeed it is equivalent to
\((K+Hjy)-jE=(1-p)(1+x)\), after canceling the displayed
positive source factors in the RSR dictionary.
The two exact testing identities
\[
 u(u^2-x)=S E_u-Su^2q,\qquad
 v(u^2-x)=S E_v-Suvq                                  \tag{SH5}
\]
now give the complete identity
\[
\begin{split}
 \sqrt S\,v(\alpha-yCHu^2)R
 ={}&[\alpha-yCHu^2-CHu^2(v^2-y)]\,(\sqrt S q)\\
 &+\sqrt S\,CH(v^2-y)E_u-\frac{d_s}{\sqrt S}R_v .
\end{split}                                                     \tag{SH6}
\]
There is no restriction on \(z\), and no normal has been set to
zero.  The first coefficient in (SH6) is infinitesimal, because
\(CH=O(1)\), \(y,u,v\) are infinitesimal, and
\(\alpha=O(Sy^2(j+1))\) is infinitesimal by \(jy\in\mathrm I\).
The second coefficient is infinitesimal and (SH3) is free.

For the last term write
\[
 R_v=-\sqrt x\,y(Q_*+VR_*),\qquad
 \frac{d_s}{\sqrt S}
   =\sqrt S\,y\,CH(1-p)(1+x)/E.                         \tag{SH7}
\]
Suppose the scalar has retained strict pivots of respective orders
\(jy\) and \(jy^2(1-p)\) on \(Q_*^2,R_*^2\), as in the actual
R2 source.  Twice Young's identity and (SH7) provide this row at
relative costs
\[
 O(S^2y^4(1-p)^2),\qquad
 O(S^2y^2(1-p)v^2),                                   \tag{SH8}
\]
respectively.  Both are infinitesimal even when \(1-p\) has
arbitrarily small positive order.  No inverse of that opening
remains.  Multiplication by the additional bounded source
factor \(\sqrt{\ell(1-\ell)}\), used to normalize the complete
second normal, preserves the estimate.

This controls the opposite normal uniformly through all R2 source ratios
once the stated scalar pivots are retained. The signed scalar allocation
remains a separate step.

\subsection{Unbounded source ratio and bounded relative current}
\label{component:pb88ccd68a4bf}

\label{pb88ccd68a4bf:sec:r2-thin-new}

\paragraph{Statement and domain.}
Use the actual original weak-C2 source RER1--RER8, normalized by
$d=1$, with $\lambda=\beta=0$ and every original strict gap.
Suppose $t,a$ are positive unbounded, while $B,aB,tB$ are
positive infinitesimals, and put $e=1-r$.  Assume
\[
             \lambda_n=n/(tB)=O(1),\qquad n\geq0.       \tag{TU1}
\]
The entire original parent is tensor SOS for every independent
projective test and every state.  Literal $n=0$ is included.
No comparison of a testing quantity with a source quantity is made.
This theorem treats precisely this source sector; it is not a
completion of the other R2 sectors or the global finite-lift proof.

The bounded-Gram and finite toric allocation lemmas used below are
the actual retained SC16--SC20 and QI1--QI10 proofs.  We use their
printed finite algebra and prove its source bounds on the present
domain.  We do not apply the physical C1 theorem to a formal source
whose event coefficient could be negative.

\subsubsection*{An exact common contact and bounded source comparison}

Set
\[
 S=1/t,\quad z=1/a,\quad q=a(tB),\quad w=ae,\qquad
 b=\sqrt S,\quad g=\sqrt z,\quad c=\sqrt q.
                                                               \tag{TU2}
\]
Thus $B=Szq$, $e=zw$, the actual squared angles are
$u_c^2=Sq$ and $v_c^2=zq$, and $S,z,q$ are positive
infinitesimals.  Keep $h$ as the original C2 source parameter.
The strict $m$ gap gives
\[
 R=1+Sqh,\quad M=R-w>0,
 \qquad 0<w<R\sim1.                                      \tag{TU3}
\]
Both original openings $h,1-h$ remain arbitrary positive quantities.

The exact RER2 identity, divided by $e>0$, gives
\[
 Q_{\rm RER}
 =\tau\{(1+z)(1+Sh)(1+Sq)+2r(z\Delta)\lambda_n\},
 \qquad \tau=q/(zw),\quad \lambda_n=n/(qz).
\]
Here $Q_{\rm RER}\sim1$, $z\Delta\sim1$ and $r\sim1$.
The bracket is a positive bounded unit, since $\lambda_n\geq0$
is bounded and its first term is a positive unit.  Consequently
\[
 \tau\text{ is a positive bounded unit},\qquad
 \overline\tau=\frac1{1+2\overline\lambda_n},\qquad
 q=zw\tau.
                                                               \tag{TU4}
\]
No division by $n$ has occurred.  This sector includes the earlier
thin condition $\rho_n=n/e^2\in\mathrm I$: in that case the same
RER identity gives $\tau\sim1$ and
$\lambda_n=\rho_nw/\tau\in\mathrm I$ directly, without first
assuming (TU1).

For the comparison only, substitute into the native C1 rational
formulas NC1--NC4
\[
 k^2=Sq,\quad \eta=1/S,\quad \zeta=z,\quad
 r_{\rm C1}=h,\quad t_{\rm C1}=w,\quad
 a_{\rm C1}=S(1-h).                                     \tag{TU5}
\]
Their $\theta$ is $e=zw$ and their $R,M$ are exactly (TU3).
Write a superscript $0$ for these formal C1 coefficients.
All their source denominators have the requisite positive signs:
\[
\begin{gathered}
 K=1+Sq,\quad J=1+Sq(1-e)-S(1-h),\quad
 E=1-q(1-e),\quad Q=1-Sqzh,
\end{gathered}
\]
and $J,E,Q$ are positive bounded units.  Their common polynomial is
\[
 P=1+qz+Sq(2+z-2e)+Sq^2(z+e-e^2)+S^2hq^2(1+z-e).
                                                               \tag{TU6}
\]
In particular $P\sim1$.  We make no sign claim for $\beta^0$.

The two sources have exactly the same state contact.  With $Z_c=1$,
\[
 u_c=bc,\quad v_c=cg,\quad W_c=bg,\quad
 Y_c=hu_cW_c,\quad X_c=-(1-e)u_c.                       \tag{TU7}
\]
This follows either by substitution in the two native dictionaries,
or from $Y_c=hu_cW_c$, $X_c=-ru_cZ_c$ and
$Z_c/W_c=\sqrt{at}$ in the actual C2 source.  In particular
the comparison does not move the contact.

Use the actual, not formal, normalized coefficients
\[
 c_b=SC,\ a_b=SA,\ D_b=D/S,\ \alpha_b=\alpha/(Sz),\quad
 h_2=mzM,\quad h_b=zw/2.
                                                               \tag{TU8}
\]
They satisfy $m=h_2/(zM)$ and $\beta+1/2=h_b/(zw)$.
Use the same normalizations with superscript $0$, and let $\delta$
mean actual minus formal.  The following explicit polynomials make
all needed cancellations finite and opening-safe.  Put $G=1+z-e$,
and define
\begin{align*}
 P_C={}&1+qz+S\{1+q(2+2z-e)+q^2z\}\\
 &+S^2\{q[h(1+z)+1-e]
          +q^2[hez-e(1+z)+2z+1]\}+S^3hq^2G,\\
 P_A={}&1+qz+S\{q(2+z-2e)+q^2z(1-e)\}
                 +S^2q^2G(G-hz),\\
 P_D={}&1+q(z+e)+S\{q[h(1+z)+1-e]+q^2(z+e-e^2)\}
                 +S^2hq^2G,\\
 P_m={}&1+qz+S\{q(2+z-2e)+q^2z\}
                 +S^2q^2\{(1-e)^2+z(1-e)+hez\},\\
 P_b={}&(1+z)(1+qz)+Sq(1+z)(2+z-2e)\\
       &+Sq^2z(1+z-e^2)+S^2q^2G^2.
\end{align*}
Then the exact rational identities are
\[
\begin{aligned}
 (1-h)\delta c_b&=-\frac{n(1-e)P_C}{zP},&
 h\delta a_b&=-\frac{n(1-e)P_A}{zP},\\
 \delta D_b&=-\frac{n(1-e)(1+qz)P_D}{zJP},&
 \delta h_2&=\frac{n(1-e)(1+qz)P_m}{zP},\\
 \delta h_b&=\frac{n(1-e)P_b}{zP},&
 \delta\alpha_b&=
 \frac{n(1-e)q^2(1+qz)[(1+S)e-S(1-h)(1+z)]}{zQP}.
\end{aligned}                                                   \tag{TU9}
\]
These identities are proved simply by substituting the six printed
RER1 weights and the six NC4 weights, clearing the indicated pure
source denominators, and multiplying.  They include $n=0$; the
right sides define no quotient by $n$.

After substituting $e=zw$, the polynomial
\[
       P_\Delta=\frac{P_b-(1+qz)P_m}{z}                 \tag{TU10}
\]
is a genuine finite polynomial in $S,z,q,w,h$.  Indeed the numerator
specialized at $z=0$ is zero identically, so ordinary monic division
by $z$ has zero remainder.  Consequently
\[
 \delta(h_b-h_2)=n(1-e)P_\Delta/P.
                                                               \tag{TU11}
\]
Every polynomial in (TU6), (TU9), (TU10) has bounded coefficients
and bounded arguments.  From (TU4) and (TU9)--(TU11) we obtain
\[
\begin{gathered}
 (1-h)\delta c_b,\ h\delta a_b,\ \delta D_b
                     =q\lambda_n\,O(1),\qquad
 \delta h_2,\delta h_b=q\lambda_n\,O(1),\\
 \delta\alpha_b/q^2=\lambda_n q\,O(S+e),\qquad
 \delta(h_b-h_2)/(qz)=\lambda_n O(1).
\end{gathered}                                                   \tag{TU12}
\]
Here $O(1)$ denotes an explicit bounded rational source array with
denominators among the units $P,J,Q,R$; it does not hide a state
or testing factor.  These conclusions require no lower bound on
$h$, $1-h$, or $M$.

In particular $c_b\asymp1/(1-h)$, $a_b\asymp1/h$,
$D_b\sim1/2$, $\alpha_b/q^2\sim1/2$, and
$h_2/q,h_b/q$ positive bounded units with common residue
$1/2+\overline\lambda_n$.  We also have
\[
 h_b/w=z/2,\qquad h_b-h_2=O(qz),\qquad
 H_0:=4h_2/R+4h_b/w\geq2z,
 \quad H_0\asymp q+z.                                  \tag{TU13}
\]
The last statements use the formal exact SC2 difference together
with (TU11), and hold also when $M$ is arbitrarily small.

\subsubsection*{The complete additional current, including its normal}

Let the testing $u,v$ be positive infinitesimals and first take
$Z=1$ on its pure-testing state chart.  Use the exact SC coordinates
\[
 Q_b=(W+uY)/b,\quad Y_b=Y/S,\quad p=w/R,\quad
 N=((X+u)/b-pvQ_b)/g,
\]
\[
 H=uQ_b/b,\quad L=vQ_b/g,\quad
 \sigma=c_b+a_b,\quad \rho_0=Kh/R,\quad
 N_y=\sqrt\sigma\{D_uY_b-\rho_0H\}.
                                                               \tag{TU14}
\]
Both $N,N_y$ vanish at (TU7).  The exact SC3 full-parent identity
is an algebraic identity for arbitrary $c_b,a_b,D_b,\alpha_b,h_2,h_b$;
no C1 coefficient relation is used in deriving it.  With (TU8),
it represents all the actual C2 parent except its additional term
$4nD_uD_vX^2/S$.  Retain its original nonnegative $m$ positive square term
$4h_2wN^2/(MR)$, incorporate the second paired positive square term into
$H_0N^2$, and retain $4h_b(vN)^2/w$ exactly as in the uniform
QI normal lemma.  Define $\widehat N=\sqrt{H_0}N$.

Put $h_t=H/R$, and introduce the same four positive separated
monomials and their geometric mean as in QI:
\[
 A_*=ch_t,\quad B_*=cL/R,\quad C_*=uv/b,\quad
 D_*=qg,\qquad \Gamma=(A_*B_*C_*D_*)^{1/4}.
                                                               \tag{TU15}
\]
Their exact relations include
$\Gamma^2=qh_tv$ and $B_*C_*=(c/g)v^2h_t$.
The extra monomial
\[
                   U_*=cgu/b
\]
has the two exact identities
\[
       U_*B_*=\Gamma^2,\qquad
       cgX/b=czN+eB_*-U_*.
                                                               \tag{TU16}
\]
The entire additional current therefore has the following square
completion, with no state normal set to zero:
\begin{align*}
 \frac{4n}{S}D_uD_vX^2
 ={}&4\lambda_nD_uD_v(U_*-B_*-czN)^2\\
 &+4\lambda_nD_uD_v
   \{2(1-e)\Gamma^2-(1-e^2)B_*^2
                         -2(1-e)czB_*N\}.
                                                               \tag{TU17}
\end{align*}
The first line is a nonnegative separated SOS and vanishes at the
actual contact, including its unrestricted state-normal coordinate.
Its coefficient is literally zero if $n=0$.  The remaining scalar
changes have bounded coefficients on $B_*^2$ and $\Gamma^2$.
Their possibly nonzero real residues will cancel exactly with
the source comparison below.  The full additional first-normal coefficient, on
$B_*\widehat N$, is infinitesimal because
\[
       \lambda_n cz/\sqrt{H_0}=O(\lambda_n c\sqrt z)
                                            \in\mathrm I.
                                                               \tag{TU18}
\]
In particular (TU17) does not estimate or discard $u/b$, which can
be arbitrarily large.  The new positive square contains that entire
row, while the exact product relation in (TU16) places its cross
term on the old angular circuit.

\subsubsection*{Why the full QI allocation is stable here}

We verify the actual finite hypotheses of QI1--QI10 rather than
an abstract continuity assertion.  Its coefficient arrays have
only the pure-testing denominators $D_u$ and positive pure-source
units after cancellation of the two $h$ openings.  The latter
cancellation is preserved exactly because
\[
 1-\rho_0=(1-h)/R,\qquad \rho_0=Kh/R.
                                                               \tag{TU19}
\]
For example all perturbations of
$c_b(1-\rho_0)^2$, $a_b\rho_0^2$,
$c_b(1-\rho_0)$ and $a_b\rho_0$ are $q\lambda_n$
times bounded rational arrays, by (TU9).  Moreover
\[
 \delta(\sigma\rho_0-c_b)
       =\rho_0\delta a_b-(1-\rho_0)\delta c_b
                             =q\lambda_n O(1).
                                                               \tag{TU20}
\]
This is the complete potentially singular harmonic first-normal
coefficient after multiplication by $\sqrt\sigma$.
Both $\sigma$ and its formal comparison are comparable to
$[h(1-h)]^{-1}$, and their ratio has residue one.
Their square-root ratios and inverses are consequently bounded
pure-source units.  Thus this change of hard normal loses no
opening bound.

In the exact generic scalar QI2, the changes of $a,b_*,c_*$ are
$q\lambda_n$ times bounded arrays in $u,v$.  This follows term by
term from (TU12), (TU19): its only $C,A$ occurrences are the first
two expressions in (TU19); its $D_b$ occurrences have bounded
factors; its $\alpha_b$ occurrences have bounded factors and the
additional small factors printed in QI2.  Consequently the
formal specialization $a_0,b_0,c_0$ and its exact $U,V$ divisions
in QI3--QI4 can be used unchanged, while the remaining $q$
coquotients change by bounded arrays.  Write
$\lambda_0=\overline\lambda_n$.  Their complete leading changes
are respectively $0,-8\lambda_0,4\lambda_0,0$: substituting
$S=z=q=u=v=0$ in the bounded coquotients (TU9) gives
$\delta a/q=-4\lambda_0$, $\delta b_*/q=0$,
$\delta c_*/q=-4\lambda_0$, and
$\delta(R^2P_L)/q=4\lambda_0$.  Thus before the extra current
the four residues are
\[
 \overline{a_1-c_1}=2,\qquad
 \overline{b_1+2c_1}=-8-8\lambda_0,\qquad
 \overline{R^2P_L/q}=2+4\lambda_0,\qquad
 \overline{4\alpha_bD_u/q^2}=2.
                                                               \tag{TU21}
\]
Equation (TU17) adds $-4\lambda_0$ to the $B_*^2$ residue
and $+8\lambda_0$ to the $\Gamma^2$ residue.  Consequently
the \emph{combined} full scalar has exactly the original four
residues $(2,-8,2,2)$ on $(A_*^2,\Gamma^2,B_*^2,D_*^2)$.
All differences from those residues are infinitesimal bounded
arrays, because $\lambda_n$ itself is bounded.  This cancellation
is why a thinness assumption is unnecessary on (TU1).
The one additional bound QI uses from physical $\beta\geq0$ is
also available directly, without any sign for $\beta^0$:
\[
  w/(c/g)=\sqrt{w/\tau}=O(1),\qquad w=O(1).
                                                               \tag{TU22}
\]
Thus the exact dangerous cross term $-4(wv/(c/g))D_uB_*C_*$
still has infinitesimal coefficient.

For the complete first $N$ normal, (TU13) gives exactly the two
bounds used in the uniform QI normal lemma:
$g/\sqrt{H_0}$ and $c/\sqrt{H_0}$ are bounded, and
\[
 \frac{h_bD_v-h_2}{\sqrt{qH_0}}=O(z+v^2).
                                                               \tag{TU23}
\]
The $H$ term and all other terms in SC14 therefore have the same
infinitesimal coefficients on the QI rows; (TU18) handles the whole
extra term.  For the harmonic normal, (TU20) adds a
$q\lambda_n$ bounded array to the left side of QI10.  It can be
absorbed into its existing $qK_q$ coquotient.  All other terms are
the same bounded rational arrays multiplied by the opening-safe
quantities in (TU12), (TU19).  The exact $D_b$ root and its two
normal terms in QI9 remain unchanged as algebraic identities.
The complete hard block therefore retains the real residue
$\operatorname{diag}(1,4)$ in $(\widehat N,N_y)$.

For clarity, here is the finite allocation being invoked.  Use
the QI roots
\[
 E_*=h_t-v,\quad F_*=vh_t-gD_*,\quad
 R_F=\frac{\sqrt{g\Gamma}}c(\Gamma-D_*),\quad
 R_E=\sqrt\Gamma(\sqrt{A_*}-\sqrt{cv}),
\]
together with the complete finite toric root list SC16.
The exact identities QI7--QI8 give strict real square pivots on
$F_*,R_F,R_E$; QI6 gives the strict $E_*$ pivot and twice the
four-monomial circuit.  Every remaining coefficient is an
infinitesimal bounded rational array in the testing $u,v$, with
the same finite denominators and divisions just verified.
All removed positive square terms, including (TU17), vanish at (TU7).
The actual source contact value, its three scalar first jets,
and both full first-normal contact values consequently vanish
exactly for this remaining expression.  Applying the two finite
monic divisions SC17--SC20 and the displayed QI7 first-jet
identities therefore puts the entire scalar correction and
both complete first normals on this strict root list with
infinitesimal coefficients.  The coefficient functions are
independent of the state ratio $Q_b/Z$; no boundedness of that
ratio is assumed.  This is precisely the constructive QI
allocation, with its source inequalities checked in
(TU12)--(TU23), not a conclusion from a scalar restriction alone.

The resulting finite Gram is bounded and has positive definite
real residue.  The strict bounded-Gram lemma supplies tensor SOS.
Restore the original positive paired positive square terms and the new full
square in (TU17).  This proves the local assertion for the entire
parent on the positive small-angle, negative-core chart.

\subsubsection*{States, signs, and all remaining tests}

Every inverse and radical used above is pure source, pure testing,
or the explicitly separated positive monomial radical in (TU15).
For $Z\ne0$ the normalization $Z=1$ is pure testing; multiplying
back by $Z^2$ restores homogeneity.  At $Z=0$ the original
pure-testing core restriction is nonnegative.  The exact core
Schur identity and sign gauges used in the retained QI proof
force $W+uY>0$ on the remaining negative-core chart, and hence
all four monomials are positive.  Otherwise the core and all
original event positive square terms already give the direct exit.  Literal
zero-angle charts have that same original exit.  Complex states
split into their real and imaginary parts, since the real
measurement matrices and full quadratic parent are unchanged.

The new R2 outer theorem R2O1--R2O24 treats every projective
testing pair outside the double-infinitesimal original chart,
with no condition on $a$ or $n/e^2$.  Its hypotheses are exactly
the original RER source conditions already retained here.
Combining that result with this local proof gives the claimed
all-tests, all-states theorem on (TU1).  No other source branch
is proved by this conclusion.

\subsection{Unbounded source ratio and relative current}
\label{component:pef021eabacde}

\label{pef021eabacde:sec:r2-unbounded-current-new}

\paragraph{Precise assertion.}
Retain the original RER source, all its strict gaps, and the notation
TU2--TU16 of the preceding bounded-relative-current proof.
Assume the native $a,t$ are positive unbounded and
\[
             \lambda=\frac{n}{tB}=\frac{n}{qz}
                   \text{ is positive unbounded}.                 \tag{UL1}
\]
Then the whole original weak-C2 parent is tensor SOS for every
independent projective test and every state. The argument below
complements TU1; it never assigns a sign to the formal C1 coefficient
$\beta^0$, and never replaces the original physical source by that
formal source.

\subsubsection*{Source bounds without a thinness assumption}

The exact RER physicality equation already used in TU4 is
\[
 Q_{\rm RER}=\tau\{A+2r(z\Delta)\lambda\},\qquad
 \tau=\frac q{zw},\qquad
 A=(1+z)(1+Sh)(1+Sq).
\]
Here $A,Q_{\rm RER},r,z\Delta$ are positive units with residue one.
Since $\lambda$ is unbounded, this proves, with no cancellation of
mixed factors,
\[
 \tau\in\mathrm I,\qquad \lambda\tau\sim\frac12,\qquad
 \lambda q\sim\frac{zw}{2},\qquad
 \frac n{z^2w}\sim\frac12.                                     \tag{UL2}
\]
All four divisions are pure source. Recall $0<w<R\sim1$.
The exact TU9 formulas therefore imply
\[
\begin{gathered}
 (1-h)\delta c_b,\ h\delta a_b,\ \delta D_b,\
            \delta h_2,\delta h_b=\lambda q\,O(1),\\
 \delta\alpha_b/q^2=\lambda q\,O(S+zw),\qquad
 h_2\asymp zw,\quad h_b=zw/2,\\
 H_0=4h_2/R+2z\asymp z,\qquad H_0\geq2z,\\
 \frac g{\sqrt{H_0}}=O(1),\quad
 \frac c{\sqrt{H_0}}=O(1),\quad
 \frac{\lambda q}{H_0}=O(1),\quad \frac n{z^2}=O(1).
\end{gathered}                                                   \tag{UL3}
\]
In particular $\lambda q$ is infinitesimal. The displayed coefficient
arrays have denominators only among the pure-source units
$P,J,Q,R$ after the original $h,1-h$ openings cancel.
The actual and formal $\sigma=c_b+a_b$ differ by a relative
infinitesimal. Their square-root ratio is a bounded unit.
No lower bound on $w,h,1-h$ or the strict gap $M$ is assumed.

Use the complete generic SC3 identity with the actual coefficients.
Retain its positive $m$ positive square term and use $\widehat N=\sqrt{H_0}N$.
The formal comparison is used only for finite scalar and normal
identities. Its $q$-coquotients and the QI coefficient arrays are
bounded here: their derivation requires $S,z,q$ infinitesimal and
$w$ bounded, but not $\beta^0\geq0$. The sole QI bound that is no
longer available, $w/(c/g)=O(1)$, is replaced explicitly below.

\subsubsection*{The entire extra current supplies a new angular pivot}

Use the exact positive-chart rows TU15--TU16, denoted here by
$A_*,B_*,C_*,D_*,U_*$ and $\Gamma$. Define
\[
 J_3=\sqrt\lambda(U_*-B_*),\qquad
 E_*=h_t-v,\qquad F_*=vh_t-gD_*.
                                                               \tag{UL4}
\]
The first line of the exact whole-current identity TU17 is
\[
 4D_uD_v\left(J_3-d_N\widehat N\right)^2,\qquad
 d_N=\frac{\sqrt\lambda\,cz}{\sqrt{H_0}}
       =\sqrt{\frac{nz}{H_0}}\in\mathrm I.                       \tag{UL5}
\]
Thus it supplies a strict real pivot on $J_3$, with only
infinitesimal changes to the hard normal Gram. Its scalar
remainder is still exactly
$8\lambda D_uD_v(1-e)\Gamma^2
-4\lambda D_uD_v(1-e^2)B_*^2$.
Its complete first-normal remainder, retained below, is
$-8\lambda D_uD_v(1-e)czB_*N$.

Write $\widetilde c,\widetilde a,\widetilde d,\widetilde\alpha,
\widetilde m,\widetilde b$ for the six differences
$\delta c_b,\delta a_b,\delta D_b,\delta\alpha_b,
\delta h_2,\delta h_b$ divided by $\lambda q$.
They are defined explicitly by the six right sides in TU9, using
$n=\lambda qz$; their apparent opening poles cancel in every
combination below. Put $U=u^2,V=v^2$, and retain
$\rho_0=Kh/R$, $a_t=1-\rho_0+(w/R)V$.
Define the exact rational arrays
\[
\begin{split}
\mathsf a={}&4D_vR^2\{\widetilde c(1-\rho_0)^2
                      +\widetilde a\rho_0^2/D_u\}
 +4\widetilde\alpha zS^2VR^2\rho_0^2/D_u
 +4\widetilde d S^2R^2a_t^2,\\
\mathsf b={}&-8\widetilde\alpha zSR\rho_0
                         -8\widetilde d SRa_tD_u,\\
\mathsf c={}&4\widetilde dD_u^2,\\
\mathsf l={}&4R\{\widetilde m(1-w/R)
                              +\widetilde b(w/R)D_v\}.
\end{split}                                                     \tag{UL6}
\]
For example $\delta a/q=\lambda\mathsf a$ in QI2,
and $\delta(R^2P_L/q)=\lambda\mathsf l$.
Direct specialization of TU9 gives, keeping $S,z,w,h$ exact,
\[
\begin{split}
 (\mathsf a-\mathsf c)|_{q=U=V=0}&=0,\\
 (\mathsf b+2\mathsf c+8D_uD_v(1-e))|_{q=U=V=0}&=0,\\
 (\mathsf l-4D_uD_v(1-e^2))|_{q=U=V=0}&=0.
\end{split}                                                     \tag{UL7}
\]
These are stronger than residue equalities. Clear the unit
denominators and divide successively by the monic variables
$q,U,V$. Each of the three expressions is consequently
$qF_q+UF_U+VF_V$, with finite bounded rational coefficient arrays.
The division does not invert $q,U,V$ in the tensor ring.

The combined full scalar correction to QI6 therefore consists of
bounded-array multiples of $\lambda q,\lambda U,\lambda V$
on $A_*^2,\Gamma^2,B_*^2$, an infinitesimal coefficient on
$E_*^2$, and the infinitesimal $\alpha$ correction on $D_*^2$.
The following exact identities allocate every apparently unbounded
testing term:
\[
\begin{aligned}
\lambda U A_*^2&=\lambda q(uE_*+bC_*)^2,\\
\lambda V A_*^2&=\lambda q(F_*+gD_*)^2,\\
\lambda U\Gamma^2&=\lambda q\,bC_*(uE_*+bC_*),\\
\lambda V\Gamma^2&=\lambda q(F_*+gD_*)
                                    (F_*+gD_*-vE_*),\\
\sqrt\lambda\,uB_*&=\sqrt{nS/z^2}\,(F_*+gD_*),\\
\sqrt\lambda\,vB_*&=\sqrt n\,C_*-vJ_3 .
\end{aligned}                                                   \tag{UL8}
\]
Every coefficient introduced on squaring the last two identities
is infinitesimal by UL2--UL3 and $S,u,v\in\mathrm I$.
Finally the formerly dangerous QI6 term has the exact replacement
\[
 \frac{wg}{c}\,vB_*C_*
       =wz\,C_*^2-\frac{wg}{c\sqrt\lambda}\,vC_*J_3,
 \qquad
 \frac{wg}{c\sqrt\lambda}
          =w\sqrt{\frac{z^2}{n}}=O(1).                          \tag{UL9}
\]
Both resulting coefficients are infinitesimal. This retains the
entire term for arbitrary independent source and testing rates.

\subsubsection*{Both unrestricted state-normal corrections}

The new complete first $N$ coefficient, including the remainder
of TU17, differs from the formal generic SC coefficient by
\[
\begin{split}
\delta\mathcal L_N={}&
 8\lambda c(1-e)\mathsf F_N B_*
 +8\lambda q\widetilde d\,ubvg(Sa_tH-D_uv),\\
\mathsf F_N={}&
 \frac{P_bD_v-(1+qz)P_m}{P}-zD_uD_v
                 =V-zUD_v+qF_{Nq}.
\end{split}                                                     \tag{UL10}
\]
The last equality is a finite monic $q$ division; its coquotient
has bounded arrays because the numerator at $q=0$ is exactly
$V-zUD_v$. After division by $\sqrt{H_0}$, its $q$ part is
infinitesimal on $B_*$ by UL3. The other parts use
\[
\begin{split}
\lambda cV B_*&=\lambda qg\,vC_*
                            -\sqrt{\lambda q}\,v^2J_3,\\
\lambda czU B_*&=\lambda qbg\,v(uE_*+bC_*).
\end{split}                                                     \tag{UL11}
\]
Their coefficients on $\widehat N$ are infinitesimal since
$g/\sqrt{H_0}$ and $\sqrt{\lambda q/H_0}$ are bounded.
The $\widetilde d$ term in UL10 uses
$uH=R(uE_*+bC_*)$ and $uv=bC_*$; it has the same
infinitesimal coefficient bounds. Thus this is a lift of the
complete first normal, with no state component put to zero.

For the second normal it is convenient to compare the coefficient
after multiplication by $\sqrt\sigma$. Differentiating the full
generic SC3 identity, including the restored $D$ positive square term, gives
\[
 \sqrt\sigma\,\delta\mathcal L_y
       =8\lambda q\{\mathsf F_HH+\mathsf F_vv\},                 \tag{UL12}
\]
where
\[
\begin{split}
\mathsf F_H={}&D_v\{\rho_0\widetilde a
                    -(1-\rho_0)\widetilde c\}
 -D_v\widetilde aU\rho_0/D_u
 +\widetilde\alpha zS^2V\rho_0/D_u
 -\widetilde dS^2a_t,\\
\mathsf F_v={}&\widetilde dSD_u-\widetilde\alpha zS.
\end{split}
\]
These arrays are bounded through both source openings. Exactly,
\[
 \mathsf F_H|_{q=U=V=0}
       =\frac{(1-e)S}{1-S(1-h)},\qquad
 \mathsf F_v|_{q=U=V=0}
       =-\frac{(1-e)S}{1-S(1-h)}.                              \tag{UL13}
\]
As $R=1$ there, finite monic division gives
$R\mathsf F_H+\mathsf F_v=qF_q+UF_U+VF_V$ with bounded
arrays. Since $H=R(E_*+v)$, the whole UL12 coefficient is
therefore an infinitesimal combination of $E_*,F_*,A_*,C_*,D_*$:
use only
\[
 qv=cA_*-qE_*,\qquad Uv=buC_*,\qquad
 Vv=v(F_*+gD_*-vE_*).                                         \tag{UL14}
\]
The factor $1/\sqrt\sigma$ is bounded.
Changing the formal harmonic-normal normalization to the actual
one multiplies its previously bounded arrays by a pure-source
unit with residue one.

For completeness, the unchanged formal first $N$ coefficient is
bounded after the actual $\sqrt{H_0}$ normalization:
$h_2^0,h_b^0=O(q)$ and $h_b^0-h_2^0=O(qz)$, while
$g/\sqrt{H_0}$ and $c/\sqrt{H_0}$ are bounded. These are exactly
the two estimates used in the QI normal allocation.
All remaining hard coefficients from the original $C,A,D,\alpha$
lines have the same residues as before; the relative $C,A$
changes and the absolute $D,\alpha$ changes are infinitesimal.
Together with UL5 the complete hard block on
$\widehat N,N_y,J_3$ has a positive definite real residue.

\subsubsection*{A finite contact lift for the additional row}

It remains to allocate the full scalar correction on roots that
vanish at contact. Merely bounding a coefficient on $A_*,B_*,C_*,D_*$
would not suffice. Keep the complete strict toric root list SC16,
the QI rows $E_*,F_*,R_F,R_E$, and add
\[
 V_3=\sqrt{\lambda\Gamma}(\sqrt{U_*}-\sqrt{B_*}).
\]
Since $U_*B_*=\Gamma^2$, the exact identities are
\[
\begin{split}
 J_3^2&=4V_3^2+\lambda(\sqrt{U_*}-\sqrt{B_*})^4,\\
 \Gamma J_3&=\frac{V_3^2}{\sqrt\lambda}
                      -2\sqrt\lambda\,\Gamma(B_*-\Gamma).
\end{split}                                                     \tag{UL15}
\]
Reserve half the $J_3^2$ coefficient in UL5 for the first
identity. It supplies strict real pivots on both $J_3,V_3$,
with the quartic square retained nonnegative. The normal
coupling in UL5 remains an infinitesimal hard Gram correction.

Here are the additional first-jet coefficient checks needed for
the second identity, so that its factor $\sqrt\lambda$ causes
no hidden loss. The only new scalar terms linear in $J_3$ in
UL8--UL9 have coefficients, up to bounded arrays, of the forms
\[
 \epsilon_B=-2\sqrt n\,v,\qquad
 \epsilon_T=\frac{wg}{c\sqrt\lambda}\,v .
                                                               \tag{UL16}
\]
Both have bounded finite testing difference arrays and are
infinitesimal. At the actual pure-source contact $v_c=cg$,
\[
 \sqrt\lambda\,\epsilon_B(u_c,v_c)=-2n\in\mathrm I,\qquad
 \sqrt\lambda\,\epsilon_T(u_c,v_c)=wz=e\in\mathrm I.             \tag{UL17}
\]
Multiplication by any of the bounded arrays produced in UL7
preserves these conclusions. No other scalar term introduces
$J_3$ linearly: the other five identities in UL8, and the
$\lambda q$ scalar corrections, use only the old QI rows.

Write a term $\epsilon J_3A_i$ as
$\epsilon J_3(A_i-\Gamma)+\epsilon\Gamma J_3$.
The first summand is already an infinitesimal quadratic on the
strict list. Divide $\epsilon-\epsilon(u_c,v_c)$ finitely by
$u-u_c,v-v_c$ and apply the exact SC18 conductors to
$\Gamma(u-u_c)$ and $\Gamma(v-v_c)$. This allocates the
nonconstant part as infinitesimal quadratic corrections involving
$J_3$. For the constant part use UL15. Its $V_3^2$
coefficient is infinitesimal, and its remaining coefficient on
$\Gamma(B_*-\Gamma)$ is infinitesimal by UL17.

There are also new bare-$\Gamma$ $E_*,F_*$ terms when the
first five identities in UL8 are expanded and their $C_*,D_*$
factors are written as $(A_i-\Gamma)+\Gamma$. Put
$L=\lambda q$ and $M_1=nS/z^2$, both infinitesimal.
The additional $E_*$ coefficients are bounded-array multiples
of $Lbu$ and $Lgv$; the additional $F_*$ coefficients are
bounded-array multiples of $Lg$ and $M_1g$.
At $u_c=bc,v_c=cg$ their required contact ratios are
\[
 \frac{Lbu_c}{c}=LS,\qquad
 \frac{Lgv_c}{c}=Lz,\qquad
 \frac{Lg}{g}=L,\qquad
 \frac{M_1g}{g}=M_1 .
                                                               \tag{UL18}
\]
All are infinitesimal, with bounded finite difference arrays
before these contact substitutions. The exact QI7 identities
therefore allocate these additional terms by the same finite
division as the old $E_*,F_*$ terms. This check is necessary:
the two new normal rows cannot be ignored when applying the
contact-jet argument.

After these explicit additional steps, the remainder has exactly
the QI form: a bounded infinitesimal quadratic on the listed
roots, the now-verified bare-$\Gamma$ $E_*,F_*$ couplings, and
a toric expression
$\Gamma\sum_i\mu_i(A_i-\Gamma)+\Gamma^2\sigma$.
The old $E_*,F_*$ contact coefficient ratios remain the bounded
ones in QI7--QI8: their unchanged $q=0$ arrays depend on bounded
$w$ and do not require $w/(c/g)$ bounded.
Every retained positive square term, including UL5 and UL15, vanishes at the
actual contact. The complete actual scalar consequently has zero
contact value and all three zero first jets. Apply the two
explicit finite divisions SC17--SC20. Their three-jet balance
makes the constant toric first-jet coefficients equal; SC19 and
the variance roots then allocate this common term. This puts
the entire scalar on a bounded Gram with strictly positive
real residue.

The full normal lifts UL10--UL14 have exact contact zero as well.
Their $J_3,E_*,F_*$ terms already vanish there. The sum of their
remaining four monomial coefficients is zero at contact, so one
finite SC18 division allocates them on the same strict list.
Every coefficient function in these steps depends only on the
testing $u,v$. No bound on the state ratio $Q_b/Z$ is used, and
the third contact jet is used only for its exact toric balance.

The strict bounded-Gram lemma now proves tensor SOS for the whole
original parent on the small-angle positive negative-core chart.
All its radicals are radicals of separated positive monomials;
every inverse is pure source or pure testing.

\subsubsection*{Projective completion and scope}

The state, sign, zero-axis and real/imaginary-state arguments are
the same exact original-core arguments at the end of TU.
In particular $Z=0$ is an original positive restriction, and the
negative-core $Z=1$ chart forces $W+uY>0$, permitting the
separated radicals above. The theorem R2O covers every testing
angle chart outside $u,v$ both infinitesimal without any
condition on $\lambda$ or the native $a$. This proves UL1 for
all tests and states.

Together with TU this covers native $a$ unbounded: the nonnegative
$\lambda=n/(tB)$ is bounded, including zero, or unbounded. The
bounded-$a$ theorem supplies the complementary source domain.

\subsection{Complementary testing angles}
\label{component:p8f8780350b13}

\label{p8f8780350b13:sec:r2-rederived-outer-join}

This proof treats every independent testing pair outside the positive
double-infinitesimal angle chart of the original R2 source.  It uses an
explicit full-parent Bell reflection and the two-current and six-root
identities that were freshly proved for R1.  Only those universal algebraic
identities and their real matrix arguments are reused.  No assertion that
the reflected source is an original R1 source is made.

\subsubsection*{Actual source and exact reflected parent}

The source is the original weak-C2 tuple, with $d=1$, $\lambda=\beta=0$,
the six coefficients in (RER1), all original strict gaps, $n\geq0$,
$a,B,t_s>0$, and $0<r,h<1$.  R2 means
\[
t_s\text{ unbounded},\qquad B, z=aB, y=t_sB\text{ infinitesimal}.
\tag{R2O1}
\]
Put $S=1/t_s$, $l=1-h$, $e=1-r$, $R_s=1+hz$,
$g_s=R_s-ae>0$, and $k=(1+ar)/(1+a)$.  The rederived exact source
lemma (RER1)--(RER7) gives
\[
k\text{ a positive bounded unit},\qquad k<1,\qquad
SC l\sim SAh\sim D/S\sim k/2,
\qquad \Delta/(1+a)\sim1.
\tag{R2O2}
\]
It also gives $ay$ infinitesimal, $Q\sim1$, and the exact expressions
\[
n=\frac{eQ-BK_s}{2r\Delta},\qquad
m=\frac{a^2e(1+y)c s_c}{2g_s\Delta},\qquad
BK_s=(1+a)y(1+hS)(1+z),
\tag{R2O3}
\]
where $c=1+zr\sim1$ and $s_c=1+B(l+ar)\sim1$.
There is no positive lower bound assumed on $r,h,l$, or $g_s$.

For arbitrary finite testing angles define the unnormalized testing
state transformation
\[
\begin{split}
G\psi=(&uvW-uX-vY+Z,\ -vW+X-uvY+uZ,\\
       &-uW-uvX+Y+vZ,\ W+vX+uY+uvZ).
\end{split}
\tag{R2O4}
\]
It has $G^TG=D_uD_v I$.  The Bell transformation on observables is
\[
(A'_0,A'_1,B'_0,B'_1)=(-B_1,-B_0,-A_1,-A_0).
\]
After the displayed real local basis gauges it sends the angles to
$(u',v')=(v,u)$ and the source coefficients to
\[
(\lambda',n',m',C',A',D',\alpha',\beta')
 =(\alpha,m,n,C,D,A,\lambda,\beta).
\tag{R2O5}
\]
Direct substitution in the whole parent gives
\[
F_{c'}(v,u;G\psi)=D_uD_v F_c(u,v;\psi).
\tag{R2O6}
\]
  The transformed
$\lambda'=\alpha$ event is kept as its whole positive square.
Because $D_uD_v>0$ is pure testing, a certificate on the left restores
the right by dividing each root by $\sqrt{D_uD_v}$ within $K_t$.
No source or testing contact is identified in this operation.

In the next two subsections write the transformed testing angles again
as $u,v$, and denote the reflected coefficients by
\[
\mathsf C=C,\quad\mathsf D=A,\quad\mathsf A=D,
\quad\mathsf m=n,\quad\mathsf n=m.
\]
Thus the symbols on the right refer to the original R2 source.
Define the pure-source quantities
\[
g_* =\frac{\mathsf C\mathsf D}{\mathsf C+\mathsf D},\qquad
h_* =\frac{\mathsf D}{\mathsf C+\mathsf D},\qquad
\gamma=\mathsf A g_*.
\tag{R2O7}
\]
If $U_C=SC l$ and $U_A=SAh$, then exactly
\[
Sg_* =\frac{U_CU_A}{hU_C+lU_A}\sim k/2,
\qquad \gamma\sim k^2/4.
\tag{R2O8}
\]
The denominator is a positive bounded unit even when $h$ or $l$ is
infinitesimal.  In particular $1/\mathsf C$, $1/\mathsf D$,
$1/g_*$ are infinitesimal,
$(\mathsf C+\mathsf D)^{-1/2}=O(\sqrt S)$, and $0<h_*<1$.
These are the complete source bounds needed for the first two outer
current identities.

\subsubsection*{Universal outer current identities}

Set
\[
R_t=uW-Y,\quad C_t=v(uX-Z),\quad B_t=X+uZ,\quad
J=R_t+C_t,\quad N=R_t+h_*C_t,
\]
and let $K$ be the exact signed quarter-core
\[
K=\tfrac12J^2+\tfrac12D_u(vY-Z)^2
 +\tfrac12D_vB_t^2-\tfrac12D_uD_vZ^2.
\tag{R2O9}
\]
All unselected original events, including $\lambda'$, are retained
as whole nonnegative positive square terms.

If $1/v$ is bounded, put
\[
K_0=2\sqrt{g_*}C_t,\quad B_0=\sqrt{2D_v/D_u}\,B_t,
\quad c_u=D_u^{-1/2},\quad t_u=uD_u^{-1/2},\quad
\eta_v=\sqrt{D_v}/v.
\]
The entire selected-parent identity is
\begin{align*}
4\mathsf C D_vR_t^2+4\mathsf D J^2+4K
={}&4(\mathsf C+\mathsf D)N^2+4\mathsf C v^2R_t^2
 +2J^2+2D_u(vY-Z)^2\\
&+(K_0,B_0)
\begin{pmatrix}
1-c_u^2\eta_v^2/(2g_*)&t_u\eta_v/\sqrt{2g_*}\\
t_u\eta_v/\sqrt{2g_*}&1
\end{pmatrix}
\binom{K_0}{B_0}.
\tag{R2O10}
\end{align*}
The finite Gram is bounded and has real residue $I$, because $c_u,t_u$
and $\eta_v$ are bounded while $1/g_*$ is infinitesimal.  The actual
separated bounded-Gram lemma makes it tensor SOS.  The roots themselves
need not be bounded.  This gives every state and every other finite angle
when transformed Bob is noninfinitesimal, equivalently when the original
Alice angle is noninfinitesimal.

For positive infinitesimal transformed $v$, put
$q=g_*^{-1/2}$ and define
\[
K_1=2\sqrt{g_*}C_t+qY,\quad
B_1=B_t-uvY,\quad V_0=vY,\quad Y_1=q\sqrt{D_u}Y,
\]
\[
L_0=\tfrac12q c_u(K_1-c_uY_1)-t_uvB_1+c_uvV_0.
\]
Then $L_0=\sqrt{D_u}\,v(vY-Z)$ exactly, and
\begin{align*}
&4\mathsf C D_vR_t^2+4\mathsf D J^2
       +4\mathsf A D_uD_vY^2+4K\\
&\quad=4(\mathsf C+\mathsf D)N^2+4\mathsf C v^2R_t^2
 +2J^2+\mathcal G,\\
\mathcal G={}&K_1^2+2D_v(B_1^2+V_0^2)
 +(4\gamma D_v-c_u^2)Y_1^2\\
&+2qv(K_1-c_uY_1)V_0-2L_0^2.
\tag{R2O11}
\end{align*}
This is the complete four-root Gram, including all normal terms.  Its
residue in the ordered frame $(K_1,B_1,V_0,Y_1)$ is
\[
\operatorname{diag}(1,2,2,k_0^2-c_{u,0}^2).
\tag{R2O12}
\]
Subscript $0$ denotes real residue.  Hence it is strict whenever its
last entry is positive.  This includes every unbounded transformed
$u$ and every bounded positive-unit transformed $u$ when $k_0=1$.
All inverse and square-root factors in these identities belong to one
field at a time.

\subsubsection*{The remaining unit angle and both coefficient openings}

It remains to treat transformed $v$ infinitesimal, transformed $u$ a
bounded positive unit, and $k_0<1$.  The source identity
\[
1-k=\frac{ae}{1+a},\qquad 0<ae<R_s\sim1
\tag{R2O13}
\]
shows that $a,e$ are positive bounded units: if $a$ were infinitesimal
or unbounded, or if $e$ were infinitesimal while $a$ is bounded, then
$1-k$ would be infinitesimal.  Thus in the current chart $a$ is a
positive unit, $e$ is a positive unit, and both $1+a$ and $\Delta$
are units.

Use the opening-safe source coordinates
\[
j=R_s/a,\qquad p=e,\qquad j-p=g_s/a>0.
\tag{R2O14}
\]
They satisfy $0<p<\min(1,j)$, $j,p$ are positive bounded units, and
the residue of $(j+1-p)/(j+1)$ equals $k_0$.  The exact relative
critical coefficient identities are
\[
\frac{\mathsf m}{jp/[2(1-p)(j+1)]}
 =\left(Q-\frac{BK_s}{e}\right)\frac{R_s+a}{R_s\Delta},
\qquad
\frac{\mathsf n}{p/[2(j-p)(j+1)]}
 =(1+y)c s_c\frac{R_s+a}{\Delta}.
\tag{R2O15}
\]
Their right sides have residue one.  Indeed $BK_s/e$ is infinitesimal
here by (R2O3), and all other normalized factors have the limits stated
in (RER).  No division by $r$, $g_s$, $n$, or an infinitesimal gap
was used to assert a bounded error.  The displayed quotients cancel
those openings exactly before the limit is taken.  In particular the
critical chart itself forces original $n>0$, even though the theorem
as a whole retains $n=0$.

Choose positive real caps $m_c\leq\mathsf m$ and
$n_c\leq\mathsf n$, with the choices specified below.  Keep both
nonnegative excesses as the original whole positive square terms.  Define six rows
\[
\widehat N=\sqrt{\mathsf C+\mathsf D}\,N,\quad
P=C_t/\sqrt S,\quad Y_0=\sqrt S\,Y,
\quad X,\quad Z,\quad V=vY.
\tag{R2O16}
\]
They can be unbounded or linearly dependent.  If
$\delta=(\mathsf C+\mathsf D)^{-1/2}$, then
$R_t=\delta\widehat N-h_*\sqrt S P$.  The selected entire
quarter-parent is exactly
\begin{align*}
\mathcal P_{\rm sel}={}&\widehat N^2+\mathsf C v^2R_t^2
 +\tfrac12J^2+(Sg_*)P^2
 +(\mathsf A/S)D_uD_vY_0^2\\
&+m_c\{(v/u)R_t+(D_u/u)V-X-uZ\}^2
 +n_cD_uD_vX^2\\
&+\tfrac12D_u(V-Z)^2+\tfrac12D_v(X+uZ)^2
 -\tfrac12D_uD_vZ^2.
\tag{R2O17}
\end{align*}
The full-state row relation is
\[
PY_0=V(uX-Z).
\tag{R2O18}
\]
Add $2b(PY_0-V(uX-Z))=0$ for a real $b$.  Keep the two displayed
positive squares $\mathsf C v^2R_t^2$ and $J^2/2$ whole.  The
remaining Gram is bounded: $u,1/u$ are bounded, $Sg_*$ and
$\mathsf A/S$ are units, and the $R_t$ contributions in the
$m_c$ square have coefficients $\delta v/u$ and
$h_*\sqrt S v/u$, both infinitesimal.  Its residue is the direct sum
of a unit pivot, the matrix
\[
\begin{pmatrix}k_0/2&b\\b&k_0(1+u_0^2)/2\end{pmatrix},
\tag{R2O19}
\]
and the following real quadratic form on $(X,Z,V)$:
\[
Q_3=m_c\{(D_u/u)V-X-uZ\}^2+n_cD_uX^2
 +\tfrac12D_uV^2-D_uVZ+\tfrac12(X+uZ)^2
 -2bV(uX-Z).
\tag{R2O20}
\]
In this formula $u=u_0>0$ is real.  We now give all finite-cap choices;
no opening of either source coefficient is suppressed.

If $p_0<1$ and $j_0>p_0$, use initially
\[
m_0=\frac{j_0p_0}{2(1-p_0)(j_0+1)},\quad
n_0=\frac{p_0}{2(j_0-p_0)(j_0+1)},\quad b=k_0/2.
\]
The three leading principal minors of $Q_3$ are
\[
m_0+n_0D_u+\tfrac12,\quad
\frac{p_0u_0^2D_u(j_0+1-p_0)}
 {4(j_0-p_0)(j_0+1)^2(1-p_0)},\quad
\frac{p_0u_0^2D_u(j_0+1-p_0)^2}
 {8(j_0-p_0)(j_0+1)^3},
\tag{R2O21}
\]
all strictly positive.  Matrix (R2O19) has determinant
$k_0^2u_0^2/4>0$.  Choose real caps slightly below $m_0,n_0$;
openness of real positive definiteness and (R2O15) ensure they remain
strict and are below the actual coefficients.

If $j_0=p_0<1$, the coefficient $\mathsf n$ is unbounded and
$m_0=j_0^2/[2(1-j_0)(j_0+1)]$ is finite positive.  Take
$b=1/[2(j_0+1)]$.  On the kernel $X=0$ of the $\mathsf n$ positive square term,
the first diagonal and determinant of $Q_3$ on $(Z,V)$ are
\[
u_0^2(m_0+1/2)>0,\qquad
u_0^2/[4(j_0+1)^2]>0.
\tag{R2O22}
\]
Choose $m_c<m_0$ sufficiently close and then a sufficiently large
real $n_c$.  A real block Schur complement makes the full matrix strict.

If $p_0=1<j_0$, the coefficient $\mathsf m$ is unbounded and
$n_0=1/[2(j_0^2-1)]$ is finite positive.  Choose a real $b$ with
\[
\frac{j_0}{2(j_0+1)}<b<
\min\left\{\frac{k_0\sqrt{D_u}}2,
                 \frac{j_0}{2(j_0-1)}\right\}.
\]
The interval is nonempty because $u_0>0$.  On the kernel
$(D_u/u)V-X-uZ=0$ of the $\mathsf m$ positive square term, the first diagonal
and determinant in $(X,V)$ are
\[
\frac{D_u}{2(j_0^2-1)}>0,\qquad
-\frac{D_u^2[2b(j_0-1)-j_0][2b(j_0+1)-j_0]}
       {4u_0^2(j_0^2-1)}>0.
\tag{R2O23}
\]
Choose $n_c<n_0$ sufficiently close, then a sufficiently large real
$m_c$.  Matrix (R2O19) and the full $Q_3$ are strict.

If $p_0=j_0=1$, both coefficients are unbounded and $k_0=1/2$.
Choose $1/4<b<\sqrt{D_u}/4$.  On the common event kernel
$X=0$, $Z=D_uV/u_0^2$, the remaining form is
\[
\frac{D_u(4b-1)}{2u_0^2}V^2>0.
\tag{R2O24}
\]
Sufficiently large real caps on both events make the full form strict.
The finite-cap step in each case is the elementary real fact that a
quadratic form positive definite on the kernel of a nonnegative event
Gram becomes positive definite after a sufficiently large multiple of
that event Gram is added; a block Schur complement proves it.

These four cases exhaust $0<p_0\leq\min(1,j_0)$ with positive $j_0,p_0$.
The separated bounded-Gram lemma now proves the entire selected parent
(R2O17), with every state normal retained.  Restore the original excesses
and all unused positive square terms.  The real caps are used only on this testing
chart: no globally physical capped source is asserted, so this argument
does not conflict with the finite-contact cap exclusion (CAP15).

\subsubsection*{Exact outer and projective join}

Return to the original testing angles.  The full-parent sign gauges
(SJ3) permit $u,v\geq0$ without changing source coefficients.  Literal
zeros and projective infinities are the exact signed-core squares
(SJ10), plus the untouched event squares.  For finite positive angles:
\begin{itemize}
\item if the original $u$ is noninfinitesimal, use (R2O10);
\item if the original $u$ is infinitesimal and $v$ is unbounded,
      use (R2O11)--(R2O12);
\item if the original $u$ is infinitesimal and $v$ is a positive
      bounded unit, use (R2O11) when $k_0=1$, and the complete
      six-root argument (R2O13)--(R2O24) when $k_0<1$.
\end{itemize}
These are exactly all finite positive angle pairs outside the case in
which both are infinitesimal.  Every argument is a whole-parent state
Gram, so no additional state chart or state-ratio restriction is present.
Equation (R2O6) restores the original frame and positive pure-testing
factor.  The zero state and the real/imaginary split for complex states
are immediate.  No source-to-testing rate comparison was made.

\section{Two small actual angles: the R3 and R4 sectors}
\subsection{All ratios at infinitesimal $e$}
\label{component:p1de8b5ffab69}

\label{p1de8b5ffab69:sec:new-r34-all-tests}

\paragraph{Statement and exact domain.}
For the present source partition it suffices to impose additionally $a$ unbounded; the identities below remain valid on this smaller domain.

Use the original common-II stationary source, its six native weights
$n,m,C,A,D,\alpha$, all its strict source gaps, $d=1$,
$\lambda=\beta=0$, and $n\geq0$. Suppose $B_s,a_sB_s,t_sB_s$
are positive infinitesimals, $t_s,h$ are positive bounded units,
and $e=1-r$ is positive infinitesimal. In addition impose either
\[
 \begin{array}{ll}
 \mathrm{R3}:&1-h\text{ is infinitesimal},\\
 \mathrm{R4}:&1-h\text{ is a positive unit and }
       1+a_sB_sr-(1-h)/t_s\text{ is infinitesimal}.
 \end{array}                                                    \tag{EN1}
\]
There is no condition on the size of $a_s$, $a_s^2B_s$,
$a_s e$, or their ratios. The conclusion is finite separated tensor
SOS for the entire original parent at every independent projective
testing pair and every state, including signs, axes, and literal $n=0$.
The bounded-$a$ parts of R3 and R4 are covered by Proposition~\ref{prop:c2-bounded-a}; the present calculation supplies their unbounded-$a$ parts.

The source and testing fields are independent real closed extensions
$K_c,K_t\supset\mathbb R$. A bounded coefficient means an element
of $B_c\otimes B_t$; its residue is the homomorphism to $\mathbb R$.
The strict bounded-Gram lemma is used only for matrices with positive
definite real residue. Row entries may be arbitrarily large.
Every inverse and radical below is pure source, pure testing, or a
product of such factors.

\subsubsection*{One exact compact source for both ends}

Use the exact compact variables of the exact actual-source
dictionary RSR2--RSR6:
\[
\begin{aligned}
c_s&=1+a_sB_sr,\quad R_s=1+ha_sB_s,\quad
j=a_sc_s/R_s,\quad B=B_sR_s/c_s,\quad T=t_sc_s/R_s,\\
p&=a_s e/R_s,\quad H=(1-h)/(t_sc_s),\qquad
x=jB=a_sB_s,\quad y=TB=t_sB_s,\quad q_z=jT .
\end{aligned}                                                    \tag{EN2}
\]
Thus $0<p<\min(1,j)$, $0<H<1$, and
\[
B,x,y\in\mathrm I_c,\quad T\in\mathrm U_c,\quad
1-HT\in\mathrm U_c,\quad p/j\in\mathrm I_c.              \tag{EN3}
\]
The last assertion is exact source entry:
$e=p(1+x)/(j+px)\sim p/j$.
In R3, $H$ is infinitesimal. In R4, $1-H$ is infinitesimal
and $0<\overline T<1$. In particular
$H(1-H)\in\mathrm I_c$ in both rows.
The native inverse is
\[
r=\frac{j-p}{j+px},\quad h=\frac{1-HT}{1+HTx},\quad
a_s=\frac{j+px}{1+HTx},\quad
B_s=\frac{B(1+HTx)}{1+pB},\quad
t_s=\frac{T(1+pB)}{1+HTx}.                              \tag{EN4}
\]
The exact six weights are RSR6 at this substitution.
These rational identities, including the contact first jets, hold at the displayed substitution.

Put $\sigma=1+(1-H)T$, $K=j+1-p$, $E=1+B(HT+p-1)$, and
use the explicit polynomials $\mathcal P,\mathcal L,\mathcal N,
\mathcal F,\mathcal K$ in RSR5, RSR15, RSR18. Then
\[
\frac{\mathcal P}{j+1}\sim1,\qquad
\frac{\mathcal L}{j+1}\sim\sigma,\qquad
\frac{\mathcal F}{j+1}\sim\sigma.                       \tag{EN5}
\]
Here $T,\sigma,(1-HT)^{-1}$ are bounded positive units.
To prove these uniform statements, divide the printed finite
polynomials by $j+1$. The corrections in $\mathcal P$ contain
$B+x$ or $TB$; those in $\mathcal L$ after
$(j+1)\sigma$ contain $p/(j+1)$ or $B+x$.
In $\mathcal F$ the only additional quotient is
$p(1-p)/(j(j+1))$, infinitesimal by (EN3).
Every other term contains $B$ or $Bj/(j+1)$.
Thus no restriction on $j$ or $jT$ is needed.

The exact physicality numerator has the grouping
\[
\mathcal N=p\{1+B(3j+1-p)+BT(j+1)\}
-Bj(j+1)\{1+(1-H)T\}+\mathcal R,\quad
|\mathcal R|\le C_0BT\,Bj(j+1).                         \tag{EN6}
\]
The bound follows termwise from RSR5 using $p<\min(1,j)$ and
bounded $T$. Since $n=\mathcal N/[2(j-p)\mathcal P]\ge0$,
\[
\begin{aligned}
&\eta_*=\frac{Bj(j+1)}p=O(1),\qquad
0\le n=O\!\left(\frac{p}{j(j+1)}\right)\in\mathrm I_c,\\
&\eta_*\in\mathrm I_c\Longrightarrow
n\sim\frac{p}{2(j-p)(j+1)}>0 .
\end{aligned}                                                    \tag{EN7}
\]
This includes $\mathcal N=0$ literally.

Write $\mu=m(1-p)$, $h_0=\mu+1/2$. The full source formulas give
\[
\begin{aligned}
&\mu=O(p),\quad h_0\asymp1,\quad
\alpha/[B^2(j+1)]\sim\sigma/2,\quad \alpha/B\in\mathrm I_c,\\
&A\sim\frac{T}{2(1-HT)},\qquad
CH\sim\frac{\sigma}{2T},\qquad
D(1-H)\sim\frac1{2T}.
\end{aligned}                                                    \tag{EN8}
\]
Consequently $A$ and the finite one of $C,D$ have positive real
lower bounds. The other coefficient is unbounded. The exact
normalized source limits are
\[
\begin{array}{c|c|c|c}
 & C&A&D\\ \hline
R3&+\infty&T/2&1/(2T)\\
R4&1/(2T)&T/[2(1-T)]&+\infty .
\end{array}                                                    \tag{EN9}
\]
Only real residues of the bounded entries are meant in this table.
The coefficient $m$ may be infinitesimal, a unit, or unbounded.

\subsubsection*{The scalar allocation in the complete two-normal frame}

Use the entire parent RFN4 and the exact full-state change RFN5--RFN9,
with $w=(W+uY)/(1+u^2)=1$ and arbitrary $z=(Z-uX)/(1+u^2)>0$.
The identities are rational identities of (EN4), so they retain
both complete normals and every state term in the present source.
The tests in this subsection satisfy $u,v>0$ infinitesimal.
For clarity, the inverse is
\[
X=N_0+pv-uz,\quad Z=z+uN_0+puv,\quad
W=1+u(M_0+Hvz),\quad Y=u-M_0-Hvz.                       \tag{EN10}
\]
Retain the actual positive square term $mpD_u^2N_0^2$, and divide the
remaining entire parent by the pure testing factor $D_u=1+u^2$.
Its exact expansion is RFN7--RFN9. In particular the hard entries are
\[
\begin{aligned}
&H_N=h_0D_u+\tfrac12v^2+\alpha u^2+nD_v,\\
&H_M=D_u(CD_v+D+\tfrac12)+AD_v+(\alpha+\tfrac12)v^2,\qquad
H_{NM}=(\alpha+\tfrac12)uv .
\end{aligned}                                                    \tag{EN11}
\]

Keep the separated rows
$U=u/\sqrt{x}$, $V=v/\sqrt{y}$, $Z_*=z/\sqrt{q_z}$,
$Q_*=U-VZ_*$, $B_*=V-UZ_*$, $R_*=Z_*-UV$, and
\[
\Psi=U^2+V^2+Z_*^2+U^2V^2Z_*^2-4UVZ_*
=(U-V)^2+Z_*^2(1-UV)^2+2UV(1-Z_*)^2.                 \tag{EN12}
\]
The last radical is separated, for example
$\sqrt{UV}=(xy)^{-1/4}\sqrt{uv}$.
The exact scalar identity is RFN12 with the same actual circuit
masses. From (EN5), the closed formulas RSR16--RSR20 now give
\[
\begin{aligned}
&f_c\sim Tx^2\sigma/2,\quad
q_c\sim\frac{Tx}{2(1-HT)},\quad
r_c\sim TjB^2(1-p)\sigma/2,\\
&g_c/f_c=Hy\,\mathcal L/\mathcal F\in\mathrm I_c,\quad
k_c<0,\quad |k_c|/(q_cy)=O(p+B).
\end{aligned}                                                    \tag{EN13}
\]
In particular all three displayed positive masses are positive.
The sign of $k_c$ follows from the same exact positive-surplus
identity RSR19: its error bound uses $B(j+1)\in\mathrm I_c$,
not $T$ infinitesimal.

Before spending any current mass, perform the exact source-only
completion
\[
\begin{aligned}
h_0N_0^2+n(N_0+jb)^2
&=(n+h_0)(N_0+c_bb)^2+M_eb^2,\\
b&=v-uz/j,\quad c_b=\frac{nj}{n+h_0},\quad
M_e=\frac{nh_0j^2}{n+h_0}.
\end{aligned}                                                    \tag{EN14}
\]
Keep $nv^2(N_0+pb)^2$ and the positive extra hard terms RFN15
as whole squares. This reserves exactly the correlated $n$ positive square terms
of the complete parent. If $n=0$, set $c_b=M_e=0$ and omit that
row. For $n>0$, $c_b^2/M_e=n/[h_0(n+h_0)]$ is bounded.

The signed scalar allocation differs from R1 in one source split.
The $e_0$ cost is unchanged: its negative part is at most $g_c$,
and RFN18b provides it at cost
$H(\mathcal L/\mathcal F)(y+v^2)\in\mathrm I$.
Also $d_c>0$. Use $Z_*(V^2-1)=-VQ_*-R_*$ for $k_c$.
The $VQ_*$ cost is $O((p+B)v^2)$, infinitesimal.
For its opposite-row part the exhaustive cases are:
\[
\begin{array}{c|l}
p\in\mathrm I_c&
 |k_c|/r_c=
 T\mathcal K/[(1-p)(1+Hy)\mathcal L]=O(p+B)\in\mathrm I_c\\
p\in\mathrm U_c,\ \eta_*\in\mathrm U_c&
 j\text{ unbounded},\ j^2B\in\mathrm U_c,\ 
 |k_c|/f_c\in\mathrm I_c,\ |k_c|/(f_cx)=O(1)\\
p\in\mathrm U_c,\ \eta_*\in\mathrm I_c&
 j\text{ unbounded},\ M_e\in\mathrm U_c,\ 
 |k_c|/f_c\in\mathrm I_c,\\
& |k_c|/(q_cy)=O(1),\ |k_c|/(M_eyx)=O(1).
\end{array}                                                    \tag{EN15}
\]
In the second case $j$ is unbounded because $p/j$ is infinitesimal;
$j^2B\asymp p\eta_*$. In the third,
$nj^2\sim p/2$ and $n$ is infinitesimal, so $M_e\sim p/2$.
The closed formulas show
$|k_c|/f_c=O((p+B)/j)$ in the last two cases and
$|k_c|/(M_eyx)=O((p+B)(1+1/n)/j^2)$ in the third.
These bounds are uniform through $p\to1$.

Case two uses RFN18a:
$R_*^2\preceq_{\rm sos}(1+U^2)\Psi$.
Case three uses RFN18c:
$R_*^2\preceq_{\rm sos}2(\Psi+V^2Q_*^2+U^2B_*^2)$.
The non-infinitesimal source ratios in (EN15) acquire the testing
factors $u^2$ or $v^2$. Thus all costs are infinitesimal.
Literal $n=0$ lies in the first or second case; it is not discarded.
After this finite allocation, fixed positive fractions of
\[
q_cQ_*^2,\quad r_cR_*^2,\quad f_c\Psi,\quad M_eb^2       \tag{EN16}
\]
remain, with the last row absent exactly at $n=0$.

\subsubsection*{Both complete normal arrays, with the new endpoint weights}

The exact derivatives RFN21--RFN22, and their exact source cancellations
$S_N=S_0=S_1=S_2=0$, hold unchanged. Let
$q=vz-Tu$, $R=z-uv/B$, $R_u=v(u^2-x)$. In particular
\[
\begin{aligned}
L_N^{\rm rem}
&=-j\alpha b+\tfrac12pv^2b+K_RR_u+(K_c/j)uvq,\\
L_M&=C_qq+\mu uvb+v(\alpha-BTCHu^2)R,\\
K_c&=p/2+j\{H(\alpha+1/2)-1/2\}+n(p-j),\\
K_R&=-\mu-(\alpha+1/2)(1-p)+(T/j)K_c,\\
K_d&=CH-(D+1/2)(1-H)=-\mu/j-BTCH,\\
C_q&=K_d+AH+1/2+
v^2\{CH(1+u^2)+AH+(\alpha+1/2)H\}.
\end{aligned}                                                    \tag{EN17}
\]
Every displayed coefficient other than a retained row is bounded:
$p/j$ and $n$ are infinitesimal, $\mu/j=O(p/(j+1))$ is bounded,
and $A,CH,D(1-H),T$ are bounded by (EN8).
In particular $K_R$, $K_c/j$, and $C_q$ are bounded.

The complete bounds for these costs are as follows; they use the actual
$T\in\mathrm I_c$ hypothesis of RFN.
The exact radial Gram RFN19--RFN20 purchases $R_u^2$ at infinitesimal
cost from (EN16). Its three cases are still $1-p$ a unit;
$p\sim1,\eta_*$ a unit; and $p\sim1,\eta_*$ infinitesimal.
Their costs are respectively
\[
\begin{array}{c|l}
1-p\in\mathrm U_c&O(B),\ O(Bu^2),\ O(u^2/(1-p))\\
p\sim1,\eta_*\in\mathrm U_c&
 O(B+u^2/j+u^4/(j^2B)),\ O(Bu^2)\\
p\sim1,\eta_*\in\mathrm I_c&
 O(u^2/j),\ O(u^4/j),\
 [8u^4+2(u^2-x)^2]/M_e .
\end{array}                                                    \tag{EN18}
\]
They require only $T$ bounded above and below, because $T$ cancels
between $y=TB$ and the masses (EN13).
The small-current budgets NSR3--NSR5 likewise apply unchanged:
both $pv^2b$ and $\sqrt T\,\mu uvb$ are free.
For $\eta_*$ infinitesimal their coefficients are
$O(v^4)$ and $O(Tu^2v^2)$, using $p^2/M_e=O(1)$.
For $\eta_*$ a unit the coefficients are
$O(v^2(y+v^2))$ and $O(Tu^2(y+v^2))$, using
$p^2y/f_c=O(1)$. Each is infinitesimal.

The row $q/\sqrt T$ is strict by (EN13).
Its $N$ coupling has coefficient $(K_c/j)uv\sqrt T$ infinitesimal.
The identity $j\alpha b=-(\alpha/B)R_u-\alpha uR$ handles
the remaining $N$ term. For $1-p$ a unit the latter square costs
$O(B^2(j+1)^2u^2/(1-p))$ against $r_cR_*^2$.
For $p\sim1$, RFN18a bounds the entire $j\alpha b$ square by
$O(B^3(j+1)^2+B^2(j+1)^2u^2/j)$ of the circuit.
These are infinitesimal. A sufficiently large fixed real radial
pivot pays the bounded $K_R$ coupling, including the bounded shift
$c_b/\sqrt{M_e}$ from (EN14).

Normalize the complete second normal by
$M_0=\sqrt{TH(1-H)}\,\widehat M$. The full hard diagonal (EN11)
then has residue
\[
\overline{TH(1-H)H_M}=
\begin{cases}
(1+\overline T)/2,&R3,\\
1/2,&R4 .
\end{cases}                                                    \tag{EN19}
\]
Every term of the complete diagonal is bounded: the only divergent
factors occur in $CH$ and $D(1-H)$.
The normalized $q$ coupling is
$T\sqrt{H(1-H)}C_q$, now infinitesimal because
$H(1-H)$ is infinitesimal. This is the replacement for the
$T$-infinitesimal argument in R1.
The normalized current coupling is a bounded multiple of the
already purchased $\sqrt T\,\mu uvb$.

For the entire remaining opposite row put $J=B(j+1)+u^2$ and
$a_R=Tq_zv^2(\alpha-BTCHu^2)^2$.
Since $\alpha-BTCHu^2=O(BJ)$, its costs are
\[
\begin{aligned}
a_R/f_c&=O((T/j)v^2J^2),&
a_R/(q_cy)&=O(v^2J^2),\\
a_R/r_c&=O(Tv^2J^2/(1-p)),&
a_R/(M_eyx)&=O(T(1+1/n)v^2J^2/j^2).
\end{aligned}                                                    \tag{EN20}
\]
The direct $r_c$ cost applies for $1-p$ a unit.
For $p\sim1,\eta_*$ a unit, RFN18a applies using $j^2B$ a unit.
For $p\sim1,\eta_*$ infinitesimal, RFN18c applies using
$(1+1/n)/j^2=O(1)$. Thus the entire square, including every
higher testing degree, is free. No state ratio occurs in a cost.

Finally the complete mixed hard term is
$2(\alpha+1/2)uv\sqrt{TH(1-H)}
(\widehat N/\sqrt{n+h_0}-c_bb)\widehat M$.
Its bounded coefficients have zero residue.
The strict bounded-Gram argument therefore applies exactly as in
the final paragraph of RFN: choose a large fixed radial pivot, provide
the finitely many free rows, and retain positive fractions of
(EN16). The two hard pivots are $1$ and the positive number
(EN19). Restore all whole positive square terms in (EN14), RFN15, and
$mpD_u^2N_0^2$, and the pure testing factor $D_u$.
This proves the full local parent for all infinitesimal $u,v>0$
and arbitrary $z>0$, with both unrestricted state normals.
The exact SJ state/sign reduction supplies every state and the
literal axes in this local chart.

\subsubsection*{Every other testing angle, with weak \(m\) retained}

This outer join uses only (EN9), positivity of the six native
weights, and exact full-parent identities. Set
$\widetilde m=m/(1+m)$, a bounded positive source coefficient below
$m$, and retain the difference as the whole original $m$ square.
Keep the entire $n$ and $\alpha$ events throughout.
For the single unbounded coefficient $C$ or $D$, retain its excess
over a sufficiently large positive real cap. Other selected
coefficients have the actual bounded positive residues (EN9).

\paragraph{Finite Bob with nonzero residue and bounded Alice.}
Let $v$ have nonzero finite residue, $u$ be bounded, and use the
five full-state rows
\[
P=uW,\quad Y,\quad X,\quad Z,\quad S=\sqrt{\widetilde m}\,W.
                                                               \tag{EN21}
\]
The selected $m$ event is the entire square
$[vS-\sqrt{\widetilde m}(X-uvY+uZ)]^2$.
All other terms use $P,Y,X,Z$ with bounded coefficients.
In R3 the strong row is $P-Y$; on its kernel the selected remaining
form on $(Y,X,Z)$ is
\[
\frac T2D_uD_vY^2+\frac{1+T}{2T}v^2(uX-Z)^2+K_0 .
\]
In R4 the strong row is $P+uvX-Y-vZ$; on its kernel it is
\[
\frac{D_vv^2}{2T}(Z-uX)^2+
\frac{T}{2(1-T)}D_uD_vY^2+K_0 ,
\quad
K_0=\frac{D_u(vY-Z)^2+D_v(X+uZ)^2-D_uD_vZ^2}{2}.
\]
The leading principal minors in the order $(Y,X,Z)$ are
\[
\begin{array}{c|c|c}
 &R3&R4\\ \hline
1&
\dfrac{D_u(TD_v+v^2)}2&
\dfrac{D_u(T+v^2)}{2(1-T)}\\[4pt]
2&
\dfrac{D_u(TD_v+v^2)(Tu^2v^2+Tv^2+T+u^2v^2)}{4T}&
\dfrac{D_uD_v(T+v^2)(T+u^2v^2)}{4T(1-T)}\\[4pt]
3&
\dfrac{v^2(1+T)D_u^2D_v(Tu^2+v^2)}{8T}&
\dfrac{v^2D_u^2D_v^2(Tu^2+v^2)}{8T(1-T)} .
\end{array}                                                    \tag{EN22}
\]
All are strictly positive at the real residues, even $u=0$.
The $S$ direction has the positive coefficient $v^2$ in the
selected $m$ event. A sufficiently large real cap on the strong
row makes the entire five-row residue positive definite.
This is the real finite-cap Schur lemma: if a bounded real
quadratic form is positive on the kernel of a linear map, adding
a sufficiently large real multiple of the squared map makes it
positive definite. The unbounded actual coefficient permits
every such cap. This argument includes $m$ infinitesimal.

\paragraph{Infinitesimal Bob and Alice a finite positive unit.}
Use $R=uW-Y$, $B_t=X+uZ$, $V_Z=vZ$, and $Y$; the testing inverse
$1/u$ is bounded. The full selected parent is a bounded quadratic
on these four rows because
$v(W+uY)=(v/u)(R+D_uY)$ and
$J=R+uvB_t-D_uV_Z$.
Its residue splits off $(\widetilde m+1/2)B_t^2$ and the form
\[
CR^2+AD_uY^2+(D+1/2)(R-D_uV_Z)^2
-D_uYV_Z-\tfrac12D_uV_Z^2.                            \tag{EN23}
\]
In R3 restrict to the strong-row kernel $R=0$.
After removing a positive factor $D_u$, the matrix on $(Y,V_Z)$ is
$\left(\begin{smallmatrix}T/2&-1/2\\-1/2&
(1+T)D_u/(2T)-1/2\end{smallmatrix}\right)$;
its determinant is $(1+T)u^2/4>0$.
In R4 the strong kernel is $R=D_uV_Z$.
On $(R,Y)$ its first diagonal is $1/(2T)-1/(2D_u)>0$,
and its determinant is $u^2/[4(1-T)]>0$.
A large real cap again makes (EN23) strict.

\paragraph{Infinite Bob and bounded Alice.}
Put $w=1/v$, divide the homogeneous parent by $v^2$, and use
the five rows (EN21). Select any fixed positive real lower caps
$c_0,a_0,d_0$ below $C,A,D$, and the actual $\widetilde m$ event.
The full quadratic is bounded; its residue is
\[
[S+\sqrt{\widetilde m}\,uY]^2+
c_0(P-Y)^2+(a_0+\tfrac12)D_uY^2+
d_0(uX-Z)^2+\tfrac12D_uX^2.                           \tag{EN24}
\]
This is positive definite on the five rows, for every bounded
residue of $u$. It includes $u=0$ and $m$ infinitesimal.

\paragraph{Infinite Alice.}
Let $s=1/u$ and keep the entire $m,n,\alpha$ events.
For finite $v$ of nonzero residue, after division by $u^2$
the selected residue on $(W,Y,X,Z)$ is
\[
c_0D_vW^2+a_0D_vY^2+
(d_0+\tfrac12)(W+vX)^2+\tfrac12(vY-Z)^2,
\]
which is positive definite. For $v$ infinitesimal use instead
$(N,W,Y,V)=(sX+Z,W,Y,vX)$.
The complete signed core after division by $u^2$ is
\[
\frac12\{W+(1+s^2)V-sY-svN\}^2
+\frac12(1+s^2)v^2Y^2
-(1+s^2)Y(vN-sV)-\frac12(1+s^2)(vN-sV)^2
+\frac12D_vN^2 .
\]
Adding the selected $C,A,D$ positive square terms gives strict residue
$c_0W^2+a_0Y^2+(d_0+1/2)(W+V)^2+N^2/2$.
Every full coefficient is bounded; no $s/v$ is used.
If also Bob is infinite, divide by $u^2v^2$ and use
$(W,X,Y,S_Z,T_Z)=(W,X,Y,sZ,wZ)$.
The complete selected residue is
\[
c_0W^2+a_0Y^2+d_0(X-S_Z)^2+\tfrac12X^2
+\tfrac12(Y-T_Z)^2,                                   \tag{EN25}
\]
again positive definite. The full quadratic follows directly from
$J/(uv)=wW+X-swY-S_Z$ and the exact signed core;
its coefficients are bounded polynomials in $s,w$.

These charts exhaust all nonnegative projective angles:
both infinitesimal are the complete local theorem, and every other
pair is one of (EN21)--(EN25) or (EN23).
All strict-matrix coefficients use only bounded pure source or
pure testing factors. Restore every selected excess, the untouched
$m,n,\alpha$ positive square terms, and the homogeneous testing squares.
The exact SJ sign gauges, literal projective endpoints and
real/imaginary state split finish the asserted all-tests theorem.

\section{Two small actual angles: the R6 sector}
\subsection{Positive residue of $a^2B$}
\label{component:pb60eba754641}

\label{pb60eba754641:sec:r6-b-unit}

\paragraph{Scope.}
All notation and all physical source constraints are the original weak-C2
ones. In particular $d=1$, $\lambda=\beta=0$, $n\geq0$, and the actual
finite stationary contact, every strict gap, the total-scale condition,
and the complement of earlier cases are retained. On R6,
$t,h,l=1-h,j,r$ are positive bounded units, $B,x=aB,y=tB,g$ are
positive infinitesimals, and here $a$ is unbounded. Define
\[
 s=a^{-1},\quad w_s=ae,\quad b=a^2B,\quad e=1-r,
 \qquad B=s^2b,\quad x=sb,\quad y=ts^2b.
 \tag{RB1}
\]
The existing exact source identity gives
\[
 2r\Delta_*n/s^2=w_sQ_*-bK_*,\qquad
 Q_*\sim\Delta_*\sim1,\quad K_*\sim h+t.
 \tag{RB2}
\]
The gap $g=1-w_s+hsb$ gives $w_s\sim1$. Hence $b$ is bounded and
\[
 \rho=n/e^2\sim\frac{1-b_0(h_0+t_0)}2,\qquad
 0\leq b_0(h_0+t_0)\leq1.                         \tag{RB3}
\]
Here and below a subscript zero means the real residue; a displayed
residue formula with bounded source letters has the same meaning.
No testing variable occurs in these source estimates.

\begin{theorem}
On the entire original R6 source intersection, assume in addition that
$b=a^2B$ is a positive bounded unit. Then the complete original parent
is a finite sum of squares in the separated source--testing tensor ring
at every independent projective testing pair and every real or complex
state, with either angle sign and every literal projective endpoint.
All positive $s,g$ opening rates are independent. Literal $n=0$ is
included. The theorem makes no assertion about $b$ infinitesimal.
\end{theorem}

\paragraph{Source bounds needed by the proof.}
Put $R=1+hx$, $K=1+x$ and $T=1+x(1-e)$. The original rational source
formulas, after the pure-source substitutions (RB1), have only unit
denominators after $mg$, $n/s^2$ and $\alpha/(s^3b^2)$ are formed.
Their joint residues are
\[
 mg\sim\tfrac12,\quad
 C\sim\frac{h+t}{2l},\quad A\sim\frac t{2h},\quad
 D\sim\frac1{2(t-l)},\quad
 \frac\alpha{s^3b^2}\sim\frac{h+t}2.               \tag{RB4}
\]
Thus $m$ is unbounded, $C,A,D$ are positive units, and all four have
positive real lower bounds. There is no $s/g$ or $g/s$ assumption.
These are rational cancellations in the actual physical source;
no formal negative-beta source is imported.

\paragraph{Two contact-zero positive square terms remove the two troublesome terms.}
Write the complete parent in the normalization $P=F/4$:
\[
\begin{split}
 P={}&nD_uD_vX^2+mM^2+C D_v(uW-Y)^2+A D_uD_vY^2\\
 &+(D+\tfrac12)J^2+(\alpha+\tfrac12)D_u(vY-Z)^2
       +\tfrac12D_v(X+uZ)^2-\tfrac12D_uD_vZ^2,\\
 M={}&vW-X+uvY-uZ,\qquad J=uW+uvX-Y-vZ,
 \quad D_u=1+u^2,\quad D_v=1+v^2.
\end{split}                                                     \tag{RB5}
\]
In the finite state chart $Z=1$, use the pure-testing coordinate
$\omega=(W+uY)/D_u$ and the exact normals
\[
\begin{gathered}
 p=\frac{w_s}{R}=1-\frac gR,\quad c=\frac{l}{tT},\quad
 q_A=D_uv\omega,\quad q_B=v(1-uX),\\
 \qquad N=X+u-pq_A,\quad M_B=uW-Y-cq_B.
\end{gathered}                                    \tag{RB6}
\]
Their full inverse is
\[
\begin{gathered}
 X=N+pD_uv\omega-u,\quad q_B=v(D_u-uN-puD_uv\omega),\\
 \quad W=\omega+\frac{u(M_B+cq_B)}{D_u},\quad
 Y=u\omega-\frac{M_B+cq_B}{D_u}.
\end{gathered}                                    \tag{RB7}
\]
The only testing denominator is the positive pure-testing unit $D_u$.
Set
\[
 \delta=\frac{R}{K\sqrt{bt}},\quad E=\frac{tK}{R},\quad
 k=\sqrt x,\quad K_c=\delta^2Ek^2=\frac{sR}{K},\quad
 \gamma=\frac r{K_c}=\frac{rK}{sR}.                \tag{RB8}
\]
All of $\delta,E,b$ are positive bounded units in this theorem.
The actual contact in these coordinates is
\[
 (u,\omega,v;N,M_B)=(k,\delta k,\delta k^2E;0,0).
 \tag{RB9}
\]
For direct verification, its original state coordinates with $Z=1$ are
\[
 u_c=\sqrt{sb},\quad v_c=s\sqrt{bt},\quad
 W_c=\sqrt{s/t},\quad X_c=-r\sqrt{sb},\quad
 Y_c=hs\sqrt{b/t},\quad
 \omega_c=\sqrt{s/t}\,R/K.
 \tag{RB9a}
\]
They give $u_cW_c-Y_c=ls\sqrt{b/t}=c\,v_c(1-u_cX_c)$
and $X_c+u_c=e\sqrt{sb}=pD_{u_c}v_c\omega_c$.
Also $\gamma v_c\omega_c=r\sqrt{sb}=-X_c$.
Thus both normals and the retained $n$ row vanish exactly, before
any residue calculation. The same contact substitution also gives
$\omega_c=\delta k$, $v_c=\delta k^2E$, and $K_c=sR/K$.
For any fixed positive real $K_m$, retain the two whole squares
\[
 (m-K_m)N^2,
 \qquad nD_uD_v(X+\gamma v\omega)^2.                \tag{RB10}
\]
The first coefficient is positive because $m$ is unbounded. The second
vanishes identically when $n=0$. Both rows vanish at (RB9), with every
first jet of their squares. No inverse of $n$ is used.

The complete remaining polynomial is exactly
\[
\begin{split}
 G={}&K_mN^2-2\frac{mg}{R}q_AN+\frac{mg^2}{R^2}q_A^2\\
 &-2n\gamma D_uD_vX(v\omega)-n\gamma^2D_uD_v(v\omega)^2\\
 &+CD_v(uW-Y)^2+AD_uD_vY^2+(D+\tfrac12)J^2\\
 &+(\alpha+\tfrac12)D_u(vY-1)^2
       +\tfrac12D_v(X+u)^2-\tfrac12D_uD_v.
\end{split}                                                     \tag{RB11}
\]
This is the entire parent minus precisely (RB10), including every
first normal and every hard quadratic term. The exact cancellation
$M=(g/R)q_A-N$ removes the whole $g^{-1}$ pole before an estimate.
Moreover $n\gamma=O(s)$ and $n\gamma^2$ is bounded. All coefficients
in (RB11), after (RB7), are therefore bounded. The only nonconstant
state-normal degree is one or two.

\paragraph{Raw support before the conductor.}
Define
\[
 A_0=\omega v,\quad B_0=\delta v,\quad C_0=\delta E u\omega,
 \quad D_0=\delta^2k^3E,\quad X_0=B_0-C_0.           \tag{RB12}
\]
At $N=M_B=0$ let $V_0=1-puv\omega$. Then
\[
 q_B=D_uvV_0,\quad uW=u\omega+u^2cvV_0,\quad
 Y=u\omega-cvV_0,\quad J=-(1-c)D_uvV_0,
 \quad X+u=pD_uA_0.                               \tag{RB13}
\]
Thus every term of the scalar of (RB11) lies, with bounded coefficients,
in $(A_0,B_0,C_0)^2+(D_0^2)$. The only constant is from $\alpha$,
and $\alpha/D_0^2$ is bounded with residue $b(h+t)/2$.
For the new negative $n$ contribution use exactly
\[
\begin{split}
 -2n\gamma D_uD_vX A_0-n\gamma^2D_uD_vA_0^2
 ={}&2n\gamma D_uD_vuA_0\\
 &-(2n\gamma pD_u+n\gamma^2)D_uD_vA_0^2
 -2n\gamma D_uD_vA_0N,
\end{split}                                                     \tag{RB14}
\]
where $uA_0=B_0C_0/(\delta^2E)$. Both complete first normals lie in
$(A_0,B_0,C_0)$ with bounded coefficients. Every remaining hard entry
is bounded. These are full rational support statements, not a scalar
contact-value test. Clearing $D_u^2$ makes them polynomial; its inverse
is a pure-testing square factor when restored.

\paragraph{The contact Gram and its exact leading calculation.}
Let
\[
 Y_0=A_0-kC_0,\quad Z_0=uB_0-kC_0,\quad H_0=kC_0-D_0.
\]
Use the separated nonnegative monomials $a_i^2=(A_0,uB_0,kC_0,D_0)$
and the two pairings $Q_1=a_1a_2-a_3a_4$,
$Q_2=a_1a_3-a_2a_4$. Their exact circuit is
\[
\begin{split}
 T_4&=A_0^2+(uB_0)^2+(kC_0)^2+D_0^2-4k^2B_0C_0\\
 &=\tfrac12\{(Y_0-Z_0)^2+H_0^2+Y_0^2+(Z_0+H_0)^2\}
                   +Q_1^2+Q_2^2.
\end{split}                                      \tag{RB15}
\]
Its matrix on $(Y_0,Z_0,H_0)$ is
$\left(\begin{smallmatrix}1&-1/2&0\\-1/2&1&1/2\\0&1/2&1\end{smallmatrix}\right)$,
which is positive definite with determinant $1/2$. All radicals
separate: $a_1$ is pure testing, $a_2,a_3$ are products of pure-source
and pure-testing radicals, and $a_4$ is pure source.

Here is enough explicit algebra to reproduce the leading collection.
At $g=0$, the canceled source expansions through the needed orders are
\[
\begin{aligned}
 n/s^2&=\tfrac12\{1-b(h+t)\}+O(s),&
 mg&=\tfrac12+\tfrac12(bh-1)s+O(s^2),\\
 C&=\frac{h+t}{2l}-\frac{t+1}{2l}s+O(s^2),&
 A&=\frac t{2h}-\frac t{2h}s+O(s^2),\\
 D&=\frac1{2(t-l)}-
       \frac{bt(h+t)+t-l}{2(t-l)^2}s+O(s^2),&
 \alpha&=\tfrac12 b^2(h+t)s^3+O(s^4).
\end{aligned}                                                     \tag{RB16}
\]
Here $O(s^j)$ means exact divisibility by $s^j$ in a rational ring with
unit denominator, not an asymptotic substitution in the proof identity.
Give $k,u,\omega$ weight one and $v$ weight two. Put $q=b(h+t)$.
Using $s=k^2/b$ and
$\delta=(1-lk^2+O(k^4))/\sqrt{bt}$, $\delta E=\sqrt{t/b}$,
the scalar $G_0=G|_{N=M_B=0}$ satisfies
\[
 G_0=\frac b{2h}(1+u^2)X_0^2+\frac q2T_4
       -\frac{k^2[1+bh(h+t)]}{2h}X_0^2+O_{\rm wt}(8)
       \quad\pmod g.                               \tag{RB17}
\]
The complete first-normal coefficients obey
\[
 L_N=O_{\rm wt}(5),\qquad
 L_{M_B}=\frac{X_0}{h\delta}+O_{\rm wt}(4)\quad\pmod g,
 \tag{RB18}
\]
and the hard block has residue
\[
 \operatorname{diag}\left(K_m+\tfrac12,\,
           A_*+C_*+D_*+\tfrac12\right),
 \quad A_*={t\over2h},\ C_*={h+t\over2l},\ D_*={1\over2(t-l)}.
 \tag{RB19}
\]

Apply the inherited monic contact--conductor identities CC1--CC6 to
the exact difference in (RB17), and to the exact first-normal
differences in (RB18). Their hypotheses are exactly the raw support in
(RB13)--(RB14), together with zero contact value and full first jet.
Both follow from the original stationary parent and (RB10); the model
$X_0^2,T_4$ also has these zeros. No contact difference is inverted.
The conductor maps preserve this grading: $X_0$ has weight two,
the five circuit roots have weight three. The scalar difference of
weight at least eight and normal differences of the displayed weights
therefore have strictly positive-weight Gram coefficients. Every such
coefficient is infinitesimal. Terms divisible by $g$ also lift with
bounded coefficients and remain infinitesimal: the contact (RB9),
$\delta$ and $E$ are independent of $g$, and the canceled coefficient
ring has unit denominators at $g=0$. Clearing and restoring its finite
unit denominators changes none of these conclusions.

Consequently the residue on the eight roots
$(X_0,Y_0,Z_0,H_0,Q_1,Q_2,N,\delta M_B)$ is the strict circuit with
coefficient $q/2>0$, the $N$ pivot $K_m+1/2>0$, and the coupled block
\[
 \begin{pmatrix}
 b/(2h)&bt/(2h)\\
 bt/(2h)&bt(A_*+C_*+D_*+1/2)
 \end{pmatrix}.                                    \tag{RB20}
\]
Its determinant is
\[
 \frac{b^2t}{2h}(C_*+D_*+1/2)>0.                    \tag{RB21}
\]
All pivots are positive bounded units. The separated strict-residue
Gram lemma therefore gives finite tensor SOS for the exact remainder.
Restore (RB10) and the pure-source/pure-testing clearing factors.
This proves the whole local contact parent, for every independent rate
and every normal direction.

\paragraph{All remaining finite state and angle charts.}
The following details make the join an all-tests result.
For bounded $u$ with nonzero residue and infinitesimal $v$, retain
$nD_uD_vX^2$ and $\alpha D_u(vY-Z)^2$, and the centered $m$ excess
$(m-K_m)N^2$ from (RB6), now homogeneous in $Z$. Put
$V_Z=vZ$ and $R_t=uW-Y$. After $X=N-uZ+pv(W+uY)$, the whole
remainder is a bounded quadratic on $(N,R_t,Y,V_Z)$. Its residue
splits off $K_m+1/2$ on $N$, with the other form
\[
 C_*R_t^2+A_*(1+u^2)Y^2
 +(D_*+1/2)\{R_t-(1+u^2)V_Z\}^2
 -(1+u^2)YV_Z-\tfrac12(1+u^2)V_Z^2.               \tag{RB22}
\]
Its first two leading principal minors are positive, and its determinant
is exactly
\[
 (1+u^2)^2A_*C_*(D_*+1/2)u^2>0.                    \tag{RB23}
\]
This direct four-root proof covers every state, including $Z=0$, with
no negative-beta or small-$\rho$ argument.

It remains to consider both angles infinitesimal. The exact testing-core
sign reduction from SD5 handles nonnegative core and literal axes by
original-event squares. On its negative-core branch one may use the
standard exact gauges $u,v>0$. Normalizing the original state by a
maximal coordinate, the case $Z$ infinitesimal has $W$ a unit and is
covered by BE4, with positive real caps of $m,A$ supplied by (RB4).
Retain the actual nonnegative $n,\alpha$ positive square terms. That exact selected
full-parent frame is $(X,Y,uW,vW)$ and has strict residue; no coefficient
bound involving $b$ enters its proof.

If $Z$ is a unit normalize $Z=1$. The same core identity gives
$\omega>0$, and $\omega$ is bounded. If $\omega$ is infinitesimal,
use (RB17)--(RB21). If $\omega$ is a positive unit, use the exact
relations $B_0=(\delta/\omega)A_0$ and $C_0=B_0-X_0$ in the raw
support (RB13)--(RB14). The inverse is pure testing. Its complete
residue on $(A_0,X_0,D_0,N,\delta M_B)$ consists of the positive
$q/2$ pivots on $A_0,D_0$, the $K_m+1/2$ normal pivot, and precisely
(RB20). Direct substitution in (RB11) gives these residues; all other
coefficients contain $k,g,u,v$ after the cancellations already recorded.
Thus this whole state chart is strict as well.

\paragraph{Every other projective test.}
Use the exact centered-event identity R5.5 for the finite remote pairs
$v_0\ne0$. All source residues are (RB4) and $n_0=\alpha_0=0$.
On the kernel of the limiting centered $m$ row, the determinant is
\[
 \frac{(1+U)^2(1+V)^2(h+t)}{8hl(t-l)}
       \{V^2+t(h+t)U^2V+hUV^2\}>0.                 \tag{RB24}
\]
The restriction to $Z=0$ is positive by the original $A,C$ positive square terms
and the nonnegative signed-core restriction, exactly as in RT10.
A sufficiently large positive real cap $K_m$ makes the full bounded
residue strict. Its choice depends on the real testing residues only;
unbounded $m$ permits every such real cap.

For unbounded $u$ and bounded unit $v$, R6.5 at $n_0=0$ is strict
because $C_*W^2$, $A_*Y^2$ and $(D_*+1/2)(W+vX)^2$ have common
kernel zero. For unbounded $u$ and infinitesimal $v$, the exact
BE $(\infty,0)$ frame $(pX+Z,W,Y,vX)$ retains every event except
positive caps of $C,A$, and has strict residue. For unbounded $v$
and bounded $u$, R6.6 at $n_0=0$ is strict; the centered correction
vanishes on its $m$-row kernel. For both angles unbounded, R6.7 at
$n_0=0$ has block determinants $D_*/2$ and
$A_*(K_m+1/2)$, both positive. These are exact homogeneous full-parent
identities. Their literal reciprocal zeros include the actual projective
endpoints; undoing a nonzero homogeneous factor uses a pure-testing
square only.

The charts exhaust the stated $b$-unit source domain. Restore all retained
squares and normalizations and apply the real identities to the real and
imaginary parts of a complex state.

\paragraph{Imported exact identities.}
The BE4 four-root state identity and the unnumbered $(\infty,0)$
identity in Section~\ref{component:p69995d829bb2} supply the two
selected full-parent frames used above. The latter is the paragraph
beginning ``At $(\infty,0)$ retain all events except $c,A$''.
The remote kernel determinant and its $Z=0$ restriction are printed
in (RB24) and the paragraph following it. The centered-event identity
is R5.5 in Section~\ref{component:p8da9ea6ed9e5}. The homogeneous
kernel and root matrices R6.5--R6.7 are printed in
Section~\ref{component:pd5fa361b4623}; they are used at $n_0=0$
with the source residues (RB4). Finally SD5 and its following sign
Schur identity in Section~\ref{component:p0acdb66bba2f} give the
pure-testing inverse and the exact sign reduction. Only these
algebraic identities are used with the present source estimates;
no C1 physical-source or negative-beta perturbation theorem is imported.

\subsection{Infinitesimal $a^2B$}
\label{component:p75b395956985}

\label{p75b395956985:sec:r6-b-inf}

\paragraph{Source scope and relation to the preceding theorem.}
Retain the complete original weak-C2 source: $d=1$, $\lambda=\beta=0$,
$n\geq0$, the actual stationary contact, every strict gap, total scale,
and the complement of prior cases. On R6, $t,h,l=1-h,r$ and
the original $j$ gap are positive bounded units, while $g,B,x=aB,y=tB$
are positive infinitesimals. This section treats $a$ unbounded and
$b=a^2B$ positive infinitesimal. The notation $s=1/a$, $w_s=ae$,
$B=s^2b$, $x=sb$, $e=s(1+hsb-g)$ is unchanged from (RB1).
The exact physical identity (RB2), equivalently RT2, remains in force.
It bounds $b$ on all unbounded-$a$ R6 and gives $\rho_0=1/2$ here,
where $\rho=n/e^2$. In particular this subdomain has $n>0$; the
literal $n=0$ source face is covered by the preceding $b$-unit theorem.

\begin{theorem}
On this entire original R6 source intersection with $b=a^2B$
infinitesimal, the complete original parent is a finite separated
source--testing tensor sum of squares at every independent projective
testing pair and every real or complex state. Both angle signs, all
literal projective endpoints, and all relative opening rates of
$s,b,g$ are included. No mixed inverse or mixed square root is used.
\end{theorem}

Together with the $b$-unit theorem and the inherited bounded-$a$
R6 theorem, this proves the entire R6 row. It does not by itself
admit another weak-C2 row or the global finite-lift theorem.

\paragraph{A regular source chart.}
Keep $R=1+hx$, $K=1+x$, $T=1+xr$, $k=\sqrt{x}$ and the contact
parameters $\delta=R/(K\sqrt{bt})$, $E=tK/R$ from (RB8). Define
\[
 \varepsilon=\delta^{-1},\qquad \jmath=\delta k,
 \qquad k=\varepsilon\jmath.
 \tag{RI1}
\]
Both $\varepsilon$ and $\jmath$ are positive infinitesimals. In this
section $\jmath$ denotes the contact scale, not the original source
gap called $j$. The exact source substitution is
\[
 \begin{gathered}
 x=\varepsilon^2\jmath^2,\quad K=1+x,\quad R=1+hx,\quad
 E=tK/R,\\
 b=\frac{\varepsilon^2R^2}{K^2t},\qquad
 s=\frac{t\jmath^2K^2}{R^2}.
 \end{gathered}                                                   \tag{RI2}
\]
All denominators in (RI2) are positive source units. There is no
comparison of a testing variable with either source infinitesimal.

Use precisely the complete contact-centered remainder $G$ in (RB11),
with the full coordinate inverse (RB7). Thus the two retained positive square terms
on the contact chart are $(m-K_m)N^2$ and
$nD_uD_v(X+\gamma v\omega)^2$, where
\[
 p=(1+hx-g)/R,\quad c=l/(tT),\quad
 \gamma=rK/(sR),\quad D_u=1+u^2,\quad D_v=1+v^2.
\]
Its exact contact is now
\[
 (u,\omega,v;N,M_B)
   =(\varepsilon\jmath,\jmath,\varepsilon\jmath^2E;0,0).
 \tag{RI3}
\]
Substitution of (RI3) in the complete contact formulas
(RB7)--(RB9a) makes both normals and the retained $n$ row vanish.

\paragraph{Two exact source cancellations.}
The quantities $mg$, $n\gamma$, $n\gamma^2$ and
$\alpha/(s^3b^2)$ have canceled denominators that are units at
$s=b=g=0$. Introduce
\[
 \mu_1=mg/R,\quad \mu_2=(mg)g/R^2,\quad
 F_0=p^2/2-n\gamma^2-2n\gamma p+\mu_2,\qquad
 A_N=p-2\mu_1-2n\gamma.
 \tag{RI4}
\]
Both $F_0$ and $A_N$ are divisible by $b$ in this rational unit ring.
Here is a short exact verification, independent of a series expansion.
At $b=0$, put $q=1-g$. Direct cancellation in the original source gives
\[
 \begin{gathered}
 p=q,\quad mg=\frac{q}{2(1+s)},\quad
 n\gamma=\frac{sq}{2(1+s)},\quad
 n\gamma^2=\frac{q(1-sq)}{2(1+s)}.
 \end{gathered}                                                    \tag{RI5}
\]
Substitution makes both expressions in (RI4) identically zero.
The denominators are units before this specialization; hence the
polynomial numerators are divisible by $b$, without introducing a
pole. Their needed residues are
\[
 \frac{F_0}{b}\sim\frac{h+t}{2},\qquad
 \left.\frac{A_N}{b}\right|_{s=g=0}=0,\qquad
 \left.\frac{A_N}{bs}\right|_{g=0}\sim h+t.
 \tag{RI6}
\]
The last quotient also has a unit canceled denominator at $s=b=0$.
Every statement in (RI5)--(RI6) is separately checked from the original
source formulas. In particular the source estimate is uniform in
$g/s$, $s/b$ and all other opening ratios.

\paragraph{Complete bounded raw support, including both first normals.}
Use the separated raw monomials
\[
 \mathsf A=\varepsilon\omega v,\quad
 \mathsf B=v,\quad \mathsf C=Eu\omega,\quad
 \mathsf D=\varepsilon^2\jmath^3E,
 \qquad \mathsf x=\mathsf B-\mathsf C.
 \tag{RI7}
\]
Set $H=\alpha+1/2$, $L=AD_v+Hv^2$, and $V_0=1-puv\omega$.
At zero normals, (RB7) gives
\[
 q_B=D_uvV_0,\quad uW-Y=cD_uvV_0,\quad
 Y=u\omega-cvV_0,\quad
 J=-(1-c)D_uvV_0,
\]
and the entire scalar has the following exact allocation:
\[
 G_0=f_A\mathsf A^2+f_{BB}\mathsf B^2+
 f_{BC}\mathsf B\mathsf C+f_{CC}\mathsf C^2+
 \varrho\mathsf D^2,                               \tag{RI8}
\]
where
\[
\begin{aligned}
 f_A={}&D_u^2D_vF_0/\varepsilon^2,\\
 f_{BB}={}&[Cc^2D_u^2D_v+(D+1/2)(1-c)^2D_u^2
                  +D_uLc^2]V_0^2\\
          &+2HD_ucV_0-D_u/2-D_u^2\mu_2v^2\omega^2,\\
 f_{BC}={}&[-2D_uLcV_0-2HD_u+2n\gamma D_uD_v]/E,\\
 f_{CC}={}&[D_uL+D_uD_vn\gamma^2v^2]/E^2,\\
 \varrho={}&\alpha D_u/\mathsf D^2.
\end{aligned}                                                      \tag{RI9}
\]
These coefficients are bounded. In particular
\[
 \frac{F_0}{\varepsilon^2}=
       \frac{R^2}{K^2t}\frac{F_0}{b},\qquad
 \frac{\alpha}{\mathsf D^2}=
       \frac1t\frac{\alpha}{s^3b^2}.
 \tag{RI10}
\]
The term that would give $u^2/\varepsilon^2$ in the first raw
coefficient has instead been placed in $f_{CC}$, using the exact
identity $u^2(v\omega)^2=(v\mathsf C/E)^2$.
The final term of $f_{BB}$ likewise allocates the whole
$-D_u^2\mu_2v^4\omega^2$ contribution without dividing by
$\varepsilon$. These are identities for the full scalar.

Write the exact normal expansion as
$G=G_0+L_NN+L_MM_B+G_{\rm hard}$, with $G_{\rm hard}$ quadratic.
The two complete linear coefficients are
\[
 L_N=R_A\mathsf A+R_B\mathsf B+R_C\mathsf C,
 \qquad L_M=S_B\mathsf B+S_C\mathsf C,               \tag{RI11}
\]
where
\[
 R_A=D_uA_N/\varepsilon,\qquad R_C=2cuvL/E,\qquad S_C=-2L/E,
\]
\[
 R_B=\Delta_{v,0}(L_N-R_A\mathsf A-R_C\mathsf C),\qquad
 S_B=\Delta_{v,0}(L_M-S_C\mathsf C).\tag{RI12}
\]
Here $L_N,L_M$ are the complete coefficient extractions from the bounded
full parent in the stated state coordinates. The displayed source
cancellations make both quotient numerators vanish identically at
$v=0$. Thus these are finite monic polynomial divisions, with the pure
testing denominator $D_u^2$ cleared first. They give all normal
coefficients, including the higher testing terms. Their coefficients
are bounded; division shifts testing degree by one and preserves the
weighted coefficient estimates used below.

Here $A_N/\varepsilon=
\varepsilon R^2(A_N/b)/(K^2t)$ is bounded. Thus no first-normal
pole has been discarded. Every hard coefficient is bounded as well:
the only testing inverse in the full change of state is $D_u^{-1}$,
and all original source coefficients in (RB11) are bounded after the
displayed cancellation. At $u=v=0$ the complete hard matrix is exactly
\[
 \operatorname{diag}(K_m+1/2,A+C+D+1/2).             \tag{RI13}
\]

\paragraph{A bounded conductor for these complete coefficients.}
Define the contact differences
\[
 \mathsf y=\mathsf A-k\mathsf C,\quad
 \mathsf z=u\mathsf B-k\mathsf C,\quad
 \mathsf h=k\mathsf C-\mathsf D.
\]
The scaled conductor identities are
\[
\begin{aligned}
 \mathsf C(u-k)&=\mathsf z-u\mathsf x,\\
 \varepsilon\mathsf C(\omega-\jmath)
       &=\mathsf y-\varepsilon\omega\mathsf x,\\
 \mathsf C(v-\varepsilon\jmath^2E)
       &=E(\omega\mathsf z+\jmath\mathsf y-k\omega\mathsf x),\\
 \mathsf Cv(\omega-\jmath)
       &=E\omega(\omega-\jmath)\mathsf z+
         E\jmath\omega\mathsf y-Ek\omega^2\mathsf x.
\end{aligned}                                                       \tag{RI14}
\]
These are SE5--SE6 with $w=\omega$, also checked directly here.
For a monic contact division
\[
 S=(u-k)T_u+(\omega-\jmath)T_\omega+
           (v-\varepsilon\jmath^2E)T_v,
 \qquad T_\omega=\varepsilon U+vV,
\]
the exact bounded lift is
\[
\begin{aligned}
 \mathsf CS={}&\mathsf x[-uT_u-\varepsilon\omega U
                    -Ek\omega^2V-Ek\omega T_v]\\
 &+\mathsf y[U+E\jmath\omega V+E\jmath T_v]\\
 &+\mathsf z[T_u+E\omega(\omega-\jmath)V+E\omega T_v].
\end{aligned}                                                       \tag{RI15}
\]
No inverse of a small source parameter or a contact difference occurs.

We verify the hypothesis of (RI15) for the full coefficients, rather
than import the different small-$e$ source assumptions of SE8.
Clear the finite source unit denominators and $D_u^2$. They can be
chosen independent of $\omega$ and even in $u,v$. In (RI9), all
$\omega$ dependence is through $uv\omega$ and $v^2\omega^2$;
the other testing dependence is even in $u,v$. The same property
holds after subtracting any model of the form used below.

For the scalar square-lift algorithm CC4, substitute
$\mathsf A=\mathsf y+k\mathsf C$,
$\mathsf B=\mathsf x+\mathsf C$,
$\mathsf D=k\mathsf C-\mathsf h$ into (RI8).
The coefficient requiring its first conductor is
\[
 S=k^2f_A+f_{BB}+f_{BC}+f_{CC}+k^2\varrho.
\]
Its zero contact value follows by domain cancellation, as in CC4.
Divide monically in the order $u$, then $\omega$, then $v$.
After $u=k=\varepsilon\jmath$, the $\omega$ divided difference
belongs to $(\varepsilon v,v^2)$. Accordingly choose
$T_\omega=\varepsilon U+vV$ with $U,V\in(v)$.
Modulo $\varepsilon$, the remaining $v$ polynomial is even, so
$T_v\in(\varepsilon,v)$; also
$T_u|_{\varepsilon=u=v=0}=0$.

After this first conductor the coefficient of $\mathsf h$ remains
$-2k\varrho$. The sole gradient relation is
$(-2k,1,1,-1)$ on
$(\mathsf x,\mathsf y,\mathsf z,\mathsf h)$; the determinant of
the first three contact gradients is
$-\varepsilon^3\jmath^4E^2\ne0$. The adjugate cancellation of
CC5 therefore removes the contact values without putting an inverse
of this determinant into any coefficient. Subtract the contact
multiple $2k\varrho_c(-2k,1,1,-1)$.
Each of the four resulting coefficient remainders $J_i$ vanishes at
contact. Its restriction at $\varepsilon=u=v=0$ is independent of
$\omega$: for the $\mathsf x$ coefficient this follows from
$2f_{BB}+f_{BC}$; the other three restrictions are zero by the
$U,V,T_u,T_v$ properties just established and the factor $k$.
Consequently the next monic $\omega$ divided difference of every
$J_i$ belongs to $(\varepsilon,v)$. Formula (RI15) applies again.

The last scalar repair uses the exact circuit relation
\[
 k\mathsf C(\mathsf y+\mathsf z-\mathsf h-2k\mathsf x)
       =Q_1^2-\mathsf y\mathsf z,
\]
where $Q_1$ is the first of the separated radical pairings below.
This completes the same restricted scalar square lift as CC6.
For (RI11), one conductor suffices: after replacing
$\mathsf A,\mathsf B$ as above, the coefficient multiplying
$\mathsf C$, restricted at $\varepsilon=u=v=0$, is independent of
$\omega$ directly from (RI12). Its divided difference therefore
lies in $(\varepsilon,v)$ and (RI15) applies. The full scalar and
both full normals have thus been lifted over the bounded coefficient
ring. The restricted raw support $\mathsf A^2,\mathsf B^2,
\mathsf B\mathsf C,\mathsf C^2,\mathsf D^2$ is essential;
an unrestricted ideal claim is neither needed nor asserted.

The original stationary parent and the two retained contact-zero
positive square terms give the exact scalar value and full first-jet identities
used here. The contact (RI3) is independent of $g$. Hence taking
the exact coefficient difference divisible by $g$ preserves these
identities and every raw-support condition.

\paragraph{Strict residues of the complete contact lift.}
Suppose now that $u,v,\omega$ are all positive infinitesimals, as
in the remaining contact state chart. Give
\[
 (\varepsilon,\jmath,u,\omega,v)
       \text{ weights }(1,1,2,1,3).
\]
Then $\mathsf x$ has weight three and the five circuit roots have
weight five. Every contact substitution and (RI14)--(RI15) preserves
the weights. Put $\kappa=(h+t)/(2t)$ and
\[
 \Lambda=\frac{1-s+s^2+u^2-[2l+h(h+t)]k^2}{2ht},\qquad
 \mu=\frac1{h(1+s)}.
\]
Direct exact collection of (RI9) through weight ten gives
\[
 G_0=\Lambda\mathsf x^2+\kappa T_4+O_{\rm wt}(12)
                   \quad\pmod g,                 \tag{RI16}
\]
where
\[
 T_4=\mathsf A^2+(u\mathsf B)^2+(k\mathsf C)^2+
          \mathsf D^2-4k^2\mathsf B\mathsf C.
\]
The two complete normal collections are
\[
 L_N=O_{\rm wt}(8),\qquad
 L_M=\mu\mathsf x+O_{\rm wt}(7)\quad\pmod g.
 \tag{RI17}
\]
For transparency, the only source terms needed for (RI16)--(RI17),
besides (RI6), are the following exact Taylor coefficients at $g=0$:
\[
\begin{aligned}
 A&=\frac t{2h}(1-s+s^2)+O_{\deg(s,b)}(3),\\
 C&=\frac{h+t}{2l}-\frac{t+1}{2l}s+
                      \frac{t+1}{2l}s^2+O_{\deg(s,b)}(3),\\
 D&=\frac{1-s+s^2}{2(t-l)}-
        \frac{sbt(h+t)}{2(t-l)^2}+O_{\deg(s,b)}(3),\\
 n\gamma&=\frac s2-\frac{s^2}2-\frac{sb(h+t)}2+
                         O_{\deg(s,b)}(3).
\end{aligned}                                                       \tag{RI18}
\]
Here the error denotes exact degree filtration in a rational ring
with unit denominators. Also $s=t\jmath^2+O_{\rm wt}(6)$,
$c=(l/t)(1-k^2)+O_{\rm wt}(6)$ and
$E=t(1+lk^2)+O_{\rm wt}(8)$. The weight eight bound for $L_N$
follows term by term from (RI6), (RI12).
At $g=0$, after the displayed unit denominators are cleared, all
scalar weighted monomials have even weight, and all $M_B$ linear
monomials have odd weight. Hence the next possible degrees are
twelve and seven, respectively.

The stronger coefficient-wise statement also holds for the displayed
raw representation: $f_A-\kappa$ and $\varrho-\kappa$ have weight at
least two, while
\[
 f_{BB}-\Lambda-\kappa u^2,\qquad
 f_{BC}+2\Lambda+4\kappa k^2,\qquad
 f_{CC}-\Lambda-\kappa k^2
\]
each has weight at least six. Similarly $S_B-\mu,S_C+\mu$ have
weight at least four separately, and $(R_A,R_B,R_C)$ have weights
at least $(3,5,5)$. Thus the chosen raw allocation of each exact
difference already has the claimed filtration; no low-weight
Gram coefficients are hidden by cancellations between raw products.

Apply the bounded conductor just proved to the exact differences in
(RI16)--(RI17). The scalar correction Gram coefficients have weight
at least two; the $M_B$ normal corrections have weight at least two;
the $N$ normal corrections have weight at least three. They are all
infinitesimal. The coefficients divisible by $g$ lift boundedly by the
same argument and are infinitesimal independently of every other
opening. Restoring the finite unit denominators preserves this fact.

For strictness use the two pairings on the separated monomials
\[
 a_i^2=(\mathsf A,u\mathsf B,k\mathsf C,\mathsf D),\qquad
 Q_1=a_1a_2-a_3a_4,\quad Q_2=a_1a_3-a_2a_4.
\]
Their product relation gives the exact two-pairing identity (RB15),
with the scaled roots. Its three-dimensional difference block has
determinant $1/2$, and both pairing pivots are positive. The complete
residue on
$(\mathsf x,\mathsf y,\mathsf z,\mathsf h,Q_1,Q_2,N,M_B)$
therefore consists of this strict circuit with coefficient
$\kappa>0$, the $N$ pivot $K_m+1/2>0$, and
\[
 \begin{pmatrix}
 1/(2ht)&1/(2h)\\
 1/(2h)&A_*+C_*+D_*+1/2
 \end{pmatrix},
 \quad A_*=t/(2h),\ C_*=(h+t)/(2l),\ D_*=1/[2(t-l)].
 \tag{RI19}
\]
Its determinant is exactly
\[
 \frac{C_*+D_*+1/2}{2ht}>0.                         \tag{RI20}
\]
All pivots are positive bounded units. The strict separated Gram
lemma WA10 yields a finite tensor-SOS identity for the complete
remainder. Restore the two actual nonnegative retained positive square terms and
all separated clearing factors. This proves the entire local contact
chart, including every normal direction and arbitrary relative rates.

\paragraph{The $\omega$-unit state chart uses the original $n$ positive square term.}
When $u,v$ are infinitesimal, $Z=1$, and $\omega$ is a positive unit,
retain $nD_uD_vX^2$ itself, together with $(m-K_m)N^2$.
The complete remainder is (RB11) with its two negative $n\gamma$
and $n\gamma^2$ contributions absent. This choice is legitimate
directly from $n\geq0$ and does not require a contact lift away from
the contact state scale.

Use the raw roots
\[
 (A_o,\mathsf x,\mathsf D,N,M_B),\qquad A_o=\omega v.
\]
The only new inverse is $1/\omega$, which is pure testing and
bounded. The relations $\mathsf B=A_o/\omega$ and
$\mathsf C=A_o/\omega-\mathsf x$ put the entire scalar and both
first normals in this frame, by the same full formulas (RI9),
(RI12) with the absent $n$ terms deleted. All coefficients are
bounded; the constant term is still $\alpha D_u=\varrho\mathsf D^2$.
Explicitly $\mathsf A=\varepsilon A_o$, so the first scalar term
becomes $D_u^2D_v(p^2/2+\mu_2)A_o^2$ and the first $N$-normal
term becomes $D_u(p-2\mu_1)A_o$. The printed $\varepsilon$
denominators therefore cancel before any coefficient estimate.
Its exact residue is the $A_o$ pivot $1/2$, the $\mathsf D$ pivot
$\kappa$, the $N$ pivot $K_m+1/2$, and the coupled block (RI19).
The first $N$ normal has zero residue and the first $M_B$ normal
has residue $\mathsf x/h$. Thus this entire state chart
also has a strict finite separated Gram.

\paragraph{All other independent tests and states.}
Every source bound used by the noncontact joins of the preceding
$b$-unit theorem is (RB4), which is valid here without a lower bound
on $b$. For bounded unit $u$ and infinitesimal $v$, retain the original
$n,\alpha$ events and use the exact full-state form (RB22), whose
determinant (RB23) is positive. For both angles infinitesimal, the
SD5 exact sign reduction disposes of the nonnegative-core branch
and literal zero axes by original-event squares. On the negative
branch choose the exact gauges $u,v>0$. If the normalized state has
$Z$ infinitesimal, BE4 gives the strict full-state frame
$(X,Y,uW,vW)$ using positive real caps of $m,A$.
Otherwise normalize $Z=1$; $\omega>0$ is bounded, and its two
possibilities are exactly the contact chart and the unit chart just
proved. No source/testing ratio split is used.

For finite remote $v_0\ne0$, use the actual centered-event identity
R5.5, the strict kernel determinant RT10 (printed as RB24), and a
sufficiently large positive real cap below unbounded $m$.
For unbounded $u$ and bounded unit $v$, use R6.5 specialized to
$n_0=0$ and the source residues (RB4). For unbounded $u$ and
infinitesimal $v$, use the exact BE $(\infty,0)$ full frame
$(p_uX+Z,W,Y,vX)$, $p_u=1/u$, with positive caps of $C,A$.
For unbounded $v$ and bounded $u$, R6.6 at $n_0=0$ is strict.
For both angles unbounded, R6.7 at $n_0=0$ has determinants
$D_*/2$ and $A_*(K_m+1/2)$ on its two nontrivial blocks.
These uses import the displayed homogeneous identities with the
source bounds (RB4), without importing a source theorem from another
stratum.
Literal reciprocal zeros are polynomial values of those charts.
All signs and complex states follow by the same exact gauges and
the real/imaginary split. This exhausts all independent tests.

\paragraph{The contact identities.}
Equations (RI8)--(RI15) establish the rational-support hypotheses of
SE5--SE7 in Section~\ref{component:pa7e744614f49} for the present
coefficients, with the bounds proved here. We then apply the monic
square-lift and gradient repair CC4--CC6 in
Section~\ref{component:pfabc29e28ab3}.

\section{Two small actual angles: the R7 sector}
\subsection{Unbounded $a$, including the physical zero face}
\label{component:p405d96527133}

\label{p405d96527133:sec:r7-unbounded-a}

\paragraph{Quantified source and conclusion.}
Keep the entire original weak-C2 source: $d=1$, $\lambda=\beta=0$,
the actual rational stationary coefficients, $n\ge0$, every strict gap,
the original finite stationary contact, total scale, and all prior
exclusions. In R7, $B,z,y$ are positive infinitesimals,
$a=z/B$, $t=y/B$ is a positive bounded unit, and $h,g$ are positive
infinitesimals. The original exclusion of simultaneously infinitesimal
$r$ and the source gap $j_{\rm gap}$ is retained.

\begin{theorem}
On the whole original R7 source intersection with $a$ unbounded,
the complete original parent is a finite separated source--testing
tensor sum of squares for every independent projective testing pair
and every real or complex state. Both angle signs, all literal zero
and infinite angles, literal $n=0$, and every relative opening rate
of the source infinitesimals are included. No inverse or radical of
a mixed source--testing expression is used.
\end{theorem}

Together with the bounded-$a$ theorem this covers R7.

\paragraph{Actual source scale and the physical inequality.}
Use
\[
 \begin{gathered}
 s=1/a,\quad b=a^2B,\quad x=sb=z,\quad B=s^2b,\\
 w_s=ae=1+hx-g,\quad e=sw_s,\quad r=1-e,\quad l=1-h.
 \end{gathered}
 \tag{RU1}
\]
Set $K=1+x$, $R=1+hx$, $T=1+xr$ and
$j_{\rm gap}=T-l/t$. The exact original source identity is
\[
\begin{gathered}
 2r\Delta n=eQ-BK_{\rm phys},\qquad
 2r\Delta_*n/s^2=w_sQ-bK_*,\qquad \Delta_*=s\Delta,\\
 Q=1+B(1+t+ah+2ar+at)
    +B^2\{a(h+t)+a^2[h(r+1)+t(1+r-r^2)]\},\\
 K_{\rm phys}=(a+1)(h+t)(1+x),\qquad
 K_*=(1+s)(h+t)(1+x).
\end{gathered}                                                     \tag{RU2}
\]
Before bounding $b$, $Q\sim\Delta_*\sim1$ since all their nonconstant
terms contain $s$ or $x$. Here $w_s\sim1$ and $K_*\sim t$ is a
positive unit. The original inequality $eQ\ge BK_{\rm phys}$
therefore gives $b$ bounded and $n=O(s^2)$, including its literal
equality case. In particular $r\sim1$, so the original R7 exclusion
is respected. The source gap gives $t_0\ge1$; $j_{\rm gap}$ is a
unit exactly when $t_0>1$, and otherwise $t_0=1$.

The original rational formulas, after cancellation, give jointly
\[
 mg\sim\tfrac12,\quad C\sim\tfrac t2,\quad
 A_h:=Ah\sim\tfrac t2,\quad
 \beta_D:=D(1-c)=Dj_{\rm gap}/T\sim\frac1{2t},
 \quad \frac\alpha{s^3b^2}\sim\tfrac t2,
 \qquad c=\frac l{tT}.
 \tag{RU3}
\]
Thus $m,A$ are unbounded; $C$ is a positive unit. If $t_0>1$,
$D\sim1/[2(t-1)]$ is a positive unit. If $t_0=1$, $D$ is unbounded.
The canceled coefficients $mg,C,A_h,\beta_D,n\gamma,n\gamma^2$
below have unit denominators through both small-gap ends. No ratio
$h/g$, $j_{\rm gap}/h$, or comparison with a testing rate is used.

\paragraph{Four whole positive square terms and a bounded full-parent remainder.}
Use the complete parent $P=F/4$:
\[
\begin{aligned}
 P={}&nD_uD_vX^2+mM^2+CD_v(uW-Y)^2+AD_uD_vY^2\\
 &+(D+1/2)J^2+(\alpha+1/2)D_u(vY-Z)^2
                   +D_v(X+uZ)^2/2-D_uD_vZ^2/2,\\
 M={}&vW-X+uvY-uZ,\qquad J=uW+uvX-Y-vZ.
\end{aligned}
\]
The exact coordinate inverse (RB6)--(RB7) is
\[
\begin{gathered}
 D_u=1+u^2,\quad D_v=1+v^2,\quad Z=1,\quad
 p=w_s/R,\quad q_A=D_uv\omega,\\
 N=X+u-pq_A,\qquad q_B=v(1-uX),\qquad M_B=uW-Y-cq_B,\\
 X=N+pD_uv\omega-u,\quad q_B=v(D_u-uN-puD_uv\omega),\\
 W=\omega+u(M_B+cq_B)/D_u,\qquad
 Y=u\omega-(M_B+cq_B)/D_u.
\end{gathered}                                                     \tag{RU4}
\]
The sole testing inverse is $D_u^{-1}$. Put
\[
 \begin{gathered}
 \gamma=\frac{rK}{sR},\quad \eta=Y-huW,\quad q=1-c=j_{\rm gap}/T,\\
 \mu_1=mg/R,\quad \mu_2=(mg)g/R^2,\quad n_1=n\gamma,\quad n_2=n\gamma^2.
 \end{gathered}
\]
Retain the four entire nonnegative squares
\[
 mN^2,\qquad nD_uD_v(X+\gamma v\omega)^2,
 \qquad AD_uD_v\eta^2,\qquad DM_B^2.
 \tag{RU5}
\]
All four rows vanish at the actual contact; (RB9a) verifies the first
two, and directly $Y_c=hu_cW_c$, $M_{B,c}=0$.
The $n$ positive square term vanishes identically on the literal $n=0$ face.

The full $A,D$ cancellations are the exact identities
\[
\begin{aligned}
 AD_uD_v(Y^2-\eta^2)&=2A_hD_uD_vuWY-A_hhD_uD_v(uW)^2,\\
 D(J^2-M_B^2)&=-2\beta_DM_Bq_B+\beta_Dq q_B^2,
 \qquad J=M_B-q q_B.
\end{aligned}                                                     \tag{RU6}
\]
Since $M=(g/R)q_A-N$, subtraction of precisely (RU5) leaves
\[
\begin{aligned}
 G={}&-2\mu_1q_AN+\mu_2q_A^2
       -2n_1D_uD_vX(v\omega)-n_2D_uD_v(v\omega)^2\\
    &+CD_v(uW-Y)^2+2A_hD_uD_vuWY-A_hhD_uD_v(uW)^2\\
    &-2\beta_DM_Bq_B+\beta_Dq q_B^2+J^2/2\\
    &+(\alpha+1/2)D_u(vY-1)^2+D_v(X+u)^2/2-D_uD_v/2.
\end{aligned}                                                     \tag{RU7}
\]
This is a whole-parent identity, including every scalar, normal and
hard coefficient. Equations (RU3), (RU6) bound all of them before
any contact restriction. In particular neither $h^{-1}$ nor
$j_{\rm gap}^{-1}$ survives.

\paragraph{One regular contact frame for the two $b$ alternatives.}
Use the pure-source parameters
\[
 \varepsilon=\frac{K\sqrt{bt}}R,\quad
 \jmath=\sqrt{s/t}\,\frac RK,\quad k=\varepsilon\jmath=\sqrt x,
 \quad E=tK/R.
 \tag{RU8}
\]
The symbol $\jmath$ is a contact scale, not $j_{\rm gap}$.
Here $\jmath$ is infinitesimal; $\varepsilon$ is a unit or
infinitesimal exactly as $b$ is. The exact inverse source chart is
\[
 x=\varepsilon^2\jmath^2,\qquad
 b=\frac{\varepsilon^2R^2}{K^2t},\qquad
 s=\frac{t\jmath^2K^2}{R^2}.
 \tag{RU9}
\]
Only source units are inverted. The contact is
$(u,\omega,v;N,M_B)=(k,\jmath,\varepsilon\jmath^2E;0,0)$.

Define
\[
\begin{gathered}
 \mathsf A=\varepsilon\omega v,\quad\mathsf B=v,
 \quad\mathsf C=Eu\omega,\quad\mathsf D=\varepsilon^2\jmath^3E,\\
 \mathsf x=\mathsf B-\mathsf C,\quad
 \mathsf y=\mathsf A-k\mathsf C,\quad
 \mathsf z=u\mathsf B-k\mathsf C,\quad
 \mathsf h=k\mathsf C-\mathsf D.
\end{gathered}                                                     \tag{RU10}
\]
The two potentially singular source combinations are
\[
 F_0=p^2/2-n_2-2n_1p+\mu_2,\qquad A_N=p-2\mu_1-2n_1.
\]
Their exact canceled quotients satisfy
\[
 F_0/b\sim t/2,\qquad A_N/(bs)\sim t.
 \tag{RU11}
\]
Both denominators are units through $s=b=h=g=0$. These assertions
follow from the original rational source, not from the inequality
$n\ge0$. In particular
\[
 F_0/\varepsilon^2=\frac{R^2}{K^2t}\frac{F_0}b,
 \qquad \alpha/\mathsf D^2=\frac1t\frac\alpha{s^3b^2},
 \qquad A_N/\varepsilon=\frac{\varepsilon R^2s}{K^2t}
                               \frac{A_N}{bs}
 \tag{RU12}
\]
are bounded without losing the zero-$n$ face.

\paragraph{Exact full scalar and both first normals.}
Put $H=\alpha+1/2$, $V_0=1-puv\omega$. At $N=M_B=0$,
$uW=\mathsf C/E+u^2c\mathsf BV_0$ and
$Y=\mathsf C/E-c\mathsf BV_0$. The complete scalar is
\[
 G_0=f_A\mathsf A^2+f_{BB}\mathsf B^2+
       f_{BC}\mathsf B\mathsf C+f_{CC}\mathsf C^2+
       \varrho\mathsf D^2,                                      \tag{RU13}
\]
with the bounded coefficients
\[
\begin{aligned}
 f_A={}&D_u^2D_vF_0/\varepsilon^2,\qquad
 \varrho=\alpha D_u/\mathsf D^2,\\
 f_{BB}={}&[Cc^2D_u^2D_v+(\beta_Dq+q^2/2)D_u^2+HD_uv^2c^2
             -A_hD_uD_v(2u^2+hu^4)c^2]V_0^2\\
          &+2HD_ucV_0-D_u/2-D_u^2\mu_2v^2\omega^2,\\
 f_{BC}={}&[2A_hD_uD_v(lu^2-1)cV_0-2HD_uv^2cV_0
                   -2HD_u+2n_1D_uD_v]/E,\\
 f_{CC}={}&[A_h(2-h)D_uD_v+HD_uv^2+D_uD_vn_2v^2]/E^2.
\end{aligned}                                                     \tag{RU14}
\]
For the full normal expansion
$G=G_0+L_NN+L_MM_B+G_{\rm hard}$, the exact coefficients are
\[
 L_N=R_A\mathsf A+R_B\mathsf B+R_C\mathsf C,
 \qquad L_M=S_B\mathsf B+S_C\mathsf C,                           \tag{RU15}
\]
where
\[
 R_A=D_uA_N/\varepsilon,\qquad R_C=2cuv\{A_hD_v(1-lu^2)+Hv^2\}/E,\qquad S_C=2\{A_hD_v(lu^2-1)-Hv^2\}/E,
\]
\[
 R_B=\Delta_{v,0}(L_N-R_A\mathsf A-R_C\mathsf C),\qquad
 S_B=\Delta_{v,0}(L_M-S_C\mathsf C).\tag{RU16}
\]
Here $L_N,L_M$ are the complete coefficient extractions from the bounded
full parent in the stated state coordinates. The displayed source
cancellations make both quotient numerators vanish identically at
$v=0$. Thus these are finite monic polynomial divisions, with the pure
testing denominator $D_u^2$ cleared first. They give all normal
coefficients, including the higher testing terms. Their coefficients
are bounded; division shifts testing degree by one and preserves the
weighted coefficient estimates used below.

At $u=v=0$ the complete hard matrix on $(N,M_B)$ is exactly
$\operatorname{diag}(1/2,C+1/2)$.

\paragraph{Bounded monic conductor, including the small-$\varepsilon$ end.}
The identities (RI14)--(RI15) apply to (RU10) without a source-shape
assumption. Their crucial inputs are checked here on the new complete
coefficients. After finite source unit denominators and $D_u^2$ are
cleared, every $\omega$ occurrence of (RU14) is through $uv\omega$
or $v^2\omega^2$, and all remaining testing powers are even in $u,v$.
The scalar raw support is exactly the five terms in (RU13).

For completeness, first replace $\mathsf A=\mathsf y+k\mathsf C$,
$\mathsf B=\mathsf x+\mathsf C$, and
$\mathsf D=k\mathsf C-\mathsf h$. The coefficient requiring a
first conductor is $S=k^2f_A+f_{BB}+f_{BC}+f_{CC}+k^2\varrho$.
Divide its contact-zero difference monically in order $u,\omega,v$.
At $u=k=\varepsilon\jmath$, the $\omega$ divided difference is in
$(\varepsilon v,v^2)$; write it as $\varepsilon U+vV$ with
$U,V\in(v)$. The last $v$ divided difference is in
$(\varepsilon,v)$ because, modulo $\varepsilon$, the remaining
$v$ dependence is even. The first $u$ divided difference restricts
to zero at $\varepsilon=u=v=0$.

These facts permit the exact bounded lift (RI15), which uses the
identity for $\mathsf Cv(\omega-\jmath)$ and never divides by
$\varepsilon$. Subtract the gradient multiple
$2k\varrho_c(-2k,1,1,-1)$. The first three contact gradients have
determinant $-\varepsilon^3\jmath^4E^2\ne0$; the CC5 adjugate
argument cancels a zero identity, and introduces no inverse of that
determinant. Each second-stage coefficient at
$\varepsilon=u=v=0$ is independent of $\omega$: the $\mathsf x$
one restricts to $2f_{BB}+f_{BC}$, and the other three to zero.
Hence every second $\omega$ divided difference is in
$(\varepsilon,v)$ and has the same bounded lift. The final repair is
\[
 k\mathsf C(\mathsf y+\mathsf z-\mathsf h-2k\mathsf x)
      =Q_1^2-\mathsf y\mathsf z.
 \tag{RU17}
\]
The pairings $Q_i$ are defined below. For each complete normal in
(RU15), one conductor suffices: after replacing $\mathsf A,\mathsf B$,
its $\mathsf C$ coefficient restricted at
$\varepsilon=u=v=0$ is independent of $\omega$, directly from
(RU16). This is precisely the normal criterion in the same proof.
When $b$ is a unit, $\varepsilon^{-1}$ is bounded and the ordinary
CC1--CC6 conductor is also available.

All contact value and first-jet hypotheses follow from the stationary
original parent and the four exact contact-zero positive square terms (RU5).
The algorithm is a bounded, source-linear map on the stated raw
coefficient class and preserves the filtrations used next. In
particular it preserves the source ideal $(h,g)$; no separate
contact-zero assertion for its individual coefficient summands is
needed. The source substitution in each monic step uses the actual
$E$, which is a bounded unit through $h=0$.

\paragraph{The complete filtered leading model.}
Set
\[
 \Lambda=\frac1{t(1+s)},\qquad \mu=\frac1{1+s},\qquad
 T_4=\mathsf A^2+(u\mathsf B)^2+(k\mathsf C)^2+
                \mathsf D^2-4k^2\mathsf B\mathsf C.
\]
The model scalar is
\[
\begin{aligned}
 G_{\rm mod}
 &=\Lambda\mathsf x^2+T_4/2-(1/t+1/2)k^2\mathsf x^2
                      +(k^2-u^2)\mathsf C\mathsf x/t\\
 &=\left[\Lambda+\frac{u(u+k)}t-(1/t+1/2)k^2\right]\mathsf x^2
                  +T_4/2-\frac{u+k}t\mathsf x\mathsf z.
\end{aligned}                                                     \tag{RU18}
\]
The second equality is the exact conductor identity
$\mathsf C(u-k)=\mathsf z-u\mathsf x$. Use the first expression
for the raw coefficient filtration, since it retains the even-$u$
class. Its angular raw coefficients are
\[
\begin{aligned}
 F_{BB}&=\Lambda+u^2/2-(1/t+1/2)k^2,\\
 F_{BC}&=-2\Lambda-u^2/t+(3/t-1)k^2,\\
 F_{CC}&=\Lambda+u^2/t-2k^2/t.
\end{aligned}                                                     \tag{RU19}
\]
At $h=g=0$, exact source collection gives the following two
coefficientwise filtrations:
\[
\begin{array}{c|c|c}
 &b\text{ a unit}&b\text{ infinitesimal}\\ \hline
 \text{variable weights}&(k,u,\omega,v)=(1,1,1,2)
     &(\varepsilon,\jmath,u,\omega,v)=(1,1,2,1,3)\\
 \operatorname{wt}(\mathsf x)&2&3\\
 \operatorname{wt}(\mathsf y,\mathsf z,\mathsf h,Q_1,Q_2)&3&5\\
 f_A-1/2,\ \varrho-1/2&\ge2&\ge2\\
 f_{BB}-F_{BB},\ f_{BC}-F_{BC},\ f_{CC}-F_{CC}&\ge4&\ge6\\
 S_B-\mu,\ S_C+\mu&\ge2&\ge4\\
 (R_A,R_B,R_C)&\ge(2,3,3)&\ge(3,5,5).
\end{array}                                                       \tag{RU20}
\]
These statements concern the chosen raw coefficients separately;
polynomial cancellation between raw products is not used to infer
them. Consequently
\[
\begin{array}{c|c|c}
 &b\text{ a unit}&b\text{ infinitesimal}\\ \hline
 G_0-G_{\rm mod}&O_{\rm wt}(8)&O_{\rm wt}(12)\\
 L_M-\mu\mathsf x&O_{\rm wt}(4)&O_{\rm wt}(7)\\
 L_N&O_{\rm wt}(5)&O_{\rm wt}(8).
\end{array}                                                       \tag{RU21}
\]
Here each equality is modulo the source ideal $(h,g)$.

To reproduce the collection, the necessary $b$-infinitesimal source
jets at $h=g=0$ are
\[
\begin{aligned}
 C&=t/2-(t+1)s/2+(t+1)s^2/2+O_{\deg(s,b)}(3),\\
 A_h&=t(1-s+s^2)/2+O_{\deg(s,b)}(3),\\
 \beta_D&=[1-s+s^2-(t+1)sb]/(2t)+O_{\deg(s,b)}(3),\\
 n_1&=(s-s^2-tsb)/2+O_{\deg(s,b)}(3).
\end{aligned}                                                     \tag{RU22}
\]
Also $c=(1-k^2)/t+O_{\rm wt}(6)$,
$E=t(1+k^2)$ at $h=0$, and
$s=t\jmath^2+O_{\rm wt}(6)$. For the unit-$b$ filtration use
$s=k^2/b$ and the respective coefficients of $s$:
$-(t+1)/2,-t/2,-[1+(t+1)b]/(2t),(1-bt)/2$.
The source cancellations (RU11)--(RU12) supply the first and last
rows of (RU20). Scalar weights are even in both
charts. The $M_B$ weights are even in the first and odd in the
second, giving exactly the next weights in (RU21).

Apply the bounded conductor to the exact differences in (RU21).
Their coefficientwise raw filtration makes every correction Gram
coefficient have strictly positive weight. Terms in $(h,g)$ lift
boundedly and remain infinitesimal at arbitrary relative rates.
The whole hard block differs infinitesimally from
$\operatorname{diag}(1/2,(t+1)/2)$, since all its coefficients are
bounded and $u,v,h,g,s$ are infinitesimal.

On the negative-core contact chart take $u,v,\omega>0$. The separated
nonnegative monomials
\[
 a_i^2=(\mathsf A,u\mathsf B,k\mathsf C,\mathsf D),\qquad
 Q_1=a_1a_2-a_3a_4,\quad Q_2=a_1a_3-a_2a_4
\]
have the exact product relation yielding
\[
 T_4=\tfrac12\{(\mathsf y-\mathsf z)^2+\mathsf h^2+
                 \mathsf y^2+(\mathsf z+\mathsf h)^2\}
                  +Q_1^2+Q_2^2.                                \tag{RU23}
\]
Its three-dimensional difference matrix has determinant $1/2>0$.
Each radical separates into source and testing factors.
The full eight-root residue on
$(\mathsf x,\mathsf y,\mathsf z,\mathsf h,Q_1,Q_2,N,M_B)$
therefore consists of half this strict circuit, the $N$ pivot $1/2$,
and the complete coupled block
\[
 \begin{pmatrix}1/t&1/2\\1/2&(t+1)/2\end{pmatrix},\qquad
 \det=\frac{t+2}{4t}>0.                          \tag{RU24}
\]
The additional cross term in the second line of (RU18) has
infinitesimal coefficient. WA10 now gives a finite separated SOS
for the exact remainder. Restore all four actual squares (RU5).
This proves the whole local contact chart in both $b$ alternatives.

\paragraph{The $\omega$-unit state chart.}
If $u,v$ are infinitesimal, $Z=1$, and $\omega$ is a positive unit,
replace the retained centered $n$ positive square term by $nD_uD_vX^2$ itself.
The complete remainder is (RU7) with the two $n_1,n_2$ terms removed.
Use the frame $(A_o,\mathsf x,\mathsf D,N,M_B)$, $A_o=\omega v$.
The exact identities
$\mathsf A=\varepsilon A_o$, $\mathsf B=A_o/\omega$,
$\mathsf C=A_o/\omega-\mathsf x$ use only a pure-testing unit
inverse. In particular the first scalar coefficient becomes
$D_u^2D_v(p^2/2+\mu_2)A_o^2$, and the first $N$ coefficient
becomes $D_u(p-2\mu_1)A_o$; the printed $\varepsilon$ denominators
cancel completely. The full residue has the $A_o,\mathsf D,N$
pivots $1/2,1/2,1/2$ and the same block (RU24).
The complete normal residues are $0,\mathsf x$. Hence this chart
also has a strict finite separated Gram.

\paragraph{A remote cap lemma and complete bounded parents.}
Away from the contact chart there is no need for centered excesses.
Retain the original nonnegative $n,\alpha$ positive square terms and the positive
excesses of $m,A$ above chosen positive real caps $m_c,A_c$.
If $D$ is unbounded, retain its original excess above a real cap
$D_c$ too; otherwise leave $D$ at its actual bounded value.
The remainder is the full original parent with these three
coefficients capped, and all its coefficients are bounded.

We use the elementary finite-dimensional fact that a real quadratic
form positive definite on the common kernel of finitely many real
linear rows becomes positive definite after sufficiently large
positive real multiples of the squared rows are added. Choose
coordinates on the common kernel and its orthogonal complement;
the assertion follows by a Schur complement. Unbounded actual
$m,A,D$, when used, exceed each chosen real cap. Thus every cap
is an exact nonnegative positive square term subtraction in the original parent,
and a strict residue gives WA10 for the entire bounded remainder.
The real caps depend only on real testing residues, never on a
source/testing rate comparison.

\paragraph{Finite remote $v$ with nonzero residue.}
Write $U=u^2,V=v^2$ for real testing residues, with $V>0$.
On the common kernel $Y=M=0$, $X=vW-uZ$, and the remaining
two-state form is exactly
\[
 CD_vU W^2+(D+1/2)(uD_vW-vD_uZ)^2
                  +VD_vW^2/2-D_uVZ^2/2.                         \tag{RU25}
\]
Its determinant is
\[
 VD_uD_v\{DV/2+U[CD-(D+1/2)/2]+C(D+1/2)U^2\}.
 \tag{RU26}
\]
For $t_0>1$, insert $C=t/2,D=1/[2(t-1)]$; it equals
$D_uD_v[V^2+t^2U^2V]/[4(t-1)]>0$.
Its $W^2$ coefficient is positive since $V>0$.
For $t_0=1$, $C=1/2$ and one may choose a real $D$ cap: the braces
are $D(V+U^2)/2+(U^2-U)/4$, positive for sufficiently large $D$.
The cap lemma then chooses $m_c,A_c$ for a strict full state matrix.

\paragraph{Finite unit $u$ with infinitesimal $v$, including zero.}
Use $N=-M=X+uZ-v(W+uY)$, $R_t=uW-Y$, and $V_Z=vZ$.
After the original $n,\alpha$ positive square terms and allowed positive excesses
are retained, the entire remainder lies in the square of the four
roots $(N,R_t,Y,V_Z)$. Only the pure-testing unit $u$ is inverted
to write $W=(R_t+Y)/u$; no inverse of $v$ is needed. Its exact
residue is
\[
\begin{aligned}
 (m_c+1/2)N^2+CR_t^2+A_cD_uY^2
 +(D+1/2)(R_t-D_uV_Z)^2-D_uYV_Z-D_uV_Z^2/2.
\end{aligned}                                                     \tag{RU27}
\]
For $t_0>1$, its restriction to $Y=0$ has determinant
$t^2D_u(D_u-1)/[4(t-1)]>0$ on $(R_t,V_Z)$.
Choose $A_c$ sufficiently large. For $t_0=1$, the common kernel
$Y=0,R_t-D_uV_Z=0$ has value $D_uu^2V_Z^2/2>0$;
large real $D_c,A_c$ give strictness. The normal pivot is positive.
The full four-root support and the residue are checked directly;
literal $v=0$ is their polynomial value.

\paragraph{Both finite angles infinitesimal, and all state charts.}
The exact SD5 pure-core sign identity deals with the nonnegative-core
branch and literal zero axes by original-event squares. On the
negative branch use its standard gauges $u,v>0$ and normalize the
state by a maximal coordinate. If $Z$ is infinitesimal, $W$ is a
unit and the exact BE4 selected four-root identity applies to
$(X,Y,uW,vW)$, with positive caps of $m,A$ provided by (RU3).
All other original events remain as nonnegative positive square terms.
If $Z$ is a unit normalize $Z=1$; the same sign identity gives
bounded positive $\omega$. Its infinitesimal and unit alternatives
are exactly (RU18)--(RU24) and the subsequent $\omega$-unit chart.

\paragraph{Every reciprocal chart and literal projective infinity.}
For unbounded $u$, put $p_u=1/u$ and use the entire homogeneous
polynomial $p_u^2P$. If bounded $v$ has nonzero residue, the full
capped four-state residue is
\[
 CD_vW^2+A_cD_vY^2+(D_c+1/2)(W+vX)^2
                          +(m_c+1/2)(vY-Z)^2\succ0.             \tag{RU28}
\]
Here $D_c$ denotes either a positive real cap or the actual positive
unit residue. The four square rows have common kernel zero because
$v$ is a unit. If instead $v$ is infinitesimal, including zero,
use the exact BE $(\infty,0)$ selected full frame
$(p_uX+Z,W,Y,vX)$ with positive caps of $C,A$. Its strict residue
is the sum of positive multiples of
$(p_uX+Z)^2,(W+vX)^2,W^2,Y^2$; the other original positive square terms are
retained. This handles the zero-$n$ degeneration at that corner.

For unbounded $v$ and bounded $u$, put $q_v=1/v$. The full capped
residue of $q_v^2P$ is
\[
 m_c(W+uY)^2+C(uW-Y)^2+(A_c+1/2)D_uY^2
                         +D_c(uX-Z)^2+D_uX^2/2\succ0.           \tag{RU29}
\]
Its common square kernel is zero for every bounded real $u$,
including zero. Finally, if both angles are unbounded, use
$p_u^2q_v^2P$ and the full frame
$(X,W,Y,P_Z=p_uZ,Q_Z=q_vZ)$. The exact residue is
\[
 CW^2+A_cY^2+D_c(X-P_Z)^2+X^2/2
                           +(m_c+1/2)(Y-Q_Z)^2\succ0.           \tag{RU30}
\]
The two nontrivial block determinants are $D_c/2$ and
$A_c(m_c+1/2)$, both positive. The entire polynomial has this
five-root support, with no inverse of $p_u$ or $q_v$ after
homogenization. All three formulas are checked from the original
parent. Their literal reciprocal zeros are polynomial values;
restoring a nonzero homogeneous factor uses a pure-testing square.

The charts exhaust independent testing pairs and states. Apply each
real identity to real and imaginary parts of a complex state, restore
every retained actual positive square term and separated clearing factor, and use
the exact sign gauges for both angle signs. This completes the stated
all-tests theorem without a restriction on any opening ratio.

\section{Two small actual angles: the R8 and R10 sectors}
\subsection{The opening $D$ gap}
\label{component:pe07041e033da}
\label{pe07041e033da:sec:r17-R8-joint-D}
Assume $t,g$ units, $h,j$ infinitesimal, and $a$ unbounded. Then
$e$ is infinitesimal, $r\sim1$, and $t\sim1$. Use the generic BA
identity above and the exact D identity in
\eqref{eq:c2-ba-D-circuit}. The following normalization and residue
formula hold before the thick/thin source alternatives are used.
Put $L_D=lE/t-rR_D-\sqrt hN_y+uv\xi$ and $M_D=\sqrt D L_D$.
Use the polynomial inverse
\[
 E={t\over l}\{D^{-1/2}M_D+rR_D+\sqrt hN_y-uv\xi\}. \tag{JD5}\label{eq:c2-jd-inverse}
\]
On R8, $t/l$ is a unit, $D^{-1/2}$ is infinitesimal, and $Dj$ is
bounded.  Every coefficient of the new remainder is bounded whenever
the remaining source arrays of R10.9 are bounded.  In particular no
ratio $h/j$, $x/j$, or testing angle divided by a source opening occurs.
The full divergent rank-one block has been normalized jointly.

\paragraph{The exact residue after the joint normalization.}
Put $\sigma=t\lambda x/\sqrt q$ and $\mu=m-a^2q$ as in R10.9.
At the R8 source/testing corner, \eqref{eq:c2-ba-D-circuit} and \eqref{eq:c2-jd-inverse} have the complete residue
\[
 \begin{split}
 M_D^2+r_0^2(C_0+n_0)R_D^2+(Ah)_0N_y^2
  +(n_0+\tfrac12-\mu_0)\xi^2
  +(c_F)_0F_v^2+\lambda_0J^2+R^2+Q^2-2\sigma_0Q\xi.
 \end{split}                                                   \tag{JD6}
\]
This formula applies when $r,n$ are bounded and $r$ a unit.  Its
derivation can retain $n,C,Ah,\mu,c_F,\lambda,\sigma$ as free source
letters: specialize $u=v=\sqrt h=D^{-1/2}=x=\alpha=0$ and $t=1$
in \eqref{eq:c2-ba-D-circuit} and \eqref{eq:c2-jd-inverse}.  In particular the $E,R_D,N_y$ terms become
$r^2(C+n)R_D^2+AhN_y^2$, with no missing $D$ contribution.
The separate $N^2,T^2$ remain intact.

\paragraph{Unbounded $a$: the thick layer.}
Put $s=1/a$, $w=ae$, $b=a^2B$.  The exact R10 source identity
still reads
\[
 n=s^2{wQ_*-bK_*\over2r\Delta_*},\qquad
 Q_*\sim1,\ K_*\sim t,\ \Delta_*\sim1.             \tag{JD11}
\]
Thus $b/w$ is bounded, $b$ is bounded, $g_0=1-w_0>0$, and
\[
\begin{gathered}
 (m/w)_0={1\over2(1-w_0)},\quad C_0=(Ah)_0=1/2,
 \quad (Dj)_0=1/2,\\ \lambda_0=1/2,
 \quad(c_F)_0=(1-w_0+b_0)/2>0.
\end{gathered}                       \tag{JD12}
\]
The source formulas for these quantities have no $j$ denominator;
only $D$ itself is unbounded, and \eqref{eq:c2-ba-D-circuit} removes that pole.
For $\rho=n/e^2$ positive noninfinitesimal, retain the exact R10.10
identities and R10.11 Schur pivot:
\[
 \bar q=a^2q=w^2\{\rho r^2+H\}+a^2\alpha,
 \quad H=m+\tfrac12,\quad \chi=m/w,
 \quad \chi-H\in\operatorname{Inf},
\]
\[
 S=\tfrac12-\mu-{(t\lambda x)^2\over q}
 =H-{m^2\over\bar q}
 \ge {H\rho r^2\over\rho r^2+H}
       +{H^2-\chi^2\over\rho r^2+H}.                 \tag{JD13}
\]
The first term has positive residue and the second is infinitesimal,
also for unbounded $\rho$.  Moreover $x/\sqrt q=b/\sqrt{\bar q}$
is bounded by JD11 and $\bar q\ge Hw^2$.  JD6 consequently has
strict residue on $(M_D,R_D,N_y,\xi,F_v,J,R,Q)$, with the only
possible nonzero cross residue in its strictly positive $Q,\xi$
block.  Restore $N^2,T^2$.  This proves the full local parent, not
just its contact-state restriction.

\paragraph{Unbounded $a$: the thin layer and the exact paired-$D$ repair.}
Assume $0<\rho=n/e^2$ infinitesimal.  JD11 gives
$b/w\sim1/t$, so $\psi=w/b$ is a unit.  The formal C1 contact is
exactly OP4, with native parameters
\[
 x_1=x,\quad\eta_1=t,\quad\zeta_1=s,\quad r_1=h,
 \quad\theta_1=e,\quad t_1=w.
\]
TL2 still gives $\beta_1=-s_n/2$ infinitesimal and
$\kappa=(1+2n)/(1-s_n)\sim1$.  The native gaps satisfy
$M_1=g$ a unit and $J_1=j>0$.  The canceled CA/CS/SE remainders
have no inverse $J_1$; their divergent paired-$D$ mass is retained
in SD3.  The formal source has $C_1,A_1h$ positive units,
$D_1j$ a positive unit, positive $m_1,\alpha_1$, and positive
$\beta_1+1/2$.  These facts follow directly from TL3 and JD12.
Its bounded opening ratio $t_1/e_{1,\mathrm{scale}}^2=w/(tb)$
is a unit.  Thus the explicit CA/CS/SE strict base and their finite
conductors apply exactly as in the R16 thin proof; the reservation
$(2+4\beta_1)\nu-\nu$ is positive because $\beta_1$ is
infinitesimal.  This is an algebraic use of those formulas, not an
application of a physical C1 theorem with negative beta.

Here is the one change needed to the R16 perturbation.  In native
notation put
\[
 T_1=c_a=1+xr,\quad a_1=l/t,\quad
 b_D={J_1\over T_1}={j\over c_a},\quad
 h_1={a_1\over T_1}=1-b_D,
\]
\[
 q_B=v(Z-uX),\quad M_B=(uW-Y)-h_1q_B,
 \quad R_{D,\mathrm{event}}=M_B-b_Dq_B.               \tag{JD14}
\]
The old retained paired-$D$ positive square term in $\kappa F_1/4$ is
$\kappa(D_1+1/2)h_1M_B^2$.  In the actual parent retain instead
the positive contact-zero positive square term $(D_2+1/2)h_1M_B^2$.
For $\delta_D=(D_2+1/2)-\kappa(D_1+1/2)$ the entire resulting
difference is
\[
 \boxed{\quad
 \delta_D\{R_{D,\mathrm{event}}^2-h_1M_B^2\}
 =\eta_D\{M_B^2-2M_Bq_B+b_Dq_B^2\},\quad
 \eta_D=\delta_D b_D.\quad}                          \tag{JD15}
\]
By the exact TL3 coefficients, with $\Theta=n/B$,
\[
 \eta_D={\Theta\{S_cP_D/t-jP_\kappa\}
                    \over K(h+t)c_a},\qquad
       \eta_D/\rho=O(w).                            \tag{JD16}
\]
Every denominator is a unit.  In particular the previously unbounded
$\rho w/j$ correction is absent; no relative $w/j$ bound is needed.

Keep the R16 positive contact-zero $n$ row TL4, the actual centered
$A$ positive square term, and every other unchanged retained positive square term.  The exact
full stress is TL6 with its $\delta_D R_{D,\mathrm{event}}^2$
line replaced by the right side of JD15.  Both parents and all
retained positive square terms have zero contact value and full first jet, so this
new entire stress has the same exact vanishing.

For completeness its new term satisfies every finite-support premise
of the inherited conductor.  The exact inverse state formulas give
\[
 q_B=v\{D_u-uN-puD_uv\mathfrak w\},\qquad p=t_1/R_1.
                                                               \tag{JD17}
\]
At zero normals $q_B=vD_uV_0$.  Hence the new scalar term is in
$(v)^2$, both complete first normals are in $(v)$, and the remaining
hard entries are bounded.  There are no terms above normal degree
two.  Every dependence on $\mathfrak w$ has a factor $p$.
In the CA chart JD16 divided by $\rho$ has bounded coefficients
on the same raw roots as TL8.  In the SE chart
\[
 \Theta/\rho=\varepsilon^2\tau^2QE/K
\]
gives the same extra $\varepsilon^2$ factor to every new coefficient.
Modulo $\varepsilon$, scalar coefficients are even in $u,v$ and
the normal coefficients are independent of $\mathfrak w$.
Thus the exact support, full-jet, and finite divisibility hypotheses
of CC3--CC6 and SE5--SE8 are unchanged.  The complete stress divided
by $\rho$ has a bounded Gram on the same eight roots.  Multiplication
by $\rho$ makes it an infinitesimal perturbation of the strict formal
C1 base.  Restoring the positive square terms proves the actual C2 parent
on the entire thin contact chart, including both full state normals.

At literal $n=0$, OP4 is a physical C1 identification with
$k_1,\zeta_1,r_1$ infinitesimal, $\eta_1=t$ a unit and
$e_{1,\mathrm{scale}}^2=tb$ bounded.  The physical C1.2/C1.3
all-tests theorems allow arbitrary positive $J_1$ and literal beta
zero, so they cover this boundary exactly.

\paragraph{The unit-$W$ entry on the unbounded-$a$ thin layer.}
The actual physical C2 critical coefficient is
\[
 C_*={(h+2Ah)(1+2D)\over2(4AhD-h)}.
\]
Use the polynomial UW3 numerator/denominator with $b=wz_*$,
$z_*=b/w$ bounded.  Its numerator is divisible by $sw$ because it
vanishes identically at $s=0$ and $w=0$.  The denominator after
canceling the $m$ numerator has residue $16t_0=16$, also when
$j_0=0$.  Consequently
\[
                   (C_*-C)/m=s\times O(1).            \tag{JD18}
\]
There is no $j$ inverse in this estimate.  SEO12 permits both $A,D$
unbounded: its entire hard block remains strict and its two radial
masses tend to $2,2$.  With $W=1$, the exact testing identity
$vZ=Z(R_m+X+uZ)/(1+uY)$ absorbs the full negative critical correction
by JD18 and the retained original $m$ square.  All denominators are
pure testing units.  This proves every local normalized unit-$W$
state, including $Z=0$; negative-core state normalization leaves
exactly the $Z=1$, $u,v,\mathfrak w>0$ infinitesimal chart proved
above.  This is the same actual-C2 state entry as UW, with its
source bound checked at the new $t_0=1$ boundary.

\paragraph{Finite remote tests when $e$ is infinitesimal.}
Here $C_*-C$ is infinitesimal by JD18 and bounded $m$.
The full critical-wall identities SEO17, SEO18, SEO20, in the
normalized SEO12 frame, are strict whenever at least one finite
angle has nonzero residue.  They explicitly allow $A,D$ unbounded.
The entire correction $-4(C_*-C)_+D_v(uW-Y)^2$ has infinitesimal
bounded entries in that frame.  Retain original $n,m,\alpha$ and
any positive excess.  This proves every finite remote test, with
all states, for the entire $e$-infinitesimal R8 domain.

\paragraph{Projective endpoints, signs, and the complete join.}
On all of R8, $C$ is a positive unit and $A,D$ are unbounded.
For infinity charts retain the whole original $m,n,\alpha$ positive square terms,
clip $C,A,D$ at suitable positive real lower bounds, and retain their
positive excesses.  The exact homogeneous infinity identities used
in SEO19 and its accompanying BE2 routing require only these lower
bounds and $d=1$; they do not require $D$ bounded or $m$ bounded
below.  In particular the delicate $(u,1/v)=(0,0)$ corner has the
strict selected residue
\[
 2X^2+4c(uW-Y)^2+(4A_*+2)Y^2+4D_*Z^2
\]
on $(X,Y,uW,Z)$ after retaining $m$.  All leftover $W$ occurrences
carry $u$.  The remaining single and double infinities are exactly
the same homogeneous event identities.  Literal endpoint values are
taken in those original homogeneous formulas; no higher-degree
contact graph is substituted at infinity.

Here $a$ is unbounded, so $e$ is infinitesimal and $r\sim1$.
The thick, positive-thin, and literal-zero alternatives for $n/e^2$
are exhaustive. The finite outer calculation above uses the exact
coefficient repair JD18 and the strict forms SEO17, SEO18, SEO20 in
Section~\ref{component:p05a86be0046c}. The displayed homogeneous
forms prove every infinite test. The sign gauges, exact axis identities,
state homogenization and real/imaginary splitting therefore give all
independent tests and states on this unbounded-$a$ R8 domain.
The bounded-$a$ part of R8 is Proposition~\ref{prop:c2-bounded-a}.

\subsection{Source scale, remote tests, and weighted local completion}
\label{component:p9ee4ba7326a7}

Use the native coefficient dictionary, the exact full-parent identity
BA4--BA7, the homogeneous identities BA11--BA15, and the physical
source identity \eqref{eq:c2-ba-physical}. The following bounds are
derived on the present unbounded-$a$ coefficient domain.

\paragraph{Scope and result.}
Retain the original weak-C2 R10 physical stationary source, every strict
source gap, $n\ge0$, the positive-infinitesimal total-scale intersection,
and every inherited exclusion. Thus $t,g,j$ are positive bounded units,
$B,h,e$ are infinitesimal, $a$ is unbounded, and $x=aB$ is infinitesimal.
The new all-tests case proved here is the source subdomain
\[
                         n/e^2\notin\operatorname{Inf}.
\tag{R10.1}
\]
Every nonlocal independent testing chart is also proved on all of R10.
The positive-$n$, infinitesimal-$n/e^2$ local complement is handled by the following thin-layer and state-entry sections; it is not covered by this weighted pivot.
The literal $n=0$ boundary is already an inherited C1 case, by the
exact map specified below. These statements retain all independent rates.

\paragraph{The decisive exact source scale.}
Use the pure-source variables
\[
 s=1/a,\qquad w=ae,\qquad b=a^2B=x/s,
 \quad B=s^2b,\quad x=sb,\quad e=sw,\quad r=1-sw.
\]
Initially $b$ is not assumed bounded. The gap $g=1-w+hx$ makes $w$
bounded and $1-w$ a positive unit in residue. In the exact positive
polynomials $Q_s,K_s$ in \eqref{eq:c2-ba-physical} put
\[
 Q_*=1+B(1+t+ah+2ar+at)
 +B^2\{a(h+t)+a^2[h(r+1)+t(1+r-r^2)]\},
\]
\[
 K_*=(1+s)(h+t)(1+x),\qquad \Delta_*=s\Delta.
\]
Before imposing a bound on $b$, $Q_*\sim1$ and $\Delta_*\sim1$:
write them using bounded $w$, infinitesimal $s,x,h$; their only
nonconstant monomials contain $s$ or $x$. Equation~\eqref{eq:c2-ba-physical} gives exactly
\[
 \boxed{\quad 2r\Delta_*\,n/s^2=wQ_*-bK_* .\quad}
\tag{R10.2}
\]
Consequently
\[
 0<b/w\le Q_*/K_*,\qquad
 0\le n/s^2\le wQ_*/(2r\Delta_*).
\tag{R10.3}
\]
In particular $b=O(w)$ is bounded, $n=O(s^2w)$, $an$ is
infinitesimal, and $ck=\sqrt{b/t}$ is bounded for
$c=\sqrt{a/t}$, $k=\sqrt{x}$. This is an original-source bound,
including equality $n=0$, not a comparison with either testing angle.

After this bound every displayed source denominator is a unit in
$(s,b,w,h,t)$ except the explicit positive monomials already canceled.
Direct finite rational cancellation in the source formulas gives
\[
 \begin{gathered}
 m/w\sim\frac1{2(1-w)},\quad
 C\sim Ah\sim t/2,\quad D\sim\frac1{2(t-1)},\quad
 \alpha/(s^3b^2)\sim t/2,\\
 \lambda=\frac{m-a^2q}{ta^2B}\sim\frac12,\qquad
 c_F\sim\frac{t-w+tb}{2t^2}>0,\qquad
 q=nr^2+(m+\tfrac12)e^2+\alpha.
 \end{gathered}
\tag{R10.4}
\]
The last strictness follows from $t_0>1$ and $w_0<1$.
These are joint residue identities in bounded letters; they do not
set $w$ or $b$ to a positive unit. In particular they hold when either
is infinitesimal. The cancellation defining $\lambda$ leaves no $b$
or $s$ denominator; the new checker verifies this explicitly.
We also have $m$ bounded, $n,\alpha,A^{-1/2}$ infinitesimal, and
\[
 d_4\sim\frac{t}{2(t-1)},\quad b_4\sim0,\quad
 a_E\sim-\frac1{t-1},\quad x_3\sim-1.
\tag{R10.5}
\]
All first-normal arrays printed in
Section~\ref{component:pd709027c7805} are bounded:
the only newly unbounded letter $a$ occurs there through $an$ or
$aq$, both infinitesimal. The factors $a^2n$ in $c_F$ are bounded.

\paragraph{All remote tests on all of R10.}
The exact full-parent identities printed with BA11--BA15 and the exact identities \eqref{eq:c2-literal-u-axis}--\eqref{eq:c2-literal-v-axis} apply algebraically. The
bounded-$a$ proof used $m$ infinitesimal; under R10.4 its entire
nonnegative event positive square term is retained. In BA11 this adds a nonnegative
rank-one form to the strict six-root residue displayed immediately
after that identity; all its coefficients are bounded. For the chart
$u$ infinitesimal, $v$ a bounded unit, keep $S=\sqrt m W$.
The complete $(S,X)$ residue is now
\[
 (v_0S-\sqrt{m_0}X)^2+\frac{1+v_0^2}{2}X^2,
 \qquad \det=\frac{v_0^2(1+v_0^2)}2>0.
\tag{R10.6}
\]
The $(uW,Z)$ block remains BA12 and the normalized $Y$ pivot remains
strict; the extra $uvY,uZ$ terms in the $m$ positive square term are infinitesimal.
For unbounded $v$ its reciprocal identity has the same diagonal
residue BA13 because the potential $\sqrt m X$ term carries $1/v$.
On the remaining remote charts retain $mR_m^2$ as a full positive
square and apply the unchanged strict residue to the remainder.
The original homogeneous identities include either literal infinity,
and the exact literal-axis identities \eqref{eq:c2-literal-u-axis}--\eqref{eq:c2-literal-v-axis} retain $n=0$ without a division. The original
sign gauges and real/imaginary split are unchanged. Hence every test
outside both angles infinitesimal is covered for the entire R10 source.

\paragraph{The exact local full-parent identity.}
Use the generic contact balances RC1, RC6 in
Section~\ref{component:pd709027c7805} and the complete identity
BA4--BA7. These are rational identities of the original source and
never require bounded $a$. Set $U=u^2,V=v^2$, $y=tB$,
\[
 E=tuW-vZ,\quad N_y=(Y-huW)/\sqrt h,\quad \xi=X+ruZ,
\]
\[
 F_v=(V-y)Z,\quad R_D=v(U-x)Z,\quad J=(uv-x/c)Z,
\]
\[
 Q=\sqrt q(avW-uZ-a\xi),\qquad
 N=\sqrt n\{r[(a/t)V-U]Z+u\xi\},
\]
\[
 R=\sqrt\lambda\,tk(cv-u)W,\qquad
 T=\sqrt{2\lambda/c}\,k\sqrt{uv}(tcW-Z),\quad
 \mu=m-a^2q=ta\lambda x.
\]
At positive nonzero angles these are exactly the state-linear roots
in BA4, with no testing-ratio inverse. The entire hard quadratic
$\mathcal H_0$ is the explicit expression BA6. It has no inverse
testing angle or $a$ denominator; after $\eta=\sqrt hN_y$, every
occurrence of $A$ is through $Ah$, $\sqrt hAh$, or $hAh$.
Its R10 local residue is
\[
 \operatorname{diag}\left(\frac1{2(t_0-1)},\frac{t_0}2,\frac12\right)
 \quad\hbox{on }(E,N_y,\xi).
\tag{R10.7}
\]
This permits bounded $m$; every occurrence of $m$ in the remaining
hard entries carries a positive testing degree or $h$.

The principal angular--$\xi$ coupling is resolved as a whole.
Writing $R_c=kv(tcW-Z)$, direct multiplication gives
\[
 -2t\lambda ck R_c\xi+2t\lambda c(R_D-uJ)\xi
 =-2t\lambda x Q\xi/\sqrt q-2\mu\xi^2.
\tag{R10.8}
\]
Thus neither $cu$ nor $cR_D$ is left as a Gram coefficient. The
complete generic identity BA7, with the source arrays just specified,
gives the entire original $F/4$ exactly as
\[
\begin{split}
 &Q^2+N^2+T^2+R^2+\lambda J^2+d_4R_D^2
 +(c_F+b_4x/t^2)F_v^2+(b_4/t^2)vR_DF_v\\
 &\quad+\mathcal H_0-\lambda xE^2-\mu\xi^2
                  -2(t\lambda x/\sqrt q)Q\xi\\
 &\quad+E\{(a_E+c_EV)R_D+(b_E+c_Ex)vF_v\}\\
 &\quad+\sqrt hN_y\{(a_Y+c_YV)R_D+(b_Y+c_Yx)vF_v\}\\
 &\quad+\xi\{x_5uvR_D+u(x_3+x_5x)F_v\}.
\end{split}
\tag{R10.9}
\]
All coefficients are the exact source arrays in
Section~\ref{component:pd709027c7805} and BA5.
This identity retains every scalar, first-normal and hard-normal term.

\paragraph{A strict weighted pivot when $n/e^2$ is noninfinitesimal.}
Put $\rho=n/e^2$, $H=m+1/2$, $\chi=m/w$, and
$\overline q=a^2q$. Then
\[
 \overline q=w^2(\rho r^2+H)+a^2\alpha,
 \qquad \chi-H\in\operatorname{Inf}.
\tag{R10.10}
\]
The $Q,\xi$ source block in R10.9 has Schur pivot
\[
 S=\frac12-\mu-\frac{(t\lambda x)^2}{q}
   =H-\frac{m^2}{\overline q}.
\tag{R10.11}
\]
The equality uses $a^2q+\mu=m$ exactly. Since $\alpha\ge0$,
\[
 S\ge H-\frac{\chi^2}{\rho r^2+H}
 =\frac{H\rho r^2}{\rho r^2+H}
    +\frac{H^2-\chi^2}{\rho r^2+H}.
\tag{R10.12}
\]
Here $H$ is a positive bounded unit. Under R10.1 the first summand
has positive residue and the second is infinitesimal, also when
$\rho$ is unbounded. Hence $S$ is a positive bounded unit.
Moreover $b/w$ is bounded and $\overline q\ge Hw^2$, so
$x/\sqrt q=b/\sqrt{\overline q}$ is bounded. All entries of this
strict source block are therefore bounded.

In the rest of R10.9 the only nonzero off-diagonal residue is
$a_E$ between $E$ and $R_D$, giving
\[
 \frac1{2(t_0-1)}
 \begin{pmatrix}1&-1\\-1&t_0\end{pmatrix}\succ0.
\]
The $N_y,F_v,J,R$ pivots have positive residues
$t_0/2,(t_0-w_0+t_0b_0)/(2t_0^2),1/2,1$.
Every other entry is infinitesimal by R10.4--R10.5 and the explicit
positive testing degrees in R10.9. Retain $N^2,T^2$ as exact positive
squares. The bounded-Gram lemma supplies finite tensor SOS for the
entire parent on R10.1. No source/testing rate comparison occurs.
All roots homogenize without $Z$ or testing-ratio inverses, so literal
axes and $Z=0$ are included directly. The original signs, complex-state
split and the remote theorem complete the all-tests case.

\paragraph{Literal $n=0$ and the exact remaining complement.}
At $n=0$, OP4 is a physical C1 identification with
\[
 x_1=x,\quad\eta_1=t,\quad\zeta_1=s,\quad r_1=h,
 \quad\theta_1=e,\quad e_{1,\mathrm{scale}}=\sqrt{tb}.
\]
Thus $k_1,\zeta_1,r_1$ are infinitesimal, $\eta_1$ a unit, and
$e_{1,\mathrm{scale}}$ bounded. Its infinitesimal type is covered by
the local calculation in Section~\ref{component:pa7e744614f49} together with the outer calculation in Section~\ref{component:p05a86be0046c};
its positive-unit type is covered by the neighboring
Section~\ref{component:pccaa793fd0bf}. Both retain $\beta=0$.
This is C1.2/C1.3 routing, not the unbounded-$\eta$ C1.1 theorem.
For $n>0$ OP5 gives formal $\beta_1<0$, so this physical C1 theorem
cannot be extended by continuity or imported into the neighborhood.
The complement left for the companion thin-layer proof is exactly
\[
       0<n/e^2\in\operatorname{Inf},\qquad u,v\in\operatorname{Inf},
\]
with every original physical and total-scale intersection retained.
In this complement R10.11 can have zero residue: at a unit $w$ and
$b_0=w_0/t_0$, its residue is zero. Thus R10.9 is a useful full-parent
identity there, but its stated strict-Gram argument does not prove that
complement. The companion thin-layer proof and direct state join complete this complement; the final join states the whole-R10 theorem.

\subsection{The positive-$n$ thin layer}
\label{component:pef37475d831e}

\paragraph{Domain.}
All polynomial identities below use the convention $F/4$, as in
R10.9; in particular the original $n$ positive square term is $nD_uD_vX^2$.
All parent and positive square term identities in this section use the common normalization $F/4$.
Keep every original R10 condition and assume
$\rho=n/e^2>0$ infinitesimal. Use $s,w,b,x,B$ from R10.2.
Dividing that exact identity by $w$ gives
\[
 Q_*-(b/w)K_*=2r\Delta_*\rho w,
 \qquad b/w\sim1/t.
\tag{TL1}
\]
Thus $\psi=w/b$ is a positive bounded unit. No testing rate occurs.
The formal native C1 map OP4 is
\[
 x_1=x,\quad\eta_1=t,\quad\zeta_1=s,\quad r_1=h,\quad
 \theta_1=e,\quad t_1=\theta_1/\zeta_1=w.
\]
Write $P_d=P_n/B$. OP3--OP5 give the exact slack
\[
 s_n=2r\Delta_*\rho w/P_d,\quad
 \beta_1=-s_n/2,\quad \pi=1-s_n=1+2\beta_1>0,\quad
 \kappa=(1+2n)/\pi.
\tag{TL2}
\]
Hence $s_n/w=O(\rho)$, $\beta_1$ is infinitesimal, and
$\kappa\sim1$. We substitute these coefficients into the C1 polynomial identities to obtain the comparison below.

\paragraph{Exact coefficient comparison.}
Let $F_1$ be the native C1 parent at this same contact, normalized by
$d_1=1$, and $F_2$ the actual C2 parent, with $d_2=1$.
For a coefficient $L$ put $\delta L=L_2-\kappa L_1$.
Set $K=1+x$, $c_a=1+xr$, $S_c=1+B(1-h)+xr$, $Q_c=1-hB$,
and $\Theta=n/B=\rho w\psi$. The following short polynomials are
bounded with the displayed unit denominators:
\[
 P_\kappa=c_a^2+sx(1+t)+sx^2(r+he+t),
\]
\[
 P_D=tx^2r^2+x\{t-1+e(1-h-2t)\}-1,\quad
 P_\alpha=1+s-2e+x\{r^2+s(r+he)\}.
\]
Exact cancellation of the original C2 and native C1 formulas gives
\[
\begin{array}{c|l}
 \delta m&0\\
 \delta C&-\Theta c_a/(1-h)\\
 \delta(Ah)&-\Theta t c_aS_c/[K(h+t)]\\
 \delta D&\Theta S_cP_D/[K(h+t)tj]\\
 \delta\alpha&-nxP_\alpha/(KQ_c)\\
 \delta d=1-\kappa&-2\Theta P_\kappa/[K(h+t)]\\
 \delta\beta+\delta d/2&-n.
\end{array}
\tag{TL3}
\]
In particular $\delta C,\delta D,\delta d,\delta(Ah)=O(\rho w)$,
while $\delta\alpha=O(nx)$. There is no surviving $1/h$ after the
large $A$ event is centered.

\paragraph{The formal C1 local identities still supply a strict base.}
The native $\eta_1$ is a unit, $k_1,\zeta_1,r_1$ are infinitesimal,
and $e_{1,\mathrm{scale}}^2=tb$ is bounded. Separate its positive-unit
and infinitesimal types. Use respectively the exact centered-$A$
CA5/CA8 construction (or its CS1 scaling) and the SE4/SE9--SE11
construction. This use concerns the printed formulas and conductor
lemmas; the physical-source theorem is not being invoked at $\beta_1<0$.
Here are all changes to their source premises:

The formulas and native stationary identities are rational identities,
valid at the formal OP4 contact. All native gaps remain positive; $M=g$
and $J=j$ are units. The other native event coefficients are positive,
with $C_1,D_1,A_1h$ positive units and $A_1$ unbounded. The opening
ratio $t_1/e_{1,\mathrm{scale}}^2=w/(tb)$ is a positive bounded unit
by TL1, which supplies the only scale bound previously obtained from
$\beta_1\ge0$. All four retained positive square terms SD3 involve the positive
paired coefficient $\beta_1+1/2=\pi/2$, never a separate
$\beta_1$ square. In SE4, the coefficient of the reserved normal positive square term
is $(2+4\beta_1)\nu$, so its subtraction of $\nu$ is still positive,
since $\beta_1$ is infinitesimal. The remaining source bounds, finite
coquotients and strict residue matrices are unchanged. A positive-unit
normalization to $C_1=2$, when used by those formulas, changes none of
these bounds. Consequently their exact local construction yields the
same strict bounded eight-root Gram and positive retained positive square terms for
$\kappa F_1$ on its contact chart.

\paragraph{A new positive contact-zero $n$ row.}
Use the inherited pure-testing coordinate and normals
\[
 \mathfrak w=(W+uY)/D_u,\quad A_t=v\mathfrak w,\quad
 X=N+pD_uA_t-uZ,\quad p=t_1/R_1,
\]
with $D_u=1+u^2,D_v=1+v^2$. In the exact CA chart the contact is
\[
 (u,\mathfrak w,v)=(k,\delta k,\delta k^2E),\quad
 K_c=\delta^2Ek^2=\zeta_1K/Q.
\]
Here $Q,E$ denote the inherited CA source letters, not the angular
weight of R10. The SE chart has $k=\varepsilon j$,
$\delta=1/\varepsilon$, so $K_c=j^2E$.
Put $\gamma=r/K_c$ and retain the entire state-linear square
\[
 \boxed{\quad \mathcal N^2=nD_uD_v(X+\gamma A_t)^2.\quad}
\tag{TL4}
\]
At the actual contact $A_t=K_ck$ and $X=-rk$, so its value and
all first derivatives vanish exactly. The square root of $D_uD_v$ is
pure testing; every other radical or inverse is pure source or pure
testing. The following entire-parent identity includes its full normal:
\[
\begin{split}
 nD_uD_vX^2-\mathcal N^2
 ={}&2n\gamma D_uD_vuA_tZ
 -(2n\gamma pD_u+n\gamma^2)D_uD_vA_t^2\\
 &-2n\gamma D_uD_vA_tN.
\end{split}
\tag{TL5}
\]
There is no remaining hard quadratic from this subtraction. This row
is necessary: without it the scalar $nu^2$ lies outside the old raw
square ideal, even if $n/k^4$ is infinitesimal.

\paragraph{Exact bounded stress after the positive square terms.}
Let $A_{*,1}$ be the inherited bounded centered-$A$ cap. In the actual
parent retain $(A_2-\kappa A_{*,1})D_uD_v(Y-huW)^2$, which is positive
because $A_2$ is unbounded and $\kappa A_{*,1}$ bounded. Comparing with
the corresponding retained positive square term of $\kappa F_1$ replaces the whole
$A$ difference by
\[
 \delta(Ah)D_uD_v(2uWY-hu^2W^2).
\]
After retaining TL4, the full difference of the bounded remainders is
exactly
\[
\begin{split}
 \mathscr E={}&\bigl[nD_uD_vX^2-\mathcal N^2\bigr]
   -nD_v(X+uZ)^2+\delta C D_v(uW-Y)^2\\
 &+(\delta D+\delta d/2)R_{D,\mathrm{event}}^2
   +\delta\alpha D_u(vY-Z)^2\\
 &+(\delta d/2)D_u\{(vY-Z)^2-D_vZ^2\}\\
 &+\delta(Ah)D_uD_v(2uWY-hu^2W^2),
\end{split}
\tag{TL6}
\]
where $R_{D,\mathrm{event}}=uW+uvX-Y-vZ$ is the original event
row. This follows from TL3; in particular the paired beta difference
is exactly $-n$. Every scalar, first-normal and hard-normal term is
present. The full expression and its first derivatives vanish at the
unchanged source contact, because both full parents and every retained
square do.

\paragraph{Finite support and bounds for the unchanged conductor.}
Work on the local chart $Z=1$, with $u,v,\mathfrak w$ positive
infinitesimal. At zero normals set
\[
 V_0=1-puv\mathfrak w,\quad q_B=D_uvV_0,\quad
 uW=u\mathfrak w+u^2h_1vV_0,\quad
 Y=u\mathfrak w-h_1vV_0,\quad h_1=a_1/T_1.
\tag{TL7}
\]
These are exact inherited state inverses. The rows $uW-Y$ and
$R_{D,\mathrm{event}}$ are in $(v)$ at zero normals;
$X+u=pD_uA_t$. The centered $A$ expression in TL6 is a quadratic
in $(v,u\mathfrak w)$ with bounded coefficients. The signed-core
line has no constant: $(vY-1)^2-D_v$ is explicitly
$(vY)^2-2vY-v^2$. The sole state-constant term is
$\delta\alpha D_u$. Finally TL5 is a combination of
$uA_t=(v)(u\mathfrak w)$ and $A_t^2$, with first normal in $(A_t)$.
Thus this is an explicit raw allocation, not an ideal-membership guess.

For the unit-$e_{1,\mathrm{scale}}$ chart use
\[
 (A,B_c,C_c,D_0)
 =(\mathfrak wv,\delta v,\delta E u\mathfrak w,\delta^2k^3E).
\]
Every scalar coefficient of $\mathscr E/\rho$ in
$(A,B_c,C_c)^2+(D_0^2)$ is bounded with unit denominators. Both complete
first normals belong to $(A,B_c,C_c)$ with bounded coefficients, and
the hard-normal entries are bounded. The decisive source coquotients
are
\[
 \Theta/\rho=\delta^2QEw^2/K,\quad
 n\gamma^2/\rho=w^2r^2Q^2/K^2,\quad
 nx/(\rho D_0^2)=Q^2w^2/K^2.
\tag{TL8}
\]
The remaining $n\gamma$ coquotient has an additional $\delta^2k^2E$.
All apparent negative $k$ powers have therefore canceled before the
conductor is used.

For the infinitesimal-$e_{1,\mathrm{scale}}$ chart use the exact SE
parameters $w=t_1=\varepsilon^2\tau$,
$\zeta_1=j^2QE/K$, $k=\varepsilon j$, and
\[
 (A,B_c,C_c,D_0)
 =(\varepsilon\mathfrak wv,v,Eu\mathfrak w,\varepsilon^2j^3E).
\]
The full raw scalar allocation of $\mathscr E/\rho$ is
\[
 f_AA^2+f_{BB}B_c^2+f_{BC}B_cC_c+f_{CC}C_c^2+f_0D_0^2.
\tag{TL9}
\]
All five coefficients are bounded; after division by $\rho$ each has a
factor $\varepsilon^2$. The $A^2$ coefficient has this factor by TL10;
all its other terms carry $p=\varepsilon^2\tau/R_1$, because
$\delta m=0$. The coefficients of $B_c^2,B_cC_c,C_c^2$ have it
from $\Theta/\rho$, and the constant coefficient from the final
identity in TL10. In particular
\[
 \Theta/\rho=\varepsilon^2\tau^2QE/K,
\]
\[
 \frac{n\gamma^2}{\rho\varepsilon^2}
 =\frac{\varepsilon^2\tau^2r^2Q^2}{K^2},\quad
 \frac{n\gamma}{\rho\varepsilon}
 =\frac{\varepsilon^3j^2E\tau^2rQ^2}{K^2},\quad
 \frac{nx}{\rho D_0^2}
 =\frac{\varepsilon^2\tau^2Q^2}{K^2}.
\tag{TL10}
\]
The first-normal coefficients are in $(A,B_c,C_c)$ with bounded
arrays, also divisible by $\varepsilon^2$ after division by $\rho$;
the new $A$ coefficient has the stronger $\varepsilon^3j^2$ factor
printed in TL10. The extra term $-nD_v(X+u)^2$ has scalar coefficient
$-np^2D_u^2D_v$ on $A_t^2$, first normal
$-2npD_uD_vA_tN$, and hard term $-nD_vN^2$; TL10 and
$p=\varepsilon^2\tau/R_1$ bound each normalized coefficient.

Exactly as in SE8, all dependence of the raw coefficients on
$\mathfrak w$ comes through $V_0$ and has a factor
$p=O(\varepsilon^2)$. The additional TL5 coefficients are independent
of $\mathfrak w$. Modulo $\varepsilon$, scalar raw coefficients are
even in $u,v$; the normal raw coefficients are independent of
$\mathfrak w$. Clearing powers of $D_u$ and the source unit
denominators preserves these properties. These are precisely the
finite divisibility premises of SE7: the first scalar conductor has
$T_{\mathfrak w}\in(\varepsilon)$, the second-stage remainders have
$T_{\mathfrak w}\in(\varepsilon,v)$, and each full normal has
$T_{\mathfrak w}\in(\varepsilon)$.

For either chart, TL6 has exact contact value and first-jet zero.
Apply the inherited monic contact lift to $\mathscr E/\rho$:
CC3--CC6 in the first chart and the division-free SE7 lift in the
second. The raw support restriction retains only $D_0^2$, not arbitrary
mixed $D_0$ products. The proofs use only finite monic division and
bounded unit denominators; extending their coefficient ring by the
bounded source letters in TL8--TL10 changes nothing. They give a
bounded Gram matrix on the same eight inherited roots for
$\mathscr E/\rho$. Hence the exact stress $\mathscr E$ has an
infinitesimal Gram matrix on that list. No inverse of $\rho$ occurs in
its final coefficients: division by $\rho$ here records the common
positive source factor before a linear finite lift.

Adding this infinitesimal matrix to the strict formal-C1 base Gram
preserves its positive residue. Restore all its retained positive square terms,
the modified positive centered-$A$ positive square term, and TL4. This proves the
full actual C2 parent on the thin-layer contact chart, with both
complete normals, all hard couplings and every independent rate.
The direct actual-C2 outer/state argument and the all-R10 remote
argument are required for the all-tests join; no negative formal-beta
positive square term is used in either join.

\subsection{The critical wall for unit $W$}
\label{component:pf1aac1dff7c7}

\paragraph{Statement.}
On the entire original R10 source, every independent test with $u,v$
infinitesimal and normalized $W$ a testing unit is tensor SOS. This
includes bounded $Z$ of every sign, infinitesimal $Z$, and $Z=0$.
The theorem uses the actual physical C2 coefficients. No continuation
of a negative-beta C1 source is involved.

Use the R10 source scale $s=1/a$, $w=ae$, $b=a^2B$, and put
$z=b/w$. The source constraint proves $z$ bounded. Write $A_h=Ah$.
The source estimates give $A_h,C,D$ positive units, $h$ infinitesimal,
$m/w$ a positive unit, and $t_0>1$, $w_0<1$.
Define the positive pure-source critical coefficient
\[
 C_*={ (h+2A_h)(1+2D)\over 2(4A_hD-h)}
     ={(1+2A)(1+2D)\over2(4AD-1)}.
 \tag{UW1}
\]
Its denominator after the first normalization is a positive unit.
We prove the opening-uniform estimate
\[
                 {C_*-C\over m}=s\,R,
                 \qquad R\text{ bounded}.             \tag{UW2}
\]

\paragraph{A short polynomial proof of the source estimate.}
All the following letters are polynomials in bounded source variables
$(s,w,z,h,t)$ after $b=wz$, with $x=sb$, $B=s^2b$, $r=1-sw$:
\[
 c_a=1+xr,\quad S_c=1+B(1-h)+sbr,\quad
 g=1-w+hx,\quad J=t c_a-(1-h),\quad \Delta_*=s\Delta,
\]
\[
 M_*=(1+x)(s+r)(h+t)-sw(1+tB)(1-h).
\]
Use the exact source numerators and denominators
\[
 \begin{array}{ll}
 N_C=c_aM_*,&D_C=2(1-h)\Delta_*,\\
 N_A=t c_a(s+r)S_c,&D_A=2\Delta_*,\\
 N_D=(1+x)(s+r)(1-txr)S_c,&D_D=2J\Delta_*,\\
 N_m=w(1+tB)c_aS_c,&D_m=2g\Delta_*.
 \end{array}
\]
Thus $C=N_C/D_C$, $A_h=N_A/D_A$, $D=N_D/D_D$,
and $m=N_m/D_m$. Clear the denominator in UW1:
\[
 \begin{split}
 P={}&(hD_A+2N_A)(D_D+2N_D)D_C\\
     &-2N_C(4N_AN_D-hD_AD_D),\\
 Q={}&2(4N_AN_D-hD_AD_D)D_C.
 \end{split}                                                  \tag{UW3}
\]
Then $C_*-C=P/Q$. At $s=0$ one has
\[
 C={t+h\over2(1-h)},\quad A_h=t/2,\quad
 D={1\over2(t-1+h)},
\]
and substitution gives $C_*=C$ identically. At $w=0$ (hence
$b=0$), the same three formulas hold for every $s$, because all
factors $1+s$ cancel. Therefore the polynomial $P$ is divisible by
both $s$ and $w$, hence $P=swP_1$. Since $N_m=wN_{m,1}$,
\[
 {C_*-C\over m}=s\,{P_1D_m\over QN_{m,1}}.
\]
The denominator $QN_{m,1}$ at $s=h=0$ is $16t>0$.
This proves UW2 by finite polynomial cancellation, without comparing
$w$ to $s$ or $h$.  Its displayed residue also gives
$R\sim t(1-w)(tz+2)$, though the proof only needs boundedness.

\paragraph{Complete parent and state allocation.}
Set $\epsilon=(C_*-C)_+$. By UW2, $\epsilon/m$ is
infinitesimal. If $C>C_*$, retain the nonnegative excess $C-C_*$.
Retain the entire original $n$ and $\alpha$ positive square terms; beta is zero.
The remaining signed core and $C_*,A,D$ positive square terms are exactly the
critical parent SEO11, with $c=C_*$ and $d=1$, because
\[
                  2C_*(4AD-1)=(1+2A)(1+2D).
\]
The exact critical-wall identities SEO12--SEO15 give a bounded strict
Gram on
\[
 B_t=X+uZ,\quad H_A=\sqrt A(Y-yV),\quad
 H_D=\sqrt{4D+2}\{uW-Y+uvB_t-(b_*+u^2)V\},
 \quad uV,\ vV,
\]
where $V=vZ$, $y=1/(2A)$ and
$b_*=(4AD-1)/(A(4D+2))$. Here $y,b_*$ are bounded;
$C_*,D$ are units and $A$ may be unbounded. The hard residue is
strict and the two radial masses have positive residues $2$ and
$2/t_0$. This is an application of those full polynomial identities,
not of their original C1 source theorem.

Normalize $W=1$ by a pure-testing factor. The entire $m$ event obeys
\[
 R_m=v(1+uY)-B_t,\qquad
 V={Z\over1+uY}(R_m+B_t).                         \tag{UW4}
\]
The denominator is a positive testing unit. From the displayed
$H_D$ root,
\[
 uW-Y=(b_*+u^2)V-uvB_t+H_D/\sqrt{4D+2}.
\]
Thus $\sqrt\epsilon(uW-Y)$ is a linear combination of the strict
critical roots and $\sqrt mR_m$ with infinitesimal coefficients.
The only potentially small denominator is accounted for by
$\sqrt{\epsilon/m}$, infinitesimal by UW2. Consequently the complete
negative correction $-4\epsilon(1+v^2)(uW-Y)^2$ has an
infinitesimal bounded Gram. The $m$ positive square term supplies a strict additional
pivot. Restore all retained original positive square terms to obtain the entire
C2 parent, including every hard state direction.

\paragraph{State entry for the remaining local work.}
Normalize a real testing state by a maximal coordinate. On the
negative pure-testing-core branch at $u,v$ infinitesimal, $X,Y$
are infinitesimal. If $W$ is a unit the theorem applies. Otherwise
$Z$ is a unit; set $Z=1$, and then
$w_{\rm test}=(W+uY)/(1+u^2)$ is infinitesimal. The inherited exact
Schur sign identity gives the positive contact chart. Literal axes
have the original event-square exits; complex states split into real
and imaginary parts. Together with the R10 remote theorem, the only
remaining task is the contact chart $u,v,w_{\rm test}>0$
infinitesimal, $Z=1$, on any source not already proved locally.

\label{component:p78dfeedee062}
\label{p78dfeedee062:sec:r17-r4-eunit}
\label{component:p1c9902e9a6f2}
\label{p1c9902e9a6f2:sec:r18-r3-eunit}
\label{component:p8da9ea6ed9e5}
\label{p8da9ea6ed9e5:sec:r17-r5}
\label{component:pd5fa361b4623}
\label{pd5fa361b4623:sec:r17-r6-bounded-a}
\label{component:p9a223c96c696}
\label{p9a223c96c696:sec:r17-r7-bounded-a}
\label{component:p482fcad1a2e2}
\label{p482fcad1a2e2:sec:r17-r8-small-r}
\label{component:p0d4eedd6c7a4}
\label{p0d4eedd6c7a4:sec:r9-weak-A-remote}
\label{component:p99372d145f17}
\label{p99372d145f17:sec:r11-small-r-local}
\label{component:p34609a6519a6}
\label{p34609a6519a6:sec:r11-small-r-remote}
